\documentclass[12pt, ukrainian]{article}
\setlength\textwidth{170mm} \setlength\hoffset{-25mm}
%\setlength\textheight{277mm} \setlength\voffset{-38mm}
%\setlength\parindent{0mm}
\pagestyle{empty}
\usepackage[T2A]{fontenc}
\usepackage[utf8]{inputenc}
\usepackage[ukrainian]{babel}
\usepackage{graphicx}
\usepackage{caption}
\DeclareCaptionFormat{nocaption}{}
\captionsetup{format=nocaption,aboveskip=0pt,belowskip=0pt}
\usepackage[Export]{adjustbox}
\adjustboxset{max size={0.9\linewidth}{0.9\paperheight}}
\begin{document}
{\Large{17.10.2024 Бета-версія: Контрольна робота \No 1, варіант  \No \textbf { 1 }\par}}
{
\begin {large}
{\begin {center}\textbf{Завдання 1}\end {center}}

Що буде виведено за виконання? (повідомлення інтерпретатора про помилку позначати error)\par
\begin{verbatim}
print(9//8*6*3)
\end{verbatim}

{\begin {center}\textbf{Завдання 2}\end {center}}

Що буде виведено за виконання? (повідомлення інтерпретатора про помилку позначати error)\par
\begin{verbatim}
print(3**-1+-1)
\end{verbatim}

{\begin {center}\textbf{Завдання 3}\end {center}}

Що буде виведено за виконання? (повідомлення інтерпретатора про помилку позначати error)\par
\begin{verbatim}
print(not True!=9.0 or 3>4 and 8>=3.0)
\end{verbatim}

{\begin {center}\textbf{Завдання 4}\end {center}}

Що буде виведено за виконання? (повідомлення інтерпретатора про помилку позначати error)\par
\begin{verbatim}
print(2.0 <= "4.0")
\end{verbatim}

{\begin {center}\textbf{Завдання 5}\end {center}}

Що буде виведено за виконання? (повідомлення інтерпретатора про помилку позначати error)\par
\begin{verbatim}
print(2 == 5.0 != 4.0 <= 9)
\end{verbatim}

{\begin {center}\textbf{Завдання 6}\end {center}}

Що буде виведено за виконання? (повідомлення інтерпретатора про помилку позначати error)\par
\begin{verbatim}

def f(n):
    print("f", end="");
    return n

print(f(8) or f(-9))
\end{verbatim}

{\begin {center}\textbf{Завдання 7}\end {center}}

Що буде виведено за виконання (повідомлення інтерпретатора про помилку позначати error)?\par
\begin{verbatim}
def g(a,b):
    c=29
    if b>=-2:
        return 6
    if b<=-5:
         c=0
    else: 
        c=7
    return c

print(g(-8,4))
\end{verbatim}

{\begin {center}\textbf{Завдання 8}\end {center}}

Що буде виведено за виконання (повідомлення інтерпретатора про помилку позначати error)?\par
\begin{verbatim}
a, b, c = 6, 7, 9
def g(a, b=8, c):
    print(a, b, c, end=" ")

g(b=2, c=0, 3)
print(a, b, c)
\end{verbatim}

{\begin {center}\textbf{Завдання 9}\end {center}}

Що буде виведено за виконання (повідомлення інтерпретатора про помилку позначати error)?\par
\begin{verbatim}
a,b,c=9,6,2
def g(a):
    a*=2
    b=1
    c=4
    return a+b+c

a,b,c=7,1,4
print(g(a),a,b,c)
\end{verbatim}

{\begin {center}\textbf{Завдання 10}\end {center}}


Написати функцiю f, що отримує щось та повертає щось. Невелика загадка все-таки має залишатися. Але нічого принципово нового відносно задач, що наявні у pdf-ках не планується.. Стиль запису умови також оцінюється.\par
\end {large}}
\newpage

{\Large{17.10.2024 Бета-версія: Контрольна робота \No 1, варіант  \No \textbf { 2 }\par}}
{
\begin {large}
{\begin {center}\textbf{Завдання 1}\end {center}}

Що буде виведено за виконання? (повідомлення інтерпретатора про помилку позначати error)\par
\begin{verbatim}
print(5/9/7*8)
\end{verbatim}

{\begin {center}\textbf{Завдання 2}\end {center}}

Що буде виведено за виконання? (повідомлення інтерпретатора про помилку позначати error)\par
\begin{verbatim}
print(1**-0.5*2)
\end{verbatim}

{\begin {center}\textbf{Завдання 3}\end {center}}

Що буде виведено за виконання? (повідомлення інтерпретатора про помилку позначати error)\par
\begin{verbatim}
print(not 5<=7 and 6<4.0 or False==6.0)
\end{verbatim}

{\begin {center}\textbf{Завдання 4}\end {center}}

Що буде виведено за виконання? (повідомлення інтерпретатора про помилку позначати error)\par
\begin{verbatim}
print("True" > 7)
\end{verbatim}

{\begin {center}\textbf{Завдання 5}\end {center}}

Що буде виведено за виконання? (повідомлення інтерпретатора про помилку позначати error)\par
\begin{verbatim}
print(True < 5 > 6.0 >= 4)
\end{verbatim}

{\begin {center}\textbf{Завдання 6}\end {center}}

Що буде виведено за виконання? (повідомлення інтерпретатора про помилку позначати error)\par
\begin{verbatim}

def f(n):
    print("f", end="");
    return n>-2

print(f(-6) and f(3))
\end{verbatim}

{\begin {center}\textbf{Завдання 7}\end {center}}

Що буде виведено за виконання (повідомлення інтерпретатора про помилку позначати error)?\par
\begin{verbatim}
def g(a):
    x=25
    if a: 
        x=4
    if a!=-5:
         return 7
    else:
         x=5
    return x

print(g(-1))
\end{verbatim}

{\begin {center}\textbf{Завдання 8}\end {center}}

Що буде виведено за виконання (повідомлення інтерпретатора про помилку позначати error)?\par
\begin{verbatim}
a, b, c = 7, 9, 8
def f(a, b=6, c=7):
    print(a, b, c, end=" ")

f(a=4, c=5, b=1)
print(a, b, c)
\end{verbatim}

{\begin {center}\textbf{Завдання 9}\end {center}}

Що буде виведено за виконання (повідомлення інтерпретатора про помилку позначати error)?\par
\begin{verbatim}
a,b,c=8,5,0
def f(b):
    a+=5
    b=3
    c=5
    return a+b+c

a,b,c=3,9,7
print(f(a),a,b,c)
\end{verbatim}

{\begin {center}\textbf{Завдання 10}\end {center}}


Написати функцiю f, що отримує щось та повертає щось. Невелика загадка все-таки має залишатися. Але нічого принципово нового відносно задач, що наявні у pdf-ках не планується.. Стиль запису умови також оцінюється.\par
\end {large}}
\newpage

{\Large{17.10.2024 Бета-версія: Контрольна робота \No 1, варіант  \No \textbf { 3 }\par}}
{
\begin {large}
{\begin {center}\textbf{Завдання 1}\end {center}}

Що буде виведено за виконання? (повідомлення інтерпретатора про помилку позначати error)\par
\begin{verbatim}
print(8//5%4*6)
\end{verbatim}

{\begin {center}\textbf{Завдання 2}\end {center}}

Що буде виведено за виконання? (повідомлення інтерпретатора про помилку позначати error)\par
\begin{verbatim}
print(-1**2--2)
\end{verbatim}

{\begin {center}\textbf{Завдання 3}\end {center}}

Що буде виведено за виконання? (повідомлення інтерпретатора про помилку позначати error)\par
\begin{verbatim}
print(5.0<=2 and not 7.0!=False or 9<7)
\end{verbatim}

{\begin {center}\textbf{Завдання 4}\end {center}}

Що буде виведено за виконання? (повідомлення інтерпретатора про помилку позначати error)\par
\begin{verbatim}
print(False == "8j")
\end{verbatim}

{\begin {center}\textbf{Завдання 5}\end {center}}

Що буде виведено за виконання? (повідомлення інтерпретатора про помилку позначати error)\par
\begin{verbatim}
print(3  is  7.0 > False <= 6)
\end{verbatim}

{\begin {center}\textbf{Завдання 6}\end {center}}

Що буде виведено за виконання? (повідомлення інтерпретатора про помилку позначати error)\par
\begin{verbatim}

def f(n):
    print("f", end="");
    return n!=1

print(f(-8) or f(-1))
\end{verbatim}

{\begin {center}\textbf{Завдання 7}\end {center}}

Що буде виведено за виконання (повідомлення інтерпретатора про помилку позначати error)?\par
\begin{verbatim}
def h(a,b):
    c=56
    if a==-3:
        return 1
    elif b<2:
         c=5
    else: 
        return 9
    return c

print(h(3,3))
\end{verbatim}

{\begin {center}\textbf{Завдання 8}\end {center}}

Що буде виведено за виконання (повідомлення інтерпретатора про помилку позначати error)?\par
\begin{verbatim}
a, b, c = 9, 8, 6
def h(a, b, c=8):
    print(a, b, c, end=" ")

h(4, 5)
print(a, b, c)
\end{verbatim}

{\begin {center}\textbf{Завдання 9}\end {center}}

Що буде виведено за виконання (повідомлення інтерпретатора про помилку позначати error)?\par
\begin{verbatim}
a,b,c=8,1,0
def h(b):
    a=4
    b-=1
    c=2
    return a+b+c

a,b,c=4,3,2
print(h(b),a,b,c)
\end{verbatim}

{\begin {center}\textbf{Завдання 10}\end {center}}


Написати функцiю f, що отримує щось та повертає щось. Невелика загадка все-таки має залишатися. Але нічого принципово нового відносно задач, що наявні у pdf-ках не планується.. Стиль запису умови також оцінюється.\par
\end {large}}
\newpage

{\Large{17.10.2024 Бета-версія: Контрольна робота \No 1, варіант  \No \textbf { 4 }\par}}
{
\begin {large}
{\begin {center}\textbf{Завдання 1}\end {center}}

Що буде виведено за виконання? (повідомлення інтерпретатора про помилку позначати error)\par
\begin{verbatim}
print(4/7//3*9)
\end{verbatim}

{\begin {center}\textbf{Завдання 2}\end {center}}

Що буде виведено за виконання? (повідомлення інтерпретатора про помилку позначати error)\par
\begin{verbatim}
print(4**0.5*False)
\end{verbatim}

{\begin {center}\textbf{Завдання 3}\end {center}}

Що буде виведено за виконання? (повідомлення інтерпретатора про помилку позначати error)\par
\begin{verbatim}
print(not 2.0==6 or True>=2 and 8.0>9)
\end{verbatim}

{\begin {center}\textbf{Завдання 4}\end {center}}

Що буде виведено за виконання? (повідомлення інтерпретатора про помилку позначати error)\par
\begin{verbatim}
print(1.0 < 6)
\end{verbatim}

{\begin {center}\textbf{Завдання 5}\end {center}}

Що буде виведено за виконання? (повідомлення інтерпретатора про помилку позначати error)\par
\begin{verbatim}
print(9.0 != 7 == 8 >= True)
\end{verbatim}

{\begin {center}\textbf{Завдання 6}\end {center}}

Що буде виведено за виконання? (повідомлення інтерпретатора про помилку позначати error)\par
\begin{verbatim}

def f(n):
    print("f", end="");
    return n

print(f(-3) and f(0))
\end{verbatim}

{\begin {center}\textbf{Завдання 7}\end {center}}

Що буде виведено за виконання (повідомлення інтерпретатора про помилку позначати error)?\par
\begin{verbatim}
def h(c):
    w=10
    if c>=4: 
        return 3
    elif c<=-1:
         w=2
    else:
         return 8
    return w

print(h(-6))
\end{verbatim}

{\begin {center}\textbf{Завдання 8}\end {center}}

Що буде виведено за виконання (повідомлення інтерпретатора про помилку позначати error)?\par
\begin{verbatim}
a, b, c = 7, 6, 9
def f(a, b, c):
    print(a, b, c, end=" ")

f(3, 0, b=2)
print(a, b, c)
\end{verbatim}

{\begin {center}\textbf{Завдання 9}\end {center}}

Що буде виведено за виконання (повідомлення інтерпретатора про помилку позначати error)?\par
\begin{verbatim}
a,b,c=4,1,5
def f(a):
    a=2
    b=1
    c=5
    return a+b+c

a,b,c=0,9,3
print(f(b),a,b,c)
\end{verbatim}

{\begin {center}\textbf{Завдання 10}\end {center}}


Написати функцiю f, що отримує щось та повертає щось. Невелика загадка все-таки має залишатися. Але нічого принципово нового відносно задач, що наявні у pdf-ках не планується.. Стиль запису умови також оцінюється.\par
\end {large}}
\newpage

{\Large{17.10.2024 Бета-версія: Контрольна робота \No 1, варіант  \No \textbf { 5 }\par}}
{
\begin {large}
{\begin {center}\textbf{Завдання 1}\end {center}}

Що буде виведено за виконання? (повідомлення інтерпретатора про помилку позначати error)\par
\begin{verbatim}
print(7//4*9*5)
\end{verbatim}

{\begin {center}\textbf{Завдання 2}\end {center}}

Що буде виведено за виконання? (повідомлення інтерпретатора про помилку позначати error)\par
\begin{verbatim}
print(4**4-True)
\end{verbatim}

{\begin {center}\textbf{Завдання 3}\end {center}}

Що буде виведено за виконання? (повідомлення інтерпретатора про помилку позначати error)\par
\begin{verbatim}
print(5.0>=3.0 or not 8>5 and 4!=True)
\end{verbatim}

{\begin {center}\textbf{Завдання 4}\end {center}}

Що буде виведено за виконання? (повідомлення інтерпретатора про помилку позначати error)\par
\begin{verbatim}
print(False != "8.0j")
\end{verbatim}

{\begin {center}\textbf{Завдання 5}\end {center}}

Що буде виведено за виконання? (повідомлення інтерпретатора про помилку позначати error)\par
\begin{verbatim}
print(6  is  8.0 < 2.0 >= 9)
\end{verbatim}

{\begin {center}\textbf{Завдання 6}\end {center}}

Що буде виведено за виконання? (повідомлення інтерпретатора про помилку позначати error)\par
\begin{verbatim}

def f(n):
    print("f", end="");
    return n<4

print(f(5) and f(4))
\end{verbatim}

{\begin {center}\textbf{Завдання 7}\end {center}}

Що буде виведено за виконання (повідомлення інтерпретатора про помилку позначати error)?\par
\begin{verbatim}
def f(b):
    y=56
    if b<0: 
        return 6
    elif b>5:
         y=0
    else:
         y=1
    return y

print(f(-5))
\end{verbatim}

{\begin {center}\textbf{Завдання 8}\end {center}}

Що буде виведено за виконання (повідомлення інтерпретатора про помилку позначати error)?\par
\begin{verbatim}
a, b, c = 8, 9, 7
def h(a, b=6, c):
    print(a, b, c, end=" ")

h(a=1, 3, c=1)
print(a, b, c)
\end{verbatim}

{\begin {center}\textbf{Завдання 9}\end {center}}

Що буде виведено за виконання (повідомлення інтерпретатора про помилку позначати error)?\par
\begin{verbatim}
a,b,c=6,5,2
def g(a):
    a=3
    b+=4
    c=1
    return a+b+c

a,b,c=7,8,6
print(g(b),a,b,c)
\end{verbatim}

{\begin {center}\textbf{Завдання 10}\end {center}}


Написати функцiю f, що отримує щось та повертає щось. Невелика загадка все-таки має залишатися. Але нічого принципово нового відносно задач, що наявні у pdf-ках не планується.. Стиль запису умови також оцінюється.\par
\end {large}}
\newpage

{\Large{17.10.2024 Бета-версія: Контрольна робота \No 1, варіант  \No \textbf { 6 }\par}}
{
\begin {large}
{\begin {center}\textbf{Завдання 1}\end {center}}

Що буде виведено за виконання? (повідомлення інтерпретатора про помилку позначати error)\par
\begin{verbatim}
print(6/6/5*4)
\end{verbatim}

{\begin {center}\textbf{Завдання 2}\end {center}}

Що буде виведено за виконання? (повідомлення інтерпретатора про помилку позначати error)\par
\begin{verbatim}
print(-3**-4+1)
\end{verbatim}

{\begin {center}\textbf{Завдання 3}\end {center}}

Що буде виведено за виконання? (повідомлення інтерпретатора про помилку позначати error)\par
\begin{verbatim}
print(not 3<4.0 and False<=8.0 or 8==9)
\end{verbatim}

{\begin {center}\textbf{Завдання 4}\end {center}}

Що буде виведено за виконання? (повідомлення інтерпретатора про помилку позначати error)\par
\begin{verbatim}
print("True" >= "2")
\end{verbatim}

{\begin {center}\textbf{Завдання 5}\end {center}}

Що буде виведено за виконання? (повідомлення інтерпретатора про помилку позначати error)\par
\begin{verbatim}
print(False == 2 > 4 != 3.0)
\end{verbatim}

{\begin {center}\textbf{Завдання 6}\end {center}}

Що буде виведено за виконання? (повідомлення інтерпретатора про помилку позначати error)\par
\begin{verbatim}

def f(n):
    print("f", end="");
    return n

print(f(6) or f(1))
\end{verbatim}

{\begin {center}\textbf{Завдання 7}\end {center}}

Що буде виведено за виконання (повідомлення інтерпретатора про помилку позначати error)?\par
\begin{verbatim}
def f(a,b):
    c=78
    if a:
        c=4
    if a!=0:
         return 3
    else: 
        return 8
    return c

print(f(9,9))
\end{verbatim}

{\begin {center}\textbf{Завдання 8}\end {center}}

Що буде виведено за виконання (повідомлення інтерпретатора про помилку позначати error)?\par
\begin{verbatim}
a, b, c = 6, 7, 9
def g(a, b=8, c=6):
    print(a, b, c, end=" ")

g(2, a=0)
print(a, b, c)
\end{verbatim}

{\begin {center}\textbf{Завдання 9}\end {center}}

Що буде виведено за виконання (повідомлення інтерпретатора про помилку позначати error)?\par
\begin{verbatim}
a,b,c=1,4,3
def f(b):
    a-=3
    b=5
    c=2
    return a+b+c

a,b,c=5,0,9
print(f(b),a,b,c)
\end{verbatim}

{\begin {center}\textbf{Завдання 10}\end {center}}


Написати функцiю f, що отримує щось та повертає щось. Невелика загадка все-таки має залишатися. Але нічого принципово нового відносно задач, що наявні у pdf-ках не планується.. Стиль запису умови також оцінюється.\par
\end {large}}
\newpage

{\Large{17.10.2024 Бета-версія: Контрольна робота \No 1, варіант  \No \textbf { 7 }\par}}
{
\begin {large}
{\begin {center}\textbf{Завдання 1}\end {center}}

Що буде виведено за виконання? (повідомлення інтерпретатора про помилку позначати error)\par
\begin{verbatim}
print(3//3//8*7)
\end{verbatim}

{\begin {center}\textbf{Завдання 2}\end {center}}

Що буде виведено за виконання? (повідомлення інтерпретатора про помилку позначати error)\par
\begin{verbatim}
print(9**-1.0*-2)
\end{verbatim}

{\begin {center}\textbf{Завдання 3}\end {center}}

Що буде виведено за виконання? (повідомлення інтерпретатора про помилку позначати error)\par
\begin{verbatim}
print(not False==9.0 and 2!=3 or 6.0<=5)
\end{verbatim}

{\begin {center}\textbf{Завдання 4}\end {center}}

Що буде виведено за виконання? (повідомлення інтерпретатора про помилку позначати error)\par
\begin{verbatim}
print(1 == "9.0j")
\end{verbatim}

{\begin {center}\textbf{Завдання 5}\end {center}}

Що буде виведено за виконання? (повідомлення інтерпретатора про помилку позначати error)\par
\begin{verbatim}
print(3.0 <= 2.0  is  5 < 3)
\end{verbatim}

{\begin {center}\textbf{Завдання 6}\end {center}}

Що буде виведено за виконання? (повідомлення інтерпретатора про помилку позначати error)\par
\begin{verbatim}

def f(n):
    print("f", end="");
    return n==0

print(f(-7) and f(2))
\end{verbatim}

{\begin {center}\textbf{Завдання 7}\end {center}}

Що буде виведено за виконання (повідомлення інтерпретатора про помилку позначати error)?\par
\begin{verbatim}
def f(d):
    v=37
    if d==2: 
        return 9
    if d>-4:
         v=6
    else:
         return 2
    return v

print(f(1))
\end{verbatim}

{\begin {center}\textbf{Завдання 8}\end {center}}

Що буде виведено за виконання (повідомлення інтерпретатора про помилку позначати error)?\par
\begin{verbatim}
a, b, c = 8, 9, 7
def f(a, b, c=8):
    print(a, b, c, end=" ")

f(4, c=5, b=1)
print(a, b, c)
\end{verbatim}

{\begin {center}\textbf{Завдання 9}\end {center}}

Що буде виведено за виконання (повідомлення інтерпретатора про помилку позначати error)?\par
\begin{verbatim}
a,b,c=8,6,7
def h(b):
    a=3
    b=4
    c=2
    return a+b+c

a,b,c=2,9,5
print(h(b),a,b,c)
\end{verbatim}

{\begin {center}\textbf{Завдання 10}\end {center}}


Написати функцiю f, що отримує щось та повертає щось. Невелика загадка все-таки має залишатися. Але нічого принципово нового відносно задач, що наявні у pdf-ках не планується.. Стиль запису умови також оцінюється.\par
\end {large}}
\newpage

{\Large{17.10.2024 Бета-версія: Контрольна робота \No 1, варіант  \No \textbf { 8 }\par}}
{
\begin {large}
{\begin {center}\textbf{Завдання 1}\end {center}}

Що буде виведено за виконання? (повідомлення інтерпретатора про помилку позначати error)\par
\begin{verbatim}
print(7/6%3*8)
\end{verbatim}

{\begin {center}\textbf{Завдання 2}\end {center}}

Що буде виведено за виконання? (повідомлення інтерпретатора про помилку позначати error)\par
\begin{verbatim}
print(9**1.0*False)
\end{verbatim}

{\begin {center}\textbf{Завдання 3}\end {center}}

Що буде виведено за виконання? (повідомлення інтерпретатора про помилку позначати error)\par
\begin{verbatim}
print(not 7.0>=6 or 4<2.0 and True>7)
\end{verbatim}

{\begin {center}\textbf{Завдання 4}\end {center}}

Що буде виведено за виконання? (повідомлення інтерпретатора про помилку позначати error)\par
\begin{verbatim}
print("4" != True)
\end{verbatim}

{\begin {center}\textbf{Завдання 5}\end {center}}

Що буде виведено за виконання? (повідомлення інтерпретатора про помилку позначати error)\par
\begin{verbatim}
print(False != 8 <= True < 7)
\end{verbatim}

{\begin {center}\textbf{Завдання 6}\end {center}}

Що буде виведено за виконання? (повідомлення інтерпретатора про помилку позначати error)\par
\begin{verbatim}

def f(n):
    print("f", end="");
    return n

print(f(-2) or f(9))
\end{verbatim}

{\begin {center}\textbf{Завдання 7}\end {center}}

Що буде виведено за виконання (повідомлення інтерпретатора про помилку позначати error)?\par
\begin{verbatim}
def f(a,b):
    c=52
    if b>5:
        c=2
    elif a!=-1:
         return 6
    else: 
        c=5
    return c

print(f(-1,-2))
\end{verbatim}

{\begin {center}\textbf{Завдання 8}\end {center}}

Що буде виведено за виконання (повідомлення інтерпретатора про помилку позначати error)?\par
\begin{verbatim}
a, b, c = 9, 7, 6
def h(a, b, c):
    print(a, b, c, end=" ")

h(4, 2, 3)
print(a, b, c)
\end{verbatim}

{\begin {center}\textbf{Завдання 9}\end {center}}

Що буде виведено за виконання (повідомлення інтерпретатора про помилку позначати error)?\par
\begin{verbatim}
a,b,c=1,2,3
def h(a):
    a*=4
    b=3
    c=5
    return a+b+c

a,b,c=4,5,8
print(h(a),a,b,c)
\end{verbatim}

{\begin {center}\textbf{Завдання 10}\end {center}}


Написати функцiю f, що отримує щось та повертає щось. Невелика загадка все-таки має залишатися. Але нічого принципово нового відносно задач, що наявні у pdf-ках не планується.. Стиль запису умови також оцінюється.\par
\end {large}}
\newpage

{\Large{17.10.2024 Бета-версія: Контрольна робота \No 1, варіант  \No \textbf { 9 }\par}}
{
\begin {large}
{\begin {center}\textbf{Завдання 1}\end {center}}

Що буде виведено за виконання? (повідомлення інтерпретатора про помилку позначати error)\par
\begin{verbatim}
print(5//4%7*5)
\end{verbatim}

{\begin {center}\textbf{Завдання 2}\end {center}}

Що буде виведено за виконання? (повідомлення інтерпретатора про помилку позначати error)\par
\begin{verbatim}
print(2**-2-True)
\end{verbatim}

{\begin {center}\textbf{Завдання 3}\end {center}}

Що буде виведено за виконання? (повідомлення інтерпретатора про помилку позначати error)\par
\begin{verbatim}
print(not True<6 or 8==9.0 and 4.0!=5)
\end{verbatim}

{\begin {center}\textbf{Завдання 4}\end {center}}

Що буде виведено за виконання? (повідомлення інтерпретатора про помилку позначати error)\par
\begin{verbatim}
print(7.0 <= "False")
\end{verbatim}

{\begin {center}\textbf{Завдання 5}\end {center}}

Що буде виведено за виконання? (повідомлення інтерпретатора про помилку позначати error)\par
\begin{verbatim}
print(9.0 >= 4  is  5 == 4.0)
\end{verbatim}

{\begin {center}\textbf{Завдання 6}\end {center}}

Що буде виведено за виконання? (повідомлення інтерпретатора про помилку позначати error)\par
\begin{verbatim}

def f(n):
    print("f", end="");
    return n

print(f(-4) and f(-5))
\end{verbatim}

{\begin {center}\textbf{Завдання 7}\end {center}}

Що буде виведено за виконання (повідомлення інтерпретатора про помилку позначати error)?\par
\begin{verbatim}
def h(a,b):
    c=65
    if b==1:
        return 8
    elif a<-4:
         c=4
    else: 
        c=0
    return c

print(h(4,-4))
\end{verbatim}

{\begin {center}\textbf{Завдання 8}\end {center}}

Що буде виведено за виконання (повідомлення інтерпретатора про помилку позначати error)?\par
\begin{verbatim}
a, b, c = 8, 7, 9
def g(a, b, c=6):
    print(a, b, c, end=" ")

g(a=5, 0, b=5)
print(a, b, c)
\end{verbatim}

{\begin {center}\textbf{Завдання 9}\end {center}}

Що буде виведено за виконання (повідомлення інтерпретатора про помилку позначати error)?\par
\begin{verbatim}
a,b,c=7,6,9
def f(a):
    a=2
    b-=1
    c=3
    return a+b+c

a,b,c=0,7,2
print(f(b),a,b,c)
\end{verbatim}

{\begin {center}\textbf{Завдання 10}\end {center}}


Написати функцiю f, що отримує щось та повертає щось. Невелика загадка все-таки має залишатися. Але нічого принципово нового відносно задач, що наявні у pdf-ках не планується.. Стиль запису умови також оцінюється.\par
\end {large}}
\newpage

{\Large{17.10.2024 Бета-версія: Контрольна робота \No 1, варіант  \No \textbf { 10 }\par}}
{
\begin {large}
{\begin {center}\textbf{Завдання 1}\end {center}}

Що буде виведено за виконання? (повідомлення інтерпретатора про помилку позначати error)\par
\begin{verbatim}
print(3/9*4*9)
\end{verbatim}

{\begin {center}\textbf{Завдання 2}\end {center}}

Що буде виведено за виконання? (повідомлення інтерпретатора про помилку позначати error)\par
\begin{verbatim}
print(3**False+-1)
\end{verbatim}

{\begin {center}\textbf{Завдання 3}\end {center}}

Що буде виведено за виконання? (повідомлення інтерпретатора про помилку позначати error)\par
\begin{verbatim}
print(not 7.0<=2 and 7>6.0 or False>=9)
\end{verbatim}

{\begin {center}\textbf{Завдання 4}\end {center}}

Що буде виведено за виконання? (повідомлення інтерпретатора про помилку позначати error)\par
\begin{verbatim}
print(True >= "False")
\end{verbatim}

{\begin {center}\textbf{Завдання 5}\end {center}}

Що буде виведено за виконання? (повідомлення інтерпретатора про помилку позначати error)\par
\begin{verbatim}
print(6.0 > 7  is  True >= 9)
\end{verbatim}

{\begin {center}\textbf{Завдання 6}\end {center}}

Що буде виведено за виконання? (повідомлення інтерпретатора про помилку позначати error)\par
\begin{verbatim}

def f(n):
    print("f", end="");
    return n<=-1

print(f(7) or f(-7))
\end{verbatim}

{\begin {center}\textbf{Завдання 7}\end {center}}

Що буде виведено за виконання (повідомлення інтерпретатора про помилку позначати error)?\par
\begin{verbatim}
def g(a,b):
    c=68
    if b:
        return 2
    if a>=4:
         return 7
    else: 
        c=1
    return c

print(g(8,-8))
\end{verbatim}

{\begin {center}\textbf{Завдання 8}\end {center}}

Що буде виведено за виконання (повідомлення інтерпретатора про помилку позначати error)?\par
\begin{verbatim}
a, b, c = 8, 7, 9
def h(a, b=6, c=9):
    print(a, b, c, end=" ")

h(2, 0, 1)
print(a, b, c)
\end{verbatim}

{\begin {center}\textbf{Завдання 9}\end {center}}

Що буде виведено за виконання (повідомлення інтерпретатора про помилку позначати error)?\par
\begin{verbatim}
a,b,c=4,6,0
def h(b):
    a+=2
    b=5
    c=1
    return a+b+c

a,b,c=8,1,5
print(h(a),a,b,c)
\end{verbatim}

{\begin {center}\textbf{Завдання 10}\end {center}}


Написати функцiю f, що отримує щось та повертає щось. Невелика загадка все-таки має залишатися. Але нічого принципово нового відносно задач, що наявні у pdf-ках не планується.. Стиль запису умови також оцінюється.\par
\end {large}}
\newpage

{\Large{17.10.2024 Бета-версія: Контрольна робота \No 1, варіант  \No \textbf { 11 }\par}}
{
\begin {large}
{\begin {center}\textbf{Завдання 1}\end {center}}

Що буде виведено за виконання? (повідомлення інтерпретатора про помилку позначати error)\par
\begin{verbatim}
print(8/3//8*3)
\end{verbatim}

{\begin {center}\textbf{Завдання 2}\end {center}}

Що буде виведено за виконання? (повідомлення інтерпретатора про помилку позначати error)\par
\begin{verbatim}
print(1**0.5*1)
\end{verbatim}

{\begin {center}\textbf{Завдання 3}\end {center}}

Що буде виведено за виконання? (повідомлення інтерпретатора про помилку позначати error)\par
\begin{verbatim}
print(not 3>False and 4!=5.0 or 8>=3.0)
\end{verbatim}

{\begin {center}\textbf{Завдання 4}\end {center}}

Що буде виведено за виконання? (повідомлення інтерпретатора про помилку позначати error)\par
\begin{verbatim}
print(3.0 < 9)
\end{verbatim}

{\begin {center}\textbf{Завдання 5}\end {center}}

Що буде виведено за виконання? (повідомлення інтерпретатора про помилку позначати error)\par
\begin{verbatim}
print(8.0 != 6 > 7.0 == 2)
\end{verbatim}

{\begin {center}\textbf{Завдання 6}\end {center}}

Що буде виведено за виконання? (повідомлення інтерпретатора про помилку позначати error)\par
\begin{verbatim}

def f(n):
    print("f", end="");
    return n>=2

print(f(-8) or f(0))
\end{verbatim}

{\begin {center}\textbf{Завдання 7}\end {center}}

Що буде виведено за виконання (повідомлення інтерпретатора про помилку позначати error)?\par
\begin{verbatim}
def h(b):
    u=46
    if b==1: 
        return 0
    if b<=3:
         u=7
    else:
         u=3
    return u

print(h(-9))
\end{verbatim}

{\begin {center}\textbf{Завдання 8}\end {center}}

Що буде виведено за виконання (повідомлення інтерпретатора про помилку позначати error)?\par
\begin{verbatim}
a, b, c = 8, 7, 6
def f(a, b, c):
    print(a, b, c, end=" ")

f(a=3, c=4, b=5)
print(a, b, c)
\end{verbatim}

{\begin {center}\textbf{Завдання 9}\end {center}}

Що буде виведено за виконання (повідомлення інтерпретатора про помилку позначати error)?\par
\begin{verbatim}
a,b,c=6,2,7
def g(b):
    a=4
    b=5
    c=1
    return a+b+c

a,b,c=8,3,0
print(g(b),a,b,c)
\end{verbatim}

{\begin {center}\textbf{Завдання 10}\end {center}}


Написати функцiю f, що отримує щось та повертає щось. Невелика загадка все-таки має залишатися. Але нічого принципово нового відносно задач, що наявні у pdf-ках не планується.. Стиль запису умови також оцінюється.\par
\end {large}}
\newpage

{\Large{17.10.2024 Бета-версія: Контрольна робота \No 1, варіант  \No \textbf { 12 }\par}}
{
\begin {large}
{\begin {center}\textbf{Завдання 1}\end {center}}

Що буде виведено за виконання? (повідомлення інтерпретатора про помилку позначати error)\par
\begin{verbatim}
print(6//7/6*6)
\end{verbatim}

{\begin {center}\textbf{Завдання 2}\end {center}}

Що буде виведено за виконання? (повідомлення інтерпретатора про помилку позначати error)\par
\begin{verbatim}
print(4**-1.0*2)
\end{verbatim}

{\begin {center}\textbf{Завдання 3}\end {center}}

Що буде виведено за виконання? (повідомлення інтерпретатора про помилку позначати error)\par
\begin{verbatim}
print(not 5<=True or 8.0<4 and 2.0==9)
\end{verbatim}

{\begin {center}\textbf{Завдання 4}\end {center}}

Що буде виведено за виконання? (повідомлення інтерпретатора про помилку позначати error)\par
\begin{verbatim}
print("5j" > "5.0")
\end{verbatim}

{\begin {center}\textbf{Завдання 5}\end {center}}

Що буде виведено за виконання? (повідомлення інтерпретатора про помилку позначати error)\par
\begin{verbatim}
print(3 < False <= 8 > 5.0)
\end{verbatim}

{\begin {center}\textbf{Завдання 6}\end {center}}

Що буде виведено за виконання? (повідомлення інтерпретатора про помилку позначати error)\par
\begin{verbatim}

def f(n):
    print("f", end="");
    return n

print(f(2) and f(8))
\end{verbatim}

{\begin {center}\textbf{Завдання 7}\end {center}}

Що буде виведено за виконання (повідомлення інтерпретатора про помилку позначати error)?\par
\begin{verbatim}
def g(d):
    z=74
    if d: 
        return 9
    elif d!=-2:
         return 4
    else:
         z=8
    return z

print(g(7))
\end{verbatim}

{\begin {center}\textbf{Завдання 8}\end {center}}

Що буде виведено за виконання (повідомлення інтерпретатора про помилку позначати error)?\par
\begin{verbatim}
a, b, c = 8, 7, 9
def g(a, b=6, c):
    print(a, b, c, end=" ")

g(b=2, c=1, 4)
print(a, b, c)
\end{verbatim}

{\begin {center}\textbf{Завдання 9}\end {center}}

Що буде виведено за виконання (повідомлення інтерпретатора про помилку позначати error)?\par
\begin{verbatim}
a,b,c=1,4,4
def h(a):
    a=3
    b=5
    c=3
    return a+b+c

a,b,c=6,9,7
print(h(a),a,b,c)
\end{verbatim}

{\begin {center}\textbf{Завдання 10}\end {center}}


Написати функцiю f, що отримує щось та повертає щось. Невелика загадка все-таки має залишатися. Але нічого принципово нового відносно задач, що наявні у pdf-ках не планується.. Стиль запису умови також оцінюється.\par
\end {large}}
\newpage

{\Large{17.10.2024 Бета-версія: Контрольна робота \No 1, варіант  \No \textbf { 13 }\par}}
{
\begin {large}
{\begin {center}\textbf{Завдання 1}\end {center}}

Що буде виведено за виконання? (повідомлення інтерпретатора про помилку позначати error)\par
\begin{verbatim}
print(9/5//5*4)
\end{verbatim}

{\begin {center}\textbf{Завдання 2}\end {center}}

Що буде виведено за виконання? (повідомлення інтерпретатора про помилку позначати error)\par
\begin{verbatim}
print(9**1.0*2)
\end{verbatim}

{\begin {center}\textbf{Завдання 3}\end {center}}

Що буде виведено за виконання? (повідомлення інтерпретатора про помилку позначати error)\par
\begin{verbatim}
print(not 2>2.0 or False<=3 or 6!=8.0)
\end{verbatim}

{\begin {center}\textbf{Завдання 4}\end {center}}

Що буде виведено за виконання? (повідомлення інтерпретатора про помилку позначати error)\par
\begin{verbatim}
print("False" == True)
\end{verbatim}

{\begin {center}\textbf{Завдання 5}\end {center}}

Що буде виведено за виконання? (повідомлення інтерпретатора про помилку позначати error)\par
\begin{verbatim}
print(3.0 >= 7.0 == 2 != 3)
\end{verbatim}

{\begin {center}\textbf{Завдання 6}\end {center}}

Що буде виведено за виконання? (повідомлення інтерпретатора про помилку позначати error)\par
\begin{verbatim}

def f(n):
    print("f", end="");
    return n<=-3

print(f(-2) and f(-4))
\end{verbatim}

{\begin {center}\textbf{Завдання 7}\end {center}}

Що буде виведено за виконання (повідомлення інтерпретатора про помилку позначати error)?\par
\begin{verbatim}
def g(c):
    w=82
    if c>=-3: 
        w=1
    if c<0:
         return 5
    else:
         return 1
    return w

print(g(9))
\end{verbatim}

{\begin {center}\textbf{Завдання 8}\end {center}}

Що буде виведено за виконання (повідомлення інтерпретатора про помилку позначати error)?\par
\begin{verbatim}
a, b, c = 9, 7, 8
def h(a, b, c=6):
    print(a, b, c, end=" ")

h(0, c=3, b=3)
print(a, b, c)
\end{verbatim}

{\begin {center}\textbf{Завдання 9}\end {center}}

Що буде виведено за виконання (повідомлення інтерпретатора про помилку позначати error)?\par
\begin{verbatim}
a,b,c=9,3,7
def g(b):
    a=4
    b*=2
    c=4
    return a+b+c

a,b,c=3,0,5
print(g(b),a,b,c)
\end{verbatim}

{\begin {center}\textbf{Завдання 10}\end {center}}


Написати функцiю f, що отримує щось та повертає щось. Невелика загадка все-таки має залишатися. Але нічого принципово нового відносно задач, що наявні у pdf-ках не планується.. Стиль запису умови також оцінюється.\par
\end {large}}
\newpage

{\Large{17.10.2024 Бета-версія: Контрольна робота \No 1, варіант  \No \textbf { 14 }\par}}
{
\begin {large}
{\begin {center}\textbf{Завдання 1}\end {center}}

Що буде виведено за виконання? (повідомлення інтерпретатора про помилку позначати error)\par
\begin{verbatim}
print(4//8/9*7)
\end{verbatim}

{\begin {center}\textbf{Завдання 2}\end {center}}

Що буде виведено за виконання? (повідомлення інтерпретатора про помилку позначати error)\par
\begin{verbatim}
print(-3**False--2)
\end{verbatim}

{\begin {center}\textbf{Завдання 3}\end {center}}

Що буде виведено за виконання? (повідомлення інтерпретатора про помилку позначати error)\par
\begin{verbatim}
print(7>=3 and not True==7 and 3.0<7.0)
\end{verbatim}

{\begin {center}\textbf{Завдання 4}\end {center}}

Що буде виведено за виконання? (повідомлення інтерпретатора про помилку позначати error)\par
\begin{verbatim}
print(6.0 != 3)
\end{verbatim}

{\begin {center}\textbf{Завдання 5}\end {center}}

Що буде виведено за виконання? (повідомлення інтерпретатора про помилку позначати error)\par
\begin{verbatim}
print(False <= 6  is  8.0 < 5)
\end{verbatim}

{\begin {center}\textbf{Завдання 6}\end {center}}

Що буде виведено за виконання? (повідомлення інтерпретатора про помилку позначати error)\par
\begin{verbatim}

def f(n):
    print("f", end="");
    return n

print(f(5) or f(-9))
\end{verbatim}

{\begin {center}\textbf{Завдання 7}\end {center}}

Що буде виведено за виконання (повідомлення інтерпретатора про помилку позначати error)?\par
\begin{verbatim}
def f(a,b):
    c=15
    if b<=3:
        c=3
    if a>-5:
         return 9
    else: 
        c=0
    return c

print(f(7,8))
\end{verbatim}

{\begin {center}\textbf{Завдання 8}\end {center}}

Що буде виведено за виконання (повідомлення інтерпретатора про помилку позначати error)?\par
\begin{verbatim}
a, b, c = 6, 9, 8
def g(a, b=7, c):
    print(a, b, c, end=" ")

g(4, 5, c=0)
print(a, b, c)
\end{verbatim}

{\begin {center}\textbf{Завдання 9}\end {center}}

Що буде виведено за виконання (повідомлення інтерпретатора про помилку позначати error)?\par
\begin{verbatim}
a,b,c=9,2,1
def g(a):
    a=3
    b+=5
    c=1
    return a+b+c

a,b,c=8,4,6
print(g(b),a,b,c)
\end{verbatim}

{\begin {center}\textbf{Завдання 10}\end {center}}


Написати функцiю f, що отримує щось та повертає щось. Невелика загадка все-таки має залишатися. Але нічого принципово нового відносно задач, що наявні у pdf-ках не планується.. Стиль запису умови також оцінюється.\par
\end {large}}
\newpage

{\Large{17.10.2024 Бета-версія: Контрольна робота \No 1, варіант  \No \textbf { 15 }\par}}
{
\begin {large}
{\begin {center}\textbf{Завдання 1}\end {center}}

Що буде виведено за виконання? (повідомлення інтерпретатора про помилку позначати error)\par
\begin{verbatim}
print(4//8*3*4)
\end{verbatim}

{\begin {center}\textbf{Завдання 2}\end {center}}

Що буде виведено за виконання? (повідомлення інтерпретатора про помилку позначати error)\par
\begin{verbatim}
print(-1**-2+1)
\end{verbatim}

{\begin {center}\textbf{Завдання 3}\end {center}}

Що буде виведено за виконання? (повідомлення інтерпретатора про помилку позначати error)\par
\begin{verbatim}
print(9>5.0 and not True>=6.0 and 4<=8)
\end{verbatim}

{\begin {center}\textbf{Завдання 4}\end {center}}

Що буде виведено за виконання? (повідомлення інтерпретатора про помилку позначати error)\par
\begin{verbatim}
print("4.0j" < "2j")
\end{verbatim}

{\begin {center}\textbf{Завдання 5}\end {center}}

Що буде виведено за виконання? (повідомлення інтерпретатора про помилку позначати error)\par
\begin{verbatim}
print(True  is  7 != 4 <= 6.0)
\end{verbatim}

{\begin {center}\textbf{Завдання 6}\end {center}}

Що буде виведено за виконання? (повідомлення інтерпретатора про помилку позначати error)\par
\begin{verbatim}

def f(n):
    print("f", end="");
    return n

print(f(9) or f(-5))
\end{verbatim}

{\begin {center}\textbf{Завдання 7}\end {center}}

Що буде виведено за виконання (повідомлення інтерпретатора про помилку позначати error)?\par
\begin{verbatim}
def f(a):
    u=60
    if a: 
        u=8
    elif a<5:
         u=6
    else:
         return 7
    return u

print(f(0))
\end{verbatim}

{\begin {center}\textbf{Завдання 8}\end {center}}

Що буде виведено за виконання (повідомлення інтерпретатора про помилку позначати error)?\par
\begin{verbatim}
a, b, c = 7, 8, 9
def f(a, b=6, c=8):
    print(a, b, c, end=" ")

f(1, a=2)
print(a, b, c)
\end{verbatim}

{\begin {center}\textbf{Завдання 9}\end {center}}

Що буде виведено за виконання (повідомлення інтерпретатора про помилку позначати error)?\par
\begin{verbatim}
a,b,c=0,7,2
def h(a):
    a=2
    b*=1
    c=4
    return a+b+c

a,b,c=4,1,5
print(h(a),a,b,c)
\end{verbatim}

{\begin {center}\textbf{Завдання 10}\end {center}}


Написати функцiю f, що отримує щось та повертає щось. Невелика загадка все-таки має залишатися. Але нічого принципово нового відносно задач, що наявні у pdf-ках не планується.. Стиль запису умови також оцінюється.\par
\end {large}}
\newpage

{\Large{17.10.2024 Бета-версія: Контрольна робота \No 1, варіант  \No \textbf { 16 }\par}}
{
\begin {large}
{\begin {center}\textbf{Завдання 1}\end {center}}

Що буде виведено за виконання? (повідомлення інтерпретатора про помилку позначати error)\par
\begin{verbatim}
print(3/6%8*9)
\end{verbatim}

{\begin {center}\textbf{Завдання 2}\end {center}}

Що буде виведено за виконання? (повідомлення інтерпретатора про помилку позначати error)\par
\begin{verbatim}
print(1**-0.5*-1)
\end{verbatim}

{\begin {center}\textbf{Завдання 3}\end {center}}

Що буде виведено за виконання? (повідомлення інтерпретатора про помилку позначати error)\par
\begin{verbatim}
print(False!=9.0 or not 2==4.0 or 5<6)
\end{verbatim}

{\begin {center}\textbf{Завдання 4}\end {center}}

Що буде виведено за виконання? (повідомлення інтерпретатора про помилку позначати error)\par
\begin{verbatim}
print("8.0" >= "9")
\end{verbatim}

{\begin {center}\textbf{Завдання 5}\end {center}}

Що буде виведено за виконання? (повідомлення інтерпретатора про помилку позначати error)\par
\begin{verbatim}
print(8 < False >= 9 == 3)
\end{verbatim}

{\begin {center}\textbf{Завдання 6}\end {center}}

Що буде виведено за виконання? (повідомлення інтерпретатора про помилку позначати error)\par
\begin{verbatim}

def f(n):
    print("f", end="");
    return n!=-4

print(f(3) and f(-1))
\end{verbatim}

{\begin {center}\textbf{Завдання 7}\end {center}}

Що буде виведено за виконання (повідомлення інтерпретатора про помилку позначати error)?\par
\begin{verbatim}
def g(a,b):
    c=36
    if a<=5:
        return 7
    elif a!=-1:
         return 6
    else: 
        c=3
    return c

print(g(1,1))
\end{verbatim}

{\begin {center}\textbf{Завдання 8}\end {center}}

Що буде виведено за виконання (повідомлення інтерпретатора про помилку позначати error)?\par
\begin{verbatim}
a, b, c = 9, 7, 6
def f(a, b, c):
    print(a, b, c, end=" ")

f(5, 1)
print(a, b, c)
\end{verbatim}

{\begin {center}\textbf{Завдання 9}\end {center}}

Що буде виведено за виконання (повідомлення інтерпретатора про помилку позначати error)?\par
\begin{verbatim}
a,b,c=3,9,8
def f(b):
    a=5
    b-=3
    c=1
    return a+b+c

a,b,c=6,6,2
print(f(a),a,b,c)
\end{verbatim}

{\begin {center}\textbf{Завдання 10}\end {center}}


Написати функцiю f, що отримує щось та повертає щось. Невелика загадка все-таки має залишатися. Але нічого принципово нового відносно задач, що наявні у pdf-ках не планується.. Стиль запису умови також оцінюється.\par
\end {large}}
\newpage

{\Large{17.10.2024 Бета-версія: Контрольна робота \No 1, варіант  \No \textbf { 17 }\par}}
{
\begin {large}
{\begin {center}\textbf{Завдання 1}\end {center}}

Що буде виведено за виконання? (повідомлення інтерпретатора про помилку позначати error)\par
\begin{verbatim}
print(8/9/9*3)
\end{verbatim}

{\begin {center}\textbf{Завдання 2}\end {center}}

Що буде виведено за виконання? (повідомлення інтерпретатора про помилку позначати error)\par
\begin{verbatim}
print(2**1.0+True)
\end{verbatim}

{\begin {center}\textbf{Завдання 3}\end {center}}

Що буде виведено за виконання? (повідомлення інтерпретатора про помилку позначати error)\par
\begin{verbatim}
print(6>=False and not 2.0>7.0 or 2<=7)
\end{verbatim}

{\begin {center}\textbf{Завдання 4}\end {center}}

Що буде виведено за виконання? (повідомлення інтерпретатора про помилку позначати error)\par
\begin{verbatim}
print(6 > 5.0)
\end{verbatim}

{\begin {center}\textbf{Завдання 5}\end {center}}

Що буде виведено за виконання? (повідомлення інтерпретатора про помилку позначати error)\par
\begin{verbatim}
print(2.0 > 9.0 >= 2 <= True)
\end{verbatim}

{\begin {center}\textbf{Завдання 6}\end {center}}

Що буде виведено за виконання? (повідомлення інтерпретатора про помилку позначати error)\par
\begin{verbatim}

def f(n):
    print("f", end="");
    return n>3

print(f(6) or f(7))
\end{verbatim}

{\begin {center}\textbf{Завдання 7}\end {center}}

Що буде виведено за виконання (повідомлення інтерпретатора про помилку позначати error)?\par
\begin{verbatim}
def h(a,b):
    c=98
    if b:
        c=9
    if b==-4:
         return 2
    else: 
        return 5
    return c

print(h(0,0))
\end{verbatim}

{\begin {center}\textbf{Завдання 8}\end {center}}

Що буде виведено за виконання (повідомлення інтерпретатора про помилку позначати error)?\par
\begin{verbatim}
a, b, c = 9, 7, 6
def g(a, b, c=8):
    print(a, b, c, end=" ")

g(4, c=3, b=2)
print(a, b, c)
\end{verbatim}

{\begin {center}\textbf{Завдання 9}\end {center}}

Що буде виведено за виконання (повідомлення інтерпретатора про помилку позначати error)?\par
\begin{verbatim}
a,b,c=1,2,8
def g(b):
    a=4
    b=2
    c=1
    return a+b+c

a,b,c=5,0,3
print(g(a),a,b,c)
\end{verbatim}

{\begin {center}\textbf{Завдання 10}\end {center}}


Написати функцiю f, що отримує щось та повертає щось. Невелика загадка все-таки має залишатися. Але нічого принципово нового відносно задач, що наявні у pdf-ках не планується.. Стиль запису умови також оцінюється.\par
\end {large}}
\newpage

{\Large{17.10.2024 Бета-версія: Контрольна робота \No 1, варіант  \No \textbf { 18 }\par}}
{
\begin {large}
{\begin {center}\textbf{Завдання 1}\end {center}}

Що буде виведено за виконання? (повідомлення інтерпретатора про помилку позначати error)\par
\begin{verbatim}
print(6//7%5*7)
\end{verbatim}

{\begin {center}\textbf{Завдання 2}\end {center}}

Що буде виведено за виконання? (повідомлення інтерпретатора про помилку позначати error)\par
\begin{verbatim}
print(4**-3-False)
\end{verbatim}

{\begin {center}\textbf{Завдання 3}\end {center}}

Що буде виведено за виконання? (повідомлення інтерпретатора про помилку позначати error)\par
\begin{verbatim}
print(not 4<5.0 or 3.0!=3 and True==9)
\end{verbatim}

{\begin {center}\textbf{Завдання 4}\end {center}}

Що буде виведено за виконання? (повідомлення інтерпретатора про помилку позначати error)\par
\begin{verbatim}
print("True" <= False)
\end{verbatim}

{\begin {center}\textbf{Завдання 5}\end {center}}

Що буде виведено за виконання? (повідомлення інтерпретатора про помилку позначати error)\par
\begin{verbatim}
print(7 > 8 != 5.0 == 4.0)
\end{verbatim}

{\begin {center}\textbf{Завдання 6}\end {center}}

Що буде виведено за виконання? (повідомлення інтерпретатора про помилку позначати error)\par
\begin{verbatim}

def f(n):
    print("f", end="");
    return n

print(f(-6) and f(1))
\end{verbatim}

{\begin {center}\textbf{Завдання 7}\end {center}}

Що буде виведено за виконання (повідомлення інтерпретатора про помилку позначати error)?\par
\begin{verbatim}
def h(c):
    z=11
    if c<=2: 
        z=9
    if c>=-5:
         return 3
    else:
         return 2
    return z

print(h(-8))
\end{verbatim}

{\begin {center}\textbf{Завдання 8}\end {center}}

Що буде виведено за виконання (повідомлення інтерпретатора про помилку позначати error)?\par
\begin{verbatim}
a, b, c = 9, 7, 6
def h(a, b=8, c=7):
    print(a, b, c, end=" ")

h(0, 0, a=1)
print(a, b, c)
\end{verbatim}

{\begin {center}\textbf{Завдання 9}\end {center}}

Що буде виведено за виконання (повідомлення інтерпретатора про помилку позначати error)?\par
\begin{verbatim}
a,b,c=4,9,0
def h(a):
    a-=5
    b=4
    c=2
    return a+b+c

a,b,c=8,7,1
print(h(b),a,b,c)
\end{verbatim}

{\begin {center}\textbf{Завдання 10}\end {center}}


Написати функцiю f, що отримує щось та повертає щось. Невелика загадка все-таки має залишатися. Але нічого принципово нового відносно задач, що наявні у pdf-ках не планується.. Стиль запису умови також оцінюється.\par
\end {large}}
\newpage

{\Large{17.10.2024 Бета-версія: Контрольна робота \No 1, варіант  \No \textbf { 19 }\par}}
{
\begin {large}
{\begin {center}\textbf{Завдання 1}\end {center}}

Що буде виведено за виконання? (повідомлення інтерпретатора про помилку позначати error)\par
\begin{verbatim}
print(9/4//6*8)
\end{verbatim}

{\begin {center}\textbf{Завдання 2}\end {center}}

Що буде виведено за виконання? (повідомлення інтерпретатора про помилку позначати error)\par
\begin{verbatim}
print(4**0.5*-2)
\end{verbatim}

{\begin {center}\textbf{Завдання 3}\end {center}}

Що буде виведено за виконання? (повідомлення інтерпретатора про помилку позначати error)\par
\begin{verbatim}
print(not 5<=8 or 9.0>True and 6.0!=8)
\end{verbatim}

{\begin {center}\textbf{Завдання 4}\end {center}}

Що буде виведено за виконання? (повідомлення інтерпретатора про помилку позначати error)\par
\begin{verbatim}
print(True > "5j")
\end{verbatim}

{\begin {center}\textbf{Завдання 5}\end {center}}

Що буде виведено за виконання? (повідомлення інтерпретатора про помилку позначати error)\par
\begin{verbatim}
print(3.0  is  5 < 6 >= True)
\end{verbatim}

{\begin {center}\textbf{Завдання 6}\end {center}}

Що буде виведено за виконання? (повідомлення інтерпретатора про помилку позначати error)\par
\begin{verbatim}

def f(n):
    print("f", end="");
    return n==3

print(f(4) or f(-3))
\end{verbatim}

{\begin {center}\textbf{Завдання 7}\end {center}}

Що буде виведено за виконання (повідомлення інтерпретатора про помилку позначати error)?\par
\begin{verbatim}
def f(d):
    x=42
    if d==-2: 
        x=5
    elif d!=-1:
         x=4
    else:
         return 0
    return x

print(f(-7))
\end{verbatim}

{\begin {center}\textbf{Завдання 8}\end {center}}

Що буде виведено за виконання (повідомлення інтерпретатора про помилку позначати error)?\par
\begin{verbatim}
a, b, c = 9, 8, 6
def f(a, b, c):
    print(a, b, c, end=" ")

f(4, 5)
print(a, b, c)
\end{verbatim}

{\begin {center}\textbf{Завдання 9}\end {center}}

Що буде виведено за виконання (повідомлення інтерпретатора про помилку позначати error)?\par
\begin{verbatim}
a,b,c=3,5,7
def f(b):
    a+=3
    b=5
    c=3
    return a+b+c

a,b,c=0,4,2
print(f(a),a,b,c)
\end{verbatim}

{\begin {center}\textbf{Завдання 10}\end {center}}


Написати функцiю f, що отримує щось та повертає щось. Невелика загадка все-таки має залишатися. Але нічого принципово нового відносно задач, що наявні у pdf-ках не планується.. Стиль запису умови також оцінюється.\par
\end {large}}
\newpage

{\Large{17.10.2024 Бета-версія: Контрольна робота \No 1, варіант  \No \textbf { 20 }\par}}
{
\begin {large}
{\begin {center}\textbf{Завдання 1}\end {center}}

Що буде виведено за виконання? (повідомлення інтерпретатора про помилку позначати error)\par
\begin{verbatim}
print(7//3*7*5)
\end{verbatim}

{\begin {center}\textbf{Завдання 2}\end {center}}

Що буде виведено за виконання? (повідомлення інтерпретатора про помилку позначати error)\par
\begin{verbatim}
print(-2**-2+-1)
\end{verbatim}

{\begin {center}\textbf{Завдання 3}\end {center}}

Що буде виведено за виконання? (повідомлення інтерпретатора про помилку позначати error)\par
\begin{verbatim}
print(not 6==7 and False>=4 or 8.0<4.0)
\end{verbatim}

{\begin {center}\textbf{Завдання 4}\end {center}}

Що буде виведено за виконання? (повідомлення інтерпретатора про помилку позначати error)\par
\begin{verbatim}
print("False" <= 3.0)
\end{verbatim}

{\begin {center}\textbf{Завдання 5}\end {center}}

Що буде виведено за виконання? (повідомлення інтерпретатора про помилку позначати error)\par
\begin{verbatim}
print(7.0  is  9 <= 5.0 == False)
\end{verbatim}

{\begin {center}\textbf{Завдання 6}\end {center}}

Що буде виведено за виконання? (повідомлення інтерпретатора про помилку позначати error)\par
\begin{verbatim}

def f(n):
    print("f", end="");
    return n

print(f(-8) and f(2))
\end{verbatim}

{\begin {center}\textbf{Завдання 7}\end {center}}

Що буде виведено за виконання (повідомлення інтерпретатора про помилку позначати error)?\par
\begin{verbatim}
def f(a,b):
    c=63
    if b>=-3:
        c=8
    elif b>1:
         c=4
    else: 
        return 1
    return c

print(f(6,-1))
\end{verbatim}

{\begin {center}\textbf{Завдання 8}\end {center}}

Що буде виведено за виконання (повідомлення інтерпретатора про помилку позначати error)?\par
\begin{verbatim}
a, b, c = 7, 8, 6
def h(a, b=9, c):
    print(a, b, c, end=" ")

h(b=3, c=2, 0)
print(a, b, c)
\end{verbatim}

{\begin {center}\textbf{Завдання 9}\end {center}}

Що буде виведено за виконання (повідомлення інтерпретатора про помилку позначати error)?\par
\begin{verbatim}
a,b,c=5,1,9
def g(a):
    a*=4
    b=1
    c=2
    return a+b+c

a,b,c=3,6,8
print(g(a),a,b,c)
\end{verbatim}

{\begin {center}\textbf{Завдання 10}\end {center}}


Написати функцiю f, що отримує щось та повертає щось. Невелика загадка все-таки має залишатися. Але нічого принципово нового відносно задач, що наявні у pdf-ках не планується.. Стиль запису умови також оцінюється.\par
\end {large}}
\newpage

{\Large{17.10.2024 Бета-версія: Контрольна робота \No 1, варіант  \No \textbf { 21 }\par}}
{
\begin {large}
{\begin {center}\textbf{Завдання 1}\end {center}}

Що буде виведено за виконання? (повідомлення інтерпретатора про помилку позначати error)\par
\begin{verbatim}
print(5//5//4*6)
\end{verbatim}

{\begin {center}\textbf{Завдання 2}\end {center}}

Що буде виведено за виконання? (повідомлення інтерпретатора про помилку позначати error)\par
\begin{verbatim}
print(1**-0.5*2)
\end{verbatim}

{\begin {center}\textbf{Завдання 3}\end {center}}

Що буде виведено за виконання? (повідомлення інтерпретатора про помилку позначати error)\par
\begin{verbatim}
print(6.0<=5 or not 9==3.0 and 3<True)
\end{verbatim}

{\begin {center}\textbf{Завдання 4}\end {center}}

Що буде виведено за виконання? (повідомлення інтерпретатора про помилку позначати error)\par
\begin{verbatim}
print("6.0j" >= 7)
\end{verbatim}

{\begin {center}\textbf{Завдання 5}\end {center}}

Що буде виведено за виконання? (повідомлення інтерпретатора про помилку позначати error)\par
\begin{verbatim}
print(4 > 8 != 3 < 2.0)
\end{verbatim}

{\begin {center}\textbf{Завдання 6}\end {center}}

Що буде виведено за виконання? (повідомлення інтерпретатора про помилку позначати error)\par
\begin{verbatim}

def f(n):
    print("f", end="");
    return n<2

print(f(-1) and f(5))
\end{verbatim}

{\begin {center}\textbf{Завдання 7}\end {center}}

Що буде виведено за виконання (повідомлення інтерпретатора про помилку позначати error)?\par
\begin{verbatim}
def h(a):
    v=38
    if a: 
        return 9
    if a>-3:
         v=4
    else:
         v=2
    return v

print(h(-4))
\end{verbatim}

{\begin {center}\textbf{Завдання 8}\end {center}}

Що буде виведено за виконання (повідомлення інтерпретатора про помилку позначати error)?\par
\begin{verbatim}
a, b, c = 9, 7, 8
def g(a, b=6, c):
    print(a, b, c, end=" ")

g(b=3, c=1, a=5)
print(a, b, c)
\end{verbatim}

{\begin {center}\textbf{Завдання 9}\end {center}}

Що буде виведено за виконання (повідомлення інтерпретатора про помилку позначати error)?\par
\begin{verbatim}
a,b,c=9,1,6
def f(a):
    a=2
    b=4
    c=1
    return a+b+c

a,b,c=4,3,7
print(f(a),a,b,c)
\end{verbatim}

{\begin {center}\textbf{Завдання 10}\end {center}}


Написати функцiю f, що отримує щось та повертає щось. Невелика загадка все-таки має залишатися. Але нічого принципово нового відносно задач, що наявні у pdf-ках не планується.. Стиль запису умови також оцінюється.\par
\end {large}}
\newpage

{\Large{17.10.2024 Бета-версія: Контрольна робота \No 1, варіант  \No \textbf { 22 }\par}}
{
\begin {large}
{\begin {center}\textbf{Завдання 1}\end {center}}

Що буде виведено за виконання? (повідомлення інтерпретатора про помилку позначати error)\par
\begin{verbatim}
print(3/6%4*7)
\end{verbatim}

{\begin {center}\textbf{Завдання 2}\end {center}}

Що буде виведено за виконання? (повідомлення інтерпретатора про помилку позначати error)\par
\begin{verbatim}
print(4**1.0*True)
\end{verbatim}

{\begin {center}\textbf{Завдання 3}\end {center}}

Що буде виведено за виконання? (повідомлення інтерпретатора про помилку позначати error)\par
\begin{verbatim}
print(4.0>=2 and not 5.0>False or 7!=3)
\end{verbatim}

{\begin {center}\textbf{Завдання 4}\end {center}}

Що буде виведено за виконання? (повідомлення інтерпретатора про помилку позначати error)\par
\begin{verbatim}
print("1" < 3)
\end{verbatim}

{\begin {center}\textbf{Завдання 5}\end {center}}

Що буде виведено за виконання? (повідомлення інтерпретатора про помилку позначати error)\par
\begin{verbatim}
print(4 <= 9 == True >= 6.0)
\end{verbatim}

{\begin {center}\textbf{Завдання 6}\end {center}}

Що буде виведено за виконання? (повідомлення інтерпретатора про помилку позначати error)\par
\begin{verbatim}

def f(n):
    print("f", end="");
    return n

print(f(-9) or f(8))
\end{verbatim}

{\begin {center}\textbf{Завдання 7}\end {center}}

Що буде виведено за виконання (повідомлення інтерпретатора про помилку позначати error)?\par
\begin{verbatim}
def h(a,b):
    c=48
    if a:
        c=7
    if b<4:
         return 6
    else: 
        return 2
    return c

print(h(-2,-6))
\end{verbatim}

{\begin {center}\textbf{Завдання 8}\end {center}}

Що буде виведено за виконання (повідомлення інтерпретатора про помилку позначати error)?\par
\begin{verbatim}
a, b, c = 8, 9, 7
def h(a, b=6, c=6):
    print(a, b, c, end=" ")

h(a=2, 4, c=0)
print(a, b, c)
\end{verbatim}

{\begin {center}\textbf{Завдання 9}\end {center}}

Що буде виведено за виконання (повідомлення інтерпретатора про помилку позначати error)?\par
\begin{verbatim}
a,b,c=3,2,9
def g(b):
    a=3
    b+=4
    c=5
    return a+b+c

a,b,c=6,4,7
print(g(b),a,b,c)
\end{verbatim}

{\begin {center}\textbf{Завдання 10}\end {center}}


Написати функцiю f, що отримує щось та повертає щось. Невелика загадка все-таки має залишатися. Але нічого принципово нового відносно задач, що наявні у pdf-ках не планується.. Стиль запису умови також оцінюється.\par
\end {large}}
\newpage

{\Large{17.10.2024 Бета-версія: Контрольна робота \No 1, варіант  \No \textbf { 23 }\par}}
{
\begin {large}
{\begin {center}\textbf{Завдання 1}\end {center}}

Що буде виведено за виконання? (повідомлення інтерпретатора про помилку позначати error)\par
\begin{verbatim}
print(6/4*3*6)
\end{verbatim}

{\begin {center}\textbf{Завдання 2}\end {center}}

Що буде виведено за виконання? (повідомлення інтерпретатора про помилку позначати error)\par
\begin{verbatim}
print(-2**True-1)
\end{verbatim}

{\begin {center}\textbf{Завдання 3}\end {center}}

Що буде виведено за виконання? (повідомлення інтерпретатора про помилку позначати error)\par
\begin{verbatim}
print(not 6>=True and 9.0<=5 and 2!=7.0)
\end{verbatim}

{\begin {center}\textbf{Завдання 4}\end {center}}

Що буде виведено за виконання? (повідомлення інтерпретатора про помилку позначати error)\par
\begin{verbatim}
print("2.0" == 9.0)
\end{verbatim}

{\begin {center}\textbf{Завдання 5}\end {center}}

Що буде виведено за виконання? (повідомлення інтерпретатора про помилку позначати error)\par
\begin{verbatim}
print(4.0  is  7 != 9.0 < False)
\end{verbatim}

{\begin {center}\textbf{Завдання 6}\end {center}}

Що буде виведено за виконання? (повідомлення інтерпретатора про помилку позначати error)\par
\begin{verbatim}

def f(n):
    print("f", end="");
    return n>=-2

print(f(0) and f(9))
\end{verbatim}

{\begin {center}\textbf{Завдання 7}\end {center}}

Що буде виведено за виконання (повідомлення інтерпретатора про помилку позначати error)?\par
\begin{verbatim}
def g(a,b):
    c=32
    if a>-2:
        c=4
    elif a<3:
         c=1
    else: 
        return 5
    return c

print(g(2,6))
\end{verbatim}

{\begin {center}\textbf{Завдання 8}\end {center}}

Що буде виведено за виконання (повідомлення інтерпретатора про помилку позначати error)?\par
\begin{verbatim}
a, b, c = 9, 7, 8
def f(a, b, c):
    print(a, b, c, end=" ")

f(5, b=2)
print(a, b, c)
\end{verbatim}

{\begin {center}\textbf{Завдання 9}\end {center}}

Що буде виведено за виконання (повідомлення інтерпретатора про помилку позначати error)?\par
\begin{verbatim}
a,b,c=8,5,0
def h(a):
    a=3
    b=5
    c=5
    return a+b+c

a,b,c=2,4,5
print(h(b),a,b,c)
\end{verbatim}

{\begin {center}\textbf{Завдання 10}\end {center}}


Написати функцiю f, що отримує щось та повертає щось. Невелика загадка все-таки має залишатися. Але нічого принципово нового відносно задач, що наявні у pdf-ках не планується.. Стиль запису умови також оцінюється.\par
\end {large}}
\newpage

{\Large{17.10.2024 Бета-версія: Контрольна робота \No 1, варіант  \No \textbf { 24 }\par}}
{
\begin {large}
{\begin {center}\textbf{Завдання 1}\end {center}}

Що буде виведено за виконання? (повідомлення інтерпретатора про помилку позначати error)\par
\begin{verbatim}
print(4//3/5*9)
\end{verbatim}

{\begin {center}\textbf{Завдання 2}\end {center}}

Що буде виведено за виконання? (повідомлення інтерпретатора про помилку позначати error)\par
\begin{verbatim}
print(9**-1.0*False)
\end{verbatim}

{\begin {center}\textbf{Завдання 3}\end {center}}

Що буде виведено за виконання? (повідомлення інтерпретатора про помилку позначати error)\par
\begin{verbatim}
print(not False<4 or 8.0>8 or 9==2.0)
\end{verbatim}

{\begin {center}\textbf{Завдання 4}\end {center}}

Що буде виведено за виконання? (повідомлення інтерпретатора про помилку позначати error)\par
\begin{verbatim}
print(True != "False")
\end{verbatim}

{\begin {center}\textbf{Завдання 5}\end {center}}

Що буде виведено за виконання? (повідомлення інтерпретатора про помилку позначати error)\par
\begin{verbatim}
print(5 > 2 < 6  is  8.0)
\end{verbatim}

{\begin {center}\textbf{Завдання 6}\end {center}}

Що буде виведено за виконання? (повідомлення інтерпретатора про помилку позначати error)\par
\begin{verbatim}

def f(n):
    print("f", end="");
    return n

print(f(-4) or f(4))
\end{verbatim}

{\begin {center}\textbf{Завдання 7}\end {center}}

Що буде виведено за виконання (повідомлення інтерпретатора про помилку позначати error)?\par
\begin{verbatim}
def g(b):
    y=41
    if b<=3: 
        return 7
    elif b<4:
         return 3
    else:
         y=8
    return y

print(g(8))
\end{verbatim}

{\begin {center}\textbf{Завдання 8}\end {center}}

Що буде виведено за виконання (повідомлення інтерпретатора про помилку позначати error)?\par
\begin{verbatim}
a, b, c = 7, 8, 6
def g(a, b, c=9):
    print(a, b, c, end=" ")

g(1, 4, 3)
print(a, b, c)
\end{verbatim}

{\begin {center}\textbf{Завдання 9}\end {center}}

Що буде виведено за виконання (повідомлення інтерпретатора про помилку позначати error)?\par
\begin{verbatim}
a,b,c=8,1,5
def h(a):
    a=1
    b*=2
    c=5
    return a+b+c

a,b,c=0,9,3
print(h(b),a,b,c)
\end{verbatim}

{\begin {center}\textbf{Завдання 10}\end {center}}


Написати функцiю f, що отримує щось та повертає щось. Невелика загадка все-таки має залишатися. Але нічого принципово нового відносно задач, що наявні у pdf-ках не планується.. Стиль запису умови також оцінюється.\par
\end {large}}
\newpage

{\Large{17.10.2024 Бета-версія: Контрольна робота \No 1, варіант  \No \textbf { 25 }\par}}
{
\begin {large}
{\begin {center}\textbf{Завдання 1}\end {center}}

Що буде виведено за виконання? (повідомлення інтерпретатора про помилку позначати error)\par
\begin{verbatim}
print(7//8%8*4)
\end{verbatim}

{\begin {center}\textbf{Завдання 2}\end {center}}

Що буде виведено за виконання? (повідомлення інтерпретатора про помилку позначати error)\par
\begin{verbatim}
print(1**-4+1)
\end{verbatim}

{\begin {center}\textbf{Завдання 3}\end {center}}

Що буде виведено за виконання? (повідомлення інтерпретатора про помилку позначати error)\par
\begin{verbatim}
print(not 5.0!=False and 6.0<9 and 6>2)
\end{verbatim}

{\begin {center}\textbf{Завдання 4}\end {center}}

Що буде виведено за виконання? (повідомлення інтерпретатора про помилку позначати error)\par
\begin{verbatim}
print("4j" == "False")
\end{verbatim}

{\begin {center}\textbf{Завдання 5}\end {center}}

Що буде виведено за виконання? (повідомлення інтерпретатора про помилку позначати error)\par
\begin{verbatim}
print(6 != True == 4 <= 3.0)
\end{verbatim}

{\begin {center}\textbf{Завдання 6}\end {center}}

Що буде виведено за виконання? (повідомлення інтерпретатора про помилку позначати error)\par
\begin{verbatim}

def f(n):
    print("f", end="");
    return n==1

print(f(3) or f(1))
\end{verbatim}

{\begin {center}\textbf{Завдання 7}\end {center}}

Що буде виведено за виконання (повідомлення інтерпретатора про помилку позначати error)?\par
\begin{verbatim}
def f(b):
    v=45
    if b>1: 
        return 0
    elif b==-4:
         v=1
    else:
         return 5
    return v

print(f(3))
\end{verbatim}

{\begin {center}\textbf{Завдання 8}\end {center}}

Що буде виведено за виконання (повідомлення інтерпретатора про помилку позначати error)?\par
\begin{verbatim}
a, b, c = 8, 6, 9
def f(a, b, c):
    print(a, b, c, end=" ")

f(4, 3)
print(a, b, c)
\end{verbatim}

{\begin {center}\textbf{Завдання 9}\end {center}}

Що буде виведено за виконання (повідомлення інтерпретатора про помилку позначати error)?\par
\begin{verbatim}
a,b,c=4,5,8
def f(b):
    a=3
    b-=4
    c=1
    return a+b+c

a,b,c=2,6,7
print(f(a),a,b,c)
\end{verbatim}

{\begin {center}\textbf{Завдання 10}\end {center}}


Написати функцiю f, що отримує щось та повертає щось. Невелика загадка все-таки має залишатися. Але нічого принципово нового відносно задач, що наявні у pdf-ках не планується.. Стиль запису умови також оцінюється.\par
\end {large}}
\newpage

{\Large{17.10.2024 Бета-версія: Контрольна робота \No 1, варіант  \No \textbf { 26 }\par}}
{
\begin {large}
{\begin {center}\textbf{Завдання 1}\end {center}}

Що буде виведено за виконання? (повідомлення інтерпретатора про помилку позначати error)\par
\begin{verbatim}
print(9/7/7*5)
\end{verbatim}

{\begin {center}\textbf{Завдання 2}\end {center}}

Що буде виведено за виконання? (повідомлення інтерпретатора про помилку позначати error)\par
\begin{verbatim}
print(1**0.5*-2)
\end{verbatim}

{\begin {center}\textbf{Завдання 3}\end {center}}

Що буде виведено за виконання? (повідомлення інтерпретатора про помилку позначати error)\par
\begin{verbatim}
print(not 5<=2.0 or 3==8 or 3.0>=True)
\end{verbatim}

{\begin {center}\textbf{Завдання 4}\end {center}}

Що буде виведено за виконання? (повідомлення інтерпретатора про помилку позначати error)\par
\begin{verbatim}
print("7.0j" > True)
\end{verbatim}

{\begin {center}\textbf{Завдання 5}\end {center}}

Що буде виведено за виконання? (повідомлення інтерпретатора про помилку позначати error)\par
\begin{verbatim}
print(2 > 9.0 >= False  is  8.0)
\end{verbatim}

{\begin {center}\textbf{Завдання 6}\end {center}}

Що буде виведено за виконання? (повідомлення інтерпретатора про помилку позначати error)\par
\begin{verbatim}

def f(n):
    print("f", end="");
    return n

print(f(-7) and f(-6))
\end{verbatim}

{\begin {center}\textbf{Завдання 7}\end {center}}

Що буде виведено за виконання (повідомлення інтерпретатора про помилку позначати error)?\par
\begin{verbatim}
def g(a):
    y=17
    if a: 
        y=6
    if a>=-3:
         y=5
    else:
         return 0
    return y

print(g(6))
\end{verbatim}

{\begin {center}\textbf{Завдання 8}\end {center}}

Що буде виведено за виконання (повідомлення інтерпретатора про помилку позначати error)?\par
\begin{verbatim}
a, b, c = 7, 9, 7
def g(a, b, c=6):
    print(a, b, c, end=" ")

g(a=2, 5, a=1)
print(a, b, c)
\end{verbatim}

{\begin {center}\textbf{Завдання 9}\end {center}}

Що буде виведено за виконання (повідомлення інтерпретатора про помилку позначати error)?\par
\begin{verbatim}
a,b,c=0,1,8
def h(a):
    a-=2
    b=5
    c=4
    return a+b+c

a,b,c=9,0,6
print(h(b),a,b,c)
\end{verbatim}

{\begin {center}\textbf{Завдання 10}\end {center}}


Написати функцiю f, що отримує щось та повертає щось. Невелика загадка все-таки має залишатися. Але нічого принципово нового відносно задач, що наявні у pdf-ках не планується.. Стиль запису умови також оцінюється.\par
\end {large}}
\newpage

{\Large{17.10.2024 Бета-версія: Контрольна робота \No 1, варіант  \No \textbf { 27 }\par}}
{
\begin {large}
{\begin {center}\textbf{Завдання 1}\end {center}}

Що буде виведено за виконання? (повідомлення інтерпретатора про помилку позначати error)\par
\begin{verbatim}
print(5/9*9*8)
\end{verbatim}

{\begin {center}\textbf{Завдання 2}\end {center}}

Що буде виведено за виконання? (повідомлення інтерпретатора про помилку позначати error)\par
\begin{verbatim}
print(3**True--1)
\end{verbatim}

{\begin {center}\textbf{Завдання 3}\end {center}}

Що буде виведено за виконання? (повідомлення інтерпретатора про помилку позначати error)\par
\begin{verbatim}
print(8.0<=7 or not 4>False and 3==7.0)
\end{verbatim}

{\begin {center}\textbf{Завдання 4}\end {center}}

Що буде виведено за виконання? (повідомлення інтерпретатора про помилку позначати error)\par
\begin{verbatim}
print(8 <= 1.0)
\end{verbatim}

{\begin {center}\textbf{Завдання 5}\end {center}}

Що буде виведено за виконання? (повідомлення інтерпретатора про помилку позначати error)\par
\begin{verbatim}
print(9 >= 5 != 7 < 2.0)
\end{verbatim}

{\begin {center}\textbf{Завдання 6}\end {center}}

Що буде виведено за виконання? (повідомлення інтерпретатора про помилку позначати error)\par
\begin{verbatim}

def f(n):
    print("f", end="");
    return n

print(f(6) or f(-3))
\end{verbatim}

{\begin {center}\textbf{Завдання 7}\end {center}}

Що буде виведено за виконання (повідомлення інтерпретатора про помилку позначати error)?\par
\begin{verbatim}
def g(a,b):
    c=39
    if b!=2:
        return 8
    if a>=0:
         c=9
    else: 
        return 0
    return c

print(g(-6,5))
\end{verbatim}

{\begin {center}\textbf{Завдання 8}\end {center}}

Що буде виведено за виконання (повідомлення інтерпретатора про помилку позначати error)?\par
\begin{verbatim}
a, b, c = 8, 8, 7
def h(a, b=6, c=9):
    print(a, b, c, end=" ")

h(b=0, c=1, 3)
print(a, b, c)
\end{verbatim}

{\begin {center}\textbf{Завдання 9}\end {center}}

Що буде виведено за виконання (повідомлення інтерпретатора про помилку позначати error)?\par
\begin{verbatim}
a,b,c=1,5,3
def f(b):
    a=3
    b+=2
    c=1
    return a+b+c

a,b,c=4,7,2
print(f(a),a,b,c)
\end{verbatim}

{\begin {center}\textbf{Завдання 10}\end {center}}


Написати функцiю f, що отримує щось та повертає щось. Невелика загадка все-таки має залишатися. Але нічого принципово нового відносно задач, що наявні у pdf-ках не планується.. Стиль запису умови також оцінюється.\par
\end {large}}
\newpage

{\Large{17.10.2024 Бета-версія: Контрольна робота \No 1, варіант  \No \textbf { 28 }\par}}
{
\begin {large}
{\begin {center}\textbf{Завдання 1}\end {center}}

Що буде виведено за виконання? (повідомлення інтерпретатора про помилку позначати error)\par
\begin{verbatim}
print(8//5//6*3)
\end{verbatim}

{\begin {center}\textbf{Завдання 2}\end {center}}

Що буде виведено за виконання? (повідомлення інтерпретатора про помилку позначати error)\par
\begin{verbatim}
print(4**-0.5*2)
\end{verbatim}

{\begin {center}\textbf{Завдання 3}\end {center}}

Що буде виведено за виконання? (повідомлення інтерпретатора про помилку позначати error)\par
\begin{verbatim}
print(9.0<7 and not 6>=8 or True!=4.0)
\end{verbatim}

{\begin {center}\textbf{Завдання 4}\end {center}}

Що буде виведено за виконання? (повідомлення інтерпретатора про помилку позначати error)\par
\begin{verbatim}
print("True" < 2)
\end{verbatim}

{\begin {center}\textbf{Завдання 5}\end {center}}

Що буде виведено за виконання? (повідомлення інтерпретатора про помилку позначати error)\par
\begin{verbatim}
print(4.0 <= 7.0 > True == 8)
\end{verbatim}

{\begin {center}\textbf{Завдання 6}\end {center}}

Що буде виведено за виконання? (повідомлення інтерпретатора про помилку позначати error)\par
\begin{verbatim}

def f(n):
    print("f", end="");
    return n<=-1

print(f(7) and f(-2))
\end{verbatim}

{\begin {center}\textbf{Завдання 7}\end {center}}

Що буде виведено за виконання (повідомлення інтерпретатора про помилку позначати error)?\par
\begin{verbatim}
def h(c):
    w=78
    if c!=-1: 
        w=3
    elif c>0:
         return 7
    else:
         return 4
    return w

print(h(4))
\end{verbatim}

{\begin {center}\textbf{Завдання 8}\end {center}}

Що буде виведено за виконання (повідомлення інтерпретатора про помилку позначати error)?\par
\begin{verbatim}
a, b, c = 9, 6, 7
def f(a, b=8, c):
    print(a, b, c, end=" ")

f(2, 5, b=4)
print(a, b, c)
\end{verbatim}

{\begin {center}\textbf{Завдання 9}\end {center}}

Що буде виведено за виконання (повідомлення інтерпретатора про помилку позначати error)?\par
\begin{verbatim}
a,b,c=2,9,8
def f(b):
    a=1
    b=4
    c=3
    return a+b+c

a,b,c=1,7,3
print(f(b),a,b,c)
\end{verbatim}

{\begin {center}\textbf{Завдання 10}\end {center}}


Написати функцiю f, що отримує щось та повертає щось. Невелика загадка все-таки має залишатися. Але нічого принципово нового відносно задач, що наявні у pdf-ках не планується.. Стиль запису умови також оцінюється.\par
\end {large}}
\newpage

{\Large{17.10.2024 Бета-версія: Контрольна робота \No 1, варіант  \No \textbf { 29 }\par}}
{
\begin {large}
{\begin {center}\textbf{Завдання 1}\end {center}}

Що буде виведено за виконання? (повідомлення інтерпретатора про помилку позначати error)\par
\begin{verbatim}
print(7/7*5*9)
\end{verbatim}

{\begin {center}\textbf{Завдання 2}\end {center}}

Що буде виведено за виконання? (повідомлення інтерпретатора про помилку позначати error)\par
\begin{verbatim}
print(9**-1.0*False)
\end{verbatim}

{\begin {center}\textbf{Завдання 3}\end {center}}

Що буде виведено за виконання? (повідомлення інтерпретатора про помилку позначати error)\par
\begin{verbatim}
print(9>5.0 and not False<2 or 5<=8.0)
\end{verbatim}

{\begin {center}\textbf{Завдання 4}\end {center}}

Що буде виведено за виконання? (повідомлення інтерпретатора про помилку позначати error)\par
\begin{verbatim}
print(False >= 5.0)
\end{verbatim}

{\begin {center}\textbf{Завдання 5}\end {center}}

Що буде виведено за виконання? (повідомлення інтерпретатора про помилку позначати error)\par
\begin{verbatim}
print(3 < 8 != 4 > 2)
\end{verbatim}

{\begin {center}\textbf{Завдання 6}\end {center}}

Що буде виведено за виконання? (повідомлення інтерпретатора про помилку позначати error)\par
\begin{verbatim}

def f(n):
    print("f", end="");
    return n!=-3

print(f(-5) or f(6))
\end{verbatim}

{\begin {center}\textbf{Завдання 7}\end {center}}

Що буде виведено за виконання (повідомлення інтерпретатора про помилку позначати error)?\par
\begin{verbatim}
def h(a,b):
    c=86
    if b:
        c=3
    elif a==-3:
         return 9
    else: 
        c=7
    return c

print(h(5,-3))
\end{verbatim}

{\begin {center}\textbf{Завдання 8}\end {center}}

Що буде виведено за виконання (повідомлення інтерпретатора про помилку позначати error)?\par
\begin{verbatim}
a, b, c = 9, 6, 7
def h(a, b=8, c):
    print(a, b, c, end=" ")

h(0, c=4)
print(a, b, c)
\end{verbatim}

{\begin {center}\textbf{Завдання 9}\end {center}}

Що буде виведено за виконання (повідомлення інтерпретатора про помилку позначати error)?\par
\begin{verbatim}
a,b,c=6,2,8
def g(b):
    a*=4
    b=5
    c=1
    return a+b+c

a,b,c=0,5,9
print(g(b),a,b,c)
\end{verbatim}

{\begin {center}\textbf{Завдання 10}\end {center}}


Написати функцiю f, що отримує щось та повертає щось. Невелика загадка все-таки має залишатися. Але нічого принципово нового відносно задач, що наявні у pdf-ках не планується.. Стиль запису умови також оцінюється.\par
\end {large}}
\newpage

{\Large{17.10.2024 Бета-версія: Контрольна робота \No 1, варіант  \No \textbf { 30 }\par}}
{
\begin {large}
{\begin {center}\textbf{Завдання 1}\end {center}}

Що буде виведено за виконання? (повідомлення інтерпретатора про помилку позначати error)\par
\begin{verbatim}
print(9//6/4*6)
\end{verbatim}

{\begin {center}\textbf{Завдання 2}\end {center}}

Що буде виведено за виконання? (повідомлення інтерпретатора про помилку позначати error)\par
\begin{verbatim}
print(3**-2-True)
\end{verbatim}

{\begin {center}\textbf{Завдання 3}\end {center}}

Що буде виведено за виконання? (повідомлення інтерпретатора про помилку позначати error)\par
\begin{verbatim}
print(4==5 or not 4.0>=6.0 and True!=2)
\end{verbatim}

{\begin {center}\textbf{Завдання 4}\end {center}}

Що буде виведено за виконання? (повідомлення інтерпретатора про помилку позначати error)\par
\begin{verbatim}
print("9.0" != "1")
\end{verbatim}

{\begin {center}\textbf{Завдання 5}\end {center}}

Що буде виведено за виконання? (повідомлення інтерпретатора про помилку позначати error)\par
\begin{verbatim}
print(5.0 == False  is  6.0 <= 9)
\end{verbatim}

{\begin {center}\textbf{Завдання 6}\end {center}}

Що буде виведено за виконання? (повідомлення інтерпретатора про помилку позначати error)\par
\begin{verbatim}

def f(n):
    print("f", end="");
    return n

print(f(9) and f(-1))
\end{verbatim}

{\begin {center}\textbf{Завдання 7}\end {center}}

Що буде виведено за виконання (повідомлення інтерпретатора про помилку позначати error)?\par
\begin{verbatim}
def f(d):
    x=21
    if d>=-5: 
        x=6
    if d<4:
         return 1
    else:
         x=2
    return x

print(f(-3))
\end{verbatim}

{\begin {center}\textbf{Завдання 8}\end {center}}

Що буде виведено за виконання (повідомлення інтерпретатора про помилку позначати error)?\par
\begin{verbatim}
a, b, c = 9, 6, 7
def g(a, b=8, c=9):
    print(a, b, c, end=" ")

g(0, c=2, b=5)
print(a, b, c)
\end{verbatim}

{\begin {center}\textbf{Завдання 9}\end {center}}

Що буде виведено за виконання (повідомлення інтерпретатора про помилку позначати error)?\par
\begin{verbatim}
a,b,c=3,4,1
def f(a):
    a+=3
    b=2
    c=4
    return a+b+c

a,b,c=7,0,8
print(f(a),a,b,c)
\end{verbatim}

{\begin {center}\textbf{Завдання 10}\end {center}}


Написати функцiю f, що отримує щось та повертає щось. Невелика загадка все-таки має залишатися. Але нічого принципово нового відносно задач, що наявні у pdf-ках не планується.. Стиль запису умови також оцінюється.\par
\end {large}}
\newpage

{\Large{17.10.2024 Бета-версія: Контрольна робота \No 1, варіант  \No \textbf { 31 }\par}}
{
\begin {large}
{\begin {center}\textbf{Завдання 1}\end {center}}

Що буде виведено за виконання? (повідомлення інтерпретатора про помилку позначати error)\par
\begin{verbatim}
print(3//4%9*5)
\end{verbatim}

{\begin {center}\textbf{Завдання 2}\end {center}}

Що буде виведено за виконання? (повідомлення інтерпретатора про помилку позначати error)\par
\begin{verbatim}
print(-3**1.0+False)
\end{verbatim}

{\begin {center}\textbf{Завдання 3}\end {center}}

Що буде виведено за виконання? (повідомлення інтерпретатора про помилку позначати error)\par
\begin{verbatim}
print(3.0!=7 or not 4>3 or 2.0<False)
\end{verbatim}

{\begin {center}\textbf{Завдання 4}\end {center}}

Що буде виведено за виконання? (повідомлення інтерпретатора про помилку позначати error)\par
\begin{verbatim}
print("8j" != True)
\end{verbatim}

{\begin {center}\textbf{Завдання 5}\end {center}}

Що буде виведено за виконання? (повідомлення інтерпретатора про помилку позначати error)\par
\begin{verbatim}
print(False >= 5.0 > 5 <= 3)
\end{verbatim}

{\begin {center}\textbf{Завдання 6}\end {center}}

Що буде виведено за виконання? (повідомлення інтерпретатора про помилку позначати error)\par
\begin{verbatim}

def f(n):
    print("f", end="");
    return n<-4

print(f(1) and f(5))
\end{verbatim}

{\begin {center}\textbf{Завдання 7}\end {center}}

Що буде виведено за виконання (повідомлення інтерпретатора про помилку позначати error)?\par
\begin{verbatim}
def f(a,b):
    c=20
    if a<=1:
        return 0
    if b>=-2:
         c=8
    else: 
        return 4
    return c

print(f(-5,-7))
\end{verbatim}

{\begin {center}\textbf{Завдання 8}\end {center}}

Що буде виведено за виконання (повідомлення інтерпретатора про помилку позначати error)?\par
\begin{verbatim}
a, b, c = 7, 8, 6
def h(a, b, c):
    print(a, b, c, end=" ")

h(c=1, b=3, a=5)
print(a, b, c)
\end{verbatim}

{\begin {center}\textbf{Завдання 9}\end {center}}

Що буде виведено за виконання (повідомлення інтерпретатора про помилку позначати error)?\par
\begin{verbatim}
a,b,c=6,4,2
def g(b):
    a*=5
    b=3
    c=2
    return a+b+c

a,b,c=5,7,9
print(g(b),a,b,c)
\end{verbatim}

{\begin {center}\textbf{Завдання 10}\end {center}}


Написати функцiю f, що отримує щось та повертає щось. Невелика загадка все-таки має залишатися. Але нічого принципово нового відносно задач, що наявні у pdf-ках не планується.. Стиль запису умови також оцінюється.\par
\end {large}}
\newpage

{\Large{17.10.2024 Бета-версія: Контрольна робота \No 1, варіант  \No \textbf { 32 }\par}}
{
\begin {large}
{\begin {center}\textbf{Завдання 1}\end {center}}

Що буде виведено за виконання? (повідомлення інтерпретатора про помилку позначати error)\par
\begin{verbatim}
print(8/8//3*8)
\end{verbatim}

{\begin {center}\textbf{Завдання 2}\end {center}}

Що буде виведено за виконання? (повідомлення інтерпретатора про помилку позначати error)\par
\begin{verbatim}
print(9**1.0*2)
\end{verbatim}

{\begin {center}\textbf{Завдання 3}\end {center}}

Що буде виведено за виконання? (повідомлення інтерпретатора про помилку позначати error)\par
\begin{verbatim}
print(not 8>=9.0 and True<=6 and 9==7.0)
\end{verbatim}

{\begin {center}\textbf{Завдання 4}\end {center}}

Що буде виведено за виконання? (повідомлення інтерпретатора про помилку позначати error)\par
\begin{verbatim}
print("2.0" < 3)
\end{verbatim}

{\begin {center}\textbf{Завдання 5}\end {center}}

Що буде виведено за виконання? (повідомлення інтерпретатора про помилку позначати error)\par
\begin{verbatim}
print(2.0 >= 6 < 7 != 5)
\end{verbatim}

{\begin {center}\textbf{Завдання 6}\end {center}}

Що буде виведено за виконання? (повідомлення інтерпретатора про помилку позначати error)\par
\begin{verbatim}

def f(n):
    print("f", end="");
    return n

print(f(-6) or f(0))
\end{verbatim}

{\begin {center}\textbf{Завдання 7}\end {center}}

Що буде виведено за виконання (повідомлення інтерпретатора про помилку позначати error)?\par
\begin{verbatim}
def h(a,b):
    c=81
    if b<=4:
        c=2
    elif a>-1:
         c=1
    else: 
        return 6
    return c

print(h(-4,2))
\end{verbatim}

{\begin {center}\textbf{Завдання 8}\end {center}}

Що буде виведено за виконання (повідомлення інтерпретатора про помилку позначати error)?\par
\begin{verbatim}
a, b, c = 9, 8, 7
def g(a, b, c=6):
    print(a, b, c, end=" ")

g(4, 1, 2)
print(a, b, c)
\end{verbatim}

{\begin {center}\textbf{Завдання 9}\end {center}}

Що буде виведено за виконання (повідомлення інтерпретатора про помилку позначати error)?\par
\begin{verbatim}
a,b,c=3,1,7
def h(a):
    a=1
    b-=3
    c=5
    return a+b+c

a,b,c=4,0,1
print(h(a),a,b,c)
\end{verbatim}

{\begin {center}\textbf{Завдання 10}\end {center}}


Написати функцiю f, що отримує щось та повертає щось. Невелика загадка все-таки має залишатися. Але нічого принципово нового відносно задач, що наявні у pdf-ках не планується.. Стиль запису умови також оцінюється.\par
\end {large}}
\newpage

{\Large{17.10.2024 Бета-версія: Контрольна робота \No 1, варіант  \No \textbf { 33 }\par}}
{
\begin {large}
{\begin {center}\textbf{Завдання 1}\end {center}}

Що буде виведено за виконання? (повідомлення інтерпретатора про помилку позначати error)\par
\begin{verbatim}
print(6/5//8*3)
\end{verbatim}

{\begin {center}\textbf{Завдання 2}\end {center}}

Що буде виведено за виконання? (повідомлення інтерпретатора про помилку позначати error)\par
\begin{verbatim}
print(1**-0.5*1)
\end{verbatim}

{\begin {center}\textbf{Завдання 3}\end {center}}

Що буде виведено за виконання? (повідомлення інтерпретатора про помилку позначати error)\par
\begin{verbatim}
print(not 2.0<=8 or 4.0>3 or 5!=False)
\end{verbatim}

{\begin {center}\textbf{Завдання 4}\end {center}}

Що буде виведено за виконання? (повідомлення інтерпретатора про помилку позначати error)\par
\begin{verbatim}
print(4.0 == "False")
\end{verbatim}

{\begin {center}\textbf{Завдання 5}\end {center}}

Що буде виведено за виконання? (повідомлення інтерпретатора про помилку позначати error)\par
\begin{verbatim}
print(9.0  is  2 == True < 3.0)
\end{verbatim}

{\begin {center}\textbf{Завдання 6}\end {center}}

Що буде виведено за виконання? (повідомлення інтерпретатора про помилку позначати error)\par
\begin{verbatim}

def f(n):
    print("f", end="");
    return n>4

print(f(-8) or f(-5))
\end{verbatim}

{\begin {center}\textbf{Завдання 7}\end {center}}

Що буде виведено за виконання (повідомлення інтерпретатора про помилку позначати error)?\par
\begin{verbatim}
def f(a,b):
    c=10
    if b:
        c=5
    if a==2:
         return 3
    else: 
        c=8
    return c

print(f(-7,-9))
\end{verbatim}

{\begin {center}\textbf{Завдання 8}\end {center}}

Що буде виведено за виконання (повідомлення інтерпретатора про помилку позначати error)?\par
\begin{verbatim}
a, b, c = 7, 8, 6
def f(a, b, c=9):
    print(a, b, c, end=" ")

f(0, 3, b=5)
print(a, b, c)
\end{verbatim}

{\begin {center}\textbf{Завдання 9}\end {center}}

Що буде виведено за виконання (повідомлення інтерпретатора про помилку позначати error)?\par
\begin{verbatim}
a,b,c=9,3,6
def g(b):
    a+=4
    b=2
    c=1
    return a+b+c

a,b,c=8,5,2
print(g(b),a,b,c)
\end{verbatim}

{\begin {center}\textbf{Завдання 10}\end {center}}


Написати функцiю f, що отримує щось та повертає щось. Невелика загадка все-таки має залишатися. Але нічого принципово нового відносно задач, що наявні у pdf-ках не планується.. Стиль запису умови також оцінюється.\par
\end {large}}
\newpage

{\Large{17.10.2024 Бета-версія: Контрольна робота \No 1, варіант  \No \textbf { 34 }\par}}
{
\begin {large}
{\begin {center}\textbf{Завдання 1}\end {center}}

Що буде виведено за виконання? (повідомлення інтерпретатора про помилку позначати error)\par
\begin{verbatim}
print(5//9/6*4)
\end{verbatim}

{\begin {center}\textbf{Завдання 2}\end {center}}

Що буде виведено за виконання? (повідомлення інтерпретатора про помилку позначати error)\par
\begin{verbatim}
print(4**-3-True)
\end{verbatim}

{\begin {center}\textbf{Завдання 3}\end {center}}

Що буде виведено за виконання? (повідомлення інтерпретатора про помилку позначати error)\par
\begin{verbatim}
print(True==2 and not 9.0<5.0 and 9>=7)
\end{verbatim}

{\begin {center}\textbf{Завдання 4}\end {center}}

Що буде виведено за виконання? (повідомлення інтерпретатора про помилку позначати error)\par
\begin{verbatim}
print(4 <= 1.0)
\end{verbatim}

{\begin {center}\textbf{Завдання 5}\end {center}}

Що буде виведено за виконання? (повідомлення інтерпретатора про помилку позначати error)\par
\begin{verbatim}
print(4.0 == 7.0 > True  is  4)
\end{verbatim}

{\begin {center}\textbf{Завдання 6}\end {center}}

Що буде виведено за виконання? (повідомлення інтерпретатора про помилку позначати error)\par
\begin{verbatim}

def f(n):
    print("f", end="");
    return n

print(f(2) and f(3))
\end{verbatim}

{\begin {center}\textbf{Завдання 7}\end {center}}

Що буде виведено за виконання (повідомлення інтерпретатора про помилку позначати error)?\par
\begin{verbatim}
def g(a,b):
    c=22
    if a!=-5:
        return 9
    elif b<5:
         c=5
    else: 
        return 7
    return c

print(g(-3,-5))
\end{verbatim}

{\begin {center}\textbf{Завдання 8}\end {center}}

Що буде виведено за виконання (повідомлення інтерпретатора про помилку позначати error)?\par
\begin{verbatim}
a, b, c = 7, 9, 6
def g(a, b, c):
    print(a, b, c, end=" ")

g(4, c=1, b=2)
print(a, b, c)
\end{verbatim}

{\begin {center}\textbf{Завдання 9}\end {center}}

Що буде виведено за виконання (повідомлення інтерпретатора про помилку позначати error)?\par
\begin{verbatim}
a,b,c=8,2,1
def f(a):
    a=3
    b*=1
    c=2
    return a+b+c

a,b,c=9,0,6
print(f(a),a,b,c)
\end{verbatim}

{\begin {center}\textbf{Завдання 10}\end {center}}


Написати функцiю f, що отримує щось та повертає щось. Невелика загадка все-таки має залишатися. Але нічого принципово нового відносно задач, що наявні у pdf-ках не планується.. Стиль запису умови також оцінюється.\par
\end {large}}
\newpage

{\Large{17.10.2024 Бета-версія: Контрольна робота \No 1, варіант  \No \textbf { 35 }\par}}
{
\begin {large}
{\begin {center}\textbf{Завдання 1}\end {center}}

Що буде виведено за виконання? (повідомлення інтерпретатора про помилку позначати error)\par
\begin{verbatim}
print(4/3%7*7)
\end{verbatim}

{\begin {center}\textbf{Завдання 2}\end {center}}

Що буде виведено за виконання? (повідомлення інтерпретатора про помилку позначати error)\par
\begin{verbatim}
print(4**-1.0*-1)
\end{verbatim}

{\begin {center}\textbf{Завдання 3}\end {center}}

Що буде виведено за виконання? (повідомлення інтерпретатора про помилку позначати error)\par
\begin{verbatim}
print(6!=4 and not 2>True or 6.0==7.0)
\end{verbatim}

{\begin {center}\textbf{Завдання 4}\end {center}}

Що буде виведено за виконання? (повідомлення інтерпретатора про помилку позначати error)\par
\begin{verbatim}
print("6" >= "True")
\end{verbatim}

{\begin {center}\textbf{Завдання 5}\end {center}}

Що буде виведено за виконання? (повідомлення інтерпретатора про помилку позначати error)\par
\begin{verbatim}
print(7 <= 8 >= False != 8.0)
\end{verbatim}

{\begin {center}\textbf{Завдання 6}\end {center}}

Що буде виведено за виконання? (повідомлення інтерпретатора про помилку позначати error)\par
\begin{verbatim}

def f(n):
    print("f", end="");
    return n

print(f(4) and f(8))
\end{verbatim}

{\begin {center}\textbf{Завдання 7}\end {center}}

Що буде виведено за виконання (повідомлення інтерпретатора про помилку позначати error)?\par
\begin{verbatim}
def h(a,b):
    c=40
    if a>3:
        c=6
    if a<=-4:
         return 0
    else: 
        return 3
    return c

print(h(-9,7))
\end{verbatim}

{\begin {center}\textbf{Завдання 8}\end {center}}

Що буде виведено за виконання (повідомлення інтерпретатора про помилку позначати error)?\par
\begin{verbatim}
a, b, c = 8, 6, 8
def f(a, b=9, c):
    print(a, b, c, end=" ")

f(b=0, c=3, 5)
print(a, b, c)
\end{verbatim}

{\begin {center}\textbf{Завдання 9}\end {center}}

Що буде виведено за виконання (повідомлення інтерпретатора про помилку позначати error)?\par
\begin{verbatim}
a,b,c=6,0,4
def g(a):
    a=2
    b=5
    c=2
    return a+b+c

a,b,c=2,3,1
print(g(a),a,b,c)
\end{verbatim}

{\begin {center}\textbf{Завдання 10}\end {center}}


Написати функцiю f, що отримує щось та повертає щось. Невелика загадка все-таки має залишатися. Але нічого принципово нового відносно задач, що наявні у pdf-ках не планується.. Стиль запису умови також оцінюється.\par
\end {large}}
\newpage

{\Large{17.10.2024 Бета-версія: Контрольна робота \No 1, варіант  \No \textbf { 36 }\par}}
{
\begin {large}
{\begin {center}\textbf{Завдання 1}\end {center}}

Що буде виведено за виконання? (повідомлення інтерпретатора про помилку позначати error)\par
\begin{verbatim}
print(4//4*6*5)
\end{verbatim}

{\begin {center}\textbf{Завдання 2}\end {center}}

Що буде виведено за виконання? (повідомлення інтерпретатора про помилку позначати error)\par
\begin{verbatim}
print(1**1.0*-2)
\end{verbatim}

{\begin {center}\textbf{Завдання 3}\end {center}}

Що буде виведено за виконання? (повідомлення інтерпретатора про помилку позначати error)\par
\begin{verbatim}
print(8.0>=False or not 4<=9 and 3.0<5)
\end{verbatim}

{\begin {center}\textbf{Завдання 4}\end {center}}

Що буде виведено за виконання? (повідомлення інтерпретатора про помилку позначати error)\par
\begin{verbatim}
print("3.0j" > False)
\end{verbatim}

{\begin {center}\textbf{Завдання 5}\end {center}}

Що буде виведено за виконання? (повідомлення інтерпретатора про помилку позначати error)\par
\begin{verbatim}
print(6.0 != 9 >= 3 == 6)
\end{verbatim}

{\begin {center}\textbf{Завдання 6}\end {center}}

Що буде виведено за виконання? (повідомлення інтерпретатора про помилку позначати error)\par
\begin{verbatim}

def f(n):
    print("f", end="");
    return n>=0

print(f(-2) or f(-3))
\end{verbatim}

{\begin {center}\textbf{Завдання 7}\end {center}}

Що буде виведено за виконання (повідомлення інтерпретатора про помилку позначати error)?\par
\begin{verbatim}
def g(a):
    z=67
    if a==3: 
        return 8
    elif a<=-2:
         z=9
    else:
         z=1
    return z

print(g(5))
\end{verbatim}

{\begin {center}\textbf{Завдання 8}\end {center}}

Що буде виведено за виконання (повідомлення інтерпретатора про помилку позначати error)?\par
\begin{verbatim}
a, b, c = 7, 7, 6
def h(a, b=9, c=8):
    print(a, b, c, end=" ")

h(a=2, c=4, b=1)
print(a, b, c)
\end{verbatim}

{\begin {center}\textbf{Завдання 9}\end {center}}

Що буде виведено за виконання (повідомлення інтерпретатора про помилку позначати error)?\par
\begin{verbatim}
a,b,c=4,5,7
def h(b):
    a=5
    b-=4
    c=5
    return a+b+c

a,b,c=3,2,7
print(h(b),a,b,c)
\end{verbatim}

{\begin {center}\textbf{Завдання 10}\end {center}}


Написати функцiю f, що отримує щось та повертає щось. Невелика загадка все-таки має залишатися. Але нічого принципово нового відносно задач, що наявні у pdf-ках не планується.. Стиль запису умови також оцінюється.\par
\end {large}}
\newpage

{\Large{17.10.2024 Бета-версія: Контрольна робота \No 1, варіант  \No \textbf { 37 }\par}}
{
\begin {large}
{\begin {center}\textbf{Завдання 1}\end {center}}

Що буде виведено за виконання? (повідомлення інтерпретатора про помилку позначати error)\par
\begin{verbatim}
print(3//7%3*8)
\end{verbatim}

{\begin {center}\textbf{Завдання 2}\end {center}}

Що буде виведено за виконання? (повідомлення інтерпретатора про помилку позначати error)\par
\begin{verbatim}
print(-1**2.0+2)
\end{verbatim}

{\begin {center}\textbf{Завдання 3}\end {center}}

Що буде виведено за виконання? (повідомлення інтерпретатора про помилку позначати error)\par
\begin{verbatim}
print(3>5.0 and not 8<=True and 6<6.0)
\end{verbatim}

{\begin {center}\textbf{Завдання 4}\end {center}}

Що буде виведено за виконання? (повідомлення інтерпретатора про помилку позначати error)\par
\begin{verbatim}
print(False > 7)
\end{verbatim}

{\begin {center}\textbf{Завдання 5}\end {center}}

Що буде виведено за виконання? (повідомлення інтерпретатора про помилку позначати error)\par
\begin{verbatim}
print(5 > True <= 7.0  is  4.0)
\end{verbatim}

{\begin {center}\textbf{Завдання 6}\end {center}}

Що буде виведено за виконання? (повідомлення інтерпретатора про помилку позначати error)\par
\begin{verbatim}

def f(n):
    print("f", end="");
    return n>=-1

print(f(-4) and f(-9))
\end{verbatim}

{\begin {center}\textbf{Завдання 7}\end {center}}

Що буде виведено за виконання (повідомлення інтерпретатора про помилку позначати error)?\par
\begin{verbatim}
def h(b):
    u=57
    if b: 
        return 0
    if b!=-4:
         return 7
    else:
         u=3
    return u

print(h(2))
\end{verbatim}

{\begin {center}\textbf{Завдання 8}\end {center}}

Що буде виведено за виконання (повідомлення інтерпретатора про помилку позначати error)?\par
\begin{verbatim}
a, b, c = 7, 6, 8
def g(a, b=9, c=6):
    print(a, b, c, end=" ")

g(a=3, 0, a=3)
print(a, b, c)
\end{verbatim}

{\begin {center}\textbf{Завдання 9}\end {center}}

Що буде виведено за виконання (повідомлення інтерпретатора про помилку позначати error)?\par
\begin{verbatim}
a,b,c=5,6,9
def h(a):
    a=3
    b*=1
    c=2
    return a+b+c

a,b,c=0,8,3
print(h(a),a,b,c)
\end{verbatim}

{\begin {center}\textbf{Завдання 10}\end {center}}


Написати функцiю f, що отримує щось та повертає щось. Невелика загадка все-таки має залишатися. Але нічого принципово нового відносно задач, що наявні у pdf-ках не планується.. Стиль запису умови також оцінюється.\par
\end {large}}
\newpage

{\Large{17.10.2024 Бета-версія: Контрольна робота \No 1, варіант  \No \textbf { 38 }\par}}
{
\begin {large}
{\begin {center}\textbf{Завдання 1}\end {center}}

Що буде виведено за виконання? (повідомлення інтерпретатора про помилку позначати error)\par
\begin{verbatim}
print(6/6*7*9)
\end{verbatim}

{\begin {center}\textbf{Завдання 2}\end {center}}

Що буде виведено за виконання? (повідомлення інтерпретатора про помилку позначати error)\par
\begin{verbatim}
print(4**2.0-True)
\end{verbatim}

{\begin {center}\textbf{Завдання 3}\end {center}}

Що буде виведено за виконання? (повідомлення інтерпретатора про помилку позначати error)\par
\begin{verbatim}
print(not 7!=6 or False==8.0 or 5>=7.0)
\end{verbatim}

{\begin {center}\textbf{Завдання 4}\end {center}}

Що буде виведено за виконання? (повідомлення інтерпретатора про помилку позначати error)\par
\begin{verbatim}
print(8.0 <= "5j")
\end{verbatim}

{\begin {center}\textbf{Завдання 5}\end {center}}

Що буде виведено за виконання? (повідомлення інтерпретатора про помилку позначати error)\par
\begin{verbatim}
print(6 < 7 == 8  is  3)
\end{verbatim}

{\begin {center}\textbf{Завдання 6}\end {center}}

Що буде виведено за виконання? (повідомлення інтерпретатора про помилку позначати error)\par
\begin{verbatim}

def f(n):
    print("f", end="");
    return n

print(f(7) or f(-7))
\end{verbatim}

{\begin {center}\textbf{Завдання 7}\end {center}}

Що буде виведено за виконання (повідомлення інтерпретатора про помилку позначати error)?\par
\begin{verbatim}
def f(c):
    z=22
    if c!=1: 
        z=5
    if c<2:
         return 2
    else:
         z=4
    return z

print(f(-2))
\end{verbatim}

{\begin {center}\textbf{Завдання 8}\end {center}}

Що буде виведено за виконання (повідомлення інтерпретатора про помилку позначати error)?\par
\begin{verbatim}
a, b, c = 9, 7, 8
def f(a, b, c=6):
    print(a, b, c, end=" ")

f(1, 4)
print(a, b, c)
\end{verbatim}

{\begin {center}\textbf{Завдання 9}\end {center}}

Що буде виведено за виконання (повідомлення інтерпретатора про помилку позначати error)?\par
\begin{verbatim}
a,b,c=4,1,7
def f(b):
    a+=4
    b=3
    c=2
    return a+b+c

a,b,c=4,6,5
print(f(b),a,b,c)
\end{verbatim}

{\begin {center}\textbf{Завдання 10}\end {center}}


Написати функцiю f, що отримує щось та повертає щось. Невелика загадка все-таки має залишатися. Але нічого принципово нового відносно задач, що наявні у pdf-ках не планується.. Стиль запису умови також оцінюється.\par
\end {large}}
\newpage

{\Large{17.10.2024 Бета-версія: Контрольна робота \No 1, варіант  \No \textbf { 39 }\par}}
{
\begin {large}
{\begin {center}\textbf{Завдання 1}\end {center}}

Що буде виведено за виконання? (повідомлення інтерпретатора про помилку позначати error)\par
\begin{verbatim}
print(7//8/9*4)
\end{verbatim}

{\begin {center}\textbf{Завдання 2}\end {center}}

Що буде виведено за виконання? (повідомлення інтерпретатора про помилку позначати error)\par
\begin{verbatim}
print(9**0.5*-1)
\end{verbatim}

{\begin {center}\textbf{Завдання 3}\end {center}}

Що буде виведено за виконання? (повідомлення інтерпретатора про помилку позначати error)\par
\begin{verbatim}
print(not 3.0>4.0 and 9<=False or 4>=8)
\end{verbatim}

{\begin {center}\textbf{Завдання 4}\end {center}}

Що буде виведено за виконання? (повідомлення інтерпретатора про помилку позначати error)\par
\begin{verbatim}
print("True" < "6.0j")
\end{verbatim}

{\begin {center}\textbf{Завдання 5}\end {center}}

Що буде виведено за виконання? (повідомлення інтерпретатора про помилку позначати error)\par
\begin{verbatim}
print(False > 4 >= 9.0 < 5.0)
\end{verbatim}

{\begin {center}\textbf{Завдання 6}\end {center}}

Що буде виведено за виконання? (повідомлення інтерпретатора про помилку позначати error)\par
\begin{verbatim}

def f(n):
    print("f", end="");
    return n==-3

print(f(-5) or f(5))
\end{verbatim}

{\begin {center}\textbf{Завдання 7}\end {center}}

Що буде виведено за виконання (повідомлення інтерпретатора про помилку позначати error)?\par
\begin{verbatim}
def g(d):
    x=52
    if d: 
        return 9
    elif d>5:
         return 6
    else:
         x=8
    return x

print(g(0))
\end{verbatim}

{\begin {center}\textbf{Завдання 8}\end {center}}

Що буде виведено за виконання (повідомлення інтерпретатора про помилку позначати error)?\par
\begin{verbatim}
a, b, c = 8, 7, 9
def h(a, b, c):
    print(a, b, c, end=" ")

h(0, 2, 5)
print(a, b, c)
\end{verbatim}

{\begin {center}\textbf{Завдання 9}\end {center}}

Що буде виведено за виконання (повідомлення інтерпретатора про помилку позначати error)?\par
\begin{verbatim}
a,b,c=5,8,6
def g(b):
    a=4
    b=3
    c=1
    return a+b+c

a,b,c=0,9,7
print(g(a),a,b,c)
\end{verbatim}

{\begin {center}\textbf{Завдання 10}\end {center}}


Написати функцiю f, що отримує щось та повертає щось. Невелика загадка все-таки має залишатися. Але нічого принципово нового відносно задач, що наявні у pdf-ках не планується.. Стиль запису умови також оцінюється.\par
\end {large}}
\newpage

{\Large{17.10.2024 Бета-версія: Контрольна робота \No 1, варіант  \No \textbf { 40 }\par}}
{
\begin {large}
{\begin {center}\textbf{Завдання 1}\end {center}}

Що буде виведено за виконання? (повідомлення інтерпретатора про помилку позначати error)\par
\begin{verbatim}
print(8/9//4*7)
\end{verbatim}

{\begin {center}\textbf{Завдання 2}\end {center}}

Що буде виведено за виконання? (повідомлення інтерпретатора про помилку позначати error)\par
\begin{verbatim}
print(-3**-4+-2)
\end{verbatim}

{\begin {center}\textbf{Завдання 3}\end {center}}

Що буде виведено за виконання? (повідомлення інтерпретатора про помилку позначати error)\par
\begin{verbatim}
print(not 7==3 or 2.0<9.0 and 2!=True)
\end{verbatim}

{\begin {center}\textbf{Завдання 4}\end {center}}

Що буде виведено за виконання? (повідомлення інтерпретатора про помилку позначати error)\par
\begin{verbatim}
print("True" >= 9)
\end{verbatim}

{\begin {center}\textbf{Завдання 5}\end {center}}

Що буде виведено за виконання? (повідомлення інтерпретатора про помилку позначати error)\par
\begin{verbatim}
print(9 != False <= 2 < 3.0)
\end{verbatim}

{\begin {center}\textbf{Завдання 6}\end {center}}

Що буде виведено за виконання? (повідомлення інтерпретатора про помилку позначати error)\par
\begin{verbatim}

def f(n):
    print("f", end="");
    return n

print(f(6) and f(3))
\end{verbatim}

{\begin {center}\textbf{Завдання 7}\end {center}}

Що буде виведено за виконання (повідомлення інтерпретатора про помилку позначати error)?\par
\begin{verbatim}
def g(a,b):
    c=34
    if b>=0:
        c=4
    elif a==-1:
         c=2
    else: 
        return 1
    return c

print(g(7,-7))
\end{verbatim}

{\begin {center}\textbf{Завдання 8}\end {center}}

Що буде виведено за виконання (повідомлення інтерпретатора про помилку позначати error)?\par
\begin{verbatim}
a, b, c = 9, 7, 6
def f(a, b=8, c):
    print(a, b, c, end=" ")

f(0, c=5)
print(a, b, c)
\end{verbatim}

{\begin {center}\textbf{Завдання 9}\end {center}}

Що буде виведено за виконання (повідомлення інтерпретатора про помилку позначати error)?\par
\begin{verbatim}
a,b,c=0,1,9
def g(a):
    a=5
    b-=1
    c=4
    return a+b+c

a,b,c=2,3,8
print(g(a),a,b,c)
\end{verbatim}

{\begin {center}\textbf{Завдання 10}\end {center}}


Написати функцiю f, що отримує щось та повертає щось. Невелика загадка все-таки має залишатися. Але нічого принципово нового відносно задач, що наявні у pdf-ках не планується.. Стиль запису умови також оцінюється.\par
\end {large}}
\newpage

{\Large{17.10.2024 Бета-версія: Контрольна робота \No 1, варіант  \No \textbf { 41 }\par}}
{
\begin {large}
{\begin {center}\textbf{Завдання 1}\end {center}}

Що буде виведено за виконання? (повідомлення інтерпретатора про помилку позначати error)\par
\begin{verbatim}
print(9//5/5*6)
\end{verbatim}

{\begin {center}\textbf{Завдання 2}\end {center}}

Що буде виведено за виконання? (повідомлення інтерпретатора про помилку позначати error)\par
\begin{verbatim}
print(4**-0.5*1)
\end{verbatim}

{\begin {center}\textbf{Завдання 3}\end {center}}

Що буде виведено за виконання? (повідомлення інтерпретатора про помилку позначати error)\par
\begin{verbatim}
print(not 7.0==2.0 or 3<5 and False!=9)
\end{verbatim}

{\begin {center}\textbf{Завдання 4}\end {center}}

Що буде виведено за виконання? (повідомлення інтерпретатора про помилку позначати error)\par
\begin{verbatim}
print("7.0" == False)
\end{verbatim}

{\begin {center}\textbf{Завдання 5}\end {center}}

Що буде виведено за виконання? (повідомлення інтерпретатора про помилку позначати error)\par
\begin{verbatim}
print(9 <= 6.0 >= 2.0 != 8.0)
\end{verbatim}

{\begin {center}\textbf{Завдання 6}\end {center}}

Що буде виведено за виконання? (повідомлення інтерпретатора про помилку позначати error)\par
\begin{verbatim}

def f(n):
    print("f", end="");
    return n

print(f(8) or f(-4))
\end{verbatim}

{\begin {center}\textbf{Завдання 7}\end {center}}

Що буде виведено за виконання (повідомлення інтерпретатора про помилку позначати error)?\par
\begin{verbatim}
def f(a,b):
    c=76
    if b:
        return 4
    if b!=-4:
         c=0
    else: 
        return 2
    return c

print(f(1,-4))
\end{verbatim}

{\begin {center}\textbf{Завдання 8}\end {center}}

Що буде виведено за виконання (повідомлення інтерпретатора про помилку позначати error)?\par
\begin{verbatim}
a, b, c = 7, 8, 6
def h(a, b=9, c):
    print(a, b, c, end=" ")

h(b=4, c=3, 1)
print(a, b, c)
\end{verbatim}

{\begin {center}\textbf{Завдання 9}\end {center}}

Що буде виведено за виконання (повідомлення інтерпретатора про помилку позначати error)?\par
\begin{verbatim}
a,b,c=2,0,1
def g(b):
    a+=4
    b=1
    c=3
    return a+b+c

a,b,c=5,9,6
print(g(a),a,b,c)
\end{verbatim}

{\begin {center}\textbf{Завдання 10}\end {center}}


Написати функцiю f, що отримує щось та повертає щось. Невелика загадка все-таки має залишатися. Але нічого принципово нового відносно задач, що наявні у pdf-ках не планується.. Стиль запису умови також оцінюється.\par
\end {large}}
\newpage

{\Large{17.10.2024 Бета-версія: Контрольна робота \No 1, варіант  \No \textbf { 42 }\par}}
{
\begin {large}
{\begin {center}\textbf{Завдання 1}\end {center}}

Що буде виведено за виконання? (повідомлення інтерпретатора про помилку позначати error)\par
\begin{verbatim}
print(5/3*8*3)
\end{verbatim}

{\begin {center}\textbf{Завдання 2}\end {center}}

Що буде виведено за виконання? (повідомлення інтерпретатора про помилку позначати error)\par
\begin{verbatim}
print(9**1.0*False)
\end{verbatim}

{\begin {center}\textbf{Завдання 3}\end {center}}

Що буде виведено за виконання? (повідомлення інтерпретатора про помилку позначати error)\par
\begin{verbatim}
print(not 9.0<=2 and 3.0>7 or 6>=True)
\end{verbatim}

{\begin {center}\textbf{Завдання 4}\end {center}}

Що буде виведено за виконання? (повідомлення інтерпретатора про помилку позначати error)\par
\begin{verbatim}
print(4.0 != "8")
\end{verbatim}

{\begin {center}\textbf{Завдання 5}\end {center}}

Що буде виведено за виконання? (повідомлення інтерпретатора про помилку позначати error)\par
\begin{verbatim}
print(5  is  6 == True > 4)
\end{verbatim}

{\begin {center}\textbf{Завдання 6}\end {center}}

Що буде виведено за виконання? (повідомлення інтерпретатора про помилку позначати error)\par
\begin{verbatim}

def f(n):
    print("f", end="");
    return n<-4

print(f(7) and f(-7))
\end{verbatim}

{\begin {center}\textbf{Завдання 7}\end {center}}

Що буде виведено за виконання (повідомлення інтерпретатора про помилку позначати error)?\par
\begin{verbatim}
def h(a):
    y=59
    if a==-3: 
        return 2
    elif a>=5:
         y=9
    else:
         return 5
    return y

print(h(-1))
\end{verbatim}

{\begin {center}\textbf{Завдання 8}\end {center}}

Що буде виведено за виконання (повідомлення інтерпретатора про помилку позначати error)?\par
\begin{verbatim}
a, b, c = 6, 9, 8
def g(a, b, c=7):
    print(a, b, c, end=" ")

g(a=2, 3, b=1)
print(a, b, c)
\end{verbatim}

{\begin {center}\textbf{Завдання 9}\end {center}}

Що буде виведено за виконання (повідомлення інтерпретатора про помилку позначати error)?\par
\begin{verbatim}
a,b,c=5,9,1
def h(a):
    a=5
    b=4
    c=3
    return a+b+c

a,b,c=6,7,4
print(h(b),a,b,c)
\end{verbatim}

{\begin {center}\textbf{Завдання 10}\end {center}}


Написати функцiю f, що отримує щось та повертає щось. Невелика загадка все-таки має залишатися. Але нічого принципово нового відносно задач, що наявні у pdf-ках не планується.. Стиль запису умови також оцінюється.\par
\end {large}}
\newpage

{\Large{17.10.2024 Бета-версія: Контрольна робота \No 1, варіант  \No \textbf { 43 }\par}}
{
\begin {large}
{\begin {center}\textbf{Завдання 1}\end {center}}

Що буде виведено за виконання? (повідомлення інтерпретатора про помилку позначати error)\par
\begin{verbatim}
print(4/6//5*4)
\end{verbatim}

{\begin {center}\textbf{Завдання 2}\end {center}}

Що буде виведено за виконання? (повідомлення інтерпретатора про помилку позначати error)\par
\begin{verbatim}
print(2**2.0+True)
\end{verbatim}

{\begin {center}\textbf{Завдання 3}\end {center}}

Що буде виведено за виконання? (повідомлення інтерпретатора про помилку позначати error)\par
\begin{verbatim}
print(8!=4 or not 8.0>=4 and 6.0>True)
\end{verbatim}

{\begin {center}\textbf{Завдання 4}\end {center}}

Що буде виведено за виконання? (повідомлення інтерпретатора про помилку позначати error)\par
\begin{verbatim}
print(3 <= "7.0j")
\end{verbatim}

{\begin {center}\textbf{Завдання 5}\end {center}}

Що буде виведено за виконання? (повідомлення інтерпретатора про помилку позначати error)\par
\begin{verbatim}
print(3  is  True > 8 >= 7)
\end{verbatim}

{\begin {center}\textbf{Завдання 6}\end {center}}

Що буде виведено за виконання? (повідомлення інтерпретатора про помилку позначати error)\par
\begin{verbatim}

def f(n):
    print("f", end="");
    return n

print(f(2) and f(0))
\end{verbatim}

{\begin {center}\textbf{Завдання 7}\end {center}}

Що буде виведено за виконання (повідомлення інтерпретатора про помилку позначати error)?\par
\begin{verbatim}
def h(a,b):
    c=24
    if a<-5:
        c=3
    elif a<=-2:
         return 1
    else: 
        c=6
    return c

print(h(9,8))
\end{verbatim}

{\begin {center}\textbf{Завдання 8}\end {center}}

Що буде виведено за виконання (повідомлення інтерпретатора про помилку позначати error)?\par
\begin{verbatim}
a, b, c = 9, 8, 7
def g(a, b, c):
    print(a, b, c, end=" ")

g(2, a=0)
print(a, b, c)
\end{verbatim}

{\begin {center}\textbf{Завдання 9}\end {center}}

Що буде виведено за виконання (повідомлення інтерпретатора про помилку позначати error)?\par
\begin{verbatim}
a,b,c=7,4,3
def f(a):
    a=5
    b-=2
    c=3
    return a+b+c

a,b,c=8,0,9
print(f(a),a,b,c)
\end{verbatim}

{\begin {center}\textbf{Завдання 10}\end {center}}


Написати функцiю f, що отримує щось та повертає щось. Невелика загадка все-таки має залишатися. Але нічого принципово нового відносно задач, що наявні у pdf-ках не планується.. Стиль запису умови також оцінюється.\par
\end {large}}
\newpage

{\Large{17.10.2024 Бета-версія: Контрольна робота \No 1, варіант  \No \textbf { 44 }\par}}
{
\begin {large}
{\begin {center}\textbf{Завдання 1}\end {center}}

Що буде виведено за виконання? (повідомлення інтерпретатора про помилку позначати error)\par
\begin{verbatim}
print(9//3%4*3)
\end{verbatim}

{\begin {center}\textbf{Завдання 2}\end {center}}

Що буде виведено за виконання? (повідомлення інтерпретатора про помилку позначати error)\par
\begin{verbatim}
print(2**-1--2)
\end{verbatim}

{\begin {center}\textbf{Завдання 3}\end {center}}

Що буде виведено за виконання? (повідомлення інтерпретатора про помилку позначати error)\par
\begin{verbatim}
print(not 2==4.0 and 5<=3 or False<5.0)
\end{verbatim}

{\begin {center}\textbf{Завдання 4}\end {center}}

Що буде виведено за виконання? (повідомлення інтерпретатора про помилку позначати error)\par
\begin{verbatim}
print("True" >= False)
\end{verbatim}

{\begin {center}\textbf{Завдання 5}\end {center}}

Що буде виведено за виконання? (повідомлення інтерпретатора про помилку позначати error)\par
\begin{verbatim}
print(4.0 == 9.0 <= 2 < False)
\end{verbatim}

{\begin {center}\textbf{Завдання 6}\end {center}}

Що буде виведено за виконання? (повідомлення інтерпретатора про помилку позначати error)\par
\begin{verbatim}

def f(n):
    print("f", end="");
    return n>2

print(f(4) or f(-1))
\end{verbatim}

{\begin {center}\textbf{Завдання 7}\end {center}}

Що буде виведено за виконання (повідомлення інтерпретатора про помилку позначати error)?\par
\begin{verbatim}
def f(c):
    w=27
    if c<=-5: 
        w=4
    if c!=-1:
         return 7
    else:
         w=8
    return w

print(f(6))
\end{verbatim}

{\begin {center}\textbf{Завдання 8}\end {center}}

Що буде виведено за виконання (повідомлення інтерпретатора про помилку позначати error)?\par
\begin{verbatim}
a, b, c = 6, 6, 7
def h(a, b=8, c=9):
    print(a, b, c, end=" ")

h(4, c=5, b=1)
print(a, b, c)
\end{verbatim}

{\begin {center}\textbf{Завдання 9}\end {center}}

Що буде виведено за виконання (повідомлення інтерпретатора про помилку позначати error)?\par
\begin{verbatim}
a,b,c=4,1,7
def h(b):
    a=5
    b*=2
    c=4
    return a+b+c

a,b,c=6,3,2
print(h(b),a,b,c)
\end{verbatim}

{\begin {center}\textbf{Завдання 10}\end {center}}


Написати функцiю f, що отримує щось та повертає щось. Невелика загадка все-таки має залишатися. Але нічого принципово нового відносно задач, що наявні у pdf-ках не планується.. Стиль запису умови також оцінюється.\par
\end {large}}
\newpage

{\Large{17.10.2024 Бета-версія: Контрольна робота \No 1, варіант  \No \textbf { 45 }\par}}
{
\begin {large}
{\begin {center}\textbf{Завдання 1}\end {center}}

Що буде виведено за виконання? (повідомлення інтерпретатора про помилку позначати error)\par
\begin{verbatim}
print(3//4//8*6)
\end{verbatim}

{\begin {center}\textbf{Завдання 2}\end {center}}

Що буде виведено за виконання? (повідомлення інтерпретатора про помилку позначати error)\par
\begin{verbatim}
print(4**0.5*-1)
\end{verbatim}

{\begin {center}\textbf{Завдання 3}\end {center}}

Що буде виведено за виконання? (повідомлення інтерпретатора про помилку позначати error)\par
\begin{verbatim}
print(not 7.0==9 or 6<=True or 9.0>=8)
\end{verbatim}

{\begin {center}\textbf{Завдання 4}\end {center}}

Що буде виведено за виконання? (повідомлення інтерпретатора про помилку позначати error)\par
\begin{verbatim}
print(9.0 == "6j")
\end{verbatim}

{\begin {center}\textbf{Завдання 5}\end {center}}

Що буде виведено за виконання? (повідомлення інтерпретатора про помилку позначати error)\par
\begin{verbatim}
print(5.0 != 3.0  is  4 >= 3)
\end{verbatim}

{\begin {center}\textbf{Завдання 6}\end {center}}

Що буде виведено за виконання? (повідомлення інтерпретатора про помилку позначати error)\par
\begin{verbatim}

def f(n):
    print("f", end="");
    return n!=3

print(f(1) or f(-2))
\end{verbatim}

{\begin {center}\textbf{Завдання 7}\end {center}}

Що буде виведено за виконання (повідомлення інтерпретатора про помилку позначати error)?\par
\begin{verbatim}
def g(a,b):
    c=77
    if a:
        return 7
    elif b!=3:
         c=5
    else: 
        return 8
    return c

print(g(-1,3))
\end{verbatim}

{\begin {center}\textbf{Завдання 8}\end {center}}

Що буде виведено за виконання (повідомлення інтерпретатора про помилку позначати error)?\par
\begin{verbatim}
a, b, c = 6, 9, 7
def f(a, b, c):
    print(a, b, c, end=" ")

f(0, 5, 3)
print(a, b, c)
\end{verbatim}

{\begin {center}\textbf{Завдання 9}\end {center}}

Що буде виведено за виконання (повідомлення інтерпретатора про помилку позначати error)?\par
\begin{verbatim}
a,b,c=8,5,5
def h(a):
    a=1
    b*=5
    c=4
    return a+b+c

a,b,c=2,9,6
print(h(a),a,b,c)
\end{verbatim}

{\begin {center}\textbf{Завдання 10}\end {center}}


Написати функцiю f, що отримує щось та повертає щось. Невелика загадка все-таки має залишатися. Але нічого принципово нового відносно задач, що наявні у pdf-ках не планується.. Стиль запису умови також оцінюється.\par
\end {large}}
\newpage

{\Large{17.10.2024 Бета-версія: Контрольна робота \No 1, варіант  \No \textbf { 46 }\par}}
{
\begin {large}
{\begin {center}\textbf{Завдання 1}\end {center}}

Що буде виведено за виконання? (повідомлення інтерпретатора про помилку позначати error)\par
\begin{verbatim}
print(5/5*6*5)
\end{verbatim}

{\begin {center}\textbf{Завдання 2}\end {center}}

Що буде виведено за виконання? (повідомлення інтерпретатора про помилку позначати error)\par
\begin{verbatim}
print(1**-1.0*1)
\end{verbatim}

{\begin {center}\textbf{Завдання 3}\end {center}}

Що буде виведено за виконання? (повідомлення інтерпретатора про помилку позначати error)\par
\begin{verbatim}
print(not False>3.0 and 7<5 and 3!=6.0)
\end{verbatim}

{\begin {center}\textbf{Завдання 4}\end {center}}

Що буде виведено за виконання? (повідомлення інтерпретатора про помилку позначати error)\par
\begin{verbatim}
print(1 < 2.0)
\end{verbatim}

{\begin {center}\textbf{Завдання 5}\end {center}}

Що буде виведено за виконання? (повідомлення інтерпретатора про помилку позначати error)\par
\begin{verbatim}
print(6 == 8 <= False < 7.0)
\end{verbatim}

{\begin {center}\textbf{Завдання 6}\end {center}}

Що буде виведено за виконання? (повідомлення інтерпретатора про помилку позначати error)\par
\begin{verbatim}

def f(n):
    print("f", end="");
    return n

print(f(-9) and f(-6))
\end{verbatim}

{\begin {center}\textbf{Завдання 7}\end {center}}

Що буде виведено за виконання (повідомлення інтерпретатора про помилку позначати error)?\par
\begin{verbatim}
def h(b):
    v=90
    if b: 
        v=3
    if b<=-2:
         return 6
    else:
         return 0
    return v

print(h(-2))
\end{verbatim}

{\begin {center}\textbf{Завдання 8}\end {center}}

Що буде виведено за виконання (повідомлення інтерпретатора про помилку позначати error)?\par
\begin{verbatim}
a, b, c = 8, 8, 7
def g(a, b=6, c):
    print(a, b, c, end=" ")

g(4, 2, c=2)
print(a, b, c)
\end{verbatim}

{\begin {center}\textbf{Завдання 9}\end {center}}

Що буде виведено за виконання (повідомлення інтерпретатора про помилку позначати error)?\par
\begin{verbatim}
a,b,c=1,0,4
def g(a):
    a-=2
    b=1
    c=3
    return a+b+c

a,b,c=3,8,7
print(g(b),a,b,c)
\end{verbatim}

{\begin {center}\textbf{Завдання 10}\end {center}}


Написати функцiю f, що отримує щось та повертає щось. Невелика загадка все-таки має залишатися. Але нічого принципово нового відносно задач, що наявні у pdf-ках не планується.. Стиль запису умови також оцінюється.\par
\end {large}}
\newpage

{\Large{17.10.2024 Бета-версія: Контрольна робота \No 1, варіант  \No \textbf { 47 }\par}}
{
\begin {large}
{\begin {center}\textbf{Завдання 1}\end {center}}

Що буде виведено за виконання? (повідомлення інтерпретатора про помилку позначати error)\par
\begin{verbatim}
print(8//9/7*8)
\end{verbatim}

{\begin {center}\textbf{Завдання 2}\end {center}}

Що буде виведено за виконання? (повідомлення інтерпретатора про помилку позначати error)\par
\begin{verbatim}
print(-1**-2+False)
\end{verbatim}

{\begin {center}\textbf{Завдання 3}\end {center}}

Що буде виведено за виконання? (повідомлення інтерпретатора про помилку позначати error)\par
\begin{verbatim}
print(8.0<False or not 7>=4.0 and 8==2)
\end{verbatim}

{\begin {center}\textbf{Завдання 4}\end {center}}

Що буде виведено за виконання? (повідомлення інтерпретатора про помилку позначати error)\par
\begin{verbatim}
print("False" != "8.0")
\end{verbatim}

{\begin {center}\textbf{Завдання 5}\end {center}}

Що буде виведено за виконання? (повідомлення інтерпретатора про помилку позначати error)\par
\begin{verbatim}
print(7 != 8.0 > 9 < True)
\end{verbatim}

{\begin {center}\textbf{Завдання 6}\end {center}}

Що буде виведено за виконання? (повідомлення інтерпретатора про помилку позначати error)\par
\begin{verbatim}

def f(n):
    print("f", end="");
    return n<=-2

print(f(-8) or f(-3))
\end{verbatim}

{\begin {center}\textbf{Завдання 7}\end {center}}

Що буде виведено за виконання (повідомлення інтерпретатора про помилку позначати error)?\par
\begin{verbatim}
def g(d):
    u=63
    if d>=-4: 
        u=1
    elif d<4:
         u=6
    else:
         return 2
    return u

print(g(-6))
\end{verbatim}

{\begin {center}\textbf{Завдання 8}\end {center}}

Що буде виведено за виконання (повідомлення інтерпретатора про помилку позначати error)?\par
\begin{verbatim}
a, b, c = 9, 6, 8
def h(a, b, c=9):
    print(a, b, c, end=" ")

h(5, 4)
print(a, b, c)
\end{verbatim}

{\begin {center}\textbf{Завдання 9}\end {center}}

Що буде виведено за виконання (повідомлення інтерпретатора про помилку позначати error)?\par
\begin{verbatim}
a,b,c=7,9,2
def f(b):
    a=4
    b+=3
    c=2
    return a+b+c

a,b,c=6,8,5
print(f(b),a,b,c)
\end{verbatim}

{\begin {center}\textbf{Завдання 10}\end {center}}


Написати функцiю f, що отримує щось та повертає щось. Невелика загадка все-таки має залишатися. Але нічого принципово нового відносно задач, що наявні у pdf-ках не планується.. Стиль запису умови також оцінюється.\par
\end {large}}
\newpage

{\Large{17.10.2024 Бета-версія: Контрольна робота \No 1, варіант  \No \textbf { 48 }\par}}
{
\begin {large}
{\begin {center}\textbf{Завдання 1}\end {center}}

Що буде виведено за виконання? (повідомлення інтерпретатора про помилку позначати error)\par
\begin{verbatim}
print(7/7%9*9)
\end{verbatim}

{\begin {center}\textbf{Завдання 2}\end {center}}

Що буде виведено за виконання? (повідомлення інтерпретатора про помилку позначати error)\par
\begin{verbatim}
print(-2**4-2)
\end{verbatim}

{\begin {center}\textbf{Завдання 3}\end {center}}

Що буде виведено за виконання? (повідомлення інтерпретатора про помилку позначати error)\par
\begin{verbatim}
print(2.0<=True and not 5.0!=6 or 9>4)
\end{verbatim}

{\begin {center}\textbf{Завдання 4}\end {center}}

Що буде виведено за виконання? (повідомлення інтерпретатора про помилку позначати error)\par
\begin{verbatim}
print(True > "9")
\end{verbatim}

{\begin {center}\textbf{Завдання 5}\end {center}}

Що буде виведено за виконання? (повідомлення інтерпретатора про помилку позначати error)\par
\begin{verbatim}
print(6.0  is  5 == 2 <= 2.0)
\end{verbatim}

{\begin {center}\textbf{Завдання 6}\end {center}}

Що буде виведено за виконання? (повідомлення інтерпретатора про помилку позначати error)\par
\begin{verbatim}

def f(n):
    print("f", end="");
    return n

print(f(9) and f(7))
\end{verbatim}

{\begin {center}\textbf{Завдання 7}\end {center}}

Що буде виведено за виконання (повідомлення інтерпретатора про помилку позначати error)?\par
\begin{verbatim}
def f(a,b):
    c=42
    if b==0:
        c=9
    if b<-3:
         c=4
    else: 
        return 2
    return c

print(f(-7,-6))
\end{verbatim}

{\begin {center}\textbf{Завдання 8}\end {center}}

Що буде виведено за виконання (повідомлення інтерпретатора про помилку позначати error)?\par
\begin{verbatim}
a, b, c = 7, 7, 8
def f(a, b=9, c=6):
    print(a, b, c, end=" ")

f(a=0, b=1, c=3)
print(a, b, c)
\end{verbatim}

{\begin {center}\textbf{Завдання 9}\end {center}}

Що буде виведено за виконання (повідомлення інтерпретатора про помилку позначати error)?\par
\begin{verbatim}
a,b,c=1,3,4
def h(b):
    a=5
    b*=1
    c=5
    return a+b+c

a,b,c=0,8,6
print(h(a),a,b,c)
\end{verbatim}

{\begin {center}\textbf{Завдання 10}\end {center}}


Написати функцiю f, що отримує щось та повертає щось. Невелика загадка все-таки має залишатися. Але нічого принципово нового відносно задач, що наявні у pdf-ках не планується.. Стиль запису умови також оцінюється.\par
\end {large}}
\newpage

{\Large{17.10.2024 Бета-версія: Контрольна робота \No 1, варіант  \No \textbf { 49 }\par}}
{
\begin {large}
{\begin {center}\textbf{Завдання 1}\end {center}}

Що буде виведено за виконання? (повідомлення інтерпретатора про помилку позначати error)\par
\begin{verbatim}
print(6/8/3*7)
\end{verbatim}

{\begin {center}\textbf{Завдання 2}\end {center}}

Що буде виведено за виконання? (повідомлення інтерпретатора про помилку позначати error)\par
\begin{verbatim}
print(-2**-2-1)
\end{verbatim}

{\begin {center}\textbf{Завдання 3}\end {center}}

Що буде виведено за виконання? (повідомлення інтерпретатора про помилку позначати error)\par
\begin{verbatim}
print(9>=3 or not 2.0<4.0 or 8>True)
\end{verbatim}

{\begin {center}\textbf{Завдання 4}\end {center}}

Що буде виведено за виконання? (повідомлення інтерпретатора про помилку позначати error)\par
\begin{verbatim}
print(True > 6.0)
\end{verbatim}

{\begin {center}\textbf{Завдання 5}\end {center}}

Що буде виведено за виконання? (повідомлення інтерпретатора про помилку позначати error)\par
\begin{verbatim}
print(7 > 4 != 2.0 >= 3.0)
\end{verbatim}

{\begin {center}\textbf{Завдання 6}\end {center}}

Що буде виведено за виконання? (повідомлення інтерпретатора про помилку позначати error)\par
\begin{verbatim}

def f(n):
    print("f", end="");
    return n<=0

print(f(-3) or f(-9))
\end{verbatim}

{\begin {center}\textbf{Завдання 7}\end {center}}

Що буде виведено за виконання (повідомлення інтерпретатора про помилку позначати error)?\par
\begin{verbatim}
def g(a,b):
    c=99
    if b>2:
        return 5
    if a>=5:
         c=6
    else: 
        return 8
    return c

print(g(5,2))
\end{verbatim}

{\begin {center}\textbf{Завдання 8}\end {center}}

Що буде виведено за виконання (повідомлення інтерпретатора про помилку позначати error)?\par
\begin{verbatim}
a, b, c = 8, 9, 7
def h(a, b=6, c=9):
    print(a, b, c, end=" ")

h(2, 1, b=3)
print(a, b, c)
\end{verbatim}

{\begin {center}\textbf{Завдання 9}\end {center}}

Що буде виведено за виконання (повідомлення інтерпретатора про помилку позначати error)?\par
\begin{verbatim}
a,b,c=7,5,9
def f(a):
    a=4
    b-=3
    c=2
    return a+b+c

a,b,c=1,2,0
print(f(b),a,b,c)
\end{verbatim}

{\begin {center}\textbf{Завдання 10}\end {center}}


Написати функцiю f, що отримує щось та повертає щось. Невелика загадка все-таки має залишатися. Але нічого принципово нового відносно задач, що наявні у pdf-ках не планується.. Стиль запису умови також оцінюється.\par
\end {large}}
\newpage

{\Large{17.10.2024 Бета-версія: Контрольна робота \No 1, варіант  \No \textbf { 50 }\par}}
{
\begin {large}
{\begin {center}\textbf{Завдання 1}\end {center}}

Що буде виведено за виконання? (повідомлення інтерпретатора про помилку позначати error)\par
\begin{verbatim}
print(5//9*7*5)
\end{verbatim}

{\begin {center}\textbf{Завдання 2}\end {center}}

Що буде виведено за виконання? (повідомлення інтерпретатора про помилку позначати error)\par
\begin{verbatim}
print(2**False+2)
\end{verbatim}

{\begin {center}\textbf{Завдання 3}\end {center}}

Що буде виведено за виконання? (повідомлення інтерпретатора про помилку позначати error)\par
\begin{verbatim}
print(2<=False and not 8.0==4 and 3.0!=5)
\end{verbatim}

{\begin {center}\textbf{Завдання 4}\end {center}}

Що буде виведено за виконання? (повідомлення інтерпретатора про помилку позначати error)\par
\begin{verbatim}
print("False" < 7)
\end{verbatim}

{\begin {center}\textbf{Завдання 5}\end {center}}

Що буде виведено за виконання? (повідомлення інтерпретатора про помилку позначати error)\par
\begin{verbatim}
print(5 == True != 8 <= 4.0)
\end{verbatim}

{\begin {center}\textbf{Завдання 6}\end {center}}

Що буде виведено за виконання? (повідомлення інтерпретатора про помилку позначати error)\par
\begin{verbatim}

def f(n):
    print("f", end="");
    return n

print(f(4) and f(2))
\end{verbatim}

{\begin {center}\textbf{Завдання 7}\end {center}}

Що буде виведено за виконання (повідомлення інтерпретатора про помилку позначати error)?\par
\begin{verbatim}
def h(a,b):
    c=33
    if a>1:
        c=9
    elif b>=4:
         c=0
    else: 
        return 1
    return c

print(h(-5,7))
\end{verbatim}

{\begin {center}\textbf{Завдання 8}\end {center}}

Що буде виведено за виконання (повідомлення інтерпретатора про помилку позначати error)?\par
\begin{verbatim}
a, b, c = 6, 7, 8
def f(a, b=6, c):
    print(a, b, c, end=" ")

f(5, c=0, b=4)
print(a, b, c)
\end{verbatim}

{\begin {center}\textbf{Завдання 9}\end {center}}

Що буде виведено за виконання (повідомлення інтерпретатора про помилку позначати error)?\par
\begin{verbatim}
a,b,c=4,3,8
def g(b):
    a=1
    b+=1
    c=2
    return a+b+c

a,b,c=0,9,4
print(g(a),a,b,c)
\end{verbatim}

{\begin {center}\textbf{Завдання 10}\end {center}}


Написати функцiю f, що отримує щось та повертає щось. Невелика загадка все-таки має залишатися. Але нічого принципово нового відносно задач, що наявні у pdf-ках не планується.. Стиль запису умови також оцінюється.\par
\end {large}}
\newpage

{\Large{17.10.2024 Бета-версія: Контрольна робота \No 1, варіант  \No \textbf { 51 }\par}}
{
\begin {large}
{\begin {center}\textbf{Завдання 1}\end {center}}

Що буде виведено за виконання? (повідомлення інтерпретатора про помилку позначати error)\par
\begin{verbatim}
print(3/5//9*7)
\end{verbatim}

{\begin {center}\textbf{Завдання 2}\end {center}}

Що буде виведено за виконання? (повідомлення інтерпретатора про помилку позначати error)\par
\begin{verbatim}
print(9**-0.5*True)
\end{verbatim}

{\begin {center}\textbf{Завдання 3}\end {center}}

Що буде виведено за виконання? (повідомлення інтерпретатора про помилку позначати error)\par
\begin{verbatim}
print(not 6==6.0 and True>5.0 or 7!=6)
\end{verbatim}

{\begin {center}\textbf{Завдання 4}\end {center}}

Що буде виведено за виконання? (повідомлення інтерпретатора про помилку позначати error)\par
\begin{verbatim}
print("1.0" == "5")
\end{verbatim}

{\begin {center}\textbf{Завдання 5}\end {center}}

Що буде виведено за виконання? (повідомлення інтерпретатора про помилку позначати error)\par
\begin{verbatim}
print(False >= 2  is  9 > 7.0)
\end{verbatim}

{\begin {center}\textbf{Завдання 6}\end {center}}

Що буде виведено за виконання? (повідомлення інтерпретатора про помилку позначати error)\par
\begin{verbatim}

def f(n):
    print("f", end="");
    return n>=4

print(f(-7) and f(8))
\end{verbatim}

{\begin {center}\textbf{Завдання 7}\end {center}}

Що буде виведено за виконання (повідомлення інтерпретатора про помилку позначати error)?\par
\begin{verbatim}
def f(a,b):
    c=30
    if a:
        return 3
    if b==-5:
         c=7
    else: 
        c=8
    return c

print(f(-9,-9))
\end{verbatim}

{\begin {center}\textbf{Завдання 8}\end {center}}

Що буде виведено за виконання (повідомлення інтерпретатора про помилку позначати error)?\par
\begin{verbatim}
a, b, c = 8, 9, 7
def g(a, b, c=6):
    print(a, b, c, end=" ")

g(a=5, 0, a=4)
print(a, b, c)
\end{verbatim}

{\begin {center}\textbf{Завдання 9}\end {center}}

Що буде виведено за виконання (повідомлення інтерпретатора про помилку позначати error)?\par
\begin{verbatim}
a,b,c=2,3,8
def f(b):
    a=1
    b=2
    c=5
    return a+b+c

a,b,c=0,1,0
print(f(a),a,b,c)
\end{verbatim}

{\begin {center}\textbf{Завдання 10}\end {center}}


Написати функцiю f, що отримує щось та повертає щось. Невелика загадка все-таки має залишатися. Але нічого принципово нового відносно задач, що наявні у pdf-ках не планується.. Стиль запису умови також оцінюється.\par
\end {large}}
\newpage

{\Large{17.10.2024 Бета-версія: Контрольна робота \No 1, варіант  \No \textbf { 52 }\par}}
{
\begin {large}
{\begin {center}\textbf{Завдання 1}\end {center}}

Що буде виведено за виконання? (повідомлення інтерпретатора про помилку позначати error)\par
\begin{verbatim}
print(9//8%3*3)
\end{verbatim}

{\begin {center}\textbf{Завдання 2}\end {center}}

Що буде виведено за виконання? (повідомлення інтерпретатора про помилку позначати error)\par
\begin{verbatim}
print(3**-2+-1)
\end{verbatim}

{\begin {center}\textbf{Завдання 3}\end {center}}

Що буде виведено за виконання? (повідомлення інтерпретатора про помилку позначати error)\par
\begin{verbatim}
print(not 5>=9.0 or False<=9 and 7.0<4)
\end{verbatim}

{\begin {center}\textbf{Завдання 4}\end {center}}

Що буде виведено за виконання? (повідомлення інтерпретатора про помилку позначати error)\par
\begin{verbatim}
print("False" >= True)
\end{verbatim}

{\begin {center}\textbf{Завдання 5}\end {center}}

Що буде виведено за виконання? (повідомлення інтерпретатора про помилку позначати error)\par
\begin{verbatim}
print(9.0 < False == 6 <= 6.0)
\end{verbatim}

{\begin {center}\textbf{Завдання 6}\end {center}}

Що буде виведено за виконання? (повідомлення інтерпретатора про помилку позначати error)\par
\begin{verbatim}

def f(n):
    print("f", end="");
    return n

print(f(9) or f(3))
\end{verbatim}

{\begin {center}\textbf{Завдання 7}\end {center}}

Що буде виведено за виконання (повідомлення інтерпретатора про помилку позначати error)?\par
\begin{verbatim}
def f(a):
    w=85
    if a>3: 
        w=4
    elif a==1:
         return 5
    else:
         w=3
    return w

print(f(-9))
\end{verbatim}

{\begin {center}\textbf{Завдання 8}\end {center}}

Що буде виведено за виконання (повідомлення інтерпретатора про помилку позначати error)?\par
\begin{verbatim}
a, b, c = 8, 7, 9
def g(a, b, c):
    print(a, b, c, end=" ")

g(2, 1, 3)
print(a, b, c)
\end{verbatim}

{\begin {center}\textbf{Завдання 9}\end {center}}

Що буде виведено за виконання (повідомлення інтерпретатора про помилку позначати error)?\par
\begin{verbatim}
a,b,c=5,2,3
def h(a):
    a=4
    b+=5
    c=3
    return a+b+c

a,b,c=6,7,1
print(h(b),a,b,c)
\end{verbatim}

{\begin {center}\textbf{Завдання 10}\end {center}}


Написати функцiю f, що отримує щось та повертає щось. Невелика загадка все-таки має залишатися. Але нічого принципово нового відносно задач, що наявні у pdf-ках не планується.. Стиль запису умови також оцінюється.\par
\end {large}}
\newpage

{\Large{17.10.2024 Бета-версія: Контрольна робота \No 1, варіант  \No \textbf { 53 }\par}}
{
\begin {large}
{\begin {center}\textbf{Завдання 1}\end {center}}

Що буде виведено за виконання? (повідомлення інтерпретатора про помилку позначати error)\par
\begin{verbatim}
print(7//7*4*8)
\end{verbatim}

{\begin {center}\textbf{Завдання 2}\end {center}}

Що буде виведено за виконання? (повідомлення інтерпретатора про помилку позначати error)\par
\begin{verbatim}
print(4**-1.0*False)
\end{verbatim}

{\begin {center}\textbf{Завдання 3}\end {center}}

Що буде виведено за виконання? (повідомлення інтерпретатора про помилку позначати error)\par
\begin{verbatim}
print(not 7>False and 2<=4.0 or 3<7.0)
\end{verbatim}

{\begin {center}\textbf{Завдання 4}\end {center}}

Що буде виведено за виконання? (повідомлення інтерпретатора про помилку позначати error)\par
\begin{verbatim}
print(4 <= 5.0)
\end{verbatim}

{\begin {center}\textbf{Завдання 5}\end {center}}

Що буде виведено за виконання? (повідомлення інтерпретатора про помилку позначати error)\par
\begin{verbatim}
print(3  is  6 >= 8.0 < 5.0)
\end{verbatim}

{\begin {center}\textbf{Завдання 6}\end {center}}

Що буде виведено за виконання? (повідомлення інтерпретатора про помилку позначати error)\par
\begin{verbatim}

def f(n):
    print("f", end="");
    return n

print(f(0) or f(-8))
\end{verbatim}

{\begin {center}\textbf{Завдання 7}\end {center}}

Що буде виведено за виконання (повідомлення інтерпретатора про помилку позначати error)?\par
\begin{verbatim}
def h(c):
    z=44
    if c: 
        return 1
    if c>2:
         return 8
    else:
         z=9
    return z

print(h(9))
\end{verbatim}

{\begin {center}\textbf{Завдання 8}\end {center}}

Що буде виведено за виконання (повідомлення інтерпретатора про помилку позначати error)?\par
\begin{verbatim}
a, b, c = 6, 9, 8
def h(a, b, c):
    print(a, b, c, end=" ")

h(b=4, c=5, 2)
print(a, b, c)
\end{verbatim}

{\begin {center}\textbf{Завдання 9}\end {center}}

Що буде виведено за виконання (повідомлення інтерпретатора про помилку позначати error)?\par
\begin{verbatim}
a,b,c=5,4,8
def f(a):
    a=4
    b*=5
    c=3
    return a+b+c

a,b,c=1,7,9
print(f(b),a,b,c)
\end{verbatim}

{\begin {center}\textbf{Завдання 10}\end {center}}


Написати функцiю f, що отримує щось та повертає щось. Невелика загадка все-таки має залишатися. Але нічого принципово нового відносно задач, що наявні у pdf-ках не планується.. Стиль запису умови також оцінюється.\par
\end {large}}
\newpage

{\Large{17.10.2024 Бета-версія: Контрольна робота \No 1, варіант  \No \textbf { 54 }\par}}
{
\begin {large}
{\begin {center}\textbf{Завдання 1}\end {center}}

Що буде виведено за виконання? (повідомлення інтерпретатора про помилку позначати error)\par
\begin{verbatim}
print(4/3//8*4)
\end{verbatim}

{\begin {center}\textbf{Завдання 2}\end {center}}

Що буде виведено за виконання? (повідомлення інтерпретатора про помилку позначати error)\par
\begin{verbatim}
print(1**0.5*-2)
\end{verbatim}

{\begin {center}\textbf{Завдання 3}\end {center}}

Що буде виведено за виконання? (повідомлення інтерпретатора про помилку позначати error)\par
\begin{verbatim}
print(not 8.0==5.0 or 8!=8 and True>=9)
\end{verbatim}

{\begin {center}\textbf{Завдання 4}\end {center}}

Що буде виведено за виконання? (повідомлення інтерпретатора про помилку позначати error)\par
\begin{verbatim}
print("2j" != "3.0j")
\end{verbatim}

{\begin {center}\textbf{Завдання 5}\end {center}}

Що буде виведено за виконання? (повідомлення інтерпретатора про помилку позначати error)\par
\begin{verbatim}
print(3 != 5 > 7 != True)
\end{verbatim}

{\begin {center}\textbf{Завдання 6}\end {center}}

Що буде виведено за виконання? (повідомлення інтерпретатора про помилку позначати error)\par
\begin{verbatim}

def f(n):
    print("f", end="");
    return n>1

print(f(-5) and f(-6))
\end{verbatim}

{\begin {center}\textbf{Завдання 7}\end {center}}

Що буде виведено за виконання (повідомлення інтерпретатора про помилку позначати error)?\par
\begin{verbatim}
def h(a,b):
    c=71
    if a!=5:
        return 4
    elif a<=0:
         return 7
    else: 
        c=0
    return c

print(h(-6,-5))
\end{verbatim}

{\begin {center}\textbf{Завдання 8}\end {center}}

Що буде виведено за виконання (повідомлення інтерпретатора про помилку позначати error)?\par
\begin{verbatim}
a, b, c = 7, 6, 9
def f(a, b=8, c=7):
    print(a, b, c, end=" ")

f(1, 0)
print(a, b, c)
\end{verbatim}

{\begin {center}\textbf{Завдання 9}\end {center}}

Що буде виведено за виконання (повідомлення інтерпретатора про помилку позначати error)?\par
\begin{verbatim}
a,b,c=2,8,5
def g(a):
    a=1
    b=4
    c=3
    return a+b+c

a,b,c=4,3,6
print(g(a),a,b,c)
\end{verbatim}

{\begin {center}\textbf{Завдання 10}\end {center}}


Написати функцiю f, що отримує щось та повертає щось. Невелика загадка все-таки має залишатися. Але нічого принципово нового відносно задач, що наявні у pdf-ках не планується.. Стиль запису умови також оцінюється.\par
\end {large}}
\newpage

{\Large{17.10.2024 Бета-версія: Контрольна робота \No 1, варіант  \No \textbf { 55 }\par}}
{
\begin {large}
{\begin {center}\textbf{Завдання 1}\end {center}}

Що буде виведено за виконання? (повідомлення інтерпретатора про помилку позначати error)\par
\begin{verbatim}
print(8/6/5*9)
\end{verbatim}

{\begin {center}\textbf{Завдання 2}\end {center}}

Що буде виведено за виконання? (повідомлення інтерпретатора про помилку позначати error)\par
\begin{verbatim}
print(1**1.0*False)
\end{verbatim}

{\begin {center}\textbf{Завдання 3}\end {center}}

Що буде виведено за виконання? (повідомлення інтерпретатора про помилку позначати error)\par
\begin{verbatim}
print(not 6.0<True or 3==4 and 7>3.0)
\end{verbatim}

{\begin {center}\textbf{Завдання 4}\end {center}}

Що буде виведено за виконання? (повідомлення інтерпретатора про помилку позначати error)\par
\begin{verbatim}
print(False < "7j")
\end{verbatim}

{\begin {center}\textbf{Завдання 5}\end {center}}

Що буде виведено за виконання? (повідомлення інтерпретатора про помилку позначати error)\par
\begin{verbatim}
print(8 < 8.0  is  4 == 6.0)
\end{verbatim}

{\begin {center}\textbf{Завдання 6}\end {center}}

Що буде виведено за виконання? (повідомлення інтерпретатора про помилку позначати error)\par
\begin{verbatim}

def f(n):
    print("f", end="");
    return n==-3

print(f(1) or f(-4))
\end{verbatim}

{\begin {center}\textbf{Завдання 7}\end {center}}

Що буде виведено за виконання (повідомлення інтерпретатора про помилку позначати error)?\par
\begin{verbatim}
def g(b):
    u=36
    if b>=0: 
        u=0
    if b!=4:
         return 7
    else:
         u=8
    return u

print(g(-4))
\end{verbatim}

{\begin {center}\textbf{Завдання 8}\end {center}}

Що буде виведено за виконання (повідомлення інтерпретатора про помилку позначати error)?\par
\begin{verbatim}
a, b, c = 7, 8, 9
def h(a, b=6, c):
    print(a, b, c, end=" ")

h(c=3, c=0, b=1)
print(a, b, c)
\end{verbatim}

{\begin {center}\textbf{Завдання 9}\end {center}}

Що буде виведено за виконання (повідомлення інтерпретатора про помилку позначати error)?\par
\begin{verbatim}
a,b,c=0,3,6
def g(b):
    a-=2
    b=1
    c=5
    return a+b+c

a,b,c=2,6,8
print(g(a),a,b,c)
\end{verbatim}

{\begin {center}\textbf{Завдання 10}\end {center}}


Написати функцiю f, що отримує щось та повертає щось. Невелика загадка все-таки має залишатися. Але нічого принципово нового відносно задач, що наявні у pdf-ках не планується.. Стиль запису умови також оцінюється.\par
\end {large}}
\newpage

{\Large{17.10.2024 Бета-версія: Контрольна робота \No 1, варіант  \No \textbf { 56 }\par}}
{
\begin {large}
{\begin {center}\textbf{Завдання 1}\end {center}}

Що буде виведено за виконання? (повідомлення інтерпретатора про помилку позначати error)\par
\begin{verbatim}
print(6//4%6*6)
\end{verbatim}

{\begin {center}\textbf{Завдання 2}\end {center}}

Що буде виведено за виконання? (повідомлення інтерпретатора про помилку позначати error)\par
\begin{verbatim}
print(-1**2-1)
\end{verbatim}

{\begin {center}\textbf{Завдання 3}\end {center}}

Що буде виведено за виконання? (повідомлення інтерпретатора про помилку позначати error)\par
\begin{verbatim}
print(not 6!=5 and 9.0>=False or 2<=2.0)
\end{verbatim}

{\begin {center}\textbf{Завдання 4}\end {center}}

Що буде виведено за виконання? (повідомлення інтерпретатора про помилку позначати error)\par
\begin{verbatim}
print("2.0" == 8.0)
\end{verbatim}

{\begin {center}\textbf{Завдання 5}\end {center}}

Що буде виведено за виконання? (повідомлення інтерпретатора про помилку позначати error)\par
\begin{verbatim}
print(2 >= False > 9 <= True)
\end{verbatim}

{\begin {center}\textbf{Завдання 6}\end {center}}

Що буде виведено за виконання? (повідомлення інтерпретатора про помилку позначати error)\par
\begin{verbatim}

def f(n):
    print("f", end="");
    return n

print(f(6) and f(-1))
\end{verbatim}

{\begin {center}\textbf{Завдання 7}\end {center}}

Що буде виведено за виконання (повідомлення інтерпретатора про помилку позначати error)?\par
\begin{verbatim}
def h(d):
    v=35
    if d<5: 
        return 9
    elif d<=-3:
         return 5
    else:
         v=1
    return v

print(h(-5))
\end{verbatim}

{\begin {center}\textbf{Завдання 8}\end {center}}

Що буде виведено за виконання (повідомлення інтерпретатора про помилку позначати error)?\par
\begin{verbatim}
a, b, c = 8, 7, 6
def f(a, b, c=9):
    print(a, b, c, end=" ")

f(2, a=4)
print(a, b, c)
\end{verbatim}

{\begin {center}\textbf{Завдання 9}\end {center}}

Що буде виведено за виконання (повідомлення інтерпретатора про помилку позначати error)?\par
\begin{verbatim}
a,b,c=0,7,3
def g(a):
    a+=1
    b=4
    c=3
    return a+b+c

a,b,c=9,2,5
print(g(b),a,b,c)
\end{verbatim}

{\begin {center}\textbf{Завдання 10}\end {center}}


Написати функцiю f, що отримує щось та повертає щось. Невелика загадка все-таки має залишатися. Але нічого принципово нового відносно задач, що наявні у pdf-ках не планується.. Стиль запису умови також оцінюється.\par
\end {large}}
\newpage

{\Large{17.10.2024 Бета-версія: Контрольна робота \No 1, варіант  \No \textbf { 57 }\par}}
{
\begin {large}
{\begin {center}\textbf{Завдання 1}\end {center}}

Що буде виведено за виконання? (повідомлення інтерпретатора про помилку позначати error)\par
\begin{verbatim}
print(9/8*6*9)
\end{verbatim}

{\begin {center}\textbf{Завдання 2}\end {center}}

Що буде виведено за виконання? (повідомлення інтерпретатора про помилку позначати error)\par
\begin{verbatim}
print(2**-1+-2)
\end{verbatim}

{\begin {center}\textbf{Завдання 3}\end {center}}

Що буде виведено за виконання? (повідомлення інтерпретатора про помилку позначати error)\par
\begin{verbatim}
print(not 3!=6.0 or 9.0>=5 or True<=8)
\end{verbatim}

{\begin {center}\textbf{Завдання 4}\end {center}}

Що буде виведено за виконання? (повідомлення інтерпретатора про помилку позначати error)\par
\begin{verbatim}
print(8 <= "True")
\end{verbatim}

{\begin {center}\textbf{Завдання 5}\end {center}}

Що буде виведено за виконання? (повідомлення інтерпретатора про помилку позначати error)\par
\begin{verbatim}
print(6 >= 7.0 < 5 == 5.0)
\end{verbatim}

{\begin {center}\textbf{Завдання 6}\end {center}}

Що буде виведено за виконання? (повідомлення інтерпретатора про помилку позначати error)\par
\begin{verbatim}

def f(n):
    print("f", end="");
    return n!=3

print(f(-2) and f(5))
\end{verbatim}

{\begin {center}\textbf{Завдання 7}\end {center}}

Що буде виведено за виконання (повідомлення інтерпретатора про помилку позначати error)?\par
\begin{verbatim}
def f(a,b):
    c=85
    if b:
        c=9
    if b<4:
         return 6
    else: 
        c=1
    return c

print(f(-2,9))
\end{verbatim}

{\begin {center}\textbf{Завдання 8}\end {center}}

Що буде виведено за виконання (повідомлення інтерпретатора про помилку позначати error)?\par
\begin{verbatim}
a, b, c = 9, 7, 6
def g(a, b=8, c):
    print(a, b, c, end=" ")

g(b=3, c=5, a=5)
print(a, b, c)
\end{verbatim}

{\begin {center}\textbf{Завдання 9}\end {center}}

Що буде виведено за виконання (повідомлення інтерпретатора про помилку позначати error)?\par
\begin{verbatim}
a,b,c=4,1,7
def h(b):
    a-=2
    b=5
    c=3
    return a+b+c

a,b,c=6,5,8
print(h(b),a,b,c)
\end{verbatim}

{\begin {center}\textbf{Завдання 10}\end {center}}


Написати функцiю f, що отримує щось та повертає щось. Невелика загадка все-таки має залишатися. Але нічого принципово нового відносно задач, що наявні у pdf-ках не планується.. Стиль запису умови також оцінюється.\par
\end {large}}
\newpage

{\Large{17.10.2024 Бета-версія: Контрольна робота \No 1, варіант  \No \textbf { 58 }\par}}
{
\begin {large}
{\begin {center}\textbf{Завдання 1}\end {center}}

Що буде виведено за виконання? (повідомлення інтерпретатора про помилку позначати error)\par
\begin{verbatim}
print(8//6/9*3)
\end{verbatim}

{\begin {center}\textbf{Завдання 2}\end {center}}

Що буде виведено за виконання? (повідомлення інтерпретатора про помилку позначати error)\par
\begin{verbatim}
print(9**-1.0*True)
\end{verbatim}

{\begin {center}\textbf{Завдання 3}\end {center}}

Що буде виведено за виконання? (повідомлення інтерпретатора про помилку позначати error)\par
\begin{verbatim}
print(not 7.0>4 and False<9 and 5.0==6)
\end{verbatim}

{\begin {center}\textbf{Завдання 4}\end {center}}

Що буде виведено за виконання? (повідомлення інтерпретатора про помилку позначати error)\par
\begin{verbatim}
print("True" != 5)
\end{verbatim}

{\begin {center}\textbf{Завдання 5}\end {center}}

Що буде виведено за виконання? (повідомлення інтерпретатора про помилку позначати error)\par
\begin{verbatim}
print(4.0 <= 4  is  True != 3)
\end{verbatim}

{\begin {center}\textbf{Завдання 6}\end {center}}

Що буде виведено за виконання? (повідомлення інтерпретатора про помилку позначати error)\par
\begin{verbatim}

def f(n):
    print("f", end="");
    return n

print(f(-3) or f(1))
\end{verbatim}

{\begin {center}\textbf{Завдання 7}\end {center}}

Що буде виведено за виконання (повідомлення інтерпретатора про помилку позначати error)?\par
\begin{verbatim}
def g(a,b):
    c=45
    if a<-4:
        return 2
    elif b<=-1:
         c=3
    else: 
        return 5
    return c

print(g(4,6))
\end{verbatim}

{\begin {center}\textbf{Завдання 8}\end {center}}

Що буде виведено за виконання (повідомлення інтерпретатора про помилку позначати error)?\par
\begin{verbatim}
a, b, c = 7, 8, 6
def h(a, b, c=9):
    print(a, b, c, end=" ")

h(b=4, c=0, 3)
print(a, b, c)
\end{verbatim}

{\begin {center}\textbf{Завдання 9}\end {center}}

Що буде виведено за виконання (повідомлення інтерпретатора про помилку позначати error)?\par
\begin{verbatim}
a,b,c=2,9,0
def f(b):
    a*=2
    b=4
    c=1
    return a+b+c

a,b,c=1,4,3
print(f(b),a,b,c)
\end{verbatim}

{\begin {center}\textbf{Завдання 10}\end {center}}


Написати функцiю f, що отримує щось та повертає щось. Невелика загадка все-таки має залишатися. Але нічого принципово нового відносно задач, що наявні у pdf-ках не планується.. Стиль запису умови також оцінюється.\par
\end {large}}
\newpage

{\Large{17.10.2024 Бета-версія: Контрольна робота \No 1, варіант  \No \textbf { 59 }\par}}
{
\begin {large}
{\begin {center}\textbf{Завдання 1}\end {center}}

Що буде виведено за виконання? (повідомлення інтерпретатора про помилку позначати error)\par
\begin{verbatim}
print(6/7%5*8)
\end{verbatim}

{\begin {center}\textbf{Завдання 2}\end {center}}

Що буде виведено за виконання? (повідомлення інтерпретатора про помилку позначати error)\par
\begin{verbatim}
print(4**True--1)
\end{verbatim}

{\begin {center}\textbf{Завдання 3}\end {center}}

Що буде виведено за виконання? (повідомлення інтерпретатора про помилку позначати error)\par
\begin{verbatim}
print(not 4.0==2 and 7>=6 and True<3.0)
\end{verbatim}

{\begin {center}\textbf{Завдання 4}\end {center}}

Що буде виведено за виконання? (повідомлення інтерпретатора про помилку позначати error)\par
\begin{verbatim}
print(1.0 >= "2")
\end{verbatim}

{\begin {center}\textbf{Завдання 5}\end {center}}

Що буде виведено за виконання? (повідомлення інтерпретатора про помилку позначати error)\par
\begin{verbatim}
print(2.0 > 7 != 9.0 >= 3.0)
\end{verbatim}

{\begin {center}\textbf{Завдання 6}\end {center}}

Що буде виведено за виконання? (повідомлення інтерпретатора про помилку позначати error)\par
\begin{verbatim}

def f(n):
    print("f", end="");
    return n

print(f(4) and f(-6))
\end{verbatim}

{\begin {center}\textbf{Завдання 7}\end {center}}

Що буде виведено за виконання (повідомлення інтерпретатора про помилку позначати error)?\par
\begin{verbatim}
def f(d):
    x=58
    if d: 
        return 3
    if d==-5:
         x=0
    else:
         return 7
    return x

print(f(5))
\end{verbatim}

{\begin {center}\textbf{Завдання 8}\end {center}}

Що буде виведено за виконання (повідомлення інтерпретатора про помилку позначати error)?\par
\begin{verbatim}
a, b, c = 7, 9, 6
def g(a, b, c):
    print(a, b, c, end=" ")

g(a=1, 2, b=4)
print(a, b, c)
\end{verbatim}

{\begin {center}\textbf{Завдання 9}\end {center}}

Що буде виведено за виконання (повідомлення інтерпретатора про помилку позначати error)?\par
\begin{verbatim}
a,b,c=7,9,7
def f(b):
    a=2
    b=1
    c=3
    return a+b+c

a,b,c=3,5,8
print(f(b),a,b,c)
\end{verbatim}

{\begin {center}\textbf{Завдання 10}\end {center}}


Написати функцiю f, що отримує щось та повертає щось. Невелика загадка все-таки має залишатися. Але нічого принципово нового відносно задач, що наявні у pdf-ках не планується.. Стиль запису умови також оцінюється.\par
\end {large}}
\newpage

{\Large{17.10.2024 Бета-версія: Контрольна робота \No 1, варіант  \No \textbf { 60 }\par}}
{
\begin {large}
{\begin {center}\textbf{Завдання 1}\end {center}}

Що буде виведено за виконання? (повідомлення інтерпретатора про помилку позначати error)\par
\begin{verbatim}
print(3//4//7*7)
\end{verbatim}

{\begin {center}\textbf{Завдання 2}\end {center}}

Що буде виведено за виконання? (повідомлення інтерпретатора про помилку позначати error)\par
\begin{verbatim}
print(4**1.0*2)
\end{verbatim}

{\begin {center}\textbf{Завдання 3}\end {center}}

Що буде виведено за виконання? (повідомлення інтерпретатора про помилку позначати error)\par
\begin{verbatim}
print(False<=8.0 or not 5>4 or 8!=2.0)
\end{verbatim}

{\begin {center}\textbf{Завдання 4}\end {center}}

Що буде виведено за виконання? (повідомлення інтерпретатора про помилку позначати error)\par
\begin{verbatim}
print("6.0j" > False)
\end{verbatim}

{\begin {center}\textbf{Завдання 5}\end {center}}

Що буде виведено за виконання? (повідомлення інтерпретатора про помилку позначати error)\par
\begin{verbatim}
print(False <= 9 == 8  is  2)
\end{verbatim}

{\begin {center}\textbf{Завдання 6}\end {center}}

Що буде виведено за виконання? (повідомлення інтерпретатора про помилку позначати error)\par
\begin{verbatim}

def f(n):
    print("f", end="");
    return n<-4

print(f(-8) or f(2))
\end{verbatim}

{\begin {center}\textbf{Завдання 7}\end {center}}

Що буде виведено за виконання (повідомлення інтерпретатора про помилку позначати error)?\par
\begin{verbatim}
def g(b):
    y=15
    if b<-2: 
        y=6
    elif b<=0:
         return 4
    else:
         y=2
    return y

print(g(1))
\end{verbatim}

{\begin {center}\textbf{Завдання 8}\end {center}}

Що буде виведено за виконання (повідомлення інтерпретатора про помилку позначати error)?\par
\begin{verbatim}
a, b, c = 8, 9, 8
def f(a, b=7, c=6):
    print(a, b, c, end=" ")

f(1, 2, a=0)
print(a, b, c)
\end{verbatim}

{\begin {center}\textbf{Завдання 9}\end {center}}

Що буде виведено за виконання (повідомлення інтерпретатора про помилку позначати error)?\par
\begin{verbatim}
a,b,c=9,6,1
def h(b):
    a=4
    b=5
    c=2
    return a+b+c

a,b,c=2,4,0
print(h(a),a,b,c)
\end{verbatim}

{\begin {center}\textbf{Завдання 10}\end {center}}


Написати функцiю f, що отримує щось та повертає щось. Невелика загадка все-таки має залишатися. Але нічого принципово нового відносно задач, що наявні у pdf-ках не планується.. Стиль запису умови також оцінюється.\par
\end {large}}
\newpage

{\Large{17.10.2024 Бета-версія: Контрольна робота \No 1, варіант  \No \textbf { 61 }\par}}
{
\begin {large}
{\begin {center}\textbf{Завдання 1}\end {center}}

Що буде виведено за виконання? (повідомлення інтерпретатора про помилку позначати error)\par
\begin{verbatim}
print(7//9*3*4)
\end{verbatim}

{\begin {center}\textbf{Завдання 2}\end {center}}

Що буде виведено за виконання? (повідомлення інтерпретатора про помилку позначати error)\par
\begin{verbatim}
print(1**-4-2)
\end{verbatim}

{\begin {center}\textbf{Завдання 3}\end {center}}

Що буде виведено за виконання? (повідомлення інтерпретатора про помилку позначати error)\par
\begin{verbatim}
print(not 2==6.0 and True<9 or 7.0>7)
\end{verbatim}

{\begin {center}\textbf{Завдання 4}\end {center}}

Що буде виведено за виконання? (повідомлення інтерпретатора про помилку позначати error)\par
\begin{verbatim}
print("3.0" >= 6)
\end{verbatim}

{\begin {center}\textbf{Завдання 5}\end {center}}

Що буде виведено за виконання? (повідомлення інтерпретатора про помилку позначати error)\par
\begin{verbatim}
print(4.0 > 8.0 < False >= 7)
\end{verbatim}

{\begin {center}\textbf{Завдання 6}\end {center}}

Що буде виведено за виконання? (повідомлення інтерпретатора про помилку позначати error)\par
\begin{verbatim}

def f(n):
    print("f", end="");
    return n

print(f(-2) and f(-7))
\end{verbatim}

{\begin {center}\textbf{Завдання 7}\end {center}}

Що буде виведено за виконання (повідомлення інтерпретатора про помилку позначати error)?\par
\begin{verbatim}
def f(a,b):
    c=25
    if a:
        c=4
    if a!=1:
         return 9
    else: 
        return 6
    return c

print(f(-3,-8))
\end{verbatim}

{\begin {center}\textbf{Завдання 8}\end {center}}

Що буде виведено за виконання (повідомлення інтерпретатора про помилку позначати error)?\par
\begin{verbatim}
a, b, c = 6, 7, 8
def g(a, b, c=9):
    print(a, b, c, end=" ")

g(3, 5)
print(a, b, c)
\end{verbatim}

{\begin {center}\textbf{Завдання 9}\end {center}}

Що буде виведено за виконання (повідомлення інтерпретатора про помилку позначати error)?\par
\begin{verbatim}
a,b,c=5,4,9
def h(a):
    a*=3
    b=4
    c=2
    return a+b+c

a,b,c=1,7,3
print(h(a),a,b,c)
\end{verbatim}

{\begin {center}\textbf{Завдання 10}\end {center}}


Написати функцiю f, що отримує щось та повертає щось. Невелика загадка все-таки має залишатися. Але нічого принципово нового відносно задач, що наявні у pdf-ках не планується.. Стиль запису умови також оцінюється.\par
\end {large}}
\newpage

{\Large{17.10.2024 Бета-версія: Контрольна робота \No 1, варіант  \No \textbf { 62 }\par}}
{
\begin {large}
{\begin {center}\textbf{Завдання 1}\end {center}}

Що буде виведено за виконання? (повідомлення інтерпретатора про помилку позначати error)\par
\begin{verbatim}
print(5/3/4*6)
\end{verbatim}

{\begin {center}\textbf{Завдання 2}\end {center}}

Що буде виведено за виконання? (повідомлення інтерпретатора про помилку позначати error)\par
\begin{verbatim}
print(1**-0.5*-1)
\end{verbatim}

{\begin {center}\textbf{Завдання 3}\end {center}}

Що буде виведено за виконання? (повідомлення інтерпретатора про помилку позначати error)\par
\begin{verbatim}
print(not 8.0>=False or 3<=9 and 3.0!=5)
\end{verbatim}

{\begin {center}\textbf{Завдання 4}\end {center}}

Що буде виведено за виконання? (повідомлення інтерпретатора про помилку позначати error)\par
\begin{verbatim}
print(True < "3")
\end{verbatim}

{\begin {center}\textbf{Завдання 5}\end {center}}

Що буде виведено за виконання? (повідомлення інтерпретатора про помилку позначати error)\par
\begin{verbatim}
print(8 == 3 <= 6.0 < 5.0)
\end{verbatim}

{\begin {center}\textbf{Завдання 6}\end {center}}

Що буде виведено за виконання? (повідомлення інтерпретатора про помилку позначати error)\par
\begin{verbatim}

def f(n):
    print("f", end="");
    return n==4

print(f(6) or f(7))
\end{verbatim}

{\begin {center}\textbf{Завдання 7}\end {center}}

Що буде виведено за виконання (повідомлення інтерпретатора про помилку позначати error)?\par
\begin{verbatim}
def g(a,b):
    c=12
    if a>3:
        c=1
    elif b==-2:
         return 2
    else: 
        c=0
    return c

print(g(2,-1))
\end{verbatim}

{\begin {center}\textbf{Завдання 8}\end {center}}

Що буде виведено за виконання (повідомлення інтерпретатора про помилку позначати error)?\par
\begin{verbatim}
a, b, c = 7, 6, 8
def f(a, b=9, c):
    print(a, b, c, end=" ")

f(2, c=1, b=0)
print(a, b, c)
\end{verbatim}

{\begin {center}\textbf{Завдання 9}\end {center}}

Що буде виведено за виконання (повідомлення інтерпретатора про помилку позначати error)?\par
\begin{verbatim}
a,b,c=8,6,2
def f(a):
    a+=1
    b=5
    c=5
    return a+b+c

a,b,c=0,6,9
print(f(a),a,b,c)
\end{verbatim}

{\begin {center}\textbf{Завдання 10}\end {center}}


Написати функцiю f, що отримує щось та повертає щось. Невелика загадка все-таки має залишатися. Але нічого принципово нового відносно задач, що наявні у pdf-ках не планується.. Стиль запису умови також оцінюється.\par
\end {large}}
\newpage

{\Large{17.10.2024 Бета-версія: Контрольна робота \No 1, варіант  \No \textbf { 63 }\par}}
{
\begin {large}
{\begin {center}\textbf{Завдання 1}\end {center}}

Що буде виведено за виконання? (повідомлення інтерпретатора про помилку позначати error)\par
\begin{verbatim}
print(4//5//8*5)
\end{verbatim}

{\begin {center}\textbf{Завдання 2}\end {center}}

Що буде виведено за виконання? (повідомлення інтерпретатора про помилку позначати error)\par
\begin{verbatim}
print(-2**1.0+True)
\end{verbatim}

{\begin {center}\textbf{Завдання 3}\end {center}}

Що буде виведено за виконання? (повідомлення інтерпретатора про помилку позначати error)\par
\begin{verbatim}
print(3!=7 or not 4.0==True or 5.0>=4)
\end{verbatim}

{\begin {center}\textbf{Завдання 4}\end {center}}

Що буде виведено за виконання? (повідомлення інтерпретатора про помилку позначати error)\par
\begin{verbatim}
print("False" == 4.0)
\end{verbatim}

{\begin {center}\textbf{Завдання 5}\end {center}}

Що буде виведено за виконання? (повідомлення інтерпретатора про помилку позначати error)\par
\begin{verbatim}
print(2 != 4 > True  is  6)
\end{verbatim}

{\begin {center}\textbf{Завдання 6}\end {center}}

Що буде виведено за виконання? (повідомлення інтерпретатора про помилку позначати error)\par
\begin{verbatim}

def f(n):
    print("f", end="");
    return n

print(f(-4) or f(-1))
\end{verbatim}

{\begin {center}\textbf{Завдання 7}\end {center}}

Що буде виведено за виконання (повідомлення інтерпретатора про помилку позначати error)?\par
\begin{verbatim}
def h(a):
    v=50
    if a!=-4: 
        return 8
    elif a==2:
         v=9
    else:
         return 4
    return v

print(h(7))
\end{verbatim}

{\begin {center}\textbf{Завдання 8}\end {center}}

Що буде виведено за виконання (повідомлення інтерпретатора про помилку позначати error)?\par
\begin{verbatim}
a, b, c = 8, 9, 6
def h(a, b, c):
    print(a, b, c, end=" ")

h(3, 5, 4)
print(a, b, c)
\end{verbatim}

{\begin {center}\textbf{Завдання 9}\end {center}}

Що буде виведено за виконання (повідомлення інтерпретатора про помилку позначати error)?\par
\begin{verbatim}
a,b,c=9,4,0
def f(a):
    a=4
    b=5
    c=2
    return a+b+c

a,b,c=8,5,6
print(f(b),a,b,c)
\end{verbatim}

{\begin {center}\textbf{Завдання 10}\end {center}}


Написати функцiю f, що отримує щось та повертає щось. Невелика загадка все-таки має залишатися. Але нічого принципово нового відносно задач, що наявні у pdf-ках не планується.. Стиль запису умови також оцінюється.\par
\end {large}}
\newpage

{\Large{17.10.2024 Бета-версія: Контрольна робота \No 1, варіант  \No \textbf { 64 }\par}}
{
\begin {large}
{\begin {center}\textbf{Завдання 1}\end {center}}

Що буде виведено за виконання? (повідомлення інтерпретатора про помилку позначати error)\par
\begin{verbatim}
print(8/4%7*5)
\end{verbatim}

{\begin {center}\textbf{Завдання 2}\end {center}}

Що буде виведено за виконання? (повідомлення інтерпретатора про помилку позначати error)\par
\begin{verbatim}
print(4**0.5*1)
\end{verbatim}

{\begin {center}\textbf{Завдання 3}\end {center}}

Що буде виведено за виконання? (повідомлення інтерпретатора про помилку позначати error)\par
\begin{verbatim}
print(False<=6 and not 2.0>2 and 8<9.0)
\end{verbatim}

{\begin {center}\textbf{Завдання 4}\end {center}}

Що буде виведено за виконання? (повідомлення інтерпретатора про помилку позначати error)\par
\begin{verbatim}
print(False <= 7.0)
\end{verbatim}

{\begin {center}\textbf{Завдання 5}\end {center}}

Що буде виведено за виконання? (повідомлення інтерпретатора про помилку позначати error)\par
\begin{verbatim}
print(5 != 3.0 >= 9.0 <= 9)
\end{verbatim}

{\begin {center}\textbf{Завдання 6}\end {center}}

Що буде виведено за виконання? (повідомлення інтерпретатора про помилку позначати error)\par
\begin{verbatim}

def f(n):
    print("f", end="");
    return n>2

print(f(-5) and f(8))
\end{verbatim}

{\begin {center}\textbf{Завдання 7}\end {center}}

Що буде виведено за виконання (повідомлення інтерпретатора про помилку позначати error)?\par
\begin{verbatim}
def g(c):
    w=24
    if c>=-1: 
        w=3
    if c>1:
         return 6
    else:
         w=5
    return w

print(g(-8))
\end{verbatim}

{\begin {center}\textbf{Завдання 8}\end {center}}

Що буде виведено за виконання (повідомлення інтерпретатора про помилку позначати error)?\par
\begin{verbatim}
a, b, c = 7, 6, 9
def h(a, b=7, c=8):
    print(a, b, c, end=" ")

h(3, c=0)
print(a, b, c)
\end{verbatim}

{\begin {center}\textbf{Завдання 9}\end {center}}

Що буде виведено за виконання (повідомлення інтерпретатора про помилку позначати error)?\par
\begin{verbatim}
a,b,c=3,4,0
def g(b):
    a-=2
    b=4
    c=3
    return a+b+c

a,b,c=2,8,1
print(g(a),a,b,c)
\end{verbatim}

{\begin {center}\textbf{Завдання 10}\end {center}}


Написати функцiю f, що отримує щось та повертає щось. Невелика загадка все-таки має залишатися. Але нічого принципово нового відносно задач, що наявні у pdf-ках не планується.. Стиль запису умови також оцінюється.\par
\end {large}}
\newpage

{\Large{17.10.2024 Бета-версія: Контрольна робота \No 1, варіант  \No \textbf { 65 }\par}}
{
\begin {large}
{\begin {center}\textbf{Завдання 1}\end {center}}

Що буде виведено за виконання? (повідомлення інтерпретатора про помилку позначати error)\par
\begin{verbatim}
print(7//5%3*8)
\end{verbatim}

{\begin {center}\textbf{Завдання 2}\end {center}}

Що буде виведено за виконання? (повідомлення інтерпретатора про помилку позначати error)\par
\begin{verbatim}
print(4**-3-False)
\end{verbatim}

{\begin {center}\textbf{Завдання 3}\end {center}}

Що буде виведено за виконання? (повідомлення інтерпретатора про помилку позначати error)\par
\begin{verbatim}
print(3>4.0 and not 9!=True and 7.0<4)
\end{verbatim}

{\begin {center}\textbf{Завдання 4}\end {center}}

Що буде виведено за виконання? (повідомлення інтерпретатора про помилку позначати error)\par
\begin{verbatim}
print("1j" > "9.0j")
\end{verbatim}

{\begin {center}\textbf{Завдання 5}\end {center}}

Що буде виведено за виконання? (повідомлення інтерпретатора про помилку позначати error)\par
\begin{verbatim}
print(False < 9 == 8  is  6)
\end{verbatim}

{\begin {center}\textbf{Завдання 6}\end {center}}

Що буде виведено за виконання? (повідомлення інтерпретатора про помилку позначати error)\par
\begin{verbatim}

def f(n):
    print("f", end="");
    return n

print(f(-9) and f(9))
\end{verbatim}

{\begin {center}\textbf{Завдання 7}\end {center}}

Що буде виведено за виконання (повідомлення інтерпретатора про помилку позначати error)?\par
\begin{verbatim}
def f(b):
    y=73
    if b: 
        return 2
    elif b<=3:
         y=0
    else:
         y=1
    return y

print(f(4))
\end{verbatim}

{\begin {center}\textbf{Завдання 8}\end {center}}

Що буде виведено за виконання (повідомлення інтерпретатора про помилку позначати error)?\par
\begin{verbatim}
a, b, c = 6, 8, 9
def g(a, b, c):
    print(a, b, c, end=" ")

g(4, c=5)
print(a, b, c)
\end{verbatim}

{\begin {center}\textbf{Завдання 9}\end {center}}

Що буде виведено за виконання (повідомлення інтерпретатора про помилку позначати error)?\par
\begin{verbatim}
a,b,c=5,7,5
def g(b):
    a=1
    b-=5
    c=1
    return a+b+c

a,b,c=9,7,0
print(g(a),a,b,c)
\end{verbatim}

{\begin {center}\textbf{Завдання 10}\end {center}}


Написати функцiю f, що отримує щось та повертає щось. Невелика загадка все-таки має залишатися. Але нічого принципово нового відносно задач, що наявні у pdf-ках не планується.. Стиль запису умови також оцінюється.\par
\end {large}}
\newpage

{\Large{17.10.2024 Бета-версія: Контрольна робота \No 1, варіант  \No \textbf { 66 }\par}}
{
\begin {large}
{\begin {center}\textbf{Завдання 1}\end {center}}

Що буде виведено за виконання? (повідомлення інтерпретатора про помилку позначати error)\par
\begin{verbatim}
print(3/3*6*3)
\end{verbatim}

{\begin {center}\textbf{Завдання 2}\end {center}}

Що буде виведено за виконання? (повідомлення інтерпретатора про помилку позначати error)\par
\begin{verbatim}
print(3**4+-2)
\end{verbatim}

{\begin {center}\textbf{Завдання 3}\end {center}}

Що буде виведено за виконання? (повідомлення інтерпретатора про помилку позначати error)\par
\begin{verbatim}
print(5<=8 or not 8.0==False or 2>=6.0)
\end{verbatim}

{\begin {center}\textbf{Завдання 4}\end {center}}

Що буде виведено за виконання? (повідомлення інтерпретатора про помилку позначати error)\par
\begin{verbatim}
print(4 != "True")
\end{verbatim}

{\begin {center}\textbf{Завдання 5}\end {center}}

Що буде виведено за виконання? (повідомлення інтерпретатора про помилку позначати error)\par
\begin{verbatim}
print(7.0 > True >= 2 <= 2.0)
\end{verbatim}

{\begin {center}\textbf{Завдання 6}\end {center}}

Що буде виведено за виконання? (повідомлення інтерпретатора про помилку позначати error)\par
\begin{verbatim}

def f(n):
    print("f", end="");
    return n<=-1

print(f(3) or f(0))
\end{verbatim}

{\begin {center}\textbf{Завдання 7}\end {center}}

Що буде виведено за виконання (повідомлення інтерпретатора про помилку позначати error)?\par
\begin{verbatim}
def f(d):
    z=71
    if d==2: 
        return 7
    if d!=0:
         z=9
    else:
         return 8
    return z

print(f(8))
\end{verbatim}

{\begin {center}\textbf{Завдання 8}\end {center}}

Що буде виведено за виконання (повідомлення інтерпретатора про помилку позначати error)?\par
\begin{verbatim}
a, b, c = 7, 8, 9
def f(a, b, c=6):
    print(a, b, c, end=" ")

f(1, 2, 4)
print(a, b, c)
\end{verbatim}

{\begin {center}\textbf{Завдання 9}\end {center}}

Що буде виведено за виконання (повідомлення інтерпретатора про помилку позначати error)?\par
\begin{verbatim}
a,b,c=1,2,7
def g(b):
    a=1
    b=3
    c=3
    return a+b+c

a,b,c=3,0,4
print(g(a),a,b,c)
\end{verbatim}

{\begin {center}\textbf{Завдання 10}\end {center}}


Написати функцiю f, що отримує щось та повертає щось. Невелика загадка все-таки має залишатися. Але нічого принципово нового відносно задач, що наявні у pdf-ках не планується.. Стиль запису умови також оцінюється.\par
\end {large}}
\newpage

{\Large{17.10.2024 Бета-версія: Контрольна робота \No 1, варіант  \No \textbf { 67 }\par}}
{
\begin {large}
{\begin {center}\textbf{Завдання 1}\end {center}}

Що буде виведено за виконання? (повідомлення інтерпретатора про помилку позначати error)\par
\begin{verbatim}
print(9/7/5*7)
\end{verbatim}

{\begin {center}\textbf{Завдання 2}\end {center}}

Що буде виведено за виконання? (повідомлення інтерпретатора про помилку позначати error)\par
\begin{verbatim}
print(9**1.0*True)
\end{verbatim}

{\begin {center}\textbf{Завдання 3}\end {center}}

Що буде виведено за виконання? (повідомлення інтерпретатора про помилку позначати error)\par
\begin{verbatim}
print(6>7 or not 5.0!=False or 3<9.0)
\end{verbatim}

{\begin {center}\textbf{Завдання 4}\end {center}}

Що буде виведено за виконання? (повідомлення інтерпретатора про помилку позначати error)\par
\begin{verbatim}
print(5.0 < "9")
\end{verbatim}

{\begin {center}\textbf{Завдання 5}\end {center}}

Що буде виведено за виконання? (повідомлення інтерпретатора про помилку позначати error)\par
\begin{verbatim}
print(False != 3 > 4 == 5.0)
\end{verbatim}

{\begin {center}\textbf{Завдання 6}\end {center}}

Що буде виведено за виконання? (повідомлення інтерпретатора про помилку позначати error)\par
\begin{verbatim}

def f(n):
    print("f", end="");
    return n>=0

print(f(5) or f(7))
\end{verbatim}

{\begin {center}\textbf{Завдання 7}\end {center}}

Що буде виведено за виконання (повідомлення інтерпретатора про помилку позначати error)?\par
\begin{verbatim}
def g(a):
    u=88
    if a>=-5: 
        u=1
    elif a>-3:
         return 0
    else:
         return 6
    return u

print(g(-7))
\end{verbatim}

{\begin {center}\textbf{Завдання 8}\end {center}}

Що буде виведено за виконання (повідомлення інтерпретатора про помилку позначати error)?\par
\begin{verbatim}
a, b, c = 7, 9, 8
def h(a, b=6, c):
    print(a, b, c, end=" ")

h(5, c=1, b=3)
print(a, b, c)
\end{verbatim}

{\begin {center}\textbf{Завдання 9}\end {center}}

Що буде виведено за виконання (повідомлення інтерпретатора про помилку позначати error)?\par
\begin{verbatim}
a,b,c=3,4,1
def h(a):
    a*=4
    b=3
    c=2
    return a+b+c

a,b,c=2,6,8
print(h(a),a,b,c)
\end{verbatim}

{\begin {center}\textbf{Завдання 10}\end {center}}


Написати функцiю f, що отримує щось та повертає щось. Невелика загадка все-таки має залишатися. Але нічого принципово нового відносно задач, що наявні у pdf-ках не планується.. Стиль запису умови також оцінюється.\par
\end {large}}
\newpage

{\Large{17.10.2024 Бета-версія: Контрольна робота \No 1, варіант  \No \textbf { 68 }\par}}
{
\begin {large}
{\begin {center}\textbf{Завдання 1}\end {center}}

Що буде виведено за виконання? (повідомлення інтерпретатора про помилку позначати error)\par
\begin{verbatim}
print(6//9//9*6)
\end{verbatim}

{\begin {center}\textbf{Завдання 2}\end {center}}

Що буде виведено за виконання? (повідомлення інтерпретатора про помилку позначати error)\par
\begin{verbatim}
print(4**-1.0*1)
\end{verbatim}

{\begin {center}\textbf{Завдання 3}\end {center}}

Що буде виведено за виконання? (повідомлення інтерпретатора про помилку позначати error)\par
\begin{verbatim}
print(8<=True and not 7>=2.0 and 4==3.0)
\end{verbatim}

{\begin {center}\textbf{Завдання 4}\end {center}}

Що буде виведено за виконання? (повідомлення інтерпретатора про помилку позначати error)\par
\begin{verbatim}
print("7.0" >= "False")
\end{verbatim}

{\begin {center}\textbf{Завдання 5}\end {center}}

Що буде виведено за виконання? (повідомлення інтерпретатора про помилку позначати error)\par
\begin{verbatim}
print(5 < 7.0  is  7  is  2.0)
\end{verbatim}

{\begin {center}\textbf{Завдання 6}\end {center}}

Що буде виведено за виконання? (повідомлення інтерпретатора про помилку позначати error)\par
\begin{verbatim}

def f(n):
    print("f", end="");
    return n

print(f(1) and f(-8))
\end{verbatim}

{\begin {center}\textbf{Завдання 7}\end {center}}

Що буде виведено за виконання (повідомлення інтерпретатора про помилку позначати error)?\par
\begin{verbatim}
def h(a,b):
    c=96
    if b>=-3:
        c=8
    if b<2:
         return 7
    else: 
        c=5
    return c

print(h(3,-3))
\end{verbatim}

{\begin {center}\textbf{Завдання 8}\end {center}}

Що буде виведено за виконання (повідомлення інтерпретатора про помилку позначати error)?\par
\begin{verbatim}
a, b, c = 7, 9, 8
def f(a, b=6, c=7):
    print(a, b, c, end=" ")

f(2, 0, a=4)
print(a, b, c)
\end{verbatim}

{\begin {center}\textbf{Завдання 9}\end {center}}

Що буде виведено за виконання (повідомлення інтерпретатора про помилку позначати error)?\par
\begin{verbatim}
a,b,c=1,5,3
def f(b):
    a+=4
    b=1
    c=5
    return a+b+c

a,b,c=8,0,9
print(f(b),a,b,c)
\end{verbatim}

{\begin {center}\textbf{Завдання 10}\end {center}}


Написати функцiю f, що отримує щось та повертає щось. Невелика загадка все-таки має залишатися. Але нічого принципово нового відносно задач, що наявні у pdf-ках не планується.. Стиль запису умови також оцінюється.\par
\end {large}}
\newpage

{\Large{17.10.2024 Бета-версія: Контрольна робота \No 1, варіант  \No \textbf { 69 }\par}}
{
\begin {large}
{\begin {center}\textbf{Завдання 1}\end {center}}

Що буде виведено за виконання? (повідомлення інтерпретатора про помилку позначати error)\par
\begin{verbatim}
print(4/6/4*9)
\end{verbatim}

{\begin {center}\textbf{Завдання 2}\end {center}}

Що буде виведено за виконання? (повідомлення інтерпретатора про помилку позначати error)\par
\begin{verbatim}
print(-3**-2--2)
\end{verbatim}

{\begin {center}\textbf{Завдання 3}\end {center}}

Що буде виведено за виконання? (повідомлення інтерпретатора про помилку позначати error)\par
\begin{verbatim}
print(2>=6 or not 6.0<=4.0 and 5==True)
\end{verbatim}

{\begin {center}\textbf{Завдання 4}\end {center}}

Що буде виведено за виконання? (повідомлення інтерпретатора про помилку позначати error)\par
\begin{verbatim}
print(7 == True)
\end{verbatim}

{\begin {center}\textbf{Завдання 5}\end {center}}

Що буде виведено за виконання? (повідомлення інтерпретатора про помилку позначати error)\par
\begin{verbatim}
print(7 > True <= 4.0 != 6)
\end{verbatim}

{\begin {center}\textbf{Завдання 6}\end {center}}

Що буде виведено за виконання? (повідомлення інтерпретатора про помилку позначати error)\par
\begin{verbatim}

def f(n):
    print("f", end="");
    return n!=1

print(f(-9) and f(9))
\end{verbatim}

{\begin {center}\textbf{Завдання 7}\end {center}}

Що буде виведено за виконання (повідомлення інтерпретатора про помилку позначати error)?\par
\begin{verbatim}
def h(c):
    x=92
    if c: 
        x=5
    if c<1:
         return 7
    else:
         x=4
    return x

print(h(3))
\end{verbatim}

{\begin {center}\textbf{Завдання 8}\end {center}}

Що буде виведено за виконання (повідомлення інтерпретатора про помилку позначати error)?\par
\begin{verbatim}
a, b, c = 9, 6, 8
def g(a, b=7, c):
    print(a, b, c, end=" ")

g(b=0, c=2, 1)
print(a, b, c)
\end{verbatim}

{\begin {center}\textbf{Завдання 9}\end {center}}

Що буде виведено за виконання (повідомлення інтерпретатора про помилку позначати error)?\par
\begin{verbatim}
a,b,c=4,7,2
def h(a):
    a=2
    b*=3
    c=2
    return a+b+c

a,b,c=6,3,4
print(h(b),a,b,c)
\end{verbatim}

{\begin {center}\textbf{Завдання 10}\end {center}}


Написати функцiю f, що отримує щось та повертає щось. Невелика загадка все-таки має залишатися. Але нічого принципово нового відносно задач, що наявні у pdf-ках не планується.. Стиль запису умови також оцінюється.\par
\end {large}}
\newpage

{\Large{17.10.2024 Бета-версія: Контрольна робота \No 1, варіант  \No \textbf { 70 }\par}}
{
\begin {large}
{\begin {center}\textbf{Завдання 1}\end {center}}

Що буде виведено за виконання? (повідомлення інтерпретатора про помилку позначати error)\par
\begin{verbatim}
print(5//8*8*4)
\end{verbatim}

{\begin {center}\textbf{Завдання 2}\end {center}}

Що буде виведено за виконання? (повідомлення інтерпретатора про помилку позначати error)\par
\begin{verbatim}
print(9**0.5*-1)
\end{verbatim}

{\begin {center}\textbf{Завдання 3}\end {center}}

Що буде виведено за виконання? (повідомлення інтерпретатора про помилку позначати error)\par
\begin{verbatim}
print(not 9>5.0 and False<9 or 7.0!=5)
\end{verbatim}

{\begin {center}\textbf{Завдання 4}\end {center}}

Що буде виведено за виконання? (повідомлення інтерпретатора про помилку позначати error)\par
\begin{verbatim}
print("9.0j" != "False")
\end{verbatim}

{\begin {center}\textbf{Завдання 5}\end {center}}

Що буде виведено за виконання? (повідомлення інтерпретатора про помилку позначати error)\par
\begin{verbatim}
print(9 == False < 4 >= 6.0)
\end{verbatim}

{\begin {center}\textbf{Завдання 6}\end {center}}

Що буде виведено за виконання? (повідомлення інтерпретатора про помилку позначати error)\par
\begin{verbatim}

def f(n):
    print("f", end="");
    return n

print(f(2) or f(6))
\end{verbatim}

{\begin {center}\textbf{Завдання 7}\end {center}}

Що буде виведено за виконання (повідомлення інтерпретатора про помилку позначати error)?\par
\begin{verbatim}
def f(a,b):
    c=16
    if b:
        return 3
    elif b==0:
         c=2
    else: 
        return 3
    return c

print(f(8,1))
\end{verbatim}

{\begin {center}\textbf{Завдання 8}\end {center}}

Що буде виведено за виконання (повідомлення інтерпретатора про помилку позначати error)?\par
\begin{verbatim}
a, b, c = 9, 8, 6
def f(a, b, c):
    print(a, b, c, end=" ")

f(b=5, b=3, c=1)
print(a, b, c)
\end{verbatim}

{\begin {center}\textbf{Завдання 9}\end {center}}

Що буде виведено за виконання (повідомлення інтерпретатора про помилку позначати error)?\par
\begin{verbatim}
a,b,c=3,1,9
def h(a):
    a=5
    b=4
    c=1
    return a+b+c

a,b,c=7,5,2
print(h(b),a,b,c)
\end{verbatim}

{\begin {center}\textbf{Завдання 10}\end {center}}


Написати функцiю f, що отримує щось та повертає щось. Невелика загадка все-таки має залишатися. Але нічого принципово нового відносно задач, що наявні у pdf-ках не планується.. Стиль запису умови також оцінюється.\par
\end {large}}
\newpage

{\Large{17.10.2024 Бета-версія: Контрольна робота \No 1, варіант  \No \textbf { 71 }\par}}
{
\begin {large}
{\begin {center}\textbf{Завдання 1}\end {center}}

Що буде виведено за виконання? (повідомлення інтерпретатора про помилку позначати error)\par
\begin{verbatim}
print(6/4//3*7)
\end{verbatim}

{\begin {center}\textbf{Завдання 2}\end {center}}

Що буде виведено за виконання? (повідомлення інтерпретатора про помилку позначати error)\par
\begin{verbatim}
print(-3**False+False)
\end{verbatim}

{\begin {center}\textbf{Завдання 3}\end {center}}

Що буде виведено за виконання? (повідомлення інтерпретатора про помилку позначати error)\par
\begin{verbatim}
print(not 2.0<2 or 9.0<=8 or 3==True)
\end{verbatim}

{\begin {center}\textbf{Завдання 4}\end {center}}

Що буде виведено за виконання? (повідомлення інтерпретатора про помилку позначати error)\par
\begin{verbatim}
print(True <= 5.0)
\end{verbatim}

{\begin {center}\textbf{Завдання 5}\end {center}}

Що буде виведено за виконання? (повідомлення інтерпретатора про помилку позначати error)\par
\begin{verbatim}
print(3 >= 9.0 == 3.0  is  8)
\end{verbatim}

{\begin {center}\textbf{Завдання 6}\end {center}}

Що буде виведено за виконання? (повідомлення інтерпретатора про помилку позначати error)\par
\begin{verbatim}

def f(n):
    print("f", end="");
    return n<-2

print(f(-2) or f(-3))
\end{verbatim}

{\begin {center}\textbf{Завдання 7}\end {center}}

Що буде виведено за виконання (повідомлення інтерпретатора про помилку позначати error)?\par
\begin{verbatim}
def g(a):
    x=12
    if a: 
        x=2
    if a!=-1:
         return 3
    else:
         return 9
    return x

print(g(-3))
\end{verbatim}

{\begin {center}\textbf{Завдання 8}\end {center}}

Що буде виведено за виконання (повідомлення інтерпретатора про помилку позначати error)?\par
\begin{verbatim}
a, b, c = 7, 6, 9
def g(a, b=8, c=7):
    print(a, b, c, end=" ")

g(2, 4)
print(a, b, c)
\end{verbatim}

{\begin {center}\textbf{Завдання 9}\end {center}}

Що буде виведено за виконання (повідомлення інтерпретатора про помилку позначати error)?\par
\begin{verbatim}
a,b,c=8,2,6
def g(a):
    a=3
    b+=4
    c=5
    return a+b+c

a,b,c=9,7,0
print(g(a),a,b,c)
\end{verbatim}

{\begin {center}\textbf{Завдання 10}\end {center}}


Написати функцiю f, що отримує щось та повертає щось. Невелика загадка все-таки має залишатися. Але нічого принципово нового відносно задач, що наявні у pdf-ках не планується.. Стиль запису умови також оцінюється.\par
\end {large}}
\newpage

{\Large{17.10.2024 Бета-версія: Контрольна робота \No 1, варіант  \No \textbf { 72 }\par}}
{
\begin {large}
{\begin {center}\textbf{Завдання 1}\end {center}}

Що буде виведено за виконання? (повідомлення інтерпретатора про помилку позначати error)\par
\begin{verbatim}
print(7//7%4*4)
\end{verbatim}

{\begin {center}\textbf{Завдання 2}\end {center}}

Що буде виведено за виконання? (повідомлення інтерпретатора про помилку позначати error)\par
\begin{verbatim}
print(1**-0.5*2)
\end{verbatim}

{\begin {center}\textbf{Завдання 3}\end {center}}

Що буде виведено за виконання? (повідомлення інтерпретатора про помилку позначати error)\par
\begin{verbatim}
print(not False>=7 and 4!=6 and 8.0>3.0)
\end{verbatim}

{\begin {center}\textbf{Завдання 4}\end {center}}

Що буде виведено за виконання? (повідомлення інтерпретатора про помилку позначати error)\par
\begin{verbatim}
print("2j" > 1)
\end{verbatim}

{\begin {center}\textbf{Завдання 5}\end {center}}

Що буде виведено за виконання? (повідомлення інтерпретатора про помилку позначати error)\par
\begin{verbatim}
print(8.0 < True > 5 <= 2)
\end{verbatim}

{\begin {center}\textbf{Завдання 6}\end {center}}

Що буде виведено за виконання? (повідомлення інтерпретатора про помилку позначати error)\par
\begin{verbatim}

def f(n):
    print("f", end="");
    return n

print(f(8) and f(-6))
\end{verbatim}

{\begin {center}\textbf{Завдання 7}\end {center}}

Що буде виведено за виконання (повідомлення інтерпретатора про помилку позначати error)?\par
\begin{verbatim}
def h(a,b):
    c=75
    if a!=-5:
        return 0
    elif b>5:
         c=8
    else: 
        return 5
    return c

print(h(-4,0))
\end{verbatim}

{\begin {center}\textbf{Завдання 8}\end {center}}

Що буде виведено за виконання (повідомлення інтерпретатора про помилку позначати error)?\par
\begin{verbatim}
a, b, c = 8, 9, 7
def h(a, b, c=6):
    print(a, b, c, end=" ")

h(a=0, 3, a=5)
print(a, b, c)
\end{verbatim}

{\begin {center}\textbf{Завдання 9}\end {center}}

Що буде виведено за виконання (повідомлення інтерпретатора про помилку позначати error)?\par
\begin{verbatim}
a,b,c=5,1,6
def f(b):
    a-=1
    b=1
    c=5
    return a+b+c

a,b,c=0,2,5
print(f(b),a,b,c)
\end{verbatim}

{\begin {center}\textbf{Завдання 10}\end {center}}


Написати функцiю f, що отримує щось та повертає щось. Невелика загадка все-таки має залишатися. Але нічого принципово нового відносно задач, що наявні у pdf-ках не планується.. Стиль запису умови також оцінюється.\par
\end {large}}
\newpage

{\Large{17.10.2024 Бета-версія: Контрольна робота \No 1, варіант  \No \textbf { 73 }\par}}
{
\begin {large}
{\begin {center}\textbf{Завдання 1}\end {center}}

Що буде виведено за виконання? (повідомлення інтерпретатора про помилку позначати error)\par
\begin{verbatim}
print(9//6/7*3)
\end{verbatim}

{\begin {center}\textbf{Завдання 2}\end {center}}

Що буде виведено за виконання? (повідомлення інтерпретатора про помилку позначати error)\par
\begin{verbatim}
print(-1**2.0--1)
\end{verbatim}

{\begin {center}\textbf{Завдання 3}\end {center}}

Що буде виведено за виконання? (повідомлення інтерпретатора про помилку позначати error)\par
\begin{verbatim}
print(not 5!=False and 9.0>7 or 2>=6.0)
\end{verbatim}

{\begin {center}\textbf{Завдання 4}\end {center}}

Що буде виведено за виконання? (повідомлення інтерпретатора про помилку позначати error)\par
\begin{verbatim}
print(True > 5)
\end{verbatim}

{\begin {center}\textbf{Завдання 5}\end {center}}

Що буде виведено за виконання? (повідомлення інтерпретатора про помилку позначати error)\par
\begin{verbatim}
print(6 != 5 <= 2.0 != 5.0)
\end{verbatim}

{\begin {center}\textbf{Завдання 6}\end {center}}

Що буде виведено за виконання? (повідомлення інтерпретатора про помилку позначати error)\par
\begin{verbatim}

def f(n):
    print("f", end="");
    return n<-3

print(f(0) or f(-7))
\end{verbatim}

{\begin {center}\textbf{Завдання 7}\end {center}}

Що буде виведено за виконання (повідомлення інтерпретатора про помилку позначати error)?\par
\begin{verbatim}
def g(a,b):
    c=89
    if a<=2:
        c=9
    if a>=-1:
         return 7
    else: 
        c=1
    return c

print(g(0,5))
\end{verbatim}

{\begin {center}\textbf{Завдання 8}\end {center}}

Що буде виведено за виконання (повідомлення інтерпретатора про помилку позначати error)?\par
\begin{verbatim}
a, b, c = 6, 8, 9
def f(a, b, c):
    print(a, b, c, end=" ")

f(a=5, c=1, b=3)
print(a, b, c)
\end{verbatim}

{\begin {center}\textbf{Завдання 9}\end {center}}

Що буде виведено за виконання (повідомлення інтерпретатора про помилку позначати error)?\par
\begin{verbatim}
a,b,c=8,9,7
def h(a):
    a-=4
    b=3
    c=2
    return a+b+c

a,b,c=4,1,3
print(h(b),a,b,c)
\end{verbatim}

{\begin {center}\textbf{Завдання 10}\end {center}}


Написати функцiю f, що отримує щось та повертає щось. Невелика загадка все-таки має залишатися. Але нічого принципово нового відносно задач, що наявні у pdf-ках не планується.. Стиль запису умови також оцінюється.\par
\end {large}}
\newpage

{\Large{17.10.2024 Бета-версія: Контрольна робота \No 1, варіант  \No \textbf { 74 }\par}}
{
\begin {large}
{\begin {center}\textbf{Завдання 1}\end {center}}

Що буде виведено за виконання? (повідомлення інтерпретатора про помилку позначати error)\par
\begin{verbatim}
print(3/8*8*5)
\end{verbatim}

{\begin {center}\textbf{Завдання 2}\end {center}}

Що буде виведено за виконання? (повідомлення інтерпретатора про помилку позначати error)\par
\begin{verbatim}
print(1**1.0*True)
\end{verbatim}

{\begin {center}\textbf{Завдання 3}\end {center}}

Що буде виведено за виконання? (повідомлення інтерпретатора про помилку позначати error)\par
\begin{verbatim}
print(4.0<=9 or not True==7.0 and 4<8)
\end{verbatim}

{\begin {center}\textbf{Завдання 4}\end {center}}

Що буде виведено за виконання? (повідомлення інтерпретатора про помилку позначати error)\par
\begin{verbatim}
print("8j" == "8.0")
\end{verbatim}

{\begin {center}\textbf{Завдання 5}\end {center}}

Що буде виведено за виконання? (повідомлення інтерпретатора про помилку позначати error)\par
\begin{verbatim}
print(False < 8 >= True  is  9)
\end{verbatim}

{\begin {center}\textbf{Завдання 6}\end {center}}

Що буде виведено за виконання? (повідомлення інтерпретатора про помилку позначати error)\par
\begin{verbatim}

def f(n):
    print("f", end="");
    return n

print(f(4) and f(-1))
\end{verbatim}

{\begin {center}\textbf{Завдання 7}\end {center}}

Що буде виведено за виконання (повідомлення інтерпретатора про помилку позначати error)?\par
\begin{verbatim}
def h(a,b):
    c=41
    if b:
        return 4
    if b>=4:
         c=6
    else: 
        c=6
    return c

print(h(6,4))
\end{verbatim}

{\begin {center}\textbf{Завдання 8}\end {center}}

Що буде виведено за виконання (повідомлення інтерпретатора про помилку позначати error)?\par
\begin{verbatim}
a, b, c = 7, 6, 8
def h(a, b=7, c):
    print(a, b, c, end=" ")

h(0, 4)
print(a, b, c)
\end{verbatim}

{\begin {center}\textbf{Завдання 9}\end {center}}

Що буде виведено за виконання (повідомлення інтерпретатора про помилку позначати error)?\par
\begin{verbatim}
a,b,c=2,5,0
def g(b):
    a+=5
    b=3
    c=4
    return a+b+c

a,b,c=8,1,9
print(g(b),a,b,c)
\end{verbatim}

{\begin {center}\textbf{Завдання 10}\end {center}}


Написати функцiю f, що отримує щось та повертає щось. Невелика загадка все-таки має залишатися. Але нічого принципово нового відносно задач, що наявні у pdf-ках не планується.. Стиль запису умови також оцінюється.\par
\end {large}}
\newpage

{\Large{17.10.2024 Бета-версія: Контрольна робота \No 1, варіант  \No \textbf { 75 }\par}}
{
\begin {large}
{\begin {center}\textbf{Завдання 1}\end {center}}

Що буде виведено за виконання? (повідомлення інтерпретатора про помилку позначати error)\par
\begin{verbatim}
print(4/5%9*8)
\end{verbatim}

{\begin {center}\textbf{Завдання 2}\end {center}}

Що буде виведено за виконання? (повідомлення інтерпретатора про помилку позначати error)\par
\begin{verbatim}
print(3**-3+False)
\end{verbatim}

{\begin {center}\textbf{Завдання 3}\end {center}}

Що буде виведено за виконання? (повідомлення інтерпретатора про помилку позначати error)\par
\begin{verbatim}
print(3>=6 and not 5.0<2.0 or 8!=False)
\end{verbatim}

{\begin {center}\textbf{Завдання 4}\end {center}}

Що буде виведено за виконання? (повідомлення інтерпретатора про помилку позначати error)\par
\begin{verbatim}
print(1.0 != "False")
\end{verbatim}

{\begin {center}\textbf{Завдання 5}\end {center}}

Що буде виведено за виконання? (повідомлення інтерпретатора про помилку позначати error)\par
\begin{verbatim}
print(4.0 > 7.0 == 2 == 3)
\end{verbatim}

{\begin {center}\textbf{Завдання 6}\end {center}}

Що буде виведено за виконання? (повідомлення інтерпретатора про помилку позначати error)\par
\begin{verbatim}

def f(n):
    print("f", end="");
    return n==4

print(f(5) and f(-4))
\end{verbatim}

{\begin {center}\textbf{Завдання 7}\end {center}}

Що буде виведено за виконання (повідомлення інтерпретатора про помилку позначати error)?\par
\begin{verbatim}
def g(a,b):
    c=58
    if b>1:
        return 3
    elif a<=-2:
         c=8
    else: 
        return 7
    return c

print(g(-8,-2))
\end{verbatim}

{\begin {center}\textbf{Завдання 8}\end {center}}

Що буде виведено за виконання (повідомлення інтерпретатора про помилку позначати error)?\par
\begin{verbatim}
a, b, c = 9, 6, 7
def g(a, b, c=8):
    print(a, b, c, end=" ")

g(2, a=5)
print(a, b, c)
\end{verbatim}

{\begin {center}\textbf{Завдання 9}\end {center}}

Що буде виведено за виконання (повідомлення інтерпретатора про помилку позначати error)?\par
\begin{verbatim}
a,b,c=8,6,5
def h(a):
    a=2
    b=5
    c=1
    return a+b+c

a,b,c=1,7,0
print(h(b),a,b,c)
\end{verbatim}

{\begin {center}\textbf{Завдання 10}\end {center}}


Написати функцiю f, що отримує щось та повертає щось. Невелика загадка все-таки має залишатися. Але нічого принципово нового відносно задач, що наявні у pdf-ках не планується.. Стиль запису умови також оцінюється.\par
\end {large}}
\newpage

{\Large{17.10.2024 Бета-версія: Контрольна робота \No 1, варіант  \No \textbf { 76 }\par}}
{
\begin {large}
{\begin {center}\textbf{Завдання 1}\end {center}}

Що буде виведено за виконання? (повідомлення інтерпретатора про помилку позначати error)\par
\begin{verbatim}
print(5//3//6*9)
\end{verbatim}

{\begin {center}\textbf{Завдання 2}\end {center}}

Що буде виведено за виконання? (повідомлення інтерпретатора про помилку позначати error)\par
\begin{verbatim}
print(-2**-4+1)
\end{verbatim}

{\begin {center}\textbf{Завдання 3}\end {center}}

Що буде виведено за виконання? (повідомлення інтерпретатора про помилку позначати error)\par
\begin{verbatim}
print(not 2>True or 8.0<=3.0 and 3==4)
\end{verbatim}

{\begin {center}\textbf{Завдання 4}\end {center}}

Що буде виведено за виконання? (повідомлення інтерпретатора про помилку позначати error)\par
\begin{verbatim}
print(6 >= "True")
\end{verbatim}

{\begin {center}\textbf{Завдання 5}\end {center}}

Що буде виведено за виконання? (повідомлення інтерпретатора про помилку позначати error)\par
\begin{verbatim}
print(True >= 9.0 > 6.0 < 7)
\end{verbatim}

{\begin {center}\textbf{Завдання 6}\end {center}}

Що буде виведено за виконання? (повідомлення інтерпретатора про помилку позначати error)\par
\begin{verbatim}

def f(n):
    print("f", end="");
    return n

print(f(-5) or f(3))
\end{verbatim}

{\begin {center}\textbf{Завдання 7}\end {center}}

Що буде виведено за виконання (повідомлення інтерпретатора про помилку позначати error)?\par
\begin{verbatim}
def h(c):
    y=16
    if c>=4: 
        y=4
    elif c==3:
         return 7
    else:
         y=1
    return y

print(h(2))
\end{verbatim}

{\begin {center}\textbf{Завдання 8}\end {center}}

Що буде виведено за виконання (повідомлення інтерпретатора про помилку позначати error)?\par
\begin{verbatim}
a, b, c = 9, 6, 9
def g(a, b=7, c=8):
    print(a, b, c, end=" ")

g(0, 2, 1)
print(a, b, c)
\end{verbatim}

{\begin {center}\textbf{Завдання 9}\end {center}}

Що буде виведено за виконання (повідомлення інтерпретатора про помилку позначати error)?\par
\begin{verbatim}
a,b,c=6,7,3
def f(b):
    a=1
    b*=2
    c=2
    return a+b+c

a,b,c=4,0,9
print(f(a),a,b,c)
\end{verbatim}

{\begin {center}\textbf{Завдання 10}\end {center}}


Написати функцiю f, що отримує щось та повертає щось. Невелика загадка все-таки має залишатися. Але нічого принципово нового відносно задач, що наявні у pdf-ках не планується.. Стиль запису умови також оцінюється.\par
\end {large}}
\newpage

{\Large{17.10.2024 Бета-версія: Контрольна робота \No 1, варіант  \No \textbf { 77 }\par}}
{
\begin {large}
{\begin {center}\textbf{Завдання 1}\end {center}}

Що буде виведено за виконання? (повідомлення інтерпретатора про помилку позначати error)\par
\begin{verbatim}
print(8/9*5*6)
\end{verbatim}

{\begin {center}\textbf{Завдання 2}\end {center}}

Що буде виведено за виконання? (повідомлення інтерпретатора про помилку позначати error)\par
\begin{verbatim}
print(4**0.5*2)
\end{verbatim}

{\begin {center}\textbf{Завдання 3}\end {center}}

Що буде виведено за виконання? (повідомлення інтерпретатора про помилку позначати error)\par
\begin{verbatim}
print(5<5.0 and not 6>=7 or 4.0<=False)
\end{verbatim}

{\begin {center}\textbf{Завдання 4}\end {center}}

Що буде виведено за виконання? (повідомлення інтерпретатора про помилку позначати error)\par
\begin{verbatim}
print(False <= "4.0j")
\end{verbatim}

{\begin {center}\textbf{Завдання 5}\end {center}}

Що буде виведено за виконання? (повідомлення інтерпретатора про помилку позначати error)\par
\begin{verbatim}
print(4 <= 6  is  3 != 3.0)
\end{verbatim}

{\begin {center}\textbf{Завдання 6}\end {center}}

Що буде виведено за виконання? (повідомлення інтерпретатора про помилку позначати error)\par
\begin{verbatim}

def f(n):
    print("f", end="");
    return n

print(f(9) and f(0))
\end{verbatim}

{\begin {center}\textbf{Завдання 7}\end {center}}

Що буде виведено за виконання (повідомлення інтерпретатора про помилку позначати error)?\par
\begin{verbatim}
def f(a,b):
    c=94
    if a!=-3:
        return 0
    elif a==-4:
         c=5
    else: 
        return 1
    return c

print(f(3,9))
\end{verbatim}

{\begin {center}\textbf{Завдання 8}\end {center}}

Що буде виведено за виконання (повідомлення інтерпретатора про помилку позначати error)?\par
\begin{verbatim}
a, b, c = 6, 9, 7
def f(a, b, c):
    print(a, b, c, end=" ")

f(3, c=4, b=5)
print(a, b, c)
\end{verbatim}

{\begin {center}\textbf{Завдання 9}\end {center}}

Що буде виведено за виконання (повідомлення інтерпретатора про помилку позначати error)?\par
\begin{verbatim}
a,b,c=9,3,2
def f(b):
    a=4
    b=3
    c=2
    return a+b+c

a,b,c=4,6,8
print(f(a),a,b,c)
\end{verbatim}

{\begin {center}\textbf{Завдання 10}\end {center}}


Написати функцiю f, що отримує щось та повертає щось. Невелика загадка все-таки має залишатися. Але нічого принципово нового відносно задач, що наявні у pdf-ках не планується.. Стиль запису умови також оцінюється.\par
\end {large}}
\newpage

{\Large{17.10.2024 Бета-версія: Контрольна робота \No 1, варіант  \No \textbf { 78 }\par}}
{
\begin {large}
{\begin {center}\textbf{Завдання 1}\end {center}}

Що буде виведено за виконання? (повідомлення інтерпретатора про помилку позначати error)\par
\begin{verbatim}
print(4//8%5*9)
\end{verbatim}

{\begin {center}\textbf{Завдання 2}\end {center}}

Що буде виведено за виконання? (повідомлення інтерпретатора про помилку позначати error)\par
\begin{verbatim}
print(9**-0.5*-2)
\end{verbatim}

{\begin {center}\textbf{Завдання 3}\end {center}}

Що буде виведено за виконання? (повідомлення інтерпретатора про помилку позначати error)\par
\begin{verbatim}
print(2.0>9 or not 9==5 and 7.0!=True)
\end{verbatim}

{\begin {center}\textbf{Завдання 4}\end {center}}

Що буде виведено за виконання? (повідомлення інтерпретатора про помилку позначати error)\par
\begin{verbatim}
print("4" < 3.0)
\end{verbatim}

{\begin {center}\textbf{Завдання 5}\end {center}}

Що буде виведено за виконання? (повідомлення інтерпретатора про помилку позначати error)\par
\begin{verbatim}
print(False >= 2 < 5 != 8.0)
\end{verbatim}

{\begin {center}\textbf{Завдання 6}\end {center}}

Що буде виведено за виконання? (повідомлення інтерпретатора про помилку позначати error)\par
\begin{verbatim}

def f(n):
    print("f", end="");
    return n>1

print(f(-1) or f(6))
\end{verbatim}

{\begin {center}\textbf{Завдання 7}\end {center}}

Що буде виведено за виконання (повідомлення інтерпретатора про помилку позначати error)?\par
\begin{verbatim}
def f(d):
    u=33
    if d<=-2: 
        return 5
    elif d>5:
         u=3
    else:
         return 8
    return u

print(f(-8))
\end{verbatim}

{\begin {center}\textbf{Завдання 8}\end {center}}

Що буде виведено за виконання (повідомлення інтерпретатора про помилку позначати error)?\par
\begin{verbatim}
a, b, c = 8, 6, 7
def h(a, b=9, c):
    print(a, b, c, end=" ")

h(a=3, 1, c=2)
print(a, b, c)
\end{verbatim}

{\begin {center}\textbf{Завдання 9}\end {center}}

Що буде виведено за виконання (повідомлення інтерпретатора про помилку позначати error)?\par
\begin{verbatim}
a,b,c=4,5,6
def f(a):
    a+=1
    b=5
    c=3
    return a+b+c

a,b,c=3,2,8
print(f(b),a,b,c)
\end{verbatim}

{\begin {center}\textbf{Завдання 10}\end {center}}


Написати функцiю f, що отримує щось та повертає щось. Невелика загадка все-таки має залишатися. Але нічого принципово нового відносно задач, що наявні у pdf-ках не планується.. Стиль запису умови також оцінюється.\par
\end {large}}
\newpage

{\Large{17.10.2024 Бета-версія: Контрольна робота \No 1, варіант  \No \textbf { 79 }\par}}
{
\begin {large}
{\begin {center}\textbf{Завдання 1}\end {center}}

Що буде виведено за виконання? (повідомлення інтерпретатора про помилку позначати error)\par
\begin{verbatim}
print(7/9/8*7)
\end{verbatim}

{\begin {center}\textbf{Завдання 2}\end {center}}

Що буде виведено за виконання? (повідомлення інтерпретатора про помилку позначати error)\par
\begin{verbatim}
print(3**1.0-2)
\end{verbatim}

{\begin {center}\textbf{Завдання 3}\end {center}}

Що буде виведено за виконання? (повідомлення інтерпретатора про помилку позначати error)\par
\begin{verbatim}
print(8.0!=3 or not 6<=3.0 and 4==True)
\end{verbatim}

{\begin {center}\textbf{Завдання 4}\end {center}}

Що буде виведено за виконання? (повідомлення інтерпретатора про помилку позначати error)\par
\begin{verbatim}
print(False >= "2.0j")
\end{verbatim}

{\begin {center}\textbf{Завдання 5}\end {center}}

Що буде виведено за виконання? (повідомлення інтерпретатора про помилку позначати error)\par
\begin{verbatim}
print(True <= 4 > 9.0  is  2.0)
\end{verbatim}

{\begin {center}\textbf{Завдання 6}\end {center}}

Що буде виведено за виконання? (повідомлення інтерпретатора про помилку позначати error)\par
\begin{verbatim}

def f(n):
    print("f", end="");
    return n

print(f(-4) or f(8))
\end{verbatim}

{\begin {center}\textbf{Завдання 7}\end {center}}

Що буде виведено за виконання (повідомлення інтерпретатора про помилку позначати error)?\par
\begin{verbatim}
def f(b):
    v=32
    if b: 
        v=2
    if b<-4:
         return 6
    else:
         v=0
    return v

print(f(-2))
\end{verbatim}

{\begin {center}\textbf{Завдання 8}\end {center}}

Що буде виведено за виконання (повідомлення інтерпретатора про помилку позначати error)?\par
\begin{verbatim}
a, b, c = 8, 7, 9
def f(a, b=6, c=8):
    print(a, b, c, end=" ")

f(4, 0, b=0)
print(a, b, c)
\end{verbatim}

{\begin {center}\textbf{Завдання 9}\end {center}}

Що буде виведено за виконання (повідомлення інтерпретатора про помилку позначати error)?\par
\begin{verbatim}
a,b,c=1,7,3
def h(b):
    a=4
    b-=4
    c=1
    return a+b+c

a,b,c=7,5,9
print(h(a),a,b,c)
\end{verbatim}

{\begin {center}\textbf{Завдання 10}\end {center}}


Написати функцiю f, що отримує щось та повертає щось. Невелика загадка все-таки має залишатися. Але нічого принципово нового відносно задач, що наявні у pdf-ках не планується.. Стиль запису умови також оцінюється.\par
\end {large}}
\newpage

{\Large{17.10.2024 Бета-версія: Контрольна робота \No 1, варіант  \No \textbf { 80 }\par}}
{
\begin {large}
{\begin {center}\textbf{Завдання 1}\end {center}}

Що буде виведено за виконання? (повідомлення інтерпретатора про помилку позначати error)\par
\begin{verbatim}
print(6//5//4*4)
\end{verbatim}

{\begin {center}\textbf{Завдання 2}\end {center}}

Що буде виведено за виконання? (повідомлення інтерпретатора про помилку позначати error)\par
\begin{verbatim}
print(4**-1.0*-1)
\end{verbatim}

{\begin {center}\textbf{Завдання 3}\end {center}}

Що буде виведено за виконання? (повідомлення інтерпретатора про помилку позначати error)\par
\begin{verbatim}
print(not 6.0<9.0 and 8>=7 or 2>False)
\end{verbatim}

{\begin {center}\textbf{Завдання 4}\end {center}}

Що буде виведено за виконання? (повідомлення інтерпретатора про помилку позначати error)\par
\begin{verbatim}
print(9 != 6.0)
\end{verbatim}

{\begin {center}\textbf{Завдання 5}\end {center}}

Що буде виведено за виконання? (повідомлення інтерпретатора про помилку позначати error)\par
\begin{verbatim}
print(7 == 9 != 8 >= 8)
\end{verbatim}

{\begin {center}\textbf{Завдання 6}\end {center}}

Що буде виведено за виконання? (повідомлення інтерпретатора про помилку позначати error)\par
\begin{verbatim}

def f(n):
    print("f", end="");
    return n>=3

print(f(-9) and f(-2))
\end{verbatim}

{\begin {center}\textbf{Завдання 7}\end {center}}

Що буде виведено за виконання (повідомлення інтерпретатора про помилку позначати error)?\par
\begin{verbatim}
def g(a,b):
    c=54
    if a:
        c=2
    if b<3:
         return 4
    else: 
        c=9
    return c

print(g(-6,-2))
\end{verbatim}

{\begin {center}\textbf{Завдання 8}\end {center}}

Що буде виведено за виконання (повідомлення інтерпретатора про помилку позначати error)?\par
\begin{verbatim}
a, b, c = 7, 9, 6
def g(a, b, c=8):
    print(a, b, c, end=" ")

g(b=5, c=3, 1)
print(a, b, c)
\end{verbatim}

{\begin {center}\textbf{Завдання 9}\end {center}}

Що буде виведено за виконання (повідомлення інтерпретатора про помилку позначати error)?\par
\begin{verbatim}
a,b,c=6,2,4
def g(a):
    a=3
    b*=2
    c=5
    return a+b+c

a,b,c=1,0,8
print(g(a),a,b,c)
\end{verbatim}

{\begin {center}\textbf{Завдання 10}\end {center}}


Написати функцiю f, що отримує щось та повертає щось. Невелика загадка все-таки має залишатися. Але нічого принципово нового відносно задач, що наявні у pdf-ках не планується.. Стиль запису умови також оцінюється.\par
\end {large}}
\newpage

{\Large{17.10.2024 Бета-версія: Контрольна робота \No 1, варіант  \No \textbf { 81 }\par}}
{
\begin {large}
{\begin {center}\textbf{Завдання 1}\end {center}}

Що буде виведено за виконання? (повідомлення інтерпретатора про помилку позначати error)\par
\begin{verbatim}
print(9/4*3*8)
\end{verbatim}

{\begin {center}\textbf{Завдання 2}\end {center}}

Що буде виведено за виконання? (повідомлення інтерпретатора про помилку позначати error)\par
\begin{verbatim}
print(1**-4+True)
\end{verbatim}

{\begin {center}\textbf{Завдання 3}\end {center}}

Що буде виведено за виконання? (повідомлення інтерпретатора про помилку позначати error)\par
\begin{verbatim}
print(not 7>6.0 and True>=6 or 8.0==2)
\end{verbatim}

{\begin {center}\textbf{Завдання 4}\end {center}}

Що буде виведено за виконання? (повідомлення інтерпретатора про помилку позначати error)\par
\begin{verbatim}
print("3j" <= "True")
\end{verbatim}

{\begin {center}\textbf{Завдання 5}\end {center}}

Що буде виведено за виконання? (повідомлення інтерпретатора про помилку позначати error)\par
\begin{verbatim}
print(3 == 4.0 < 3.0 <= False)
\end{verbatim}

{\begin {center}\textbf{Завдання 6}\end {center}}

Що буде виведено за виконання? (повідомлення інтерпретатора про помилку позначати error)\par
\begin{verbatim}

def f(n):
    print("f", end="");
    return n!=-2

print(f(2) and f(-8))
\end{verbatim}

{\begin {center}\textbf{Завдання 7}\end {center}}

Що буде виведено за виконання (повідомлення інтерпретатора про помилку позначати error)?\par
\begin{verbatim}
def h(a):
    w=75
    if a>=2: 
        return 5
    if a!=-4:
         w=0
    else:
         return 7
    return w

print(h(-4))
\end{verbatim}

{\begin {center}\textbf{Завдання 8}\end {center}}

Що буде виведено за виконання (повідомлення інтерпретатора про помилку позначати error)?\par
\begin{verbatim}
a, b, c = 7, 6, 8
def h(a, b=9, c):
    print(a, b, c, end=" ")

h(4, 2, 3)
print(a, b, c)
\end{verbatim}

{\begin {center}\textbf{Завдання 9}\end {center}}

Що буде виведено за виконання (повідомлення інтерпретатора про помилку позначати error)?\par
\begin{verbatim}
a,b,c=7,8,1
def g(a):
    a=3
    b-=1
    c=5
    return a+b+c

a,b,c=0,4,5
print(g(b),a,b,c)
\end{verbatim}

{\begin {center}\textbf{Завдання 10}\end {center}}


Написати функцiю f, що отримує щось та повертає щось. Невелика загадка все-таки має залишатися. Але нічого принципово нового відносно задач, що наявні у pdf-ках не планується.. Стиль запису умови також оцінюється.\par
\end {large}}
\newpage

{\Large{17.10.2024 Бета-версія: Контрольна робота \No 1, варіант  \No \textbf { 82 }\par}}
{
\begin {large}
{\begin {center}\textbf{Завдання 1}\end {center}}

Що буде виведено за виконання? (повідомлення інтерпретатора про помилку позначати error)\par
\begin{verbatim}
print(3//7%6*6)
\end{verbatim}

{\begin {center}\textbf{Завдання 2}\end {center}}

Що буде виведено за виконання? (повідомлення інтерпретатора про помилку позначати error)\par
\begin{verbatim}
print(9**-0.5*-2)
\end{verbatim}

{\begin {center}\textbf{Завдання 3}\end {center}}

Що буде виведено за виконання? (повідомлення інтерпретатора про помилку позначати error)\par
\begin{verbatim}
print(not 2.0<False or 5.0!=9 and 3<=8)
\end{verbatim}

{\begin {center}\textbf{Завдання 4}\end {center}}

Що буде виведено за виконання? (повідомлення інтерпретатора про помилку позначати error)\par
\begin{verbatim}
print("1" == "False")
\end{verbatim}

{\begin {center}\textbf{Завдання 5}\end {center}}

Що буде виведено за виконання? (повідомлення інтерпретатора про помилку позначати error)\par
\begin{verbatim}
print(8.0 > 6.0  is  7 >= 6)
\end{verbatim}

{\begin {center}\textbf{Завдання 6}\end {center}}

Що буде виведено за виконання? (повідомлення інтерпретатора про помилку позначати error)\par
\begin{verbatim}

def f(n):
    print("f", end="");
    return n

print(f(-7) or f(3))
\end{verbatim}

{\begin {center}\textbf{Завдання 7}\end {center}}

Що буде виведено за виконання (повідомлення інтерпретатора про помилку позначати error)?\par
\begin{verbatim}
def h(a,b):
    c=90
    if a<3:
        return 8
    if b>5:
         c=1
    else: 
        return 2
    return c

print(h(-4,-3))
\end{verbatim}

{\begin {center}\textbf{Завдання 8}\end {center}}

Що буде виведено за виконання (повідомлення інтерпретатора про помилку позначати error)?\par
\begin{verbatim}
a, b, c = 6, 8, 7
def h(a, b=9, c=6):
    print(a, b, c, end=" ")

h(1, c=2, b=4)
print(a, b, c)
\end{verbatim}

{\begin {center}\textbf{Завдання 9}\end {center}}

Що буде виведено за виконання (повідомлення інтерпретатора про помилку позначати error)?\par
\begin{verbatim}
a,b,c=2,7,8
def g(b):
    a=3
    b=4
    c=5
    return a+b+c

a,b,c=9,6,0
print(g(a),a,b,c)
\end{verbatim}

{\begin {center}\textbf{Завдання 10}\end {center}}


Написати функцiю f, що отримує щось та повертає щось. Невелика загадка все-таки має залишатися. Але нічого принципово нового відносно задач, що наявні у pdf-ках не планується.. Стиль запису умови також оцінюється.\par
\end {large}}
\newpage

{\Large{17.10.2024 Бета-версія: Контрольна робота \No 1, варіант  \No \textbf { 83 }\par}}
{
\begin {large}
{\begin {center}\textbf{Завдання 1}\end {center}}

Що буде виведено за виконання? (повідомлення інтерпретатора про помилку позначати error)\par
\begin{verbatim}
print(8/6//9*3)
\end{verbatim}

{\begin {center}\textbf{Завдання 2}\end {center}}

Що буде виведено за виконання? (повідомлення інтерпретатора про помилку позначати error)\par
\begin{verbatim}
print(1**0.5*False)
\end{verbatim}

{\begin {center}\textbf{Завдання 3}\end {center}}

Що буде виведено за виконання? (повідомлення інтерпретатора про помилку позначати error)\par
\begin{verbatim}
print(5>4 or not 4!=9.0 or False<=7.0)
\end{verbatim}

{\begin {center}\textbf{Завдання 4}\end {center}}

Що буде виведено за виконання? (повідомлення інтерпретатора про помилку позначати error)\par
\begin{verbatim}
print(8 < True)
\end{verbatim}

{\begin {center}\textbf{Завдання 5}\end {center}}

Що буде виведено за виконання? (повідомлення інтерпретатора про помилку позначати error)\par
\begin{verbatim}
print(9  is  False == 5 < 2)
\end{verbatim}

{\begin {center}\textbf{Завдання 6}\end {center}}

Що буде виведено за виконання? (повідомлення інтерпретатора про помилку позначати error)\par
\begin{verbatim}

def f(n):
    print("f", end="");
    return n<=-4

print(f(1) and f(-5))
\end{verbatim}

{\begin {center}\textbf{Завдання 7}\end {center}}

Що буде виведено за виконання (повідомлення інтерпретатора про помилку позначати error)?\par
\begin{verbatim}
def f(a,b):
    c=91
    if b==4:
        c=5
    elif a<=2:
         return 6
    else: 
        c=4
    return c

print(f(9,-6))
\end{verbatim}

{\begin {center}\textbf{Завдання 8}\end {center}}

Що буде виведено за виконання (повідомлення інтерпретатора про помилку позначати error)?\par
\begin{verbatim}
a, b, c = 9, 8, 7
def f(a, b, c=7):
    print(a, b, c, end=" ")

f(5, 0)
print(a, b, c)
\end{verbatim}

{\begin {center}\textbf{Завдання 9}\end {center}}

Що буде виведено за виконання (повідомлення інтерпретатора про помилку позначати error)?\par
\begin{verbatim}
a,b,c=2,3,6
def f(b):
    a=2
    b*=4
    c=3
    return a+b+c

a,b,c=9,7,5
print(f(a),a,b,c)
\end{verbatim}

{\begin {center}\textbf{Завдання 10}\end {center}}


Написати функцiю f, що отримує щось та повертає щось. Невелика загадка все-таки має залишатися. Але нічого принципово нового відносно задач, що наявні у pdf-ках не планується.. Стиль запису умови також оцінюється.\par
\end {large}}
\newpage

{\Large{17.10.2024 Бета-версія: Контрольна робота \No 1, варіант  \No \textbf { 84 }\par}}
{
\begin {large}
{\begin {center}\textbf{Завдання 1}\end {center}}

Що буде виведено за виконання? (повідомлення інтерпретатора про помилку позначати error)\par
\begin{verbatim}
print(5//3/7*5)
\end{verbatim}

{\begin {center}\textbf{Завдання 2}\end {center}}

Що буде виведено за виконання? (повідомлення інтерпретатора про помилку позначати error)\par
\begin{verbatim}
print(9**-1.0*1)
\end{verbatim}

{\begin {center}\textbf{Завдання 3}\end {center}}

Що буде виведено за виконання? (повідомлення інтерпретатора про помилку позначати error)\par
\begin{verbatim}
print(7==3 and not True>=4.0 and 3.0<5)
\end{verbatim}

{\begin {center}\textbf{Завдання 4}\end {center}}

Що буде виведено за виконання? (повідомлення інтерпретатора про помилку позначати error)\par
\begin{verbatim}
print("1.0" > 9.0)
\end{verbatim}

{\begin {center}\textbf{Завдання 5}\end {center}}

Що буде виведено за виконання? (повідомлення інтерпретатора про помилку позначати error)\par
\begin{verbatim}
print(True != 4 <= 5.0 > 7.0)
\end{verbatim}

{\begin {center}\textbf{Завдання 6}\end {center}}

Що буде виведено за виконання? (повідомлення інтерпретатора про помилку позначати error)\par
\begin{verbatim}

def f(n):
    print("f", end="");
    return n

print(f(-6) or f(-3))
\end{verbatim}

{\begin {center}\textbf{Завдання 7}\end {center}}

Що буде виведено за виконання (повідомлення інтерпретатора про помилку позначати error)?\par
\begin{verbatim}
def f(a,b):
    c=18
    if b!=-2:
        return 0
    if a>=-4:
         c=3
    else: 
        c=7
    return c

print(f(-1,-9))
\end{verbatim}

{\begin {center}\textbf{Завдання 8}\end {center}}

Що буде виведено за виконання (повідомлення інтерпретатора про помилку позначати error)?\par
\begin{verbatim}
a, b, c = 6, 9, 8
def g(a, b, c):
    print(a, b, c, end=" ")

g(a=2, 4, c=5)
print(a, b, c)
\end{verbatim}

{\begin {center}\textbf{Завдання 9}\end {center}}

Що буде виведено за виконання (повідомлення інтерпретатора про помилку позначати error)?\par
\begin{verbatim}
a,b,c=2,0,3
def h(a):
    a=5
    b+=2
    c=4
    return a+b+c

a,b,c=1,9,6
print(h(b),a,b,c)
\end{verbatim}

{\begin {center}\textbf{Завдання 10}\end {center}}


Написати функцiю f, що отримує щось та повертає щось. Невелика загадка все-таки має залишатися. Але нічого принципово нового відносно задач, що наявні у pdf-ках не планується.. Стиль запису умови також оцінюється.\par
\end {large}}
\newpage

{\Large{17.10.2024 Бета-версія: Контрольна робота \No 1, варіант  \No \textbf { 85 }\par}}
{
\begin {large}
{\begin {center}\textbf{Завдання 1}\end {center}}

Що буде виведено за виконання? (повідомлення інтерпретатора про помилку позначати error)\par
\begin{verbatim}
print(8//9*8*9)
\end{verbatim}

{\begin {center}\textbf{Завдання 2}\end {center}}

Що буде виведено за виконання? (повідомлення інтерпретатора про помилку позначати error)\par
\begin{verbatim}
print(-2**4-True)
\end{verbatim}

{\begin {center}\textbf{Завдання 3}\end {center}}

Що буде виведено за виконання? (повідомлення інтерпретатора про помилку позначати error)\par
\begin{verbatim}
print(False<8 or not 9!=2 or 3.0<=7.0)
\end{verbatim}

{\begin {center}\textbf{Завдання 4}\end {center}}

Що буде виведено за виконання? (повідомлення інтерпретатора про помилку позначати error)\par
\begin{verbatim}
print(5 <= True)
\end{verbatim}

{\begin {center}\textbf{Завдання 5}\end {center}}

Що буде виведено за виконання? (повідомлення інтерпретатора про помилку позначати error)\par
\begin{verbatim}
print(2.0 <= 2 < False == 3)
\end{verbatim}

{\begin {center}\textbf{Завдання 6}\end {center}}

Що буде виведено за виконання? (повідомлення інтерпретатора про помилку позначати error)\par
\begin{verbatim}

def f(n):
    print("f", end="");
    return n>0

print(f(5) and f(4))
\end{verbatim}

{\begin {center}\textbf{Завдання 7}\end {center}}

Що буде виведено за виконання (повідомлення інтерпретатора про помилку позначати error)?\par
\begin{verbatim}
def g(a,b):
    c=53
    if a:
        return 9
    elif a>-5:
         return 8
    else: 
        c=7
    return c

print(g(-3,-5))
\end{verbatim}

{\begin {center}\textbf{Завдання 8}\end {center}}

Що буде виведено за виконання (повідомлення інтерпретатора про помилку позначати error)?\par
\begin{verbatim}
a, b, c = 8, 6, 7
def g(a, b=9, c):
    print(a, b, c, end=" ")

g(b=1, c=0, 3)
print(a, b, c)
\end{verbatim}

{\begin {center}\textbf{Завдання 9}\end {center}}

Що буде виведено за виконання (повідомлення інтерпретатора про помилку позначати error)?\par
\begin{verbatim}
a,b,c=5,4,3
def h(a):
    a=1
    b=2
    c=5
    return a+b+c

a,b,c=1,2,3
print(h(a),a,b,c)
\end{verbatim}

{\begin {center}\textbf{Завдання 10}\end {center}}


Написати функцiю f, що отримує щось та повертає щось. Невелика загадка все-таки має залишатися. Але нічого принципово нового відносно задач, що наявні у pdf-ках не планується.. Стиль запису умови також оцінюється.\par
\end {large}}
\newpage

{\Large{17.10.2024 Бета-версія: Контрольна робота \No 1, варіант  \No \textbf { 86 }\par}}
{
\begin {large}
{\begin {center}\textbf{Завдання 1}\end {center}}

Що буде виведено за виконання? (повідомлення інтерпретатора про помилку позначати error)\par
\begin{verbatim}
print(4/4%6*3)
\end{verbatim}

{\begin {center}\textbf{Завдання 2}\end {center}}

Що буде виведено за виконання? (повідомлення інтерпретатора про помилку позначати error)\par
\begin{verbatim}
print(4**1.0*1)
\end{verbatim}

{\begin {center}\textbf{Завдання 3}\end {center}}

Що буде виведено за виконання? (повідомлення інтерпретатора про помилку позначати error)\par
\begin{verbatim}
print(6>True and not 8==4 and 5.0>=2.0)
\end{verbatim}

{\begin {center}\textbf{Завдання 4}\end {center}}

Що буде виведено за виконання? (повідомлення інтерпретатора про помилку позначати error)\par
\begin{verbatim}
print("6j" != "3.0")
\end{verbatim}

{\begin {center}\textbf{Завдання 5}\end {center}}

Що буде виведено за виконання? (повідомлення інтерпретатора про помилку позначати error)\par
\begin{verbatim}
print(6 > 3.0 >= 4 != 6.0)
\end{verbatim}

{\begin {center}\textbf{Завдання 6}\end {center}}

Що буде виведено за виконання? (повідомлення інтерпретатора про помилку позначати error)\par
\begin{verbatim}

def f(n):
    print("f", end="");
    return n

print(f(7) or f(3))
\end{verbatim}

{\begin {center}\textbf{Завдання 7}\end {center}}

Що буде виведено за виконання (повідомлення інтерпретатора про помилку позначати error)?\par
\begin{verbatim}
def g(d):
    z=31
    if d<=-2: 
        z=2
    elif d>0:
         z=6
    else:
         return 9
    return z

print(g(-1))
\end{verbatim}

{\begin {center}\textbf{Завдання 8}\end {center}}

Що буде виведено за виконання (повідомлення інтерпретатора про помилку позначати error)?\par
\begin{verbatim}
a, b, c = 8, 9, 6
def f(a, b, c=7):
    print(a, b, c, end=" ")

f(2, b=0)
print(a, b, c)
\end{verbatim}

{\begin {center}\textbf{Завдання 9}\end {center}}

Що буде виведено за виконання (повідомлення інтерпретатора про помилку позначати error)?\par
\begin{verbatim}
a,b,c=4,8,9
def f(b):
    a=1
    b+=2
    c=3
    return a+b+c

a,b,c=2,5,8
print(f(b),a,b,c)
\end{verbatim}

{\begin {center}\textbf{Завдання 10}\end {center}}


Написати функцiю f, що отримує щось та повертає щось. Невелика загадка все-таки має залишатися. Але нічого принципово нового відносно задач, що наявні у pdf-ках не планується.. Стиль запису умови також оцінюється.\par
\end {large}}
\newpage

{\Large{17.10.2024 Бета-версія: Контрольна робота \No 1, варіант  \No \textbf { 87 }\par}}
{
\begin {large}
{\begin {center}\textbf{Завдання 1}\end {center}}

Що буде виведено за виконання? (повідомлення інтерпретатора про помилку позначати error)\par
\begin{verbatim}
print(5//3//3*8)
\end{verbatim}

{\begin {center}\textbf{Завдання 2}\end {center}}

Що буде виведено за виконання? (повідомлення інтерпретатора про помилку позначати error)\par
\begin{verbatim}
print(4**4-2)
\end{verbatim}

{\begin {center}\textbf{Завдання 3}\end {center}}

Що буде виведено за виконання? (повідомлення інтерпретатора про помилку позначати error)\par
\begin{verbatim}
print(3>=6 and not 9.0!=True or 8.0==5)
\end{verbatim}

{\begin {center}\textbf{Завдання 4}\end {center}}

Що буде виведено за виконання? (повідомлення інтерпретатора про помилку позначати error)\par
\begin{verbatim}
print(6.0 < "False")
\end{verbatim}

{\begin {center}\textbf{Завдання 5}\end {center}}

Що буде виведено за виконання? (повідомлення інтерпретатора про помилку позначати error)\par
\begin{verbatim}
print(7  is  5  is  True == 9.0)
\end{verbatim}

{\begin {center}\textbf{Завдання 6}\end {center}}

Що буде виведено за виконання? (повідомлення інтерпретатора про помилку позначати error)\par
\begin{verbatim}

def f(n):
    print("f", end="");
    return n

print(f(-4) and f(4))
\end{verbatim}

{\begin {center}\textbf{Завдання 7}\end {center}}

Що буде виведено за виконання (повідомлення інтерпретатора про помилку позначати error)?\par
\begin{verbatim}
def h(a,b):
    c=97
    if b>=0:
        return 0
    if b<=-1:
         c=4
    else: 
        c=1
    return c

print(h(-7,-7))
\end{verbatim}

{\begin {center}\textbf{Завдання 8}\end {center}}

Що буде виведено за виконання (повідомлення інтерпретатора про помилку позначати error)?\par
\begin{verbatim}
a, b, c = 6, 9, 8
def h(a, b=7, c=6):
    print(a, b, c, end=" ")

h(a=5, b=1, c=3)
print(a, b, c)
\end{verbatim}

{\begin {center}\textbf{Завдання 9}\end {center}}

Що буде виведено за виконання (повідомлення інтерпретатора про помилку позначати error)?\par
\begin{verbatim}
a,b,c=7,8,4
def f(a):
    a=3
    b=1
    c=4
    return a+b+c

a,b,c=1,0,5
print(f(b),a,b,c)
\end{verbatim}

{\begin {center}\textbf{Завдання 10}\end {center}}


Написати функцiю f, що отримує щось та повертає щось. Невелика загадка все-таки має залишатися. Але нічого принципово нового відносно задач, що наявні у pdf-ках не планується.. Стиль запису умови також оцінюється.\par
\end {large}}
\newpage

{\Large{17.10.2024 Бета-версія: Контрольна робота \No 1, варіант  \No \textbf { 88 }\par}}
{
\begin {large}
{\begin {center}\textbf{Завдання 1}\end {center}}

Що буде виведено за виконання? (повідомлення інтерпретатора про помилку позначати error)\par
\begin{verbatim}
print(7/6/9*4)
\end{verbatim}

{\begin {center}\textbf{Завдання 2}\end {center}}

Що буде виведено за виконання? (повідомлення інтерпретатора про помилку позначати error)\par
\begin{verbatim}
print(2**-2+-1)
\end{verbatim}

{\begin {center}\textbf{Завдання 3}\end {center}}

Що буде виведено за виконання? (повідомлення інтерпретатора про помилку позначати error)\par
\begin{verbatim}
print(2<=7 or not 9<4.0 and False>6.0)
\end{verbatim}

{\begin {center}\textbf{Завдання 4}\end {center}}

Що буде виведено за виконання? (повідомлення інтерпретатора про помилку позначати error)\par
\begin{verbatim}
print(4 > "False")
\end{verbatim}

{\begin {center}\textbf{Завдання 5}\end {center}}

Що буде виведено за виконання? (повідомлення інтерпретатора про помилку позначати error)\par
\begin{verbatim}
print(8 < True != 7.0 > 9)
\end{verbatim}

{\begin {center}\textbf{Завдання 6}\end {center}}

Що буде виведено за виконання? (повідомлення інтерпретатора про помилку позначати error)\par
\begin{verbatim}

def f(n):
    print("f", end="");
    return n!=2

print(f(5) or f(-2))
\end{verbatim}

{\begin {center}\textbf{Завдання 7}\end {center}}

Що буде виведено за виконання (повідомлення інтерпретатора про помилку позначати error)?\par
\begin{verbatim}
def g(b):
    z=79
    if b<-5: 
        z=8
    if b==5:
         return 3
    else:
         return 4
    return z

print(g(8))
\end{verbatim}

{\begin {center}\textbf{Завдання 8}\end {center}}

Що буде виведено за виконання (повідомлення інтерпретатора про помилку позначати error)?\par
\begin{verbatim}
a, b, c = 9, 8, 7
def f(a, b, c):
    print(a, b, c, end=" ")

f(4, 1, a=2)
print(a, b, c)
\end{verbatim}

{\begin {center}\textbf{Завдання 9}\end {center}}

Що буде виведено за виконання (повідомлення інтерпретатора про помилку позначати error)?\par
\begin{verbatim}
a,b,c=6,9,4
def g(b):
    a=2
    b=1
    c=3
    return a+b+c

a,b,c=2,8,9
print(g(b),a,b,c)
\end{verbatim}

{\begin {center}\textbf{Завдання 10}\end {center}}


Написати функцiю f, що отримує щось та повертає щось. Невелика загадка все-таки має залишатися. Але нічого принципово нового відносно задач, що наявні у pdf-ках не планується.. Стиль запису умови також оцінюється.\par
\end {large}}
\newpage

{\Large{17.10.2024 Бета-версія: Контрольна робота \No 1, варіант  \No \textbf { 89 }\par}}
{
\begin {large}
{\begin {center}\textbf{Завдання 1}\end {center}}

Що буде виведено за виконання? (повідомлення інтерпретатора про помилку позначати error)\par
\begin{verbatim}
print(3//7/7*6)
\end{verbatim}

{\begin {center}\textbf{Завдання 2}\end {center}}

Що буде виведено за виконання? (повідомлення інтерпретатора про помилку позначати error)\par
\begin{verbatim}
print(1**0.5*False)
\end{verbatim}

{\begin {center}\textbf{Завдання 3}\end {center}}

Що буде виведено за виконання? (повідомлення інтерпретатора про помилку позначати error)\par
\begin{verbatim}
print(not 4.0>True and 9.0!=2 and 4>=5)
\end{verbatim}

{\begin {center}\textbf{Завдання 4}\end {center}}

Що буде виведено за виконання? (повідомлення інтерпретатора про помилку позначати error)\par
\begin{verbatim}
print(True == "8.0j")
\end{verbatim}

{\begin {center}\textbf{Завдання 5}\end {center}}

Що буде виведено за виконання? (повідомлення інтерпретатора про помилку позначати error)\par
\begin{verbatim}
print(4.0 >= 5 <= 2 > False)
\end{verbatim}

{\begin {center}\textbf{Завдання 6}\end {center}}

Що буде виведено за виконання? (повідомлення інтерпретатора про помилку позначати error)\par
\begin{verbatim}

def f(n):
    print("f", end="");
    return n<=-1

print(f(-8) or f(9))
\end{verbatim}

{\begin {center}\textbf{Завдання 7}\end {center}}

Що буде виведено за виконання (повідомлення інтерпретатора про помилку позначати error)?\par
\begin{verbatim}
def g(a,b):
    c=70
    if a<-3:
        return 6
    elif a!=1:
         return 2
    else: 
        c=3
    return c

print(g(-2,1))
\end{verbatim}

{\begin {center}\textbf{Завдання 8}\end {center}}

Що буде виведено за виконання (повідомлення інтерпретатора про помилку позначати error)?\par
\begin{verbatim}
a, b, c = 6, 7, 9
def h(a, b=8, c):
    print(a, b, c, end=" ")

h(3, 5, 4)
print(a, b, c)
\end{verbatim}

{\begin {center}\textbf{Завдання 9}\end {center}}

Що буде виведено за виконання (повідомлення інтерпретатора про помилку позначати error)?\par
\begin{verbatim}
a,b,c=4,7,6
def h(b):
    a=5
    b*=1
    c=4
    return a+b+c

a,b,c=3,0,1
print(h(a),a,b,c)
\end{verbatim}

{\begin {center}\textbf{Завдання 10}\end {center}}


Написати функцiю f, що отримує щось та повертає щось. Невелика загадка все-таки має залишатися. Але нічого принципово нового відносно задач, що наявні у pdf-ках не планується.. Стиль запису умови також оцінюється.\par
\end {large}}
\newpage

{\Large{17.10.2024 Бета-версія: Контрольна робота \No 1, варіант  \No \textbf { 90 }\par}}
{
\begin {large}
{\begin {center}\textbf{Завдання 1}\end {center}}

Що буде виведено за виконання? (повідомлення інтерпретатора про помилку позначати error)\par
\begin{verbatim}
print(6/5%5*5)
\end{verbatim}

{\begin {center}\textbf{Завдання 2}\end {center}}

Що буде виведено за виконання? (повідомлення інтерпретатора про помилку позначати error)\par
\begin{verbatim}
print(4**-1+-2)
\end{verbatim}

{\begin {center}\textbf{Завдання 3}\end {center}}

Що буде виведено за виконання? (повідомлення інтерпретатора про помилку позначати error)\par
\begin{verbatim}
print(not 3<9 or False==6 or 7.0<=3.0)
\end{verbatim}

{\begin {center}\textbf{Завдання 4}\end {center}}

Що буде виведено за виконання? (повідомлення інтерпретатора про помилку позначати error)\par
\begin{verbatim}
print("2" >= 2.0)
\end{verbatim}

{\begin {center}\textbf{Завдання 5}\end {center}}

Що буде виведено за виконання? (повідомлення інтерпретатора про помилку позначати error)\par
\begin{verbatim}
print(9  is  8.0 <= 4 >= 5.0)
\end{verbatim}

{\begin {center}\textbf{Завдання 6}\end {center}}

Що буде виведено за виконання? (повідомлення інтерпретатора про помилку позначати error)\par
\begin{verbatim}

def f(n):
    print("f", end="");
    return n

print(f(-3) and f(-9))
\end{verbatim}

{\begin {center}\textbf{Завдання 7}\end {center}}

Що буде виведено за виконання (повідомлення інтерпретатора про помилку позначати error)?\par
\begin{verbatim}
def f(c):
    w=83
    if c<1: 
        w=1
    elif c==3:
         w=1
    else:
         return 0
    return w

print(f(-3))
\end{verbatim}

{\begin {center}\textbf{Завдання 8}\end {center}}

Що буде виведено за виконання (повідомлення інтерпретатора про помилку позначати error)?\par
\begin{verbatim}
a, b, c = 9, 7, 6
def g(a, b, c):
    print(a, b, c, end=" ")

g(0, c=3, b=1)
print(a, b, c)
\end{verbatim}

{\begin {center}\textbf{Завдання 9}\end {center}}

Що буде виведено за виконання (повідомлення інтерпретатора про помилку позначати error)?\par
\begin{verbatim}
a,b,c=9,6,2
def g(a):
    a-=3
    b=2
    c=5
    return a+b+c

a,b,c=8,1,7
print(g(a),a,b,c)
\end{verbatim}

{\begin {center}\textbf{Завдання 10}\end {center}}


Написати функцiю f, що отримує щось та повертає щось. Невелика загадка все-таки має залишатися. Але нічого принципово нового відносно задач, що наявні у pdf-ках не планується.. Стиль запису умови також оцінюється.\par
\end {large}}
\newpage

{\Large{17.10.2024 Бета-версія: Контрольна робота \No 1, варіант  \No \textbf { 91 }\par}}
{
\begin {large}
{\begin {center}\textbf{Завдання 1}\end {center}}

Що буде виведено за виконання? (повідомлення інтерпретатора про помилку позначати error)\par
\begin{verbatim}
print(9/8*4*7)
\end{verbatim}

{\begin {center}\textbf{Завдання 2}\end {center}}

Що буде виведено за виконання? (повідомлення інтерпретатора про помилку позначати error)\par
\begin{verbatim}
print(2**False-2)
\end{verbatim}

{\begin {center}\textbf{Завдання 3}\end {center}}

Що буде виведено за виконання? (повідомлення інтерпретатора про помилку позначати error)\par
\begin{verbatim}
print(not 8<True or 7==5.0 or 2.0!=9)
\end{verbatim}

{\begin {center}\textbf{Завдання 4}\end {center}}

Що буде виведено за виконання? (повідомлення інтерпретатора про помилку позначати error)\par
\begin{verbatim}
print("False" != 9)
\end{verbatim}

{\begin {center}\textbf{Завдання 5}\end {center}}

Що буде виведено за виконання? (повідомлення інтерпретатора про помилку позначати error)\par
\begin{verbatim}
print(6 != 4.0 == 2.0 < True)
\end{verbatim}

{\begin {center}\textbf{Завдання 6}\end {center}}

Що буде виведено за виконання? (повідомлення інтерпретатора про помилку позначати error)\par
\begin{verbatim}

def f(n):
    print("f", end="");
    return n

print(f(-5) and f(7))
\end{verbatim}

{\begin {center}\textbf{Завдання 7}\end {center}}

Що буде виведено за виконання (повідомлення інтерпретатора про помилку позначати error)?\par
\begin{verbatim}
def h(d):
    v=80
    if d: 
        return 7
    if d>4:
         v=9
    else:
         return 4
    return v

print(h(-6))
\end{verbatim}

{\begin {center}\textbf{Завдання 8}\end {center}}

Що буде виведено за виконання (повідомлення інтерпретатора про помилку позначати error)?\par
\begin{verbatim}
a, b, c = 8, 8, 7
def h(a, b, c=9):
    print(a, b, c, end=" ")

h(5, 0, a=2)
print(a, b, c)
\end{verbatim}

{\begin {center}\textbf{Завдання 9}\end {center}}

Що буде виведено за виконання (повідомлення інтерпретатора про помилку позначати error)?\par
\begin{verbatim}
a,b,c=5,0,3
def f(b):
    a+=1
    b=4
    c=2
    return a+b+c

a,b,c=4,1,7
print(f(b),a,b,c)
\end{verbatim}

{\begin {center}\textbf{Завдання 10}\end {center}}


Написати функцiю f, що отримує щось та повертає щось. Невелика загадка все-таки має залишатися. Але нічого принципово нового відносно задач, що наявні у pdf-ках не планується.. Стиль запису умови також оцінюється.\par
\end {large}}
\newpage

{\Large{17.10.2024 Бета-версія: Контрольна робота \No 1, варіант  \No \textbf { 92 }\par}}
{
\begin {large}
{\begin {center}\textbf{Завдання 1}\end {center}}

Що буде виведено за виконання? (повідомлення інтерпретатора про помилку позначати error)\par
\begin{verbatim}
print(8//8//6*3)
\end{verbatim}

{\begin {center}\textbf{Завдання 2}\end {center}}

Що буде виведено за виконання? (повідомлення інтерпретатора про помилку позначати error)\par
\begin{verbatim}
print(1**1.0*-2)
\end{verbatim}

{\begin {center}\textbf{Завдання 3}\end {center}}

Що буде виведено за виконання? (повідомлення інтерпретатора про помилку позначати error)\par
\begin{verbatim}
print(not 6<=8.0 and False>=6.0 and 5>8)
\end{verbatim}

{\begin {center}\textbf{Завдання 4}\end {center}}

Що буде виведено за виконання? (повідомлення інтерпретатора про помилку позначати error)\par
\begin{verbatim}
print(True <= 4.0)
\end{verbatim}

{\begin {center}\textbf{Завдання 5}\end {center}}

Що буде виведено за виконання? (повідомлення інтерпретатора про помилку позначати error)\par
\begin{verbatim}
print(8 >= 3 == 7  is  False)
\end{verbatim}

{\begin {center}\textbf{Завдання 6}\end {center}}

Що буде виведено за виконання? (повідомлення інтерпретатора про помилку позначати error)\par
\begin{verbatim}

def f(n):
    print("f", end="");
    return n==-3

print(f(0) or f(-6))
\end{verbatim}

{\begin {center}\textbf{Завдання 7}\end {center}}

Що буде виведено за виконання (повідомлення інтерпретатора про помилку позначати error)?\par
\begin{verbatim}
def f(a,b):
    c=64
    if b:
        return 5
    if b==3:
         c=9
    else: 
        return 8
    return c

print(f(5,2))
\end{verbatim}

{\begin {center}\textbf{Завдання 8}\end {center}}

Що буде виведено за виконання (повідомлення інтерпретатора про помилку позначати error)?\par
\begin{verbatim}
a, b, c = 6, 8, 9
def g(a, b=7, c=6):
    print(a, b, c, end=" ")

g(4, b=4)
print(a, b, c)
\end{verbatim}

{\begin {center}\textbf{Завдання 9}\end {center}}

Що буде виведено за виконання (повідомлення інтерпретатора про помилку позначати error)?\par
\begin{verbatim}
a,b,c=5,8,0
def h(a):
    a-=4
    b=5
    c=3
    return a+b+c

a,b,c=4,6,9
print(h(b),a,b,c)
\end{verbatim}

{\begin {center}\textbf{Завдання 10}\end {center}}


Написати функцiю f, що отримує щось та повертає щось. Невелика загадка все-таки має залишатися. Але нічого принципово нового відносно задач, що наявні у pdf-ках не планується.. Стиль запису умови також оцінюється.\par
\end {large}}
\newpage

{\Large{17.10.2024 Бета-версія: Контрольна робота \No 1, варіант  \No \textbf { 93 }\par}}
{
\begin {large}
{\begin {center}\textbf{Завдання 1}\end {center}}

Що буде виведено за виконання? (повідомлення інтерпретатора про помилку позначати error)\par
\begin{verbatim}
print(5//9/7*9)
\end{verbatim}

{\begin {center}\textbf{Завдання 2}\end {center}}

Що буде виведено за виконання? (повідомлення інтерпретатора про помилку позначати error)\par
\begin{verbatim}
print(4**-1.0*False)
\end{verbatim}

{\begin {center}\textbf{Завдання 3}\end {center}}

Що буде виведено за виконання? (повідомлення інтерпретатора про помилку позначати error)\par
\begin{verbatim}
print(not 7<=3 and 9.0!=False and 2>2.0)
\end{verbatim}

{\begin {center}\textbf{Завдання 4}\end {center}}

Що буде виведено за виконання? (повідомлення інтерпретатора про помилку позначати error)\par
\begin{verbatim}
print("3" > "5.0")
\end{verbatim}

{\begin {center}\textbf{Завдання 5}\end {center}}

Що буде виведено за виконання? (повідомлення інтерпретатора про помилку позначати error)\par
\begin{verbatim}
print(7 != 8.0 <= 9 > 7.0)
\end{verbatim}

{\begin {center}\textbf{Завдання 6}\end {center}}

Що буде виведено за виконання? (повідомлення інтерпретатора про помилку позначати error)\par
\begin{verbatim}

def f(n):
    print("f", end="");
    return n

print(f(-7) and f(8))
\end{verbatim}

{\begin {center}\textbf{Завдання 7}\end {center}}

Що буде виведено за виконання (повідомлення інтерпретатора про помилку позначати error)?\par
\begin{verbatim}
def h(a,b):
    c=23
    if b<=5:
        c=4
    elif a!=-1:
         c=5
    else: 
        return 2
    return c

print(h(7,4))
\end{verbatim}

{\begin {center}\textbf{Завдання 8}\end {center}}

Що буде виведено за виконання (повідомлення інтерпретатора про помилку позначати error)?\par
\begin{verbatim}
a, b, c = 7, 6, 8
def f(a, b=9, c):
    print(a, b, c, end=" ")

f(a=5, 3, c=1)
print(a, b, c)
\end{verbatim}

{\begin {center}\textbf{Завдання 9}\end {center}}

Що буде виведено за виконання (повідомлення інтерпретатора про помилку позначати error)?\par
\begin{verbatim}
a,b,c=3,2,5
def g(a):
    a*=1
    b=2
    c=1
    return a+b+c

a,b,c=8,7,2
print(g(a),a,b,c)
\end{verbatim}

{\begin {center}\textbf{Завдання 10}\end {center}}


Написати функцiю f, що отримує щось та повертає щось. Невелика загадка все-таки має залишатися. Але нічого принципово нового відносно задач, що наявні у pdf-ках не планується.. Стиль запису умови також оцінюється.\par
\end {large}}
\newpage

{\Large{17.10.2024 Бета-версія: Контрольна робота \No 1, варіант  \No \textbf { 94 }\par}}
{
\begin {large}
{\begin {center}\textbf{Завдання 1}\end {center}}

Що буде виведено за виконання? (повідомлення інтерпретатора про помилку позначати error)\par
\begin{verbatim}
print(3/4%5*5)
\end{verbatim}

{\begin {center}\textbf{Завдання 2}\end {center}}

Що буде виведено за виконання? (повідомлення інтерпретатора про помилку позначати error)\par
\begin{verbatim}
print(-3**True+1)
\end{verbatim}

{\begin {center}\textbf{Завдання 3}\end {center}}

Що буде виведено за виконання? (повідомлення інтерпретатора про помилку позначати error)\par
\begin{verbatim}
print(True>=4 or not 3.0<8 or 5.0==6)
\end{verbatim}

{\begin {center}\textbf{Завдання 4}\end {center}}

Що буде виведено за виконання? (повідомлення інтерпретатора про помилку позначати error)\par
\begin{verbatim}
print(7.0 >= 7)
\end{verbatim}

{\begin {center}\textbf{Завдання 5}\end {center}}

Що буде виведено за виконання? (повідомлення інтерпретатора про помилку позначати error)\par
\begin{verbatim}
print(False < 8  is  6.0 < 3.0)
\end{verbatim}

{\begin {center}\textbf{Завдання 6}\end {center}}

Що буде виведено за виконання? (повідомлення інтерпретатора про помилку позначати error)\par
\begin{verbatim}

def f(n):
    print("f", end="");
    return n>=2

print(f(-1) or f(6))
\end{verbatim}

{\begin {center}\textbf{Завдання 7}\end {center}}

Що буде виведено за виконання (повідомлення інтерпретатора про помилку позначати error)?\par
\begin{verbatim}
def h(c):
    x=72
    if c!=-3: 
        x=8
    elif c>=-1:
         return 3
    else:
         x=5
    return x

print(h(5))
\end{verbatim}

{\begin {center}\textbf{Завдання 8}\end {center}}

Що буде виведено за виконання (повідомлення інтерпретатора про помилку позначати error)?\par
\begin{verbatim}
a, b, c = 9, 8, 7
def g(a, b, c=6):
    print(a, b, c, end=" ")

g(0, 2)
print(a, b, c)
\end{verbatim}

{\begin {center}\textbf{Завдання 9}\end {center}}

Що буде виведено за виконання (повідомлення інтерпретатора про помилку позначати error)?\par
\begin{verbatim}
a,b,c=1,4,0
def g(b):
    a*=4
    b=5
    c=3
    return a+b+c

a,b,c=3,9,6
print(g(b),a,b,c)
\end{verbatim}

{\begin {center}\textbf{Завдання 10}\end {center}}


Написати функцiю f, що отримує щось та повертає щось. Невелика загадка все-таки має залишатися. Але нічого принципово нового відносно задач, що наявні у pdf-ках не планується.. Стиль запису умови також оцінюється.\par
\end {large}}
\newpage

{\Large{17.10.2024 Бета-версія: Контрольна робота \No 1, варіант  \No \textbf { 95 }\par}}
{
\begin {large}
{\begin {center}\textbf{Завдання 1}\end {center}}

Що буде виведено за виконання? (повідомлення інтерпретатора про помилку позначати error)\par
\begin{verbatim}
print(4//3*9*8)
\end{verbatim}

{\begin {center}\textbf{Завдання 2}\end {center}}

Що буде виведено за виконання? (повідомлення інтерпретатора про помилку позначати error)\par
\begin{verbatim}
print(9**-0.5*True)
\end{verbatim}

{\begin {center}\textbf{Завдання 3}\end {center}}

Що буде виведено за виконання? (повідомлення інтерпретатора про помилку позначати error)\par
\begin{verbatim}
print(7>9 or not 4<4.0 and 8.0!=True)
\end{verbatim}

{\begin {center}\textbf{Завдання 4}\end {center}}

Що буде виведено за виконання? (повідомлення інтерпретатора про помилку позначати error)\par
\begin{verbatim}
print("False" == True)
\end{verbatim}

{\begin {center}\textbf{Завдання 5}\end {center}}

Що буде виведено за виконання? (повідомлення інтерпретатора про помилку позначати error)\par
\begin{verbatim}
print(4 >= 5 != True == 3)
\end{verbatim}

{\begin {center}\textbf{Завдання 6}\end {center}}

Що буде виведено за виконання? (повідомлення інтерпретатора про помилку позначати error)\par
\begin{verbatim}

def f(n):
    print("f", end="");
    return n

print(f(2) and f(1))
\end{verbatim}

{\begin {center}\textbf{Завдання 7}\end {center}}

Що буде виведено за виконання (повідомлення інтерпретатора про помилку позначати error)?\par
\begin{verbatim}
def f(a):
    u=65
    if a<=-4: 
        return 6
    if a<=0:
         u=2
    else:
         return 1
    return u

print(f(-7))
\end{verbatim}

{\begin {center}\textbf{Завдання 8}\end {center}}

Що буде виведено за виконання (повідомлення інтерпретатора про помилку позначати error)?\par
\begin{verbatim}
a, b, c = 6, 9, 8
def f(a, b=7, c=8):
    print(a, b, c, end=" ")

f(a=4, b=2, c=5)
print(a, b, c)
\end{verbatim}

{\begin {center}\textbf{Завдання 9}\end {center}}

Що буде виведено за виконання (повідомлення інтерпретатора про помилку позначати error)?\par
\begin{verbatim}
a,b,c=2,8,4
def h(a):
    a=1
    b-=3
    c=4
    return a+b+c

a,b,c=3,9,0
print(h(a),a,b,c)
\end{verbatim}

{\begin {center}\textbf{Завдання 10}\end {center}}


Написати функцiю f, що отримує щось та повертає щось. Невелика загадка все-таки має залишатися. Але нічого принципово нового відносно задач, що наявні у pdf-ках не планується.. Стиль запису умови також оцінюється.\par
\end {large}}
\newpage

{\Large{17.10.2024 Бета-версія: Контрольна робота \No 1, варіант  \No \textbf { 96 }\par}}
{
\begin {large}
{\begin {center}\textbf{Завдання 1}\end {center}}

Що буде виведено за виконання? (повідомлення інтерпретатора про помилку позначати error)\par
\begin{verbatim}
print(6/7//4*7)
\end{verbatim}

{\begin {center}\textbf{Завдання 2}\end {center}}

Що буде виведено за виконання? (повідомлення інтерпретатора про помилку позначати error)\par
\begin{verbatim}
print(-1**-2--1)
\end{verbatim}

{\begin {center}\textbf{Завдання 3}\end {center}}

Що буде виведено за виконання? (повідомлення інтерпретатора про помилку позначати error)\par
\begin{verbatim}
print(not 3>=False and 6.0==2 or 5<=7.0)
\end{verbatim}

{\begin {center}\textbf{Завдання 4}\end {center}}

Що буде виведено за виконання? (повідомлення інтерпретатора про помилку позначати error)\par
\begin{verbatim}
print("8.0j" < "4j")
\end{verbatim}

{\begin {center}\textbf{Завдання 5}\end {center}}

Що буде виведено за виконання? (повідомлення інтерпретатора про помилку позначати error)\par
\begin{verbatim}
print(5.0 <= 2 > 6  is  9.0)
\end{verbatim}

{\begin {center}\textbf{Завдання 6}\end {center}}

Що буде виведено за виконання? (повідомлення інтерпретатора про помилку позначати error)\par
\begin{verbatim}

def f(n):
    print("f", end="");
    return n<4

print(f(-3) or f(8))
\end{verbatim}

{\begin {center}\textbf{Завдання 7}\end {center}}

Що буде виведено за виконання (повідомлення інтерпретатора про помилку позначати error)?\par
\begin{verbatim}
def g(a,b):
    c=28
    if b:
        c=9
    if a>=1:
         return 6
    else: 
        return 3
    return c

print(g(2,-4))
\end{verbatim}

{\begin {center}\textbf{Завдання 8}\end {center}}

Що буде виведено за виконання (повідомлення інтерпретатора про помилку позначати error)?\par
\begin{verbatim}
a, b, c = 9, 7, 6
def h(a, b, c):
    print(a, b, c, end=" ")

h(b=0, c=1, 3)
print(a, b, c)
\end{verbatim}

{\begin {center}\textbf{Завдання 9}\end {center}}

Що буде виведено за виконання (повідомлення інтерпретатора про помилку позначати error)?\par
\begin{verbatim}
a,b,c=6,5,1
def f(b):
    a+=2
    b=5
    c=1
    return a+b+c

a,b,c=7,2,8
print(f(a),a,b,c)
\end{verbatim}

{\begin {center}\textbf{Завдання 10}\end {center}}


Написати функцiю f, що отримує щось та повертає щось. Невелика загадка все-таки має залишатися. Але нічого принципово нового відносно задач, що наявні у pdf-ках не планується.. Стиль запису умови також оцінюється.\par
\end {large}}
\newpage

{\Large{17.10.2024 Бета-версія: Контрольна робота \No 1, варіант  \No \textbf { 97 }\par}}
{
\begin {large}
{\begin {center}\textbf{Завдання 1}\end {center}}

Що буде виведено за виконання? (повідомлення інтерпретатора про помилку позначати error)\par
\begin{verbatim}
print(9/5/3*4)
\end{verbatim}

{\begin {center}\textbf{Завдання 2}\end {center}}

Що буде виведено за виконання? (повідомлення інтерпретатора про помилку позначати error)\par
\begin{verbatim}
print(1**1.0*1)
\end{verbatim}

{\begin {center}\textbf{Завдання 3}\end {center}}

Що буде виведено за виконання? (повідомлення інтерпретатора про помилку позначати error)\par
\begin{verbatim}
print(not 4.0!=7 and 5.0<3 or True==2)
\end{verbatim}

{\begin {center}\textbf{Завдання 4}\end {center}}

Що буде виведено за виконання? (повідомлення інтерпретатора про помилку позначати error)\par
\begin{verbatim}
print("4.0" != 9)
\end{verbatim}

{\begin {center}\textbf{Завдання 5}\end {center}}

Що буде виведено за виконання? (повідомлення інтерпретатора про помилку позначати error)\par
\begin{verbatim}
print(8 >= 4 > 2.0 <= 5)
\end{verbatim}

{\begin {center}\textbf{Завдання 6}\end {center}}

Що буде виведено за виконання? (повідомлення інтерпретатора про помилку позначати error)\par
\begin{verbatim}

def f(n):
    print("f", end="");
    return n<0

print(f(-4) and f(5))
\end{verbatim}

{\begin {center}\textbf{Завдання 7}\end {center}}

Що буде виведено за виконання (повідомлення інтерпретатора про помилку позначати error)?\par
\begin{verbatim}
def g(b):
    y=68
    if b: 
        y=5
    elif b==-2:
         return 2
    else:
         y=4
    return y

print(g(4))
\end{verbatim}

{\begin {center}\textbf{Завдання 8}\end {center}}

Що буде виведено за виконання (повідомлення інтерпретатора про помилку позначати error)?\par
\begin{verbatim}
a, b, c = 7, 8, 6
def h(a, b=9, c=8):
    print(a, b, c, end=" ")

h(a=0, 4, a=5)
print(a, b, c)
\end{verbatim}

{\begin {center}\textbf{Завдання 9}\end {center}}

Що буде виведено за виконання (повідомлення інтерпретатора про помилку позначати error)?\par
\begin{verbatim}
a,b,c=9,1,5
def g(b):
    a+=3
    b=4
    c=5
    return a+b+c

a,b,c=0,7,4
print(g(b),a,b,c)
\end{verbatim}

{\begin {center}\textbf{Завдання 10}\end {center}}


Написати функцiю f, що отримує щось та повертає щось. Невелика загадка все-таки має залишатися. Але нічого принципово нового відносно задач, що наявні у pdf-ках не планується.. Стиль запису умови також оцінюється.\par
\end {large}}
\newpage

{\Large{17.10.2024 Бета-версія: Контрольна робота \No 1, варіант  \No \textbf { 98 }\par}}
{
\begin {large}
{\begin {center}\textbf{Завдання 1}\end {center}}

Що буде виведено за виконання? (повідомлення інтерпретатора про помилку позначати error)\par
\begin{verbatim}
print(7//6//8*6)
\end{verbatim}

{\begin {center}\textbf{Завдання 2}\end {center}}

Що буде виведено за виконання? (повідомлення інтерпретатора про помилку позначати error)\par
\begin{verbatim}
print(9**-0.5*-1)
\end{verbatim}

{\begin {center}\textbf{Завдання 3}\end {center}}

Що буде виведено за виконання? (повідомлення інтерпретатора про помилку позначати error)\par
\begin{verbatim}
print(not False>5 or 6.0>=9 and 4<=3.0)
\end{verbatim}

{\begin {center}\textbf{Завдання 4}\end {center}}

Що буде виведено за виконання? (повідомлення інтерпретатора про помилку позначати error)\par
\begin{verbatim}
print("5j" <= "True")
\end{verbatim}

{\begin {center}\textbf{Завдання 5}\end {center}}

Що буде виведено за виконання? (повідомлення інтерпретатора про помилку позначати error)\par
\begin{verbatim}
print(3.0 != True == 8.0 < 4.0)
\end{verbatim}

{\begin {center}\textbf{Завдання 6}\end {center}}

Що буде виведено за виконання? (повідомлення інтерпретатора про помилку позначати error)\par
\begin{verbatim}

def f(n):
    print("f", end="");
    return n

print(f(-9) or f(3))
\end{verbatim}

{\begin {center}\textbf{Завдання 7}\end {center}}

Що буде виведено за виконання (повідомлення інтерпретатора про помилку позначати error)?\par
\begin{verbatim}
def f(d):
    w=18
    if d: 
        w=8
    elif d<2:
         return 7
    else:
         w=0
    return w

print(f(1))
\end{verbatim}

{\begin {center}\textbf{Завдання 8}\end {center}}

Що буде виведено за виконання (повідомлення інтерпретатора про помилку позначати error)?\par
\begin{verbatim}
a, b, c = 9, 7, 6
def f(a, b, c=8):
    print(a, b, c, end=" ")

f(2, c=3, b=1)
print(a, b, c)
\end{verbatim}

{\begin {center}\textbf{Завдання 9}\end {center}}

Що буде виведено за виконання (повідомлення інтерпретатора про помилку позначати error)?\par
\begin{verbatim}
a,b,c=6,3,8
def f(a):
    a=2
    b*=3
    c=2
    return a+b+c

a,b,c=5,4,6
print(f(a),a,b,c)
\end{verbatim}

{\begin {center}\textbf{Завдання 10}\end {center}}


Написати функцiю f, що отримує щось та повертає щось. Невелика загадка все-таки має залишатися. Але нічого принципово нового відносно задач, що наявні у pdf-ках не планується.. Стиль запису умови також оцінюється.\par
\end {large}}
\newpage

{\Large{17.10.2024 Бета-версія: Контрольна робота \No 1, варіант  \No \textbf { 99 }\par}}
{
\begin {large}
{\begin {center}\textbf{Завдання 1}\end {center}}

Що буде виведено за виконання? (повідомлення інтерпретатора про помилку позначати error)\par
\begin{verbatim}
print(5/6%6*7)
\end{verbatim}

{\begin {center}\textbf{Завдання 2}\end {center}}

Що буде виведено за виконання? (повідомлення інтерпретатора про помилку позначати error)\par
\begin{verbatim}
print(4**-1.0*-2)
\end{verbatim}

{\begin {center}\textbf{Завдання 3}\end {center}}

Що буде виведено за виконання? (повідомлення інтерпретатора про помилку позначати error)\par
\begin{verbatim}
print(not 6<=8 and 9.0>True or 2.0==2)
\end{verbatim}

{\begin {center}\textbf{Завдання 4}\end {center}}

Що буде виведено за виконання? (повідомлення інтерпретатора про помилку позначати error)\par
\begin{verbatim}
print(9.0 >= False)
\end{verbatim}

{\begin {center}\textbf{Завдання 5}\end {center}}

Що буде виведено за виконання? (повідомлення інтерпретатора про помилку позначати error)\par
\begin{verbatim}
print(False >= 6 > 7 < 3)
\end{verbatim}

{\begin {center}\textbf{Завдання 6}\end {center}}

Що буде виведено за виконання? (повідомлення інтерпретатора про помилку позначати error)\par
\begin{verbatim}

def f(n):
    print("f", end="");
    return n==1

print(f(1) and f(-5))
\end{verbatim}

{\begin {center}\textbf{Завдання 7}\end {center}}

Що буде виведено за виконання (повідомлення інтерпретатора про помилку позначати error)?\par
\begin{verbatim}
def g(c):
    z=70
    if c>-1: 
        return 6
    if c>=5:
         z=3
    else:
         return 9
    return z

print(g(9))
\end{verbatim}

{\begin {center}\textbf{Завдання 8}\end {center}}

Що буде виведено за виконання (повідомлення інтерпретатора про помилку позначати error)?\par
\begin{verbatim}
a, b, c = 9, 7, 6
def g(a, b, c):
    print(a, b, c, end=" ")

g(0, 1, b=2)
print(a, b, c)
\end{verbatim}

{\begin {center}\textbf{Завдання 9}\end {center}}

Що буде виведено за виконання (повідомлення інтерпретатора про помилку позначати error)?\par
\begin{verbatim}
a,b,c=0,5,1
def h(a):
    a=2
    b=4
    c=5
    return a+b+c

a,b,c=7,3,6
print(h(a),a,b,c)
\end{verbatim}

{\begin {center}\textbf{Завдання 10}\end {center}}


Написати функцiю f, що отримує щось та повертає щось. Невелика загадка все-таки має залишатися. Але нічого принципово нового відносно задач, що наявні у pdf-ках не планується.. Стиль запису умови також оцінюється.\par
\end {large}}
\newpage

{\Large{17.10.2024 Бета-версія: Контрольна робота \No 1, варіант  \No \textbf { 100 }\par}}
{
\begin {large}
{\begin {center}\textbf{Завдання 1}\end {center}}

Що буде виведено за виконання? (повідомлення інтерпретатора про помилку позначати error)\par
\begin{verbatim}
print(6//7*3*5)
\end{verbatim}

{\begin {center}\textbf{Завдання 2}\end {center}}

Що буде виведено за виконання? (повідомлення інтерпретатора про помилку позначати error)\par
\begin{verbatim}
print(1**-4-True)
\end{verbatim}

{\begin {center}\textbf{Завдання 3}\end {center}}

Що буде виведено за виконання? (повідомлення інтерпретатора про помилку позначати error)\par
\begin{verbatim}
print(7.0>=7 or not 8.0<8 and False!=9)
\end{verbatim}

{\begin {center}\textbf{Завдання 4}\end {center}}

Що буде виведено за виконання? (повідомлення інтерпретатора про помилку позначати error)\par
\begin{verbatim}
print(2 > "True")
\end{verbatim}

{\begin {center}\textbf{Завдання 5}\end {center}}

Що буде виведено за виконання? (повідомлення інтерпретатора про помилку позначати error)\par
\begin{verbatim}
print(2 != 6.0 == True <= 7.0)
\end{verbatim}

{\begin {center}\textbf{Завдання 6}\end {center}}

Що буде виведено за виконання? (повідомлення інтерпретатора про помилку позначати error)\par
\begin{verbatim}

def f(n):
    print("f", end="");
    return n

print(f(9) or f(2))
\end{verbatim}

{\begin {center}\textbf{Завдання 7}\end {center}}

Що буде виведено за виконання (повідомлення інтерпретатора про помилку позначати error)?\par
\begin{verbatim}
def f(a,b):
    c=61
    if b==-3:
        c=7
    elif a<-4:
         return 1
    else: 
        c=0
    return c

print(f(1,-1))
\end{verbatim}

{\begin {center}\textbf{Завдання 8}\end {center}}

Що буде виведено за виконання (повідомлення інтерпретатора про помилку позначати error)?\par
\begin{verbatim}
a, b, c = 8, 7, 6
def f(a, b=9, c):
    print(a, b, c, end=" ")

f(c=5, b=3, c=4)
print(a, b, c)
\end{verbatim}

{\begin {center}\textbf{Завдання 9}\end {center}}

Що буде виведено за виконання (повідомлення інтерпретатора про помилку позначати error)?\par
\begin{verbatim}
a,b,c=9,7,1
def h(b):
    a=1
    b-=4
    c=5
    return a+b+c

a,b,c=0,3,2
print(h(b),a,b,c)
\end{verbatim}

{\begin {center}\textbf{Завдання 10}\end {center}}


Написати функцiю f, що отримує щось та повертає щось. Невелика загадка все-таки має залишатися. Але нічого принципово нового відносно задач, що наявні у pdf-ках не планується.. Стиль запису умови також оцінюється.\par
\end {large}}
\newpage

{\Large{17.10.2024 Бета-версія: Контрольна робота \No 1, варіант  \No \textbf { 101 }\par}}
{
\begin {large}
{\begin {center}\textbf{Завдання 1}\end {center}}

Що буде виведено за виконання? (повідомлення інтерпретатора про помилку позначати error)\par
\begin{verbatim}
print(3/4%9*6)
\end{verbatim}

{\begin {center}\textbf{Завдання 2}\end {center}}

Що буде виведено за виконання? (повідомлення інтерпретатора про помилку позначати error)\par
\begin{verbatim}
print(1**0.5*False)
\end{verbatim}

{\begin {center}\textbf{Завдання 3}\end {center}}

Що буде виведено за виконання? (повідомлення інтерпретатора про помилку позначати error)\par
\begin{verbatim}
print(5>=7.0 and not True>3 or 6<=4.0)
\end{verbatim}

{\begin {center}\textbf{Завдання 4}\end {center}}

Що буде виведено за виконання? (повідомлення інтерпретатора про помилку позначати error)\par
\begin{verbatim}
print(3.0 < "2.0j")
\end{verbatim}

{\begin {center}\textbf{Завдання 5}\end {center}}

Що буде виведено за виконання? (повідомлення інтерпретатора про помилку позначати error)\par
\begin{verbatim}
print(9  is  6 >= 8 == 5.0)
\end{verbatim}

{\begin {center}\textbf{Завдання 6}\end {center}}

Що буде виведено за виконання? (повідомлення інтерпретатора про помилку позначати error)\par
\begin{verbatim}

def f(n):
    print("f", end="");
    return n

print(f(7) and f(-2))
\end{verbatim}

{\begin {center}\textbf{Завдання 7}\end {center}}

Що буде виведено за виконання (повідомлення інтерпретатора про помилку позначати error)?\par
\begin{verbatim}
def h(b):
    x=81
    if b!=-3: 
        return 9
    elif b>3:
         x=0
    else:
         x=4
    return x

print(h(3))
\end{verbatim}

{\begin {center}\textbf{Завдання 8}\end {center}}

Що буде виведено за виконання (повідомлення інтерпретатора про помилку позначати error)?\par
\begin{verbatim}
a, b, c = 8, 6, 9
def g(a, b, c=7):
    print(a, b, c, end=" ")

g(b=0, c=5, 3)
print(a, b, c)
\end{verbatim}

{\begin {center}\textbf{Завдання 9}\end {center}}

Що буде виведено за виконання (повідомлення інтерпретатора про помилку позначати error)?\par
\begin{verbatim}
a,b,c=9,5,2
def g(b):
    a=2
    b=3
    c=4
    return a+b+c

a,b,c=8,7,3
print(g(a),a,b,c)
\end{verbatim}

{\begin {center}\textbf{Завдання 10}\end {center}}


Написати функцiю f, що отримує щось та повертає щось. Невелика загадка все-таки має залишатися. Але нічого принципово нового відносно задач, що наявні у pdf-ках не планується.. Стиль запису умови також оцінюється.\par
\end {large}}
\newpage

{\Large{17.10.2024 Бета-версія: Контрольна робота \No 1, варіант  \No \textbf { 102 }\par}}
{
\begin {large}
{\begin {center}\textbf{Завдання 1}\end {center}}

Що буде виведено за виконання? (повідомлення інтерпретатора про помилку позначати error)\par
\begin{verbatim}
print(9//9//4*4)
\end{verbatim}

{\begin {center}\textbf{Завдання 2}\end {center}}

Що буде виведено за виконання? (повідомлення інтерпретатора про помилку позначати error)\par
\begin{verbatim}
print(3**-3+2)
\end{verbatim}

{\begin {center}\textbf{Завдання 3}\end {center}}

Що буде виведено за виконання? (повідомлення інтерпретатора про помилку позначати error)\par
\begin{verbatim}
print(9.0<False or not 6.0!=4 and 9==3)
\end{verbatim}

{\begin {center}\textbf{Завдання 4}\end {center}}

Що буде виведено за виконання? (повідомлення інтерпретатора про помилку позначати error)\par
\begin{verbatim}
print("6" == False)
\end{verbatim}

{\begin {center}\textbf{Завдання 5}\end {center}}

Що буде виведено за виконання? (повідомлення інтерпретатора про помилку позначати error)\par
\begin{verbatim}
print(2 != 3  is  9.0 < False)
\end{verbatim}

{\begin {center}\textbf{Завдання 6}\end {center}}

Що буде виведено за виконання? (повідомлення інтерпретатора про помилку позначати error)\par
\begin{verbatim}

def f(n):
    print("f", end="");
    return n>=-2

print(f(-1) or f(-8))
\end{verbatim}

{\begin {center}\textbf{Завдання 7}\end {center}}

Що буде виведено за виконання (повідомлення інтерпретатора про помилку позначати error)?\par
\begin{verbatim}
def h(a,b):
    c=19
    if b>0:
        return 2
    elif a>-5:
         c=3
    else: 
        return 1
    return c

print(h(-5,5))
\end{verbatim}

{\begin {center}\textbf{Завдання 8}\end {center}}

Що буде виведено за виконання (повідомлення інтерпретатора про помилку позначати error)?\par
\begin{verbatim}
a, b, c = 7, 6, 9
def h(a, b=8, c):
    print(a, b, c, end=" ")

h(2, a=4)
print(a, b, c)
\end{verbatim}

{\begin {center}\textbf{Завдання 9}\end {center}}

Що буде виведено за виконання (повідомлення інтерпретатора про помилку позначати error)?\par
\begin{verbatim}
a,b,c=9,0,4
def h(a):
    a*=1
    b=3
    c=5
    return a+b+c

a,b,c=7,5,6
print(h(a),a,b,c)
\end{verbatim}

{\begin {center}\textbf{Завдання 10}\end {center}}


Написати функцiю f, що отримує щось та повертає щось. Невелика загадка все-таки має залишатися. Але нічого принципово нового відносно задач, що наявні у pdf-ках не планується.. Стиль запису умови також оцінюється.\par
\end {large}}
\newpage

{\Large{17.10.2024 Бета-версія: Контрольна робота \No 1, варіант  \No \textbf { 103 }\par}}
{
\begin {large}
{\begin {center}\textbf{Завдання 1}\end {center}}

Що буде виведено за виконання? (повідомлення інтерпретатора про помилку позначати error)\par
\begin{verbatim}
print(4/8/7*8)
\end{verbatim}

{\begin {center}\textbf{Завдання 2}\end {center}}

Що буде виведено за виконання? (повідомлення інтерпретатора про помилку позначати error)\par
\begin{verbatim}
print(-3**True+2)
\end{verbatim}

{\begin {center}\textbf{Завдання 3}\end {center}}

Що буде виведено за виконання? (повідомлення інтерпретатора про помилку позначати error)\par
\begin{verbatim}
print(4!=2.0 and not False<=8.0 and 5<6)
\end{verbatim}

{\begin {center}\textbf{Завдання 4}\end {center}}

Що буде виведено за виконання? (повідомлення інтерпретатора про помилку позначати error)\par
\begin{verbatim}
print("True" == 7.0)
\end{verbatim}

{\begin {center}\textbf{Завдання 5}\end {center}}

Що буде виведено за виконання? (повідомлення інтерпретатора про помилку позначати error)\par
\begin{verbatim}
print(7 <= 5.0 > 3.0 > 4)
\end{verbatim}

{\begin {center}\textbf{Завдання 6}\end {center}}

Що буде виведено за виконання? (повідомлення інтерпретатора про помилку позначати error)\par
\begin{verbatim}

def f(n):
    print("f", end="");
    return n!=-4

print(f(-7) and f(0))
\end{verbatim}

{\begin {center}\textbf{Завдання 7}\end {center}}

Що буде виведено за виконання (повідомлення інтерпретатора про помилку позначати error)?\par
\begin{verbatim}
def h(a):
    y=86
    if a<4: 
        return 5
    if a==-5:
         y=7
    else:
         return 6
    return y

print(h(0))
\end{verbatim}

{\begin {center}\textbf{Завдання 8}\end {center}}

Що буде виведено за виконання (повідомлення інтерпретатора про помилку позначати error)?\par
\begin{verbatim}
a, b, c = 7, 8, 6
def f(a, b=9, c=8):
    print(a, b, c, end=" ")

f(1, 2)
print(a, b, c)
\end{verbatim}

{\begin {center}\textbf{Завдання 9}\end {center}}

Що буде виведено за виконання (повідомлення інтерпретатора про помилку позначати error)?\par
\begin{verbatim}
a,b,c=1,3,2
def f(b):
    a=2
    b+=4
    c=3
    return a+b+c

a,b,c=8,9,5
print(f(b),a,b,c)
\end{verbatim}

{\begin {center}\textbf{Завдання 10}\end {center}}


Написати функцiю f, що отримує щось та повертає щось. Невелика загадка все-таки має залишатися. Але нічого принципово нового відносно задач, що наявні у pdf-ках не планується.. Стиль запису умови також оцінюється.\par
\end {large}}
\newpage

{\Large{17.10.2024 Бета-версія: Контрольна робота \No 1, варіант  \No \textbf { 104 }\par}}
{
\begin {large}
{\begin {center}\textbf{Завдання 1}\end {center}}

Що буде виведено за виконання? (повідомлення інтерпретатора про помилку позначати error)\par
\begin{verbatim}
print(8//5*5*3)
\end{verbatim}

{\begin {center}\textbf{Завдання 2}\end {center}}

Що буде виведено за виконання? (повідомлення інтерпретатора про помилку позначати error)\par
\begin{verbatim}
print(2**True--1)
\end{verbatim}

{\begin {center}\textbf{Завдання 3}\end {center}}

Що буде виведено за виконання? (повідомлення інтерпретатора про помилку позначати error)\par
\begin{verbatim}
print(not 3.0>2 or 5.0==8 or True>=7)
\end{verbatim}

{\begin {center}\textbf{Завдання 4}\end {center}}

Що буде виведено за виконання? (повідомлення інтерпретатора про помилку позначати error)\par
\begin{verbatim}
print("3j" >= False)
\end{verbatim}

{\begin {center}\textbf{Завдання 5}\end {center}}

Що буде виведено за виконання? (повідомлення інтерпретатора про помилку позначати error)\par
\begin{verbatim}
print(False == 9 <= 5 >= 6.0)
\end{verbatim}

{\begin {center}\textbf{Завдання 6}\end {center}}

Що буде виведено за виконання? (повідомлення інтерпретатора про помилку позначати error)\par
\begin{verbatim}

def f(n):
    print("f", end="");
    return n

print(f(6) or f(4))
\end{verbatim}

{\begin {center}\textbf{Завдання 7}\end {center}}

Що буде виведено за виконання (повідомлення інтерпретатора про помилку позначати error)?\par
\begin{verbatim}
def h(a,b):
    c=44
    if a==4:
        c=6
    if b!=-2:
         return 7
    else: 
        c=5
    return c

print(h(4,-8))
\end{verbatim}

{\begin {center}\textbf{Завдання 8}\end {center}}

Що буде виведено за виконання (повідомлення інтерпретатора про помилку позначати error)?\par
\begin{verbatim}
a, b, c = 9, 7, 6
def g(a, b, c):
    print(a, b, c, end=" ")

g(0, 4, 1)
print(a, b, c)
\end{verbatim}

{\begin {center}\textbf{Завдання 9}\end {center}}

Що буде виведено за виконання (повідомлення інтерпретатора про помилку позначати error)?\par
\begin{verbatim}
a,b,c=0,4,6
def f(b):
    a=5
    b=1
    c=5
    return a+b+c

a,b,c=1,4,0
print(f(b),a,b,c)
\end{verbatim}

{\begin {center}\textbf{Завдання 10}\end {center}}


Написати функцiю f, що отримує щось та повертає щось. Невелика загадка все-таки має залишатися. Але нічого принципово нового відносно задач, що наявні у pdf-ках не планується.. Стиль запису умови також оцінюється.\par
\end {large}}
\newpage

{\Large{17.10.2024 Бета-версія: Контрольна робота \No 1, варіант  \No \textbf { 105 }\par}}
{
\begin {large}
{\begin {center}\textbf{Завдання 1}\end {center}}

Що буде виведено за виконання? (повідомлення інтерпретатора про помилку позначати error)\par
\begin{verbatim}
print(7//3%8*9)
\end{verbatim}

{\begin {center}\textbf{Завдання 2}\end {center}}

Що буде виведено за виконання? (повідомлення інтерпретатора про помилку позначати error)\par
\begin{verbatim}
print(-1**2+True)
\end{verbatim}

{\begin {center}\textbf{Завдання 3}\end {center}}

Що буде виведено за виконання? (повідомлення інтерпретатора про помилку позначати error)\par
\begin{verbatim}
print(2.0>False or not 5.0<=5 and 3==6)
\end{verbatim}

{\begin {center}\textbf{Завдання 4}\end {center}}

Що буде виведено за виконання? (повідомлення інтерпретатора про помилку позначати error)\par
\begin{verbatim}
print("6.0j" < 1)
\end{verbatim}

{\begin {center}\textbf{Завдання 5}\end {center}}

Що буде виведено за виконання? (повідомлення інтерпретатора про помилку позначати error)\par
\begin{verbatim}
print(4 < 7.0 != True  is  5)
\end{verbatim}

{\begin {center}\textbf{Завдання 6}\end {center}}

Що буде виведено за виконання? (повідомлення інтерпретатора про помилку позначати error)\par
\begin{verbatim}

def f(n):
    print("f", end="");
    return n>3

print(f(-6) or f(6))
\end{verbatim}

{\begin {center}\textbf{Завдання 7}\end {center}}

Що буде виведено за виконання (повідомлення інтерпретатора про помилку позначати error)?\par
\begin{verbatim}
def g(a):
    v=39
    if a>=1: 
        return 1
    if a!=0:
         v=8
    else:
         v=3
    return v

print(g(-5))
\end{verbatim}

{\begin {center}\textbf{Завдання 8}\end {center}}

Що буде виведено за виконання (повідомлення інтерпретатора про помилку позначати error)?\par
\begin{verbatim}
a, b, c = 9, 7, 6
def h(a, b=8, c=9):
    print(a, b, c, end=" ")

h(5, 3)
print(a, b, c)
\end{verbatim}

{\begin {center}\textbf{Завдання 9}\end {center}}

Що буде виведено за виконання (повідомлення інтерпретатора про помилку позначати error)?\par
\begin{verbatim}
a,b,c=0,1,4
def g(a):
    a-=2
    b=5
    c=1
    return a+b+c

a,b,c=3,8,7
print(g(a),a,b,c)
\end{verbatim}

{\begin {center}\textbf{Завдання 10}\end {center}}


Написати функцiю f, що отримує щось та повертає щось. Невелика загадка все-таки має залишатися. Але нічого принципово нового відносно задач, що наявні у pdf-ках не планується.. Стиль запису умови також оцінюється.\par
\end {large}}
\newpage

{\Large{17.10.2024 Бета-версія: Контрольна робота \No 1, варіант  \No \textbf { 106 }\par}}
{
\begin {large}
{\begin {center}\textbf{Завдання 1}\end {center}}

Що буде виведено за виконання? (повідомлення інтерпретатора про помилку позначати error)\par
\begin{verbatim}
print(3/8*4*7)
\end{verbatim}

{\begin {center}\textbf{Завдання 2}\end {center}}

Що буде виведено за виконання? (повідомлення інтерпретатора про помилку позначати error)\par
\begin{verbatim}
print(9**-0.5*-2)
\end{verbatim}

{\begin {center}\textbf{Завдання 3}\end {center}}

Що буде виведено за виконання? (повідомлення інтерпретатора про помилку позначати error)\par
\begin{verbatim}
print(2!=True and not 7.0<4 or 8>=4.0)
\end{verbatim}

{\begin {center}\textbf{Завдання 4}\end {center}}

Що буде виведено за виконання? (повідомлення інтерпретатора про помилку позначати error)\par
\begin{verbatim}
print(1.0 != 7)
\end{verbatim}

{\begin {center}\textbf{Завдання 5}\end {center}}

Що буде виведено за виконання? (повідомлення інтерпретатора про помилку позначати error)\par
\begin{verbatim}
print(4.0 <= 8 != 9  is  False)
\end{verbatim}

{\begin {center}\textbf{Завдання 6}\end {center}}

Що буде виведено за виконання? (повідомлення інтерпретатора про помилку позначати error)\par
\begin{verbatim}

def f(n):
    print("f", end="");
    return n

print(f(1) and f(7))
\end{verbatim}

{\begin {center}\textbf{Завдання 7}\end {center}}

Що буде виведено за виконання (повідомлення інтерпретатора про помилку позначати error)?\par
\begin{verbatim}
def f(c):
    u=54
    if c: 
        return 2
    elif c<=-2:
         u=1
    else:
         return 7
    return u

print(f(-9))
\end{verbatim}

{\begin {center}\textbf{Завдання 8}\end {center}}

Що буде виведено за виконання (повідомлення інтерпретатора про помилку позначати error)?\par
\begin{verbatim}
a, b, c = 8, 7, 6
def f(a, b, c):
    print(a, b, c, end=" ")

f(4, a=5)
print(a, b, c)
\end{verbatim}

{\begin {center}\textbf{Завдання 9}\end {center}}

Що буде виведено за виконання (повідомлення інтерпретатора про помилку позначати error)?\par
\begin{verbatim}
a,b,c=6,2,3
def g(a):
    a=4
    b+=5
    c=4
    return a+b+c

a,b,c=5,1,2
print(g(a),a,b,c)
\end{verbatim}

{\begin {center}\textbf{Завдання 10}\end {center}}


Написати функцiю f, що отримує щось та повертає щось. Невелика загадка все-таки має залишатися. Але нічого принципово нового відносно задач, що наявні у pdf-ках не планується.. Стиль запису умови також оцінюється.\par
\end {large}}
\newpage

{\Large{17.10.2024 Бета-версія: Контрольна робота \No 1, варіант  \No \textbf { 107 }\par}}
{
\begin {large}
{\begin {center}\textbf{Завдання 1}\end {center}}

Що буде виведено за виконання? (повідомлення інтерпретатора про помилку позначати error)\par
\begin{verbatim}
print(4/3//3*8)
\end{verbatim}

{\begin {center}\textbf{Завдання 2}\end {center}}

Що буде виведено за виконання? (повідомлення інтерпретатора про помилку позначати error)\par
\begin{verbatim}
print(4**2.0-False)
\end{verbatim}

{\begin {center}\textbf{Завдання 3}\end {center}}

Що буде виведено за виконання? (повідомлення інтерпретатора про помилку позначати error)\par
\begin{verbatim}
print(not 9>7 or 3==9.0 or 8.0<False)
\end{verbatim}

{\begin {center}\textbf{Завдання 4}\end {center}}

Що буде виведено за виконання? (повідомлення інтерпретатора про помилку позначати error)\par
\begin{verbatim}
print("False" <= "5.0")
\end{verbatim}

{\begin {center}\textbf{Завдання 5}\end {center}}

Що буде виведено за виконання? (повідомлення інтерпретатора про помилку позначати error)\par
\begin{verbatim}
print(2 < 2.0 >= 3 == 9.0)
\end{verbatim}

{\begin {center}\textbf{Завдання 6}\end {center}}

Що буде виведено за виконання? (повідомлення інтерпретатора про помилку позначати error)\par
\begin{verbatim}

def f(n):
    print("f", end="");
    return n<=-1

print(f(8) or f(-4))
\end{verbatim}

{\begin {center}\textbf{Завдання 7}\end {center}}

Що буде виведено за виконання (повідомлення інтерпретатора про помилку позначати error)?\par
\begin{verbatim}
def f(d):
    w=30
    if d<=3: 
        return 8
    if d>=4:
         w=4
    else:
         return 5
    return w

print(f(7))
\end{verbatim}

{\begin {center}\textbf{Завдання 8}\end {center}}

Що буде виведено за виконання (повідомлення інтерпретатора про помилку позначати error)?\par
\begin{verbatim}
a, b, c = 7, 6, 8
def g(a, b, c=9):
    print(a, b, c, end=" ")

g(3, c=0, b=2)
print(a, b, c)
\end{verbatim}

{\begin {center}\textbf{Завдання 9}\end {center}}

Що буде виведено за виконання (повідомлення інтерпретатора про помилку позначати error)?\par
\begin{verbatim}
a,b,c=9,6,8
def h(b):
    a=3
    b*=2
    c=1
    return a+b+c

a,b,c=4,0,7
print(h(b),a,b,c)
\end{verbatim}

{\begin {center}\textbf{Завдання 10}\end {center}}


Написати функцiю f, що отримує щось та повертає щось. Невелика загадка все-таки має залишатися. Але нічого принципово нового відносно задач, що наявні у pdf-ках не планується.. Стиль запису умови також оцінюється.\par
\end {large}}
\newpage

{\Large{17.10.2024 Бета-версія: Контрольна робота \No 1, варіант  \No \textbf { 108 }\par}}
{
\begin {large}
{\begin {center}\textbf{Завдання 1}\end {center}}

Що буде виведено за виконання? (повідомлення інтерпретатора про помилку позначати error)\par
\begin{verbatim}
print(6//6/5*4)
\end{verbatim}

{\begin {center}\textbf{Завдання 2}\end {center}}

Що буде виведено за виконання? (повідомлення інтерпретатора про помилку позначати error)\par
\begin{verbatim}
print(2**-2-1)
\end{verbatim}

{\begin {center}\textbf{Завдання 3}\end {center}}

Що буде виведено за виконання? (повідомлення інтерпретатора про помилку позначати error)\par
\begin{verbatim}
print(not 6.0<=9 and 4>=True and 7!=3.0)
\end{verbatim}

{\begin {center}\textbf{Завдання 4}\end {center}}

Що буде виведено за виконання? (повідомлення інтерпретатора про помилку позначати error)\par
\begin{verbatim}
print(True > "8")
\end{verbatim}

{\begin {center}\textbf{Завдання 5}\end {center}}

Що буде виведено за виконання? (повідомлення інтерпретатора про помилку позначати error)\par
\begin{verbatim}
print(7 > 6 < 8.0 == True)
\end{verbatim}

{\begin {center}\textbf{Завдання 6}\end {center}}

Що буде виведено за виконання? (повідомлення інтерпретатора про помилку позначати error)\par
\begin{verbatim}

def f(n):
    print("f", end="");
    return n

print(f(-9) and f(-1))
\end{verbatim}

{\begin {center}\textbf{Завдання 7}\end {center}}

Що буде виведено за виконання (повідомлення інтерпретатора про помилку позначати error)?\par
\begin{verbatim}
def g(b):
    y=96
    if b==-5: 
        y=3
    elif b<-4:
         y=6
    else:
         return 9
    return y

print(g(6))
\end{verbatim}

{\begin {center}\textbf{Завдання 8}\end {center}}

Що буде виведено за виконання (повідомлення інтерпретатора про помилку позначати error)?\par
\begin{verbatim}
a, b, c = 9, 7, 8
def h(a, b=6, c):
    print(a, b, c, end=" ")

h(a=1, 4, b=3)
print(a, b, c)
\end{verbatim}

{\begin {center}\textbf{Завдання 9}\end {center}}

Що буде виведено за виконання (повідомлення інтерпретатора про помилку позначати error)?\par
\begin{verbatim}
a,b,c=1,0,6
def f(b):
    a=4
    b-=5
    c=1
    return a+b+c

a,b,c=9,8,3
print(f(b),a,b,c)
\end{verbatim}

{\begin {center}\textbf{Завдання 10}\end {center}}


Написати функцiю f, що отримує щось та повертає щось. Невелика загадка все-таки має залишатися. Але нічого принципово нового відносно задач, що наявні у pdf-ках не планується.. Стиль запису умови також оцінюється.\par
\end {large}}
\newpage

{\Large{17.10.2024 Бета-версія: Контрольна робота \No 1, варіант  \No \textbf { 109 }\par}}
{
\begin {large}
{\begin {center}\textbf{Завдання 1}\end {center}}

Що буде виведено за виконання? (повідомлення інтерпретатора про помилку позначати error)\par
\begin{verbatim}
print(8//4%7*6)
\end{verbatim}

{\begin {center}\textbf{Завдання 2}\end {center}}

Що буде виведено за виконання? (повідомлення інтерпретатора про помилку позначати error)\par
\begin{verbatim}
print(4**0.5*-2)
\end{verbatim}

{\begin {center}\textbf{Завдання 3}\end {center}}

Що буде виведено за виконання? (повідомлення інтерпретатора про помилку позначати error)\par
\begin{verbatim}
print(not 5>=9.0 and True>6 or 2!=3.0)
\end{verbatim}

{\begin {center}\textbf{Завдання 4}\end {center}}

Що буде виведено за виконання? (повідомлення інтерпретатора про помилку позначати error)\par
\begin{verbatim}
print("7.0" > "3j")
\end{verbatim}

{\begin {center}\textbf{Завдання 5}\end {center}}

Що буде виведено за виконання? (повідомлення інтерпретатора про помилку позначати error)\par
\begin{verbatim}
print(9 != 6.0 > False >= 6)
\end{verbatim}

{\begin {center}\textbf{Завдання 6}\end {center}}

Що буде виведено за виконання? (повідомлення інтерпретатора про помилку позначати error)\par
\begin{verbatim}

def f(n):
    print("f", end="");
    return n!=1

print(f(3) and f(-7))
\end{verbatim}

{\begin {center}\textbf{Завдання 7}\end {center}}

Що буде виведено за виконання (повідомлення інтерпретатора про помилку позначати error)?\par
\begin{verbatim}
def f(a,b):
    c=74
    if a<=2:
        c=9
    elif b>=-5:
         return 0
    else: 
        c=8
    return c

print(f(-8,8))
\end{verbatim}

{\begin {center}\textbf{Завдання 8}\end {center}}

Що буде виведено за виконання (повідомлення інтерпретатора про помилку позначати error)?\par
\begin{verbatim}
a, b, c = 8, 7, 6
def h(a, b, c):
    print(a, b, c, end=" ")

h(2, 1, c=5)
print(a, b, c)
\end{verbatim}

{\begin {center}\textbf{Завдання 9}\end {center}}

Що буде виведено за виконання (повідомлення інтерпретатора про помилку позначати error)?\par
\begin{verbatim}
a,b,c=2,7,5
def g(a):
    a+=2
    b=3
    c=1
    return a+b+c

a,b,c=4,0,5
print(g(b),a,b,c)
\end{verbatim}

{\begin {center}\textbf{Завдання 10}\end {center}}


Написати функцiю f, що отримує щось та повертає щось. Невелика загадка все-таки має залишатися. Але нічого принципово нового відносно задач, що наявні у pdf-ках не планується.. Стиль запису умови також оцінюється.\par
\end {large}}
\newpage

{\Large{17.10.2024 Бета-версія: Контрольна робота \No 1, варіант  \No \textbf { 110 }\par}}
{
\begin {large}
{\begin {center}\textbf{Завдання 1}\end {center}}

Що буде виведено за виконання? (повідомлення інтерпретатора про помилку позначати error)\par
\begin{verbatim}
print(7/9//8*5)
\end{verbatim}

{\begin {center}\textbf{Завдання 2}\end {center}}

Що буде виведено за виконання? (повідомлення інтерпретатора про помилку позначати error)\par
\begin{verbatim}
print(1**-1.0*-1)
\end{verbatim}

{\begin {center}\textbf{Завдання 3}\end {center}}

Що буде виведено за виконання? (повідомлення інтерпретатора про помилку позначати error)\par
\begin{verbatim}
print(not False==5.0 or 8<=6 and 6.0<8)
\end{verbatim}

{\begin {center}\textbf{Завдання 4}\end {center}}

Що буде виведено за виконання? (повідомлення інтерпретатора про помилку позначати error)\par
\begin{verbatim}
print(False == "True")
\end{verbatim}

{\begin {center}\textbf{Завдання 5}\end {center}}

Що буде виведено за виконання? (повідомлення інтерпретатора про помилку позначати error)\par
\begin{verbatim}
print(7 <= 2.0  is  True <= 5)
\end{verbatim}

{\begin {center}\textbf{Завдання 6}\end {center}}

Що буде виведено за виконання? (повідомлення інтерпретатора про помилку позначати error)\par
\begin{verbatim}

def f(n):
    print("f", end="");
    return n

print(f(-8) or f(-3))
\end{verbatim}

{\begin {center}\textbf{Завдання 7}\end {center}}

Що буде виведено за виконання (повідомлення інтерпретатора про помилку позначати error)?\par
\begin{verbatim}
def h(d):
    z=84
    if d: 
        z=0
    elif d!=5:
         return 2
    else:
         z=8
    return z

print(h(2))
\end{verbatim}

{\begin {center}\textbf{Завдання 8}\end {center}}

Що буде виведено за виконання (повідомлення інтерпретатора про помилку позначати error)?\par
\begin{verbatim}
a, b, c = 9, 7, 8
def f(a, b, c=6):
    print(a, b, c, end=" ")

f(0, 0, 3)
print(a, b, c)
\end{verbatim}

{\begin {center}\textbf{Завдання 9}\end {center}}

Що буде виведено за виконання (повідомлення інтерпретатора про помилку позначати error)?\par
\begin{verbatim}
a,b,c=9,8,6
def h(a):
    a*=3
    b=5
    c=4
    return a+b+c

a,b,c=3,2,4
print(h(a),a,b,c)
\end{verbatim}

{\begin {center}\textbf{Завдання 10}\end {center}}


Написати функцiю f, що отримує щось та повертає щось. Невелика загадка все-таки має залишатися. Але нічого принципово нового відносно задач, що наявні у pdf-ках не планується.. Стиль запису умови також оцінюється.\par
\end {large}}
\newpage

{\Large{17.10.2024 Бета-версія: Контрольна робота \No 1, варіант  \No \textbf { 111 }\par}}
{
\begin {large}
{\begin {center}\textbf{Завдання 1}\end {center}}

Що буде виведено за виконання? (повідомлення інтерпретатора про помилку позначати error)\par
\begin{verbatim}
print(9/7/6*9)
\end{verbatim}

{\begin {center}\textbf{Завдання 2}\end {center}}

Що буде виведено за виконання? (повідомлення інтерпретатора про помилку позначати error)\par
\begin{verbatim}
print(9**1.0*True)
\end{verbatim}

{\begin {center}\textbf{Завдання 3}\end {center}}

Що буде виведено за виконання? (повідомлення інтерпретатора про помилку позначати error)\par
\begin{verbatim}
print(not 3<7.0 and 8.0<=4 or 5>=False)
\end{verbatim}

{\begin {center}\textbf{Завдання 4}\end {center}}

Що буде виведено за виконання? (повідомлення інтерпретатора про помилку позначати error)\par
\begin{verbatim}
print(1 != 5.0)
\end{verbatim}

{\begin {center}\textbf{Завдання 5}\end {center}}

Що буде виведено за виконання? (повідомлення інтерпретатора про помилку позначати error)\par
\begin{verbatim}
print(5.0 == 9.0 >= 4 != 8)
\end{verbatim}

{\begin {center}\textbf{Завдання 6}\end {center}}

Що буде виведено за виконання? (повідомлення інтерпретатора про помилку позначати error)\par
\begin{verbatim}

def f(n):
    print("f", end="");
    return n<=3

print(f(5) or f(-6))
\end{verbatim}

{\begin {center}\textbf{Завдання 7}\end {center}}

Що буде виведено за виконання (повідомлення інтерпретатора про помилку позначати error)?\par
\begin{verbatim}
def g(b):
    v=69
    if b>-3: 
        return 4
    if b!=-1:
         return 9
    else:
         v=0
    return v

print(g(0))
\end{verbatim}

{\begin {center}\textbf{Завдання 8}\end {center}}

Що буде виведено за виконання (повідомлення інтерпретатора про помилку позначати error)?\par
\begin{verbatim}
a, b, c = 9, 8, 6
def g(a, b=7, c=9):
    print(a, b, c, end=" ")

g(b=5, c=2, 4)
print(a, b, c)
\end{verbatim}

{\begin {center}\textbf{Завдання 9}\end {center}}

Що буде виведено за виконання (повідомлення інтерпретатора про помилку позначати error)?\par
\begin{verbatim}
a,b,c=7,1,3
def f(b):
    a-=2
    b=5
    c=4
    return a+b+c

a,b,c=0,7,1
print(f(a),a,b,c)
\end{verbatim}

{\begin {center}\textbf{Завдання 10}\end {center}}


Написати функцiю f, що отримує щось та повертає щось. Невелика загадка все-таки має залишатися. Але нічого принципово нового відносно задач, що наявні у pdf-ках не планується.. Стиль запису умови також оцінюється.\par
\end {large}}
\newpage

{\Large{17.10.2024 Бета-версія: Контрольна робота \No 1, варіант  \No \textbf { 112 }\par}}
{
\begin {large}
{\begin {center}\textbf{Завдання 1}\end {center}}

Що буде виведено за виконання? (повідомлення інтерпретатора про помилку позначати error)\par
\begin{verbatim}
print(5//5*9*3)
\end{verbatim}

{\begin {center}\textbf{Завдання 2}\end {center}}

Що буде виведено за виконання? (повідомлення інтерпретатора про помилку позначати error)\par
\begin{verbatim}
print(-3**-2+1)
\end{verbatim}

{\begin {center}\textbf{Завдання 3}\end {center}}

Що буде виведено за виконання? (повідомлення інтерпретатора про помилку позначати error)\par
\begin{verbatim}
print(2==True or not 7>4.0 and 9!=2.0)
\end{verbatim}

{\begin {center}\textbf{Завдання 4}\end {center}}

Що буде виведено за виконання? (повідомлення інтерпретатора про помилку позначати error)\par
\begin{verbatim}
print("7" >= "True")
\end{verbatim}

{\begin {center}\textbf{Завдання 5}\end {center}}

Що буде виведено за виконання? (повідомлення інтерпретатора про помилку позначати error)\par
\begin{verbatim}
print(8.0 < 3 > False  is  3.0)
\end{verbatim}

{\begin {center}\textbf{Завдання 6}\end {center}}

Що буде виведено за виконання? (повідомлення інтерпретатора про помилку позначати error)\par
\begin{verbatim}

def f(n):
    print("f", end="");
    return n

print(f(-2) and f(4))
\end{verbatim}

{\begin {center}\textbf{Завдання 7}\end {center}}

Що буде виведено за виконання (повідомлення інтерпретатора про помилку позначати error)?\par
\begin{verbatim}
def g(a,b):
    c=51
    if b:
        return 4
    if a<-3:
         c=3
    else: 
        return 6
    return c

print(g(8,0))
\end{verbatim}

{\begin {center}\textbf{Завдання 8}\end {center}}

Що буде виведено за виконання (повідомлення інтерпретатора про помилку позначати error)?\par
\begin{verbatim}
a, b, c = 8, 6, 9
def f(a, b=7, c):
    print(a, b, c, end=" ")

f(b=1, a=2, c=0)
print(a, b, c)
\end{verbatim}

{\begin {center}\textbf{Завдання 9}\end {center}}

Що буде виведено за виконання (повідомлення інтерпретатора про помилку позначати error)?\par
\begin{verbatim}
a,b,c=9,8,6
def g(a):
    a=1
    b=3
    c=4
    return a+b+c

a,b,c=3,1,7
print(g(b),a,b,c)
\end{verbatim}

{\begin {center}\textbf{Завдання 10}\end {center}}


Написати функцiю f, що отримує щось та повертає щось. Невелика загадка все-таки має залишатися. Але нічого принципово нового відносно задач, що наявні у pdf-ках не планується.. Стиль запису умови також оцінюється.\par
\end {large}}
\newpage

{\Large{17.10.2024 Бета-версія: Контрольна робота \No 1, варіант  \No \textbf { 113 }\par}}
{
\begin {large}
{\begin {center}\textbf{Завдання 1}\end {center}}

Що буде виведено за виконання? (повідомлення інтерпретатора про помилку позначати error)\par
\begin{verbatim}
print(9/7/9*5)
\end{verbatim}

{\begin {center}\textbf{Завдання 2}\end {center}}

Що буде виведено за виконання? (повідомлення інтерпретатора про помилку позначати error)\par
\begin{verbatim}
print(4**-1.0*2)
\end{verbatim}

{\begin {center}\textbf{Завдання 3}\end {center}}

Що буде виведено за виконання? (повідомлення інтерпретатора про помилку позначати error)\par
\begin{verbatim}
print(not True>=7.0 or 5<=8 or 3.0!=3)
\end{verbatim}

{\begin {center}\textbf{Завдання 4}\end {center}}

Що буде виведено за виконання? (повідомлення інтерпретатора про помилку позначати error)\par
\begin{verbatim}
print(6 < 1.0)
\end{verbatim}

{\begin {center}\textbf{Завдання 5}\end {center}}

Що буде виведено за виконання? (повідомлення інтерпретатора про помилку позначати error)\par
\begin{verbatim}
print(2 >= 8 > 4.0 < 6)
\end{verbatim}

{\begin {center}\textbf{Завдання 6}\end {center}}

Що буде виведено за виконання? (повідомлення інтерпретатора про помилку позначати error)\par
\begin{verbatim}

def f(n):
    print("f", end="");
    return n>=0

print(f(9) and f(-5))
\end{verbatim}

{\begin {center}\textbf{Завдання 7}\end {center}}

Що буде виведено за виконання (повідомлення інтерпретатора про помилку позначати error)?\par
\begin{verbatim}
def h(a):
    u=29
    if a: 
        u=1
    elif a>=2:
         return 7
    else:
         u=6
    return u

print(h(-8))
\end{verbatim}

{\begin {center}\textbf{Завдання 8}\end {center}}

Що буде виведено за виконання (повідомлення інтерпретатора про помилку позначати error)?\par
\begin{verbatim}
a, b, c = 7, 6, 9
def h(a, b, c):
    print(a, b, c, end=" ")

h(4, b=5)
print(a, b, c)
\end{verbatim}

{\begin {center}\textbf{Завдання 9}\end {center}}

Що буде виведено за виконання (повідомлення інтерпретатора про помилку позначати error)?\par
\begin{verbatim}
a,b,c=6,4,9
def h(b):
    a-=2
    b=1
    c=3
    return a+b+c

a,b,c=8,5,2
print(h(a),a,b,c)
\end{verbatim}

{\begin {center}\textbf{Завдання 10}\end {center}}


Написати функцiю f, що отримує щось та повертає щось. Невелика загадка все-таки має залишатися. Але нічого принципово нового відносно задач, що наявні у pdf-ках не планується.. Стиль запису умови також оцінюється.\par
\end {large}}
\newpage

{\Large{17.10.2024 Бета-версія: Контрольна робота \No 1, варіант  \No \textbf { 114 }\par}}
{
\begin {large}
{\begin {center}\textbf{Завдання 1}\end {center}}

Що буде виведено за виконання? (повідомлення інтерпретатора про помилку позначати error)\par
\begin{verbatim}
print(7//3%8*4)
\end{verbatim}

{\begin {center}\textbf{Завдання 2}\end {center}}

Що буде виведено за виконання? (повідомлення інтерпретатора про помилку позначати error)\par
\begin{verbatim}
print(-2**1.0-False)
\end{verbatim}

{\begin {center}\textbf{Завдання 3}\end {center}}

Що буде виведено за виконання? (повідомлення інтерпретатора про помилку позначати error)\par
\begin{verbatim}
print(not 7<False and 6==5.0 and 9.0>4)
\end{verbatim}

{\begin {center}\textbf{Завдання 4}\end {center}}

Що буде виведено за виконання? (повідомлення інтерпретатора про помилку позначати error)\par
\begin{verbatim}
print(False <= "4.0j")
\end{verbatim}

{\begin {center}\textbf{Завдання 5}\end {center}}

Що буде виведено за виконання? (повідомлення інтерпретатора про помилку позначати error)\par
\begin{verbatim}
print(7.0 == 5 != True  is  2)
\end{verbatim}

{\begin {center}\textbf{Завдання 6}\end {center}}

Що буде виведено за виконання? (повідомлення інтерпретатора про помилку позначати error)\par
\begin{verbatim}

def f(n):
    print("f", end="");
    return n

print(f(0) or f(2))
\end{verbatim}

{\begin {center}\textbf{Завдання 7}\end {center}}

Що буде виведено за виконання (повідомлення інтерпретатора про помилку позначати error)?\par
\begin{verbatim}
def h(a,b):
    c=80
    if a<5:
        return 4
    elif b<=-1:
         c=5
    else: 
        c=0
    return c

print(h(6,6))
\end{verbatim}

{\begin {center}\textbf{Завдання 8}\end {center}}

Що буде виведено за виконання (повідомлення інтерпретатора про помилку позначати error)?\par
\begin{verbatim}
a, b, c = 8, 7, 8
def g(a, b=9, c):
    print(a, b, c, end=" ")

g(1, 3, a=0)
print(a, b, c)
\end{verbatim}

{\begin {center}\textbf{Завдання 9}\end {center}}

Що буде виведено за виконання (повідомлення інтерпретатора про помилку позначати error)?\par
\begin{verbatim}
a,b,c=2,5,5
def f(b):
    a=2
    b=4
    c=1
    return a+b+c

a,b,c=0,3,6
print(f(a),a,b,c)
\end{verbatim}

{\begin {center}\textbf{Завдання 10}\end {center}}


Написати функцiю f, що отримує щось та повертає щось. Невелика загадка все-таки має залишатися. Але нічого принципово нового відносно задач, що наявні у pdf-ках не планується.. Стиль запису умови також оцінюється.\par
\end {large}}
\newpage

{\Large{17.10.2024 Бета-версія: Контрольна робота \No 1, варіант  \No \textbf { 115 }\par}}
{
\begin {large}
{\begin {center}\textbf{Завдання 1}\end {center}}

Що буде виведено за виконання? (повідомлення інтерпретатора про помилку позначати error)\par
\begin{verbatim}
print(6/4//6*7)
\end{verbatim}

{\begin {center}\textbf{Завдання 2}\end {center}}

Що буде виведено за виконання? (повідомлення інтерпретатора про помилку позначати error)\par
\begin{verbatim}
print(-1**-4+-1)
\end{verbatim}

{\begin {center}\textbf{Завдання 3}\end {center}}

Що буде виведено за виконання? (повідомлення інтерпретатора про помилку позначати error)\par
\begin{verbatim}
print(False>=9 or not 4.0>6.0 and 2==3)
\end{verbatim}

{\begin {center}\textbf{Завдання 4}\end {center}}

Що буде виведено за виконання? (повідомлення інтерпретатора про помилку позначати error)\par
\begin{verbatim}
print("False" >= True)
\end{verbatim}

{\begin {center}\textbf{Завдання 5}\end {center}}

Що буде виведено за виконання? (повідомлення інтерпретатора про помилку позначати error)\par
\begin{verbatim}
print(True <= 9.0 < 4 >= 7)
\end{verbatim}

{\begin {center}\textbf{Завдання 6}\end {center}}

Що буде виведено за виконання? (повідомлення інтерпретатора про помилку позначати error)\par
\begin{verbatim}

def f(n):
    print("f", end="");
    return n>2

print(f(-5) and f(5))
\end{verbatim}

{\begin {center}\textbf{Завдання 7}\end {center}}

Що буде виведено за виконання (повідомлення інтерпретатора про помилку позначати error)?\par
\begin{verbatim}
def g(a,b):
    c=13
    if a:
        return 1
    if a>-4:
         c=9
    else: 
        return 7
    return c

print(g(-9,3))
\end{verbatim}

{\begin {center}\textbf{Завдання 8}\end {center}}

Що буде виведено за виконання (повідомлення інтерпретатора про помилку позначати error)?\par
\begin{verbatim}
a, b, c = 6, 8, 9
def f(a, b=7, c=6):
    print(a, b, c, end=" ")

f(a=3, 1, c=2)
print(a, b, c)
\end{verbatim}

{\begin {center}\textbf{Завдання 9}\end {center}}

Що буде виведено за виконання (повідомлення інтерпретатора про помилку позначати error)?\par
\begin{verbatim}
a,b,c=2,8,4
def h(a):
    a=5
    b=2
    c=3
    return a+b+c

a,b,c=1,9,7
print(h(b),a,b,c)
\end{verbatim}

{\begin {center}\textbf{Завдання 10}\end {center}}


Написати функцiю f, що отримує щось та повертає щось. Невелика загадка все-таки має залишатися. Але нічого принципово нового відносно задач, що наявні у pdf-ках не планується.. Стиль запису умови також оцінюється.\par
\end {large}}
\newpage

{\Large{17.10.2024 Бета-версія: Контрольна робота \No 1, варіант  \No \textbf { 116 }\par}}
{
\begin {large}
{\begin {center}\textbf{Завдання 1}\end {center}}

Що буде виведено за виконання? (повідомлення інтерпретатора про помилку позначати error)\par
\begin{verbatim}
print(3//8*7*8)
\end{verbatim}

{\begin {center}\textbf{Завдання 2}\end {center}}

Що буде виведено за виконання? (повідомлення інтерпретатора про помилку позначати error)\par
\begin{verbatim}
print(3**1.0+1)
\end{verbatim}

{\begin {center}\textbf{Завдання 3}\end {center}}

Що буде виведено за виконання? (повідомлення інтерпретатора про помилку позначати error)\par
\begin{verbatim}
print(2.0<=4 and not 8.0!=9 or 5<True)
\end{verbatim}

{\begin {center}\textbf{Завдання 4}\end {center}}

Що буде виведено за виконання? (повідомлення інтерпретатора про помилку позначати error)\par
\begin{verbatim}
print(8 == "9.0")
\end{verbatim}

{\begin {center}\textbf{Завдання 5}\end {center}}

Що буде виведено за виконання? (повідомлення інтерпретатора про помилку позначати error)\par
\begin{verbatim}
print(8.0 <= 3 != False == 9)
\end{verbatim}

{\begin {center}\textbf{Завдання 6}\end {center}}

Що буде виведено за виконання? (повідомлення інтерпретатора про помилку позначати error)\par
\begin{verbatim}

def f(n):
    print("f", end="");
    return n

print(f(3) or f(1))
\end{verbatim}

{\begin {center}\textbf{Завдання 7}\end {center}}

Що буде виведено за виконання (повідомлення інтерпретатора про помилку позначати error)?\par
\begin{verbatim}
def f(c):
    x=28
    if c<1: 
        return 3
    if c==4:
         x=5
    else:
         return 2
    return x

print(f(-7))
\end{verbatim}

{\begin {center}\textbf{Завдання 8}\end {center}}

Що буде виведено за виконання (повідомлення інтерпретатора про помилку позначати error)?\par
\begin{verbatim}
a, b, c = 6, 7, 8
def h(a, b, c=9):
    print(a, b, c, end=" ")

h(4, 5)
print(a, b, c)
\end{verbatim}

{\begin {center}\textbf{Завдання 9}\end {center}}

Що буде виведено за виконання (повідомлення інтерпретатора про помилку позначати error)?\par
\begin{verbatim}
a,b,c=1,6,5
def g(a):
    a+=2
    b=4
    c=3
    return a+b+c

a,b,c=8,2,4
print(g(b),a,b,c)
\end{verbatim}

{\begin {center}\textbf{Завдання 10}\end {center}}


Написати функцiю f, що отримує щось та повертає щось. Невелика загадка все-таки має залишатися. Але нічого принципово нового відносно задач, що наявні у pdf-ках не планується.. Стиль запису умови також оцінюється.\par
\end {large}}
\newpage

{\Large{17.10.2024 Бета-версія: Контрольна робота \No 1, варіант  \No \textbf { 117 }\par}}
{
\begin {large}
{\begin {center}\textbf{Завдання 1}\end {center}}

Що буде виведено за виконання? (повідомлення інтерпретатора про помилку позначати error)\par
\begin{verbatim}
print(5//6*5*9)
\end{verbatim}

{\begin {center}\textbf{Завдання 2}\end {center}}

Що буде виведено за виконання? (повідомлення інтерпретатора про помилку позначати error)\par
\begin{verbatim}
print(1**0.5*True)
\end{verbatim}

{\begin {center}\textbf{Завдання 3}\end {center}}

Що буде виведено за виконання? (повідомлення інтерпретатора про помилку позначати error)\par
\begin{verbatim}
print(2!=7 or not 9.0<7.0 and 6<=False)
\end{verbatim}

{\begin {center}\textbf{Завдання 4}\end {center}}

Що буде виведено за виконання? (повідомлення інтерпретатора про помилку позначати error)\par
\begin{verbatim}
print(3.0 <= "4j")
\end{verbatim}

{\begin {center}\textbf{Завдання 5}\end {center}}

Що буде виведено за виконання? (повідомлення інтерпретатора про помилку позначати error)\par
\begin{verbatim}
print(3.0 > 5.0  is  5 < 7)
\end{verbatim}

{\begin {center}\textbf{Завдання 6}\end {center}}

Що буде виведено за виконання? (повідомлення інтерпретатора про помилку позначати error)\par
\begin{verbatim}

def f(n):
    print("f", end="");
    return n==-1

print(f(-8) and f(8))
\end{verbatim}

{\begin {center}\textbf{Завдання 7}\end {center}}

Що буде виведено за виконання (повідомлення інтерпретатора про помилку позначати error)?\par
\begin{verbatim}
def f(a,b):
    c=35
    if b==3:
        return 8
    elif b>=0:
         c=2
    else: 
        return 5
    return c

print(f(0,7))
\end{verbatim}

{\begin {center}\textbf{Завдання 8}\end {center}}

Що буде виведено за виконання (повідомлення інтерпретатора про помилку позначати error)?\par
\begin{verbatim}
a, b, c = 9, 6, 7
def g(a, b=8, c=7):
    print(a, b, c, end=" ")

g(b=1, c=4, 2)
print(a, b, c)
\end{verbatim}

{\begin {center}\textbf{Завдання 9}\end {center}}

Що буде виведено за виконання (повідомлення інтерпретатора про помилку позначати error)?\par
\begin{verbatim}
a,b,c=3,0,9
def f(b):
    a=5
    b*=1
    c=3
    return a+b+c

a,b,c=7,0,3
print(f(a),a,b,c)
\end{verbatim}

{\begin {center}\textbf{Завдання 10}\end {center}}


Написати функцiю f, що отримує щось та повертає щось. Невелика загадка все-таки має залишатися. Але нічого принципово нового відносно задач, що наявні у pdf-ках не планується.. Стиль запису умови також оцінюється.\par
\end {large}}
\newpage

{\Large{17.10.2024 Бета-версія: Контрольна робота \No 1, варіант  \No \textbf { 118 }\par}}
{
\begin {large}
{\begin {center}\textbf{Завдання 1}\end {center}}

Що буде виведено за виконання? (повідомлення інтерпретатора про помилку позначати error)\par
\begin{verbatim}
print(8/5//4*6)
\end{verbatim}

{\begin {center}\textbf{Завдання 2}\end {center}}

Що буде виведено за виконання? (повідомлення інтерпретатора про помилку позначати error)\par
\begin{verbatim}
print(9**-0.5*2)
\end{verbatim}

{\begin {center}\textbf{Завдання 3}\end {center}}

Що буде виведено за виконання? (повідомлення інтерпретатора про помилку позначати error)\par
\begin{verbatim}
print(not 8.0>8 and True==2.0 or 3>=5)
\end{verbatim}

{\begin {center}\textbf{Завдання 4}\end {center}}

Що буде виведено за виконання? (повідомлення інтерпретатора про помилку позначати error)\par
\begin{verbatim}
print("8.0j" < 5)
\end{verbatim}

{\begin {center}\textbf{Завдання 5}\end {center}}

Що буде виведено за виконання? (повідомлення інтерпретатора про помилку позначати error)\par
\begin{verbatim}
print(8 <= True  is  2 != 4)
\end{verbatim}

{\begin {center}\textbf{Завдання 6}\end {center}}

Що буде виведено за виконання? (повідомлення інтерпретатора про помилку позначати error)\par
\begin{verbatim}

def f(n):
    print("f", end="");
    return n

print(f(4) or f(0))
\end{verbatim}

{\begin {center}\textbf{Завдання 7}\end {center}}

Що буде виведено за виконання (повідомлення інтерпретатора про помилку позначати error)?\par
\begin{verbatim}
def h(a,b):
    c=50
    if b!=-2:
        c=3
    if a!=4:
         c=6
    else: 
        return 1
    return c

print(h(-7,3))
\end{verbatim}

{\begin {center}\textbf{Завдання 8}\end {center}}

Що буде виведено за виконання (повідомлення інтерпретатора про помилку позначати error)?\par
\begin{verbatim}
a, b, c = 9, 6, 8
def f(a, b=6, c):
    print(a, b, c, end=" ")

f(a=5, b=3, c=0)
print(a, b, c)
\end{verbatim}

{\begin {center}\textbf{Завдання 9}\end {center}}

Що буде виведено за виконання (повідомлення інтерпретатора про помилку позначати error)?\par
\begin{verbatim}
a,b,c=6,8,9
def f(a):
    a-=4
    b=2
    c=5
    return a+b+c

a,b,c=2,4,7
print(f(b),a,b,c)
\end{verbatim}

{\begin {center}\textbf{Завдання 10}\end {center}}


Написати функцiю f, що отримує щось та повертає щось. Невелика загадка все-таки має залишатися. Але нічого принципово нового відносно задач, що наявні у pdf-ках не планується.. Стиль запису умови також оцінюється.\par
\end {large}}
\newpage

{\Large{17.10.2024 Бета-версія: Контрольна робота \No 1, варіант  \No \textbf { 119 }\par}}
{
\begin {large}
{\begin {center}\textbf{Завдання 1}\end {center}}

Що буде виведено за виконання? (повідомлення інтерпретатора про помилку позначати error)\par
\begin{verbatim}
print(4/9/3*3)
\end{verbatim}

{\begin {center}\textbf{Завдання 2}\end {center}}

Що буде виведено за виконання? (повідомлення інтерпретатора про помилку позначати error)\par
\begin{verbatim}
print(4**-1--2)
\end{verbatim}

{\begin {center}\textbf{Завдання 3}\end {center}}

Що буде виведено за виконання? (повідомлення інтерпретатора про помилку позначати error)\par
\begin{verbatim}
print(not 7!=3.0 or 6.0<True or 2==4)
\end{verbatim}

{\begin {center}\textbf{Завдання 4}\end {center}}

Що буде виведено за виконання? (повідомлення інтерпретатора про помилку позначати error)\par
\begin{verbatim}
print(True != "9")
\end{verbatim}

{\begin {center}\textbf{Завдання 5}\end {center}}

Що буде виведено за виконання? (повідомлення інтерпретатора про помилку позначати error)\par
\begin{verbatim}
print(2.0 > 6.0 >= 9 == 6)
\end{verbatim}

{\begin {center}\textbf{Завдання 6}\end {center}}

Що буде виведено за виконання? (повідомлення інтерпретатора про помилку позначати error)\par
\begin{verbatim}

def f(n):
    print("f", end="");
    return n<4

print(f(-6) and f(-2))
\end{verbatim}

{\begin {center}\textbf{Завдання 7}\end {center}}

Що буде виведено за виконання (повідомлення інтерпретатора про помилку позначати error)?\par
\begin{verbatim}
def f(a,b):
    c=38
    if a:
        c=2
    elif b>2:
         return 8
    else: 
        return 7
    return c

print(f(0,-4))
\end{verbatim}

{\begin {center}\textbf{Завдання 8}\end {center}}

Що буде виведено за виконання (повідомлення інтерпретатора про помилку позначати error)?\par
\begin{verbatim}
a, b, c = 7, 8, 9
def h(a, b, c=6):
    print(a, b, c, end=" ")

h(5, 4, 3)
print(a, b, c)
\end{verbatim}

{\begin {center}\textbf{Завдання 9}\end {center}}

Що буде виведено за виконання (повідомлення інтерпретатора про помилку позначати error)?\par
\begin{verbatim}
a,b,c=5,1,7
def h(a):
    a*=1
    b=1
    c=2
    return a+b+c

a,b,c=8,0,6
print(h(a),a,b,c)
\end{verbatim}

{\begin {center}\textbf{Завдання 10}\end {center}}


Написати функцiю f, що отримує щось та повертає щось. Невелика загадка все-таки має залишатися. Але нічого принципово нового відносно задач, що наявні у pdf-ках не планується.. Стиль запису умови також оцінюється.\par
\end {large}}
\newpage

{\Large{17.10.2024 Бета-версія: Контрольна робота \No 1, варіант  \No \textbf { 120 }\par}}
{
\begin {large}
{\begin {center}\textbf{Завдання 1}\end {center}}

Що буде виведено за виконання? (повідомлення інтерпретатора про помилку позначати error)\par
\begin{verbatim}
print(3//5%4*3)
\end{verbatim}

{\begin {center}\textbf{Завдання 2}\end {center}}

Що буде виведено за виконання? (повідомлення інтерпретатора про помилку позначати error)\par
\begin{verbatim}
print(4**1.0*False)
\end{verbatim}

{\begin {center}\textbf{Завдання 3}\end {center}}

Що буде виведено за виконання? (повідомлення інтерпретатора про помилку позначати error)\par
\begin{verbatim}
print(5.0>=6 and not 4.0>9 and False<=8)
\end{verbatim}

{\begin {center}\textbf{Завдання 4}\end {center}}

Що буде виведено за виконання? (повідомлення інтерпретатора про помилку позначати error)\par
\begin{verbatim}
print("False" > 6.0)
\end{verbatim}

{\begin {center}\textbf{Завдання 5}\end {center}}

Що буде виведено за виконання? (повідомлення інтерпретатора про помилку позначати error)\par
\begin{verbatim}
print(3 >= 4.0  is  7.0 < False)
\end{verbatim}

{\begin {center}\textbf{Завдання 6}\end {center}}

Що буде виведено за виконання? (повідомлення інтерпретатора про помилку позначати error)\par
\begin{verbatim}

def f(n):
    print("f", end="");
    return n

print(f(-3) or f(-9))
\end{verbatim}

{\begin {center}\textbf{Завдання 7}\end {center}}

Що буде виведено за виконання (повідомлення інтерпретатора про помилку позначати error)?\par
\begin{verbatim}
def g(a,b):
    c=31
    if b==1:
        c=9
    if a<-4:
         c=0
    else: 
        return 4
    return c

print(g(8,0))
\end{verbatim}

{\begin {center}\textbf{Завдання 8}\end {center}}

Що буде виведено за виконання (повідомлення інтерпретатора про помилку позначати error)?\par
\begin{verbatim}
a, b, c = 7, 8, 9
def g(a, b, c):
    print(a, b, c, end=" ")

g(0, c=2, b=1)
print(a, b, c)
\end{verbatim}

{\begin {center}\textbf{Завдання 9}\end {center}}

Що буде виведено за виконання (повідомлення інтерпретатора про помилку позначати error)?\par
\begin{verbatim}
a,b,c=2,3,5
def g(b):
    a=4
    b+=3
    c=5
    return a+b+c

a,b,c=1,4,9
print(g(b),a,b,c)
\end{verbatim}

{\begin {center}\textbf{Завдання 10}\end {center}}


Написати функцiю f, що отримує щось та повертає щось. Невелика загадка все-таки має залишатися. Але нічого принципово нового відносно задач, що наявні у pdf-ках не планується.. Стиль запису умови також оцінюється.\par
\end {large}}
\newpage

{\Large{17.10.2024 Бета-версія: Контрольна робота \No 1, варіант  \No \textbf { 121 }\par}}
{
\begin {large}
{\begin {center}\textbf{Завдання 1}\end {center}}

Що буде виведено за виконання? (повідомлення інтерпретатора про помилку позначати error)\par
\begin{verbatim}
print(8/3//5*7)
\end{verbatim}

{\begin {center}\textbf{Завдання 2}\end {center}}

Що буде виведено за виконання? (повідомлення інтерпретатора про помилку позначати error)\par
\begin{verbatim}
print(9**0.5*-1)
\end{verbatim}

{\begin {center}\textbf{Завдання 3}\end {center}}

Що буде виведено за виконання? (повідомлення інтерпретатора про помилку позначати error)\par
\begin{verbatim}
print(not 6<=5 and 3.0>6.0 and 2==True)
\end{verbatim}

{\begin {center}\textbf{Завдання 4}\end {center}}

Що буде виведено за виконання? (повідомлення інтерпретатора про помилку позначати error)\par
\begin{verbatim}
print("2.0j" <= 6.0)
\end{verbatim}

{\begin {center}\textbf{Завдання 5}\end {center}}

Що буде виведено за виконання? (повідомлення інтерпретатора про помилку позначати error)\par
\begin{verbatim}
print(3 == 2.0 > 9 <= 8)
\end{verbatim}

{\begin {center}\textbf{Завдання 6}\end {center}}

Що буде виведено за виконання? (повідомлення інтерпретатора про помилку позначати error)\par
\begin{verbatim}

def f(n):
    print("f", end="");
    return n

print(f(-1) and f(7))
\end{verbatim}

{\begin {center}\textbf{Завдання 7}\end {center}}

Що буде виведено за виконання (повідомлення інтерпретатора про помилку позначати error)?\par
\begin{verbatim}
def g(d):
    y=99
    if d: 
        y=8
    if d<=-2:
         return 9
    else:
         return 2
    return y

print(g(-5))
\end{verbatim}

{\begin {center}\textbf{Завдання 8}\end {center}}

Що буде виведено за виконання (повідомлення інтерпретатора про помилку позначати error)?\par
\begin{verbatim}
a, b, c = 6, 9, 7
def h(a, b, c=8):
    print(a, b, c, end=" ")

h(1, 5)
print(a, b, c)
\end{verbatim}

{\begin {center}\textbf{Завдання 9}\end {center}}

Що буде виведено за виконання (повідомлення інтерпретатора про помилку позначати error)?\par
\begin{verbatim}
a,b,c=6,3,5
def g(a):
    a+=3
    b=4
    c=5
    return a+b+c

a,b,c=8,2,7
print(g(a),a,b,c)
\end{verbatim}

{\begin {center}\textbf{Завдання 10}\end {center}}


Написати функцiю f, що отримує щось та повертає щось. Невелика загадка все-таки має залишатися. Але нічого принципово нового відносно задач, що наявні у pdf-ках не планується.. Стиль запису умови також оцінюється.\par
\end {large}}
\newpage

{\Large{17.10.2024 Бета-версія: Контрольна робота \No 1, варіант  \No \textbf { 122 }\par}}
{
\begin {large}
{\begin {center}\textbf{Завдання 1}\end {center}}

Що буде виведено за виконання? (повідомлення інтерпретатора про помилку позначати error)\par
\begin{verbatim}
print(5//6%7*6)
\end{verbatim}

{\begin {center}\textbf{Завдання 2}\end {center}}

Що буде виведено за виконання? (повідомлення інтерпретатора про помилку позначати error)\par
\begin{verbatim}
print(4**-0.5*2)
\end{verbatim}

{\begin {center}\textbf{Завдання 3}\end {center}}

Що буде виведено за виконання? (повідомлення інтерпретатора про помилку позначати error)\par
\begin{verbatim}
print(not False>=3 or 8<9.0 or 7!=4.0)
\end{verbatim}

{\begin {center}\textbf{Завдання 4}\end {center}}

Що буде виведено за виконання? (повідомлення інтерпретатора про помилку позначати error)\par
\begin{verbatim}
print("False" < 2)
\end{verbatim}

{\begin {center}\textbf{Завдання 5}\end {center}}

Що буде виведено за виконання? (повідомлення інтерпретатора про помилку позначати error)\par
\begin{verbatim}
print(9.0 != False != 7 == 6)
\end{verbatim}

{\begin {center}\textbf{Завдання 6}\end {center}}

Що буде виведено за виконання? (повідомлення інтерпретатора про помилку позначати error)\par
\begin{verbatim}

def f(n):
    print("f", end="");
    return n>-4

print(f(-4) or f(2))
\end{verbatim}

{\begin {center}\textbf{Завдання 7}\end {center}}

Що буде виведено за виконання (повідомлення інтерпретатора про помилку позначати error)?\par
\begin{verbatim}
def h(a):
    z=62
    if a>2: 
        z=1
    elif a!=-5:
         z=6
    else:
         return 4
    return z

print(h(5))
\end{verbatim}

{\begin {center}\textbf{Завдання 8}\end {center}}

Що буде виведено за виконання (повідомлення інтерпретатора про помилку позначати error)?\par
\begin{verbatim}
a, b, c = 6, 9, 8
def g(a, b, c):
    print(a, b, c, end=" ")

g(4, c=0, b=2)
print(a, b, c)
\end{verbatim}

{\begin {center}\textbf{Завдання 9}\end {center}}

Що буде виведено за виконання (повідомлення інтерпретатора про помилку позначати error)?\par
\begin{verbatim}
a,b,c=9,0,4
def h(b):
    a=1
    b*=2
    c=1
    return a+b+c

a,b,c=1,9,8
print(h(b),a,b,c)
\end{verbatim}

{\begin {center}\textbf{Завдання 10}\end {center}}


Написати функцiю f, що отримує щось та повертає щось. Невелика загадка все-таки має залишатися. Але нічого принципово нового відносно задач, що наявні у pdf-ках не планується.. Стиль запису умови також оцінюється.\par
\end {large}}
\newpage

{\Large{17.10.2024 Бета-версія: Контрольна робота \No 1, варіант  \No \textbf { 123 }\par}}
{
\begin {large}
{\begin {center}\textbf{Завдання 1}\end {center}}

Що буде виведено за виконання? (повідомлення інтерпретатора про помилку позначати error)\par
\begin{verbatim}
print(7/7*8*5)
\end{verbatim}

{\begin {center}\textbf{Завдання 2}\end {center}}

Що буде виведено за виконання? (повідомлення інтерпретатора про помилку позначати error)\par
\begin{verbatim}
print(4**2.0+True)
\end{verbatim}

{\begin {center}\textbf{Завдання 3}\end {center}}

Що буде виведено за виконання? (повідомлення інтерпретатора про помилку позначати error)\par
\begin{verbatim}
print(8.0==9 and not 4<True or 2.0!=5)
\end{verbatim}

{\begin {center}\textbf{Завдання 4}\end {center}}

Що буде виведено за виконання? (повідомлення інтерпретатора про помилку позначати error)\par
\begin{verbatim}
print(True >= "1j")
\end{verbatim}

{\begin {center}\textbf{Завдання 5}\end {center}}

Що буде виведено за виконання? (повідомлення інтерпретатора про помилку позначати error)\par
\begin{verbatim}
print(3.0 < True  is  4 >= 5.0)
\end{verbatim}

{\begin {center}\textbf{Завдання 6}\end {center}}

Що буде виведено за виконання? (повідомлення інтерпретатора про помилку позначати error)\par
\begin{verbatim}

def f(n):
    print("f", end="");
    return n==-3

print(f(-7) and f(6))
\end{verbatim}

{\begin {center}\textbf{Завдання 7}\end {center}}

Що буде виведено за виконання (повідомлення інтерпретатора про помилку позначати error)?\par
\begin{verbatim}
def g(a,b):
    c=82
    if a>=-3:
        c=2
    elif b<=3:
         return 3
    else: 
        c=5
    return c

print(g(4,6))
\end{verbatim}

{\begin {center}\textbf{Завдання 8}\end {center}}

Що буде виведено за виконання (повідомлення інтерпретатора про помилку позначати error)?\par
\begin{verbatim}
a, b, c = 7, 7, 8
def f(a, b=9, c=6):
    print(a, b, c, end=" ")

f(3, b=1)
print(a, b, c)
\end{verbatim}

{\begin {center}\textbf{Завдання 9}\end {center}}

Що буде виведено за виконання (повідомлення інтерпретатора про помилку позначати error)?\par
\begin{verbatim}
a,b,c=0,8,1
def f(a):
    a=2
    b=1
    c=5
    return a+b+c

a,b,c=9,2,3
print(f(b),a,b,c)
\end{verbatim}

{\begin {center}\textbf{Завдання 10}\end {center}}


Написати функцiю f, що отримує щось та повертає щось. Невелика загадка все-таки має залишатися. Але нічого принципово нового відносно задач, що наявні у pdf-ках не планується.. Стиль запису умови також оцінюється.\par
\end {large}}
\newpage

{\Large{17.10.2024 Бета-версія: Контрольна робота \No 1, варіант  \No \textbf { 124 }\par}}
{
\begin {large}
{\begin {center}\textbf{Завдання 1}\end {center}}

Що буде виведено за виконання? (повідомлення інтерпретатора про помилку позначати error)\par
\begin{verbatim}
print(6//8/3*4)
\end{verbatim}

{\begin {center}\textbf{Завдання 2}\end {center}}

Що буде виведено за виконання? (повідомлення інтерпретатора про помилку позначати error)\par
\begin{verbatim}
print(1**1.0*False)
\end{verbatim}

{\begin {center}\textbf{Завдання 3}\end {center}}

Що буде виведено за виконання? (повідомлення інтерпретатора про помилку позначати error)\par
\begin{verbatim}
print(3>=5.0 or not False<=6 and 4>7.0)
\end{verbatim}

{\begin {center}\textbf{Завдання 4}\end {center}}

Що буде виведено за виконання? (повідомлення інтерпретатора про помилку позначати error)\par
\begin{verbatim}
print(3 > False)
\end{verbatim}

{\begin {center}\textbf{Завдання 5}\end {center}}

Що буде виведено за виконання? (повідомлення інтерпретатора про помилку позначати error)\par
\begin{verbatim}
print(5 <= 2 > 8 == 8.0)
\end{verbatim}

{\begin {center}\textbf{Завдання 6}\end {center}}

Що буде виведено за виконання? (повідомлення інтерпретатора про помилку позначати error)\par
\begin{verbatim}

def f(n):
    print("f", end="");
    return n

print(f(9) or f(-8))
\end{verbatim}

{\begin {center}\textbf{Завдання 7}\end {center}}

Що буде виведено за виконання (повідомлення інтерпретатора про помилку позначати error)?\par
\begin{verbatim}
def f(c):
    v=26
    if c<1: 
        v=5
    elif c<=-3:
         return 0
    else:
         v=3
    return v

print(f(-9))
\end{verbatim}

{\begin {center}\textbf{Завдання 8}\end {center}}

Що буде виведено за виконання (повідомлення інтерпретатора про помилку позначати error)?\par
\begin{verbatim}
a, b, c = 6, 8, 7
def f(a, b=9, c):
    print(a, b, c, end=" ")

f(a=0, c=5, b=4)
print(a, b, c)
\end{verbatim}

{\begin {center}\textbf{Завдання 9}\end {center}}

Що буде виведено за виконання (повідомлення інтерпретатора про помилку позначати error)?\par
\begin{verbatim}
a,b,c=3,5,6
def f(b):
    a-=3
    b=5
    c=4
    return a+b+c

a,b,c=2,1,0
print(f(a),a,b,c)
\end{verbatim}

{\begin {center}\textbf{Завдання 10}\end {center}}


Написати функцiю f, що отримує щось та повертає щось. Невелика загадка все-таки має залишатися. Але нічого принципово нового відносно задач, що наявні у pdf-ках не планується.. Стиль запису умови також оцінюється.\par
\end {large}}
\newpage

{\Large{17.10.2024 Бета-версія: Контрольна робота \No 1, варіант  \No \textbf { 125 }\par}}
{
\begin {large}
{\begin {center}\textbf{Завдання 1}\end {center}}

Що буде виведено за виконання? (повідомлення інтерпретатора про помилку позначати error)\par
\begin{verbatim}
print(9//4//6*8)
\end{verbatim}

{\begin {center}\textbf{Завдання 2}\end {center}}

Що буде виведено за виконання? (повідомлення інтерпретатора про помилку позначати error)\par
\begin{verbatim}
print(-3**4-1)
\end{verbatim}

{\begin {center}\textbf{Завдання 3}\end {center}}

Що буде виведено за виконання? (повідомлення інтерпретатора про помилку позначати error)\par
\begin{verbatim}
print(not False<=7 and 2>=9 and 7.0==4.0)
\end{verbatim}

{\begin {center}\textbf{Завдання 4}\end {center}}

Що буде виведено за виконання? (повідомлення інтерпретатора про помилку позначати error)\par
\begin{verbatim}
print("3.0" != 8.0)
\end{verbatim}

{\begin {center}\textbf{Завдання 5}\end {center}}

Що буде виведено за виконання? (повідомлення інтерпретатора про помилку позначати error)\par
\begin{verbatim}
print(True != 7.0 >= 6 < 2)
\end{verbatim}

{\begin {center}\textbf{Завдання 6}\end {center}}

Що буде виведено за виконання? (повідомлення інтерпретатора про помилку позначати error)\par
\begin{verbatim}

def f(n):
    print("f", end="");
    return n

print(f(7) and f(-6))
\end{verbatim}

{\begin {center}\textbf{Завдання 7}\end {center}}

Що буде виведено за виконання (повідомлення інтерпретатора про помилку позначати error)?\par
\begin{verbatim}
def f(a,b):
    c=79
    if b:
        return 9
    if b>-5:
         return 6
    else: 
        c=1
    return c

print(f(6,-1))
\end{verbatim}

{\begin {center}\textbf{Завдання 8}\end {center}}

Що буде виведено за виконання (повідомлення інтерпретатора про помилку позначати error)?\par
\begin{verbatim}
a, b, c = 9, 7, 8
def g(a, b=6, c):
    print(a, b, c, end=" ")

g(3, 2, 0)
print(a, b, c)
\end{verbatim}

{\begin {center}\textbf{Завдання 9}\end {center}}

Що буде виведено за виконання (повідомлення інтерпретатора про помилку позначати error)?\par
\begin{verbatim}
a,b,c=4,7,5
def g(b):
    a=3
    b=4
    c=4
    return a+b+c

a,b,c=6,5,2
print(g(a),a,b,c)
\end{verbatim}

{\begin {center}\textbf{Завдання 10}\end {center}}


Написати функцiю f, що отримує щось та повертає щось. Невелика загадка все-таки має залишатися. Але нічого принципово нового відносно задач, що наявні у pdf-ках не планується.. Стиль запису умови також оцінюється.\par
\end {large}}
\newpage

{\Large{17.10.2024 Бета-версія: Контрольна робота \No 1, варіант  \No \textbf { 126 }\par}}
{
\begin {large}
{\begin {center}\textbf{Завдання 1}\end {center}}

Що буде виведено за виконання? (повідомлення інтерпретатора про помилку позначати error)\par
\begin{verbatim}
print(4/9*9*9)
\end{verbatim}

{\begin {center}\textbf{Завдання 2}\end {center}}

Що буде виведено за виконання? (повідомлення інтерпретатора про помилку позначати error)\par
\begin{verbatim}
print(-2**-2--2)
\end{verbatim}

{\begin {center}\textbf{Завдання 3}\end {center}}

Що буде виведено за виконання? (повідомлення інтерпретатора про помилку позначати error)\par
\begin{verbatim}
print(8!=3.0 or not 7<True or 5.0>5)
\end{verbatim}

{\begin {center}\textbf{Завдання 4}\end {center}}

Що буде виведено за виконання? (повідомлення інтерпретатора про помилку позначати error)\par
\begin{verbatim}
print("7" == "True")
\end{verbatim}

{\begin {center}\textbf{Завдання 5}\end {center}}

Що буде виведено за виконання? (повідомлення інтерпретатора про помилку позначати error)\par
\begin{verbatim}
print(6.0  is  5 <= False > 4.0)
\end{verbatim}

{\begin {center}\textbf{Завдання 6}\end {center}}

Що буде виведено за виконання? (повідомлення інтерпретатора про помилку позначати error)\par
\begin{verbatim}

def f(n):
    print("f", end="");
    return n>=-2

print(f(2) or f(9))
\end{verbatim}

{\begin {center}\textbf{Завдання 7}\end {center}}

Що буде виведено за виконання (повідомлення інтерпретатора про помилку позначати error)?\par
\begin{verbatim}
def h(a,b):
    c=17
    if a>=2:
        c=8
    elif a<-2:
         return 4
    else: 
        c=7
    return c

print(h(1,5))
\end{verbatim}

{\begin {center}\textbf{Завдання 8}\end {center}}

Що буде виведено за виконання (повідомлення інтерпретатора про помилку позначати error)?\par
\begin{verbatim}
a, b, c = 9, 7, 8
def h(a, b, c=6):
    print(a, b, c, end=" ")

h(a=5, 3, c=1)
print(a, b, c)
\end{verbatim}

{\begin {center}\textbf{Завдання 9}\end {center}}

Що буде виведено за виконання (повідомлення інтерпретатора про помилку позначати error)?\par
\begin{verbatim}
a,b,c=4,7,9
def f(a):
    a-=2
    b=1
    c=4
    return a+b+c

a,b,c=1,4,0
print(f(b),a,b,c)
\end{verbatim}

{\begin {center}\textbf{Завдання 10}\end {center}}


Написати функцiю f, що отримує щось та повертає щось. Невелика загадка все-таки має залишатися. Але нічого принципово нового відносно задач, що наявні у pdf-ках не планується.. Стиль запису умови також оцінюється.\par
\end {large}}
\newpage

{\Large{17.10.2024 Бета-версія: Контрольна робота \No 1, варіант  \No \textbf { 127 }\par}}
{
\begin {large}
{\begin {center}\textbf{Завдання 1}\end {center}}

Що буде виведено за виконання? (повідомлення інтерпретатора про помилку позначати error)\par
\begin{verbatim}
print(7/5/8*3)
\end{verbatim}

{\begin {center}\textbf{Завдання 2}\end {center}}

Що буде виведено за виконання? (повідомлення інтерпретатора про помилку позначати error)\par
\begin{verbatim}
print(2**-1+1)
\end{verbatim}

{\begin {center}\textbf{Завдання 3}\end {center}}

Що буде виведено за виконання? (повідомлення інтерпретатора про помилку позначати error)\par
\begin{verbatim}
print(8>=8.0 or not 3==2.0 or True<9)
\end{verbatim}

{\begin {center}\textbf{Завдання 4}\end {center}}

Що буде виведено за виконання? (повідомлення інтерпретатора про помилку позначати error)\par
\begin{verbatim}
print("8" <= 9)
\end{verbatim}

{\begin {center}\textbf{Завдання 5}\end {center}}

Що буде виведено за виконання? (повідомлення інтерпретатора про помилку позначати error)\par
\begin{verbatim}
print(4 <= 7.0 != 9 >= 8.0)
\end{verbatim}

{\begin {center}\textbf{Завдання 6}\end {center}}

Що буде виведено за виконання? (повідомлення інтерпретатора про помилку позначати error)\par
\begin{verbatim}

def f(n):
    print("f", end="");
    return n

print(f(-2) or f(-1))
\end{verbatim}

{\begin {center}\textbf{Завдання 7}\end {center}}

Що буде виведено за виконання (повідомлення інтерпретатора про помилку позначати error)?\par
\begin{verbatim}
def g(b):
    u=97
    if b==-1: 
        return 7
    if b>=3:
         return 6
    else:
         u=1
    return u

print(g(4))
\end{verbatim}

{\begin {center}\textbf{Завдання 8}\end {center}}

Що буде виведено за виконання (повідомлення інтерпретатора про помилку позначати error)?\par
\begin{verbatim}
a, b, c = 6, 7, 9
def h(a, b, c):
    print(a, b, c, end=" ")

h(4, 2, a=5)
print(a, b, c)
\end{verbatim}

{\begin {center}\textbf{Завдання 9}\end {center}}

Що буде виведено за виконання (повідомлення інтерпретатора про помилку позначати error)?\par
\begin{verbatim}
a,b,c=0,8,9
def h(a):
    a=2
    b=1
    c=3
    return a+b+c

a,b,c=7,1,6
print(h(a),a,b,c)
\end{verbatim}

{\begin {center}\textbf{Завдання 10}\end {center}}


Написати функцiю f, що отримує щось та повертає щось. Невелика загадка все-таки має залишатися. Але нічого принципово нового відносно задач, що наявні у pdf-ках не планується.. Стиль запису умови також оцінюється.\par
\end {large}}
\newpage

{\Large{17.10.2024 Бета-версія: Контрольна робота \No 1, варіант  \No \textbf { 128 }\par}}
{
\begin {large}
{\begin {center}\textbf{Завдання 1}\end {center}}

Що буде виведено за виконання? (повідомлення інтерпретатора про помилку позначати error)\par
\begin{verbatim}
print(9//7%4*5)
\end{verbatim}

{\begin {center}\textbf{Завдання 2}\end {center}}

Що буде виведено за виконання? (повідомлення інтерпретатора про помилку позначати error)\par
\begin{verbatim}
print(9**-1.0*2)
\end{verbatim}

{\begin {center}\textbf{Завдання 3}\end {center}}

Що буде виведено за виконання? (повідомлення інтерпретатора про помилку позначати error)\par
\begin{verbatim}
print(not 2!=9.0 and 6.0>False and 4<=6)
\end{verbatim}

{\begin {center}\textbf{Завдання 4}\end {center}}

Що буде виведено за виконання? (повідомлення інтерпретатора про помилку позначати error)\par
\begin{verbatim}
print("False" >= "4.0j")
\end{verbatim}

{\begin {center}\textbf{Завдання 5}\end {center}}

Що буде виведено за виконання? (повідомлення інтерпретатора про помилку позначати error)\par
\begin{verbatim}
print(True == 3 < 7 > 6.0)
\end{verbatim}

{\begin {center}\textbf{Завдання 6}\end {center}}

Що буде виведено за виконання? (повідомлення інтерпретатора про помилку позначати error)\par
\begin{verbatim}

def f(n):
    print("f", end="");
    return n<-4

print(f(-7) and f(5))
\end{verbatim}

{\begin {center}\textbf{Завдання 7}\end {center}}

Що буде виведено за виконання (повідомлення інтерпретатора про помилку позначати error)?\par
\begin{verbatim}
def f(b):
    w=19
    if b>0: 
        return 4
    if b>-4:
         w=7
    else:
         w=2
    return w

print(f(-1))
\end{verbatim}

{\begin {center}\textbf{Завдання 8}\end {center}}

Що буде виведено за виконання (повідомлення інтерпретатора про помилку позначати error)?\par
\begin{verbatim}
a, b, c = 8, 8, 7
def g(a, b=6, c=9):
    print(a, b, c, end=" ")

g(b=2, c=4, 3)
print(a, b, c)
\end{verbatim}

{\begin {center}\textbf{Завдання 9}\end {center}}

Що буде виведено за виконання (повідомлення інтерпретатора про помилку позначати error)?\par
\begin{verbatim}
a,b,c=2,7,3
def h(b):
    a=5
    b+=2
    c=3
    return a+b+c

a,b,c=5,8,6
print(h(a),a,b,c)
\end{verbatim}

{\begin {center}\textbf{Завдання 10}\end {center}}


Написати функцiю f, що отримує щось та повертає щось. Невелика загадка все-таки має залишатися. Але нічого принципово нового відносно задач, що наявні у pdf-ках не планується.. Стиль запису умови також оцінюється.\par
\end {large}}
\newpage

{\Large{17.10.2024 Бета-версія: Контрольна робота \No 1, варіант  \No \textbf { 129 }\par}}
{
\begin {large}
{\begin {center}\textbf{Завдання 1}\end {center}}

Що буде виведено за виконання? (повідомлення інтерпретатора про помилку позначати error)\par
\begin{verbatim}
print(5/6%9*9)
\end{verbatim}

{\begin {center}\textbf{Завдання 2}\end {center}}

Що буде виведено за виконання? (повідомлення інтерпретатора про помилку позначати error)\par
\begin{verbatim}
print(-2**2.0+-1)
\end{verbatim}

{\begin {center}\textbf{Завдання 3}\end {center}}

Що буде виведено за виконання? (повідомлення інтерпретатора про помилку позначати error)\par
\begin{verbatim}
print(not 6!=5.0 and 2.0>=5 or False<3)
\end{verbatim}

{\begin {center}\textbf{Завдання 4}\end {center}}

Що буде виведено за виконання? (повідомлення інтерпретатора про помилку позначати error)\par
\begin{verbatim}
print(True > 2.0)
\end{verbatim}

{\begin {center}\textbf{Завдання 5}\end {center}}

Що буде виведено за виконання? (повідомлення інтерпретатора про помилку позначати error)\par
\begin{verbatim}
print(4  is  9 == 2.0 >= False)
\end{verbatim}

{\begin {center}\textbf{Завдання 6}\end {center}}

Що буде виведено за виконання? (повідомлення інтерпретатора про помилку позначати error)\par
\begin{verbatim}

def f(n):
    print("f", end="");
    return n<=0

print(f(4) or f(6))
\end{verbatim}

{\begin {center}\textbf{Завдання 7}\end {center}}

Що буде виведено за виконання (повідомлення інтерпретатора про помилку позначати error)?\par
\begin{verbatim}
def h(a,b):
    c=21
    if b!=0:
        return 0
    if a<=-1:
         c=5
    else: 
        return 9
    return c

print(h(5,-6))
\end{verbatim}

{\begin {center}\textbf{Завдання 8}\end {center}}

Що буде виведено за виконання (повідомлення інтерпретатора про помилку позначати error)?\par
\begin{verbatim}
a, b, c = 9, 7, 8
def f(a, b=6, c):
    print(a, b, c, end=" ")

f(1, 0)
print(a, b, c)
\end{verbatim}

{\begin {center}\textbf{Завдання 9}\end {center}}

Що буде виведено за виконання (повідомлення інтерпретатора про помилку позначати error)?\par
\begin{verbatim}
a,b,c=3,9,1
def g(a):
    a*=2
    b=4
    c=3
    return a+b+c

a,b,c=8,7,2
print(g(b),a,b,c)
\end{verbatim}

{\begin {center}\textbf{Завдання 10}\end {center}}


Написати функцiю f, що отримує щось та повертає щось. Невелика загадка все-таки має залишатися. Але нічого принципово нового відносно задач, що наявні у pdf-ках не планується.. Стиль запису умови також оцінюється.\par
\end {large}}
\newpage

{\Large{17.10.2024 Бета-версія: Контрольна робота \No 1, варіант  \No \textbf { 130 }\par}}
{
\begin {large}
{\begin {center}\textbf{Завдання 1}\end {center}}

Що буде виведено за виконання? (повідомлення інтерпретатора про помилку позначати error)\par
\begin{verbatim}
print(6//4/5*8)
\end{verbatim}

{\begin {center}\textbf{Завдання 2}\end {center}}

Що буде виведено за виконання? (повідомлення інтерпретатора про помилку позначати error)\par
\begin{verbatim}
print(4**1.0*-2)
\end{verbatim}

{\begin {center}\textbf{Завдання 3}\end {center}}

Що буде виведено за виконання? (повідомлення інтерпретатора про помилку позначати error)\par
\begin{verbatim}
print(not 9.0<=7 or True>4.0 and 8==4)
\end{verbatim}

{\begin {center}\textbf{Завдання 4}\end {center}}

Що буде виведено за виконання? (повідомлення інтерпретатора про помилку позначати error)\par
\begin{verbatim}
print(False != "5.0")
\end{verbatim}

{\begin {center}\textbf{Завдання 5}\end {center}}

Що буде виведено за виконання? (повідомлення інтерпретатора про помилку позначати error)\par
\begin{verbatim}
print(8 != True <= 5.0 > 4.0)
\end{verbatim}

{\begin {center}\textbf{Завдання 6}\end {center}}

Що буде виведено за виконання? (повідомлення інтерпретатора про помилку позначати error)\par
\begin{verbatim}

def f(n):
    print("f", end="");
    return n

print(f(3) and f(-9))
\end{verbatim}

{\begin {center}\textbf{Завдання 7}\end {center}}

Що буде виведено за виконання (повідомлення інтерпретатора про помилку позначати error)?\par
\begin{verbatim}
def g(a,b):
    c=26
    if a==1:
        c=1
    elif a!=5:
         return 2
    else: 
        c=0
    return c

print(g(7,2))
\end{verbatim}

{\begin {center}\textbf{Завдання 8}\end {center}}

Що буде виведено за виконання (повідомлення інтерпретатора про помилку позначати error)?\par
\begin{verbatim}
a, b, c = 6, 9, 8
def h(a, b, c):
    print(a, b, c, end=" ")

h(1, c=5, b=4)
print(a, b, c)
\end{verbatim}

{\begin {center}\textbf{Завдання 9}\end {center}}

Що буде виведено за виконання (повідомлення інтерпретатора про помилку позначати error)?\par
\begin{verbatim}
a,b,c=0,5,4
def h(a):
    a*=1
    b=5
    c=5
    return a+b+c

a,b,c=6,4,6
print(h(b),a,b,c)
\end{verbatim}

{\begin {center}\textbf{Завдання 10}\end {center}}


Написати функцiю f, що отримує щось та повертає щось. Невелика загадка все-таки має залишатися. Але нічого принципово нового відносно задач, що наявні у pdf-ках не планується.. Стиль запису умови також оцінюється.\par
\end {large}}
\newpage

{\Large{17.10.2024 Бета-версія: Контрольна робота \No 1, варіант  \No \textbf { 131 }\par}}
{
\begin {large}
{\begin {center}\textbf{Завдання 1}\end {center}}

Що буде виведено за виконання? (повідомлення інтерпретатора про помилку позначати error)\par
\begin{verbatim}
print(3/3//7*4)
\end{verbatim}

{\begin {center}\textbf{Завдання 2}\end {center}}

Що буде виведено за виконання? (повідомлення інтерпретатора про помилку позначати error)\par
\begin{verbatim}
print(1**0.5*True)
\end{verbatim}

{\begin {center}\textbf{Завдання 3}\end {center}}

Що буде виведено за виконання? (повідомлення інтерпретатора про помилку позначати error)\par
\begin{verbatim}
print(not False!=3.0 and 6.0<=9 or 2==3)
\end{verbatim}

{\begin {center}\textbf{Завдання 4}\end {center}}

Що буде виведено за виконання? (повідомлення інтерпретатора про помилку позначати error)\par
\begin{verbatim}
print("True" < 4)
\end{verbatim}

{\begin {center}\textbf{Завдання 5}\end {center}}

Що буде виведено за виконання? (повідомлення інтерпретатора про помилку позначати error)\par
\begin{verbatim}
print(7 < 2  is  False != 6)
\end{verbatim}

{\begin {center}\textbf{Завдання 6}\end {center}}

Що буде виведено за виконання? (повідомлення інтерпретатора про помилку позначати error)\par
\begin{verbatim}

def f(n):
    print("f", end="");
    return n

print(f(-5) and f(-3))
\end{verbatim}

{\begin {center}\textbf{Завдання 7}\end {center}}

Що буде виведено за виконання (повідомлення інтерпретатора про помилку позначати error)?\par
\begin{verbatim}
def h(a):
    x=61
    if a: 
        return 9
    elif a<5:
         return 8
    else:
         x=5
    return x

print(h(1))
\end{verbatim}

{\begin {center}\textbf{Завдання 8}\end {center}}

Що буде виведено за виконання (повідомлення інтерпретатора про помилку позначати error)?\par
\begin{verbatim}
a, b, c = 7, 9, 6
def f(a, b=7, c=8):
    print(a, b, c, end=" ")

f(a=0, c=3, b=2)
print(a, b, c)
\end{verbatim}

{\begin {center}\textbf{Завдання 9}\end {center}}

Що буде виведено за виконання (повідомлення інтерпретатора про помилку позначати error)?\par
\begin{verbatim}
a,b,c=7,5,3
def g(b):
    a=2
    b-=1
    c=4
    return a+b+c

a,b,c=2,1,9
print(g(b),a,b,c)
\end{verbatim}

{\begin {center}\textbf{Завдання 10}\end {center}}


Написати функцiю f, що отримує щось та повертає щось. Невелика загадка все-таки має залишатися. Але нічого принципово нового відносно задач, що наявні у pdf-ках не планується.. Стиль запису умови також оцінюється.\par
\end {large}}
\newpage

{\Large{17.10.2024 Бета-версія: Контрольна робота \No 1, варіант  \No \textbf { 132 }\par}}
{
\begin {large}
{\begin {center}\textbf{Завдання 1}\end {center}}

Що буде виведено за виконання? (повідомлення інтерпретатора про помилку позначати error)\par
\begin{verbatim}
print(4//8*3*6)
\end{verbatim}

{\begin {center}\textbf{Завдання 2}\end {center}}

Що буде виведено за виконання? (повідомлення інтерпретатора про помилку позначати error)\par
\begin{verbatim}
print(1**-4-False)
\end{verbatim}

{\begin {center}\textbf{Завдання 3}\end {center}}

Що буде виведено за виконання? (повідомлення інтерпретатора про помилку позначати error)\par
\begin{verbatim}
print(not 8.0<6 or 5>4 and True>=7.0)
\end{verbatim}

{\begin {center}\textbf{Завдання 4}\end {center}}

Що буде виведено за виконання? (повідомлення інтерпретатора про помилку позначати error)\par
\begin{verbatim}
print("6j" == 9.0)
\end{verbatim}

{\begin {center}\textbf{Завдання 5}\end {center}}

Що буде виведено за виконання? (повідомлення інтерпретатора про помилку позначати error)\par
\begin{verbatim}
print(9.0 > 5 == 3.0 >= 3)
\end{verbatim}

{\begin {center}\textbf{Завдання 6}\end {center}}

Що буде виведено за виконання? (повідомлення інтерпретатора про помилку позначати error)\par
\begin{verbatim}

def f(n):
    print("f", end="");
    return n!=-1

print(f(-4) or f(0))
\end{verbatim}

{\begin {center}\textbf{Завдання 7}\end {center}}

Що буде виведено за виконання (повідомлення інтерпретатора про помилку позначати error)?\par
\begin{verbatim}
def g(d):
    x=47
    if d<=0: 
        x=0
    if d==-2:
         return 3
    else:
         x=8
    return x

print(g(9))
\end{verbatim}

{\begin {center}\textbf{Завдання 8}\end {center}}

Що буде виведено за виконання (повідомлення інтерпретатора про помилку позначати error)?\par
\begin{verbatim}
a, b, c = 6, 8, 9
def g(a, b, c=7):
    print(a, b, c, end=" ")

g(b=3, c=5, 2)
print(a, b, c)
\end{verbatim}

{\begin {center}\textbf{Завдання 9}\end {center}}

Що буде виведено за виконання (повідомлення інтерпретатора про помилку позначати error)?\par
\begin{verbatim}
a,b,c=4,3,2
def f(b):
    a=5
    b=2
    c=3
    return a+b+c

a,b,c=9,3,8
print(f(b),a,b,c)
\end{verbatim}

{\begin {center}\textbf{Завдання 10}\end {center}}


Написати функцiю f, що отримує щось та повертає щось. Невелика загадка все-таки має залишатися. Але нічого принципово нового відносно задач, що наявні у pdf-ках не планується.. Стиль запису умови також оцінюється.\par
\end {large}}
\newpage

{\Large{17.10.2024 Бета-версія: Контрольна робота \No 1, варіант  \No \textbf { 133 }\par}}
{
\begin {large}
{\begin {center}\textbf{Завдання 1}\end {center}}

Що буде виведено за виконання? (повідомлення інтерпретатора про помилку позначати error)\par
\begin{verbatim}
print(8/9%6*7)
\end{verbatim}

{\begin {center}\textbf{Завдання 2}\end {center}}

Що буде виведено за виконання? (повідомлення інтерпретатора про помилку позначати error)\par
\begin{verbatim}
print(1**-0.5*-2)
\end{verbatim}

{\begin {center}\textbf{Завдання 3}\end {center}}

Що буде виведено за виконання? (повідомлення інтерпретатора про помилку позначати error)\par
\begin{verbatim}
print(2<False and not 8!=9 or 3.0<=5.0)
\end{verbatim}

{\begin {center}\textbf{Завдання 4}\end {center}}

Що буде виведено за виконання? (повідомлення інтерпретатора про помилку позначати error)\par
\begin{verbatim}
print(1.0 < "7.0")
\end{verbatim}

{\begin {center}\textbf{Завдання 5}\end {center}}

Що буде виведено за виконання? (повідомлення інтерпретатора про помилку позначати error)\par
\begin{verbatim}
print(3.0 < 5  is  4 <= True)
\end{verbatim}

{\begin {center}\textbf{Завдання 6}\end {center}}

Що буде виведено за виконання? (повідомлення інтерпретатора про помилку позначати error)\par
\begin{verbatim}

def f(n):
    print("f", end="");
    return n

print(f(8) or f(1))
\end{verbatim}

{\begin {center}\textbf{Завдання 7}\end {center}}

Що буде виведено за виконання (повідомлення інтерпретатора про помилку позначати error)?\par
\begin{verbatim}
def f(a,b):
    c=84
    if b:
        return 7
    if b<=4:
         c=8
    else: 
        return 6
    return c

print(f(-1,-5))
\end{verbatim}

{\begin {center}\textbf{Завдання 8}\end {center}}

Що буде виведено за виконання (повідомлення інтерпретатора про помилку позначати error)?\par
\begin{verbatim}
a, b, c = 9, 7, 8
def g(a, b=6, c=7):
    print(a, b, c, end=" ")

g(0, 4, 1)
print(a, b, c)
\end{verbatim}

{\begin {center}\textbf{Завдання 9}\end {center}}

Що буде виведено за виконання (повідомлення інтерпретатора про помилку позначати error)?\par
\begin{verbatim}
a,b,c=8,0,8
def f(b):
    a=3
    b+=4
    c=1
    return a+b+c

a,b,c=2,4,0
print(f(b),a,b,c)
\end{verbatim}

{\begin {center}\textbf{Завдання 10}\end {center}}


Написати функцiю f, що отримує щось та повертає щось. Невелика загадка все-таки має залишатися. Але нічого принципово нового відносно задач, що наявні у pdf-ках не планується.. Стиль запису умови також оцінюється.\par
\end {large}}
\newpage

{\Large{17.10.2024 Бета-версія: Контрольна робота \No 1, варіант  \No \textbf { 134 }\par}}
{
\begin {large}
{\begin {center}\textbf{Завдання 1}\end {center}}

Що буде виведено за виконання? (повідомлення інтерпретатора про помилку позначати error)\par
\begin{verbatim}
print(4//7/9*7)
\end{verbatim}

{\begin {center}\textbf{Завдання 2}\end {center}}

Що буде виведено за виконання? (повідомлення інтерпретатора про помилку позначати error)\par
\begin{verbatim}
print(4**-2--1)
\end{verbatim}

{\begin {center}\textbf{Завдання 3}\end {center}}

Що буде виведено за виконання? (повідомлення інтерпретатора про помилку позначати error)\par
\begin{verbatim}
print(not True>=7 or 4==3 and 9.0>7.0)
\end{verbatim}

{\begin {center}\textbf{Завдання 4}\end {center}}

Що буде виведено за виконання? (повідомлення інтерпретатора про помилку позначати error)\par
\begin{verbatim}
print("True" != False)
\end{verbatim}

{\begin {center}\textbf{Завдання 5}\end {center}}

Що буде виведено за виконання? (повідомлення інтерпретатора про помилку позначати error)\par
\begin{verbatim}
print(7.0 < 3 <= 8 >= 9)
\end{verbatim}

{\begin {center}\textbf{Завдання 6}\end {center}}

Що буде виведено за виконання? (повідомлення інтерпретатора про помилку позначати error)\par
\begin{verbatim}

def f(n):
    print("f", end="");
    return n>=-3

print(f(9) and f(8))
\end{verbatim}

{\begin {center}\textbf{Завдання 7}\end {center}}

Що буде виведено за виконання (повідомлення інтерпретатора про помилку позначати error)?\par
\begin{verbatim}
def g(a,b):
    c=55
    if a:
        c=3
    if a>=2:
         return 4
    else: 
        c=6
    return c

print(g(-5,7))
\end{verbatim}

{\begin {center}\textbf{Завдання 8}\end {center}}

Що буде виведено за виконання (повідомлення інтерпретатора про помилку позначати error)?\par
\begin{verbatim}
a, b, c = 9, 8, 6
def f(a, b, c):
    print(a, b, c, end=" ")

f(5, a=1)
print(a, b, c)
\end{verbatim}

{\begin {center}\textbf{Завдання 9}\end {center}}

Що буде виведено за виконання (повідомлення інтерпретатора про помилку позначати error)?\par
\begin{verbatim}
a,b,c=6,5,1
def f(a):
    a=5
    b+=3
    c=2
    return a+b+c

a,b,c=7,9,3
print(f(a),a,b,c)
\end{verbatim}

{\begin {center}\textbf{Завдання 10}\end {center}}


Написати функцiю f, що отримує щось та повертає щось. Невелика загадка все-таки має залишатися. Але нічого принципово нового відносно задач, що наявні у pdf-ках не планується.. Стиль запису умови також оцінюється.\par
\end {large}}
\newpage

{\Large{17.10.2024 Бета-версія: Контрольна робота \No 1, варіант  \No \textbf { 135 }\par}}
{
\begin {large}
{\begin {center}\textbf{Завдання 1}\end {center}}

Що буде виведено за виконання? (повідомлення інтерпретатора про помилку позначати error)\par
\begin{verbatim}
print(9/5//6*3)
\end{verbatim}

{\begin {center}\textbf{Завдання 2}\end {center}}

Що буде виведено за виконання? (повідомлення інтерпретатора про помилку позначати error)\par
\begin{verbatim}
print(9**-1.0*True)
\end{verbatim}

{\begin {center}\textbf{Завдання 3}\end {center}}

Що буде виведено за виконання? (повідомлення інтерпретатора про помилку позначати error)\par
\begin{verbatim}
print(4.0<=9 or not 6.0!=True and 6==8)
\end{verbatim}

{\begin {center}\textbf{Завдання 4}\end {center}}

Що буде виведено за виконання? (повідомлення інтерпретатора про помилку позначати error)\par
\begin{verbatim}
print(5 <= "2j")
\end{verbatim}

{\begin {center}\textbf{Завдання 5}\end {center}}

Що буде виведено за виконання? (повідомлення інтерпретатора про помилку позначати error)\par
\begin{verbatim}
print(8.0 > False == 9.0 != 7)
\end{verbatim}

{\begin {center}\textbf{Завдання 6}\end {center}}

Що буде виведено за виконання? (повідомлення інтерпретатора про помилку позначати error)\par
\begin{verbatim}

def f(n):
    print("f", end="");
    return n

print(f(-3) and f(7))
\end{verbatim}

{\begin {center}\textbf{Завдання 7}\end {center}}

Що буде виведено за виконання (повідомлення інтерпретатора про помилку позначати error)?\par
\begin{verbatim}
def h(c):
    u=64
    if c: 
        return 5
    elif c!=1:
         u=4
    else:
         return 3
    return u

print(h(7))
\end{verbatim}

{\begin {center}\textbf{Завдання 8}\end {center}}

Що буде виведено за виконання (повідомлення інтерпретатора про помилку позначати error)?\par
\begin{verbatim}
a, b, c = 7, 9, 8
def h(a, b=6, c):
    print(a, b, c, end=" ")

h(2, 0, c=3)
print(a, b, c)
\end{verbatim}

{\begin {center}\textbf{Завдання 9}\end {center}}

Що буде виведено за виконання (повідомлення інтерпретатора про помилку позначати error)?\par
\begin{verbatim}
a,b,c=0,6,7
def g(a):
    a-=5
    b=3
    c=4
    return a+b+c

a,b,c=4,2,9
print(g(a),a,b,c)
\end{verbatim}

{\begin {center}\textbf{Завдання 10}\end {center}}


Написати функцiю f, що отримує щось та повертає щось. Невелика загадка все-таки має залишатися. Але нічого принципово нового відносно задач, що наявні у pdf-ках не планується.. Стиль запису умови також оцінюється.\par
\end {large}}
\newpage

{\Large{17.10.2024 Бета-версія: Контрольна робота \No 1, варіант  \No \textbf { 136 }\par}}
{
\begin {large}
{\begin {center}\textbf{Завдання 1}\end {center}}

Що буде виведено за виконання? (повідомлення інтерпретатора про помилку позначати error)\par
\begin{verbatim}
print(3//8*7*6)
\end{verbatim}

{\begin {center}\textbf{Завдання 2}\end {center}}

Що буде виведено за виконання? (повідомлення інтерпретатора про помилку позначати error)\par
\begin{verbatim}
print(2**4+False)
\end{verbatim}

{\begin {center}\textbf{Завдання 3}\end {center}}

Що буде виведено за виконання? (повідомлення інтерпретатора про помилку позначати error)\par
\begin{verbatim}
print(5<False and not 8.0>2 or 2.0>=7)
\end{verbatim}

{\begin {center}\textbf{Завдання 4}\end {center}}

Що буде виведено за виконання? (повідомлення інтерпретатора про помилку позначати error)\par
\begin{verbatim}
print(2 == 9.0)
\end{verbatim}

{\begin {center}\textbf{Завдання 5}\end {center}}

Що буде виведено за виконання? (повідомлення інтерпретатора про помилку позначати error)\par
\begin{verbatim}
print(False  is  2 <= 4.0 != 6.0)
\end{verbatim}

{\begin {center}\textbf{Завдання 6}\end {center}}

Що буде виведено за виконання? (повідомлення інтерпретатора про помилку позначати error)\par
\begin{verbatim}

def f(n):
    print("f", end="");
    return n>-2

print(f(-5) or f(-4))
\end{verbatim}

{\begin {center}\textbf{Завдання 7}\end {center}}

Що буде виведено за виконання (повідомлення інтерпретатора про помилку позначати error)?\par
\begin{verbatim}
def f(a,b):
    c=47
    if b>5:
        return 9
    elif a<-5:
         c=5
    else: 
        return 1
    return c

print(f(-4,9))
\end{verbatim}

{\begin {center}\textbf{Завдання 8}\end {center}}

Що буде виведено за виконання (повідомлення інтерпретатора про помилку позначати error)?\par
\begin{verbatim}
a, b, c = 7, 9, 8
def f(a, b, c=6):
    print(a, b, c, end=" ")

f(a=4, 5, b=3)
print(a, b, c)
\end{verbatim}

{\begin {center}\textbf{Завдання 9}\end {center}}

Що буде виведено за виконання (повідомлення інтерпретатора про помилку позначати error)?\par
\begin{verbatim}
a,b,c=5,8,1
def h(b):
    a=2
    b*=1
    c=1
    return a+b+c

a,b,c=3,7,5
print(h(a),a,b,c)
\end{verbatim}

{\begin {center}\textbf{Завдання 10}\end {center}}


Написати функцiю f, що отримує щось та повертає щось. Невелика загадка все-таки має залишатися. Але нічого принципово нового відносно задач, що наявні у pdf-ках не планується.. Стиль запису умови також оцінюється.\par
\end {large}}
\newpage

{\Large{17.10.2024 Бета-версія: Контрольна робота \No 1, варіант  \No \textbf { 137 }\par}}
{
\begin {large}
{\begin {center}\textbf{Завдання 1}\end {center}}

Що буде виведено за виконання? (повідомлення інтерпретатора про помилку позначати error)\par
\begin{verbatim}
print(5//3/5*5)
\end{verbatim}

{\begin {center}\textbf{Завдання 2}\end {center}}

Що буде виведено за виконання? (повідомлення інтерпретатора про помилку позначати error)\par
\begin{verbatim}
print(4**-1.0*1)
\end{verbatim}

{\begin {center}\textbf{Завдання 3}\end {center}}

Що буде виведено за виконання? (повідомлення інтерпретатора про помилку позначати error)\par
\begin{verbatim}
print(False==3 or not 5!=4 and 4.0>3.0)
\end{verbatim}

{\begin {center}\textbf{Завдання 4}\end {center}}

Що буде виведено за виконання? (повідомлення інтерпретатора про помилку позначати error)\par
\begin{verbatim}
print("4" > "7.0j")
\end{verbatim}

{\begin {center}\textbf{Завдання 5}\end {center}}

Що буде виведено за виконання? (повідомлення інтерпретатора про помилку позначати error)\par
\begin{verbatim}
print(6 >= 6 == 5.0 > 2)
\end{verbatim}

{\begin {center}\textbf{Завдання 6}\end {center}}

Що буде виведено за виконання? (повідомлення інтерпретатора про помилку позначати error)\par
\begin{verbatim}

def f(n):
    print("f", end="");
    return n

print(f(-6) and f(-2))
\end{verbatim}

{\begin {center}\textbf{Завдання 7}\end {center}}

Що буде виведено за виконання (повідомлення інтерпретатора про помилку позначати error)?\par
\begin{verbatim}
def f(a):
    z=53
    if a>=-5: 
        z=6
    elif a<=-4:
         return 0
    else:
         z=1
    return z

print(f(-6))
\end{verbatim}

{\begin {center}\textbf{Завдання 8}\end {center}}

Що буде виведено за виконання (повідомлення інтерпретатора про помилку позначати error)?\par
\begin{verbatim}
a, b, c = 9, 6, 7
def h(a, b, c=8):
    print(a, b, c, end=" ")

h(0, c=4, b=1)
print(a, b, c)
\end{verbatim}

{\begin {center}\textbf{Завдання 9}\end {center}}

Що буде виведено за виконання (повідомлення інтерпретатора про помилку позначати error)?\par
\begin{verbatim}
a,b,c=2,3,4
def g(a):
    a=2
    b+=3
    c=4
    return a+b+c

a,b,c=8,6,0
print(g(a),a,b,c)
\end{verbatim}

{\begin {center}\textbf{Завдання 10}\end {center}}


Написати функцiю f, що отримує щось та повертає щось. Невелика загадка все-таки має залишатися. Але нічого принципово нового відносно задач, що наявні у pdf-ках не планується.. Стиль запису умови також оцінюється.\par
\end {large}}
\newpage

{\Large{17.10.2024 Бета-версія: Контрольна робота \No 1, варіант  \No \textbf { 138 }\par}}
{
\begin {large}
{\begin {center}\textbf{Завдання 1}\end {center}}

Що буде виведено за виконання? (повідомлення інтерпретатора про помилку позначати error)\par
\begin{verbatim}
print(6/9*4*9)
\end{verbatim}

{\begin {center}\textbf{Завдання 2}\end {center}}

Що буде виведено за виконання? (повідомлення інтерпретатора про помилку позначати error)\par
\begin{verbatim}
print(-1**-2+2)
\end{verbatim}

{\begin {center}\textbf{Завдання 3}\end {center}}

Що буде виведено за виконання? (повідомлення інтерпретатора про помилку позначати error)\par
\begin{verbatim}
print(7.0<6 and not True<=9 or 7>=5.0)
\end{verbatim}

{\begin {center}\textbf{Завдання 4}\end {center}}

Що буде виведено за виконання? (повідомлення інтерпретатора про помилку позначати error)\par
\begin{verbatim}
print("True" >= False)
\end{verbatim}

{\begin {center}\textbf{Завдання 5}\end {center}}

Що буде виведено за виконання? (повідомлення інтерпретатора про помилку позначати error)\par
\begin{verbatim}
print(8 < 4  is  True  is  2.0)
\end{verbatim}

{\begin {center}\textbf{Завдання 6}\end {center}}

Що буде виведено за виконання? (повідомлення інтерпретатора про помилку позначати error)\par
\begin{verbatim}

def f(n):
    print("f", end="");
    return n<4

print(f(2) or f(-8))
\end{verbatim}

{\begin {center}\textbf{Завдання 7}\end {center}}

Що буде виведено за виконання (повідомлення інтерпретатора про помилку позначати error)?\par
\begin{verbatim}
def f(d):
    w=91
    if d==3: 
        return 2
    if d<5:
         return 7
    else:
         w=9
    return w

print(f(2))
\end{verbatim}

{\begin {center}\textbf{Завдання 8}\end {center}}

Що буде виведено за виконання (повідомлення інтерпретатора про помилку позначати error)?\par
\begin{verbatim}
a, b, c = 7, 9, 8
def g(a, b=6, c):
    print(a, b, c, end=" ")

g(2, 3, 1)
print(a, b, c)
\end{verbatim}

{\begin {center}\textbf{Завдання 9}\end {center}}

Що буде виведено за виконання (повідомлення інтерпретатора про помилку позначати error)?\par
\begin{verbatim}
a,b,c=1,9,9
def f(b):
    a*=5
    b=1
    c=2
    return a+b+c

a,b,c=3,1,8
print(f(a),a,b,c)
\end{verbatim}

{\begin {center}\textbf{Завдання 10}\end {center}}


Написати функцiю f, що отримує щось та повертає щось. Невелика загадка все-таки має залишатися. Але нічого принципово нового відносно задач, що наявні у pdf-ках не планується.. Стиль запису умови також оцінюється.\par
\end {large}}
\newpage

{\Large{17.10.2024 Бета-версія: Контрольна робота \No 1, варіант  \No \textbf { 139 }\par}}
{
\begin {large}
{\begin {center}\textbf{Завдання 1}\end {center}}

Що буде виведено за виконання? (повідомлення інтерпретатора про помилку позначати error)\par
\begin{verbatim}
print(7/6//3*4)
\end{verbatim}

{\begin {center}\textbf{Завдання 2}\end {center}}

Що буде виведено за виконання? (повідомлення інтерпретатора про помилку позначати error)\par
\begin{verbatim}
print(3**-3-True)
\end{verbatim}

{\begin {center}\textbf{Завдання 3}\end {center}}

Що буде виведено за виконання? (повідомлення інтерпретатора про помилку позначати error)\par
\begin{verbatim}
print(2<6.0 and not 8==8 and True!=2.0)
\end{verbatim}

{\begin {center}\textbf{Завдання 4}\end {center}}

Що буде виведено за виконання? (повідомлення інтерпретатора про помилку позначати error)\par
\begin{verbatim}
print(4.0 <= "6.0j")
\end{verbatim}

{\begin {center}\textbf{Завдання 5}\end {center}}

Що буде виведено за виконання? (повідомлення інтерпретатора про помилку позначати error)\par
\begin{verbatim}
print(7 > 9.0 != 9 < 2.0)
\end{verbatim}

{\begin {center}\textbf{Завдання 6}\end {center}}

Що буде виведено за виконання? (повідомлення інтерпретатора про помилку позначати error)\par
\begin{verbatim}

def f(n):
    print("f", end="");
    return n

print(f(-9) and f(-1))
\end{verbatim}

{\begin {center}\textbf{Завдання 7}\end {center}}

Що буде виведено за виконання (повідомлення інтерпретатора про помилку позначати error)?\par
\begin{verbatim}
def h(a,b):
    c=60
    if b==-1:
        return 7
    elif b==3:
         c=2
    else: 
        c=8
    return c

print(h(9,-8))
\end{verbatim}

{\begin {center}\textbf{Завдання 8}\end {center}}

Що буде виведено за виконання (повідомлення інтерпретатора про помилку позначати error)?\par
\begin{verbatim}
a, b, c = 8, 6, 9
def h(a, b, c):
    print(a, b, c, end=" ")

h(a=0, 5, c=4)
print(a, b, c)
\end{verbatim}

{\begin {center}\textbf{Завдання 9}\end {center}}

Що буде виведено за виконання (повідомлення інтерпретатора про помилку позначати error)?\par
\begin{verbatim}
a,b,c=0,6,2
def h(b):
    a=5
    b-=4
    c=3
    return a+b+c

a,b,c=7,5,4
print(h(b),a,b,c)
\end{verbatim}

{\begin {center}\textbf{Завдання 10}\end {center}}


Написати функцiю f, що отримує щось та повертає щось. Невелика загадка все-таки має залишатися. Але нічого принципово нового відносно задач, що наявні у pdf-ках не планується.. Стиль запису умови також оцінюється.\par
\end {large}}
\newpage

{\Large{17.10.2024 Бета-версія: Контрольна робота \No 1, варіант  \No \textbf { 140 }\par}}
{
\begin {large}
{\begin {center}\textbf{Завдання 1}\end {center}}

Що буде виведено за виконання? (повідомлення інтерпретатора про помилку позначати error)\par
\begin{verbatim}
print(8//4%8*8)
\end{verbatim}

{\begin {center}\textbf{Завдання 2}\end {center}}

Що буде виведено за виконання? (повідомлення інтерпретатора про помилку позначати error)\par
\begin{verbatim}
print(1**-0.5*False)
\end{verbatim}

{\begin {center}\textbf{Завдання 3}\end {center}}

Що буде виведено за виконання? (повідомлення інтерпретатора про помилку позначати error)\par
\begin{verbatim}
print(not False>=8.0 or 9<=2 or 9.0>5)
\end{verbatim}

{\begin {center}\textbf{Завдання 4}\end {center}}

Що буде виведено за виконання? (повідомлення інтерпретатора про помилку позначати error)\par
\begin{verbatim}
print(9 == False)
\end{verbatim}

{\begin {center}\textbf{Завдання 5}\end {center}}

Що буде виведено за виконання? (повідомлення інтерпретатора про помилку позначати error)\par
\begin{verbatim}
print(False >= 5 == 7.0 <= 3)
\end{verbatim}

{\begin {center}\textbf{Завдання 6}\end {center}}

Що буде виведено за виконання? (повідомлення інтерпретатора про помилку позначати error)\par
\begin{verbatim}

def f(n):
    print("f", end="");
    return n<=1

print(f(-7) or f(4))
\end{verbatim}

{\begin {center}\textbf{Завдання 7}\end {center}}

Що буде виведено за виконання (повідомлення інтерпретатора про помилку позначати error)?\par
\begin{verbatim}
def f(a,b):
    c=72
    if a>=1:
        return 0
    if b>-3:
         c=4
    else: 
        return 3
    return c

print(f(2,-9))
\end{verbatim}

{\begin {center}\textbf{Завдання 8}\end {center}}

Що буде виведено за виконання (повідомлення інтерпретатора про помилку позначати error)?\par
\begin{verbatim}
a, b, c = 7, 7, 9
def g(a, b=8, c=6):
    print(a, b, c, end=" ")

g(2, 4, a=3)
print(a, b, c)
\end{verbatim}

{\begin {center}\textbf{Завдання 9}\end {center}}

Що буде виведено за виконання (повідомлення інтерпретатора про помилку позначати error)?\par
\begin{verbatim}
a,b,c=4,5,1
def h(b):
    a=4
    b=5
    c=1
    return a+b+c

a,b,c=6,7,0
print(h(b),a,b,c)
\end{verbatim}

{\begin {center}\textbf{Завдання 10}\end {center}}


Написати функцiю f, що отримує щось та повертає щось. Невелика загадка все-таки має залишатися. Але нічого принципово нового відносно задач, що наявні у pdf-ках не планується.. Стиль запису умови також оцінюється.\par
\end {large}}
\newpage

{\Large{17.10.2024 Бета-версія: Контрольна робота \No 1, варіант  \No \textbf { 141 }\par}}
{
\begin {large}
{\begin {center}\textbf{Завдання 1}\end {center}}

Що буде виведено за виконання? (повідомлення інтерпретатора про помилку позначати error)\par
\begin{verbatim}
print(8/5//4*8)
\end{verbatim}

{\begin {center}\textbf{Завдання 2}\end {center}}

Що буде виведено за виконання? (повідомлення інтерпретатора про помилку позначати error)\par
\begin{verbatim}
print(9**0.5*-2)
\end{verbatim}

{\begin {center}\textbf{Завдання 3}\end {center}}

Що буде виведено за виконання? (повідомлення інтерпретатора про помилку позначати error)\par
\begin{verbatim}
print(not 6>=9.0 and 7>False or 7.0<3)
\end{verbatim}

{\begin {center}\textbf{Завдання 4}\end {center}}

Що буде виведено за виконання? (повідомлення інтерпретатора про помилку позначати error)\par
\begin{verbatim}
print("6" >= "True")
\end{verbatim}

{\begin {center}\textbf{Завдання 5}\end {center}}

Що буде виведено за виконання? (повідомлення інтерпретатора про помилку позначати error)\par
\begin{verbatim}
print(3 <= True < 8.0 != 7)
\end{verbatim}

{\begin {center}\textbf{Завдання 6}\end {center}}

Що буде виведено за виконання? (повідомлення інтерпретатора про помилку позначати error)\par
\begin{verbatim}

def f(n):
    print("f", end="");
    return n!=2

print(f(3) and f(6))
\end{verbatim}

{\begin {center}\textbf{Завдання 7}\end {center}}

Що буде виведено за виконання (повідомлення інтерпретатора про помилку позначати error)?\par
\begin{verbatim}
def h(a,b):
    c=69
    if b<0:
        return 4
    elif a<=-2:
         c=2
    else: 
        return 5
    return c

print(h(3,1))
\end{verbatim}

{\begin {center}\textbf{Завдання 8}\end {center}}

Що буде виведено за виконання (повідомлення інтерпретатора про помилку позначати error)?\par
\begin{verbatim}
a, b, c = 6, 7, 9
def f(a, b=8, c=6):
    print(a, b, c, end=" ")

f(b=1, c=5, a=0)
print(a, b, c)
\end{verbatim}

{\begin {center}\textbf{Завдання 9}\end {center}}

Що буде виведено за виконання (повідомлення інтерпретатора про помилку позначати error)?\par
\begin{verbatim}
a,b,c=5,7,2
def f(a):
    a*=5
    b=1
    c=4
    return a+b+c

a,b,c=8,4,3
print(f(b),a,b,c)
\end{verbatim}

{\begin {center}\textbf{Завдання 10}\end {center}}


Написати функцiю f, що отримує щось та повертає щось. Невелика загадка все-таки має залишатися. Але нічого принципово нового відносно задач, що наявні у pdf-ках не планується.. Стиль запису умови також оцінюється.\par
\end {large}}
\newpage

{\Large{17.10.2024 Бета-версія: Контрольна робота \No 1, варіант  \No \textbf { 142 }\par}}
{
\begin {large}
{\begin {center}\textbf{Завдання 1}\end {center}}

Що буде виведено за виконання? (повідомлення інтерпретатора про помилку позначати error)\par
\begin{verbatim}
print(7//7%6*6)
\end{verbatim}

{\begin {center}\textbf{Завдання 2}\end {center}}

Що буде виведено за виконання? (повідомлення інтерпретатора про помилку позначати error)\par
\begin{verbatim}
print(4**1.0*-1)
\end{verbatim}

{\begin {center}\textbf{Завдання 3}\end {center}}

Що буде виведено за виконання? (повідомлення інтерпретатора про помилку позначати error)\par
\begin{verbatim}
print(5.0<=4 or not 7==3.0 and True!=9)
\end{verbatim}

{\begin {center}\textbf{Завдання 4}\end {center}}

Що буде виведено за виконання? (повідомлення інтерпретатора про помилку позначати error)\par
\begin{verbatim}
print(7 < "5.0")
\end{verbatim}

{\begin {center}\textbf{Завдання 5}\end {center}}

Що буде виведено за виконання? (повідомлення інтерпретатора про помилку позначати error)\par
\begin{verbatim}
print(4 >= 3.0 == 8 > 6.0)
\end{verbatim}

{\begin {center}\textbf{Завдання 6}\end {center}}

Що буде виведено за виконання? (повідомлення інтерпретатора про помилку позначати error)\par
\begin{verbatim}

def f(n):
    print("f", end="");
    return n

print(f(1) or f(5))
\end{verbatim}

{\begin {center}\textbf{Завдання 7}\end {center}}

Що буде виведено за виконання (повідомлення інтерпретатора про помилку позначати error)?\par
\begin{verbatim}
def g(a,b):
    c=66
    if a:
        c=3
    if b!=4:
         c=9
    else: 
        return 0
    return c

print(g(-2,-2))
\end{verbatim}

{\begin {center}\textbf{Завдання 8}\end {center}}

Що буде виведено за виконання (повідомлення інтерпретатора про помилку позначати error)?\par
\begin{verbatim}
a, b, c = 8, 7, 9
def h(a, b, c):
    print(a, b, c, end=" ")

h(2, 3)
print(a, b, c)
\end{verbatim}

{\begin {center}\textbf{Завдання 9}\end {center}}

Що буде виведено за виконання (повідомлення інтерпретатора про помилку позначати error)?\par
\begin{verbatim}
a,b,c=9,0,1
def h(b):
    a=2
    b-=3
    c=3
    return a+b+c

a,b,c=6,1,9
print(h(b),a,b,c)
\end{verbatim}

{\begin {center}\textbf{Завдання 10}\end {center}}


Написати функцiю f, що отримує щось та повертає щось. Невелика загадка все-таки має залишатися. Але нічого принципово нового відносно задач, що наявні у pdf-ках не планується.. Стиль запису умови також оцінюється.\par
\end {large}}
\newpage

{\Large{17.10.2024 Бета-версія: Контрольна робота \No 1, варіант  \No \textbf { 143 }\par}}
{
\begin {large}
{\begin {center}\textbf{Завдання 1}\end {center}}

Що буде виведено за виконання? (повідомлення інтерпретатора про помилку позначати error)\par
\begin{verbatim}
print(4/4/3*7)
\end{verbatim}

{\begin {center}\textbf{Завдання 2}\end {center}}

Що буде виведено за виконання? (повідомлення інтерпретатора про помилку позначати error)\par
\begin{verbatim}
print(3**True+1)
\end{verbatim}

{\begin {center}\textbf{Завдання 3}\end {center}}

Що буде виведено за виконання? (повідомлення інтерпретатора про помилку позначати error)\par
\begin{verbatim}
print(not 8==5 and 4>=4.0 or True<=6.0)
\end{verbatim}

{\begin {center}\textbf{Завдання 4}\end {center}}

Що буде виведено за виконання? (повідомлення інтерпретатора про помилку позначати error)\par
\begin{verbatim}
print("True" > False)
\end{verbatim}

{\begin {center}\textbf{Завдання 5}\end {center}}

Що буде виведено за виконання? (повідомлення інтерпретатора про помилку позначати error)\par
\begin{verbatim}
print(2  is  True < 6  is  False)
\end{verbatim}

{\begin {center}\textbf{Завдання 6}\end {center}}

Що буде виведено за виконання? (повідомлення інтерпретатора про помилку позначати error)\par
\begin{verbatim}

def f(n):
    print("f", end="");
    return n

print(f(0) or f(-8))
\end{verbatim}

{\begin {center}\textbf{Завдання 7}\end {center}}

Що буде виведено за виконання (повідомлення інтерпретатора про помилку позначати error)?\par
\begin{verbatim}
def g(c):
    y=76
    if c>=2: 
        y=4
    if c>4:
         return 8
    else:
         y=7
    return y

print(g(-3))
\end{verbatim}

{\begin {center}\textbf{Завдання 8}\end {center}}

Що буде виведено за виконання (повідомлення інтерпретатора про помилку позначати error)?\par
\begin{verbatim}
a, b, c = 8, 7, 6
def f(a, b, c=9):
    print(a, b, c, end=" ")

f(b=2, c=1, 5)
print(a, b, c)
\end{verbatim}

{\begin {center}\textbf{Завдання 9}\end {center}}

Що буде виведено за виконання (повідомлення інтерпретатора про помилку позначати error)?\par
\begin{verbatim}
a,b,c=7,0,3
def g(a):
    a=3
    b=4
    c=5
    return a+b+c

a,b,c=1,6,2
print(g(a),a,b,c)
\end{verbatim}

{\begin {center}\textbf{Завдання 10}\end {center}}


Написати функцiю f, що отримує щось та повертає щось. Невелика загадка все-таки має залишатися. Але нічого принципово нового відносно задач, що наявні у pdf-ках не планується.. Стиль запису умови також оцінюється.\par
\end {large}}
\newpage

{\Large{17.10.2024 Бета-версія: Контрольна робота \No 1, варіант  \No \textbf { 144 }\par}}
{
\begin {large}
{\begin {center}\textbf{Завдання 1}\end {center}}

Що буде виведено за виконання? (повідомлення інтерпретатора про помилку позначати error)\par
\begin{verbatim}
print(6//8*7*5)
\end{verbatim}

{\begin {center}\textbf{Завдання 2}\end {center}}

Що буде виведено за виконання? (повідомлення інтерпретатора про помилку позначати error)\par
\begin{verbatim}
print(-1**False-2)
\end{verbatim}

{\begin {center}\textbf{Завдання 3}\end {center}}

Що буде виведено за виконання? (повідомлення інтерпретатора про помилку позначати error)\par
\begin{verbatim}
print(8.0<6 or not 2.0>3 and 2!=False)
\end{verbatim}

{\begin {center}\textbf{Завдання 4}\end {center}}

Що буде виведено за виконання? (повідомлення інтерпретатора про помилку позначати error)\par
\begin{verbatim}
print("8j" != 8.0)
\end{verbatim}

{\begin {center}\textbf{Завдання 5}\end {center}}

Що буде виведено за виконання? (повідомлення інтерпретатора про помилку позначати error)\par
\begin{verbatim}
print(9 == 5 <= 4.0 != 5.0)
\end{verbatim}

{\begin {center}\textbf{Завдання 6}\end {center}}

Що буде виведено за виконання? (повідомлення інтерпретатора про помилку позначати error)\par
\begin{verbatim}

def f(n):
    print("f", end="");
    return n==3

print(f(-7) and f(3))
\end{verbatim}

{\begin {center}\textbf{Завдання 7}\end {center}}

Що буде виведено за виконання (повідомлення інтерпретатора про помилку позначати error)?\par
\begin{verbatim}
def h(b):
    v=23
    if b: 
        return 1
    elif b!=-1:
         return 2
    else:
         v=5
    return v

print(h(6))
\end{verbatim}

{\begin {center}\textbf{Завдання 8}\end {center}}

Що буде виведено за виконання (повідомлення інтерпретатора про помилку позначати error)?\par
\begin{verbatim}
a, b, c = 8, 9, 6
def g(a, b=7, c):
    print(a, b, c, end=" ")

g(0, b=4)
print(a, b, c)
\end{verbatim}

{\begin {center}\textbf{Завдання 9}\end {center}}

Що буде виведено за виконання (повідомлення інтерпретатора про помилку позначати error)?\par
\begin{verbatim}
a,b,c=7,0,2
def g(a):
    a+=5
    b=4
    c=1
    return a+b+c

a,b,c=4,5,3
print(g(a),a,b,c)
\end{verbatim}

{\begin {center}\textbf{Завдання 10}\end {center}}


Написати функцiю f, що отримує щось та повертає щось. Невелика загадка все-таки має залишатися. Але нічого принципово нового відносно задач, що наявні у pdf-ках не планується.. Стиль запису умови також оцінюється.\par
\end {large}}
\newpage

{\Large{17.10.2024 Бета-версія: Контрольна робота \No 1, варіант  \No \textbf { 145 }\par}}
{
\begin {large}
{\begin {center}\textbf{Завдання 1}\end {center}}

Що буде виведено за виконання? (повідомлення інтерпретатора про помилку позначати error)\par
\begin{verbatim}
print(3/3*9*3)
\end{verbatim}

{\begin {center}\textbf{Завдання 2}\end {center}}

Що буде виведено за виконання? (повідомлення інтерпретатора про помилку позначати error)\par
\begin{verbatim}
print(-3**4--1)
\end{verbatim}

{\begin {center}\textbf{Завдання 3}\end {center}}

Що буде виведено за виконання? (повідомлення інтерпретатора про помилку позначати error)\par
\begin{verbatim}
print(9.0>5 and not 4<7 and False!=7.0)
\end{verbatim}

{\begin {center}\textbf{Завдання 4}\end {center}}

Що буде виведено за виконання? (повідомлення інтерпретатора про помилку позначати error)\par
\begin{verbatim}
print("3.0" == "False")
\end{verbatim}

{\begin {center}\textbf{Завдання 5}\end {center}}

Що буде виведено за виконання? (повідомлення інтерпретатора про помилку позначати error)\par
\begin{verbatim}
print(4 >= True > 5 == 9.0)
\end{verbatim}

{\begin {center}\textbf{Завдання 6}\end {center}}

Що буде виведено за виконання? (повідомлення інтерпретатора про помилку позначати error)\par
\begin{verbatim}

def f(n):
    print("f", end="");
    return n<1

print(f(6) or f(-4))
\end{verbatim}

{\begin {center}\textbf{Завдання 7}\end {center}}

Що буде виведено за виконання (повідомлення інтерпретатора про помилку позначати error)?\par
\begin{verbatim}
def h(a,b):
    c=11
    if a>-4:
        c=8
    elif b!=-2:
         return 6
    else: 
        c=7
    return c

print(h(-3,4))
\end{verbatim}

{\begin {center}\textbf{Завдання 8}\end {center}}

Що буде виведено за виконання (повідомлення інтерпретатора про помилку позначати error)?\par
\begin{verbatim}
a, b, c = 9, 7, 6
def f(a, b=8, c=6):
    print(a, b, c, end=" ")

f(4, 0)
print(a, b, c)
\end{verbatim}

{\begin {center}\textbf{Завдання 9}\end {center}}

Що буде виведено за виконання (повідомлення інтерпретатора про помилку позначати error)?\par
\begin{verbatim}
a,b,c=6,8,2
def f(b):
    a=2
    b*=2
    c=4
    return a+b+c

a,b,c=4,6,7
print(f(b),a,b,c)
\end{verbatim}

{\begin {center}\textbf{Завдання 10}\end {center}}


Написати функцiю f, що отримує щось та повертає щось. Невелика загадка все-таки має залишатися. Але нічого принципово нового відносно задач, що наявні у pdf-ках не планується.. Стиль запису умови також оцінюється.\par
\end {large}}
\newpage

{\Large{17.10.2024 Бета-версія: Контрольна робота \No 1, варіант  \No \textbf { 146 }\par}}
{
\begin {large}
{\begin {center}\textbf{Завдання 1}\end {center}}

Що буде виведено за виконання? (повідомлення інтерпретатора про помилку позначати error)\par
\begin{verbatim}
print(5//6%8*9)
\end{verbatim}

{\begin {center}\textbf{Завдання 2}\end {center}}

Що буде виведено за виконання? (повідомлення інтерпретатора про помилку позначати error)\par
\begin{verbatim}
print(9**-0.5*2)
\end{verbatim}

{\begin {center}\textbf{Завдання 3}\end {center}}

Що буде виведено за виконання? (повідомлення інтерпретатора про помилку позначати error)\par
\begin{verbatim}
print(not 2==2.0 or True>=9 or 8.0<=6)
\end{verbatim}

{\begin {center}\textbf{Завдання 4}\end {center}}

Що буде виведено за виконання? (повідомлення інтерпретатора про помилку позначати error)\par
\begin{verbatim}
print(5 != True)
\end{verbatim}

{\begin {center}\textbf{Завдання 5}\end {center}}

Що буде виведено за виконання? (повідомлення інтерпретатора про помилку позначати error)\par
\begin{verbatim}
print(5.0 > 8 <= 9  is  2)
\end{verbatim}

{\begin {center}\textbf{Завдання 6}\end {center}}

Що буде виведено за виконання? (повідомлення інтерпретатора про помилку позначати error)\par
\begin{verbatim}

def f(n):
    print("f", end="");
    return n

print(f(-3) and f(1))
\end{verbatim}

{\begin {center}\textbf{Завдання 7}\end {center}}

Що буде виведено за виконання (повідомлення інтерпретатора про помилку позначати error)?\par
\begin{verbatim}
def g(a,b):
    c=57
    if a:
        return 1
    if b<-4:
         return 3
    else: 
        c=2
    return c

print(g(-6,-7))
\end{verbatim}

{\begin {center}\textbf{Завдання 8}\end {center}}

Що буде виведено за виконання (повідомлення інтерпретатора про помилку позначати error)?\par
\begin{verbatim}
a, b, c = 9, 8, 7
def g(a, b, c):
    print(a, b, c, end=" ")

g(5, c=2, b=3)
print(a, b, c)
\end{verbatim}

{\begin {center}\textbf{Завдання 9}\end {center}}

Що буде виведено за виконання (повідомлення інтерпретатора про помилку позначати error)?\par
\begin{verbatim}
a,b,c=0,8,9
def g(a):
    a-=3
    b=1
    c=5
    return a+b+c

a,b,c=3,5,1
print(g(b),a,b,c)
\end{verbatim}

{\begin {center}\textbf{Завдання 10}\end {center}}


Написати функцiю f, що отримує щось та повертає щось. Невелика загадка все-таки має залишатися. Але нічого принципово нового відносно задач, що наявні у pdf-ках не планується.. Стиль запису умови також оцінюється.\par
\end {large}}
\newpage

{\Large{17.10.2024 Бета-версія: Контрольна робота \No 1, варіант  \No \textbf { 147 }\par}}
{
\begin {large}
{\begin {center}\textbf{Завдання 1}\end {center}}

Що буде виведено за виконання? (повідомлення інтерпретатора про помилку позначати error)\par
\begin{verbatim}
print(9/9/5*4)
\end{verbatim}

{\begin {center}\textbf{Завдання 2}\end {center}}

Що буде виведено за виконання? (повідомлення інтерпретатора про помилку позначати error)\par
\begin{verbatim}
print(3**-3+1)
\end{verbatim}

{\begin {center}\textbf{Завдання 3}\end {center}}

Що буде виведено за виконання? (повідомлення інтерпретатора про помилку позначати error)\par
\begin{verbatim}
print(not 3>=8 or 9==5.0 or True<4.0)
\end{verbatim}

{\begin {center}\textbf{Завдання 4}\end {center}}

Що буде виведено за виконання? (повідомлення інтерпретатора про помилку позначати error)\par
\begin{verbatim}
print(1.0 <= "3j")
\end{verbatim}

{\begin {center}\textbf{Завдання 5}\end {center}}

Що буде виведено за виконання? (повідомлення інтерпретатора про помилку позначати error)\par
\begin{verbatim}
print(False < 3 != 7.0 >= 6.0)
\end{verbatim}

{\begin {center}\textbf{Завдання 6}\end {center}}

Що буде виведено за виконання? (повідомлення інтерпретатора про помилку позначати error)\par
\begin{verbatim}

def f(n):
    print("f", end="");
    return n==-4

print(f(4) or f(2))
\end{verbatim}

{\begin {center}\textbf{Завдання 7}\end {center}}

Що буде виведено за виконання (повідомлення інтерпретатора про помилку позначати error)?\par
\begin{verbatim}
def g(c):
    y=48
    if c>-3: 
        y=3
    elif c<=5:
         return 0
    else:
         return 9
    return y

print(g(3))
\end{verbatim}

{\begin {center}\textbf{Завдання 8}\end {center}}

Що буде виведено за виконання (повідомлення інтерпретатора про помилку позначати error)?\par
\begin{verbatim}
a, b, c = 6, 7, 9
def h(a, b=8, c):
    print(a, b, c, end=" ")

h(b=1, c=4, 3)
print(a, b, c)
\end{verbatim}

{\begin {center}\textbf{Завдання 9}\end {center}}

Що буде виведено за виконання (повідомлення інтерпретатора про помилку позначати error)?\par
\begin{verbatim}
a,b,c=3,9,7
def h(a):
    a=2
    b+=4
    c=1
    return a+b+c

a,b,c=2,1,5
print(h(a),a,b,c)
\end{verbatim}

{\begin {center}\textbf{Завдання 10}\end {center}}


Написати функцiю f, що отримує щось та повертає щось. Невелика загадка все-таки має залишатися. Але нічого принципово нового відносно задач, що наявні у pdf-ках не планується.. Стиль запису умови також оцінюється.\par
\end {large}}
\newpage

{\Large{17.10.2024 Бета-версія: Контрольна робота \No 1, варіант  \No \textbf { 148 }\par}}
{
\begin {large}
{\begin {center}\textbf{Завдання 1}\end {center}}

Що буде виведено за виконання? (повідомлення інтерпретатора про помилку позначати error)\par
\begin{verbatim}
print(8//6//7*4)
\end{verbatim}

{\begin {center}\textbf{Завдання 2}\end {center}}

Що буде виведено за виконання? (повідомлення інтерпретатора про помилку позначати error)\par
\begin{verbatim}
print(4**False--2)
\end{verbatim}

{\begin {center}\textbf{Завдання 3}\end {center}}

Що буде виведено за виконання? (повідомлення інтерпретатора про помилку позначати error)\par
\begin{verbatim}
print(3>False and not 5<=4 and 6.0!=3.0)
\end{verbatim}

{\begin {center}\textbf{Завдання 4}\end {center}}

Що буде виведено за виконання? (повідомлення інтерпретатора про помилку позначати error)\par
\begin{verbatim}
print(2.0 > "False")
\end{verbatim}

{\begin {center}\textbf{Завдання 5}\end {center}}

Що буде виведено за виконання? (повідомлення інтерпретатора про помилку позначати error)\par
\begin{verbatim}
print(4.0  is  7 < 6 >= True)
\end{verbatim}

{\begin {center}\textbf{Завдання 6}\end {center}}

Що буде виведено за виконання? (повідомлення інтерпретатора про помилку позначати error)\par
\begin{verbatim}

def f(n):
    print("f", end="");
    return n

print(f(7) and f(-5))
\end{verbatim}

{\begin {center}\textbf{Завдання 7}\end {center}}

Що буде виведено за виконання (повідомлення інтерпретатора про помилку позначати error)?\par
\begin{verbatim}
def f(a,b):
    c=14
    if a>=4:
        return 0
    if b<=0:
         c=5
    else: 
        c=4
    return c

print(f(-8,-3))
\end{verbatim}

{\begin {center}\textbf{Завдання 8}\end {center}}

Що буде виведено за виконання (повідомлення інтерпретатора про помилку позначати error)?\par
\begin{verbatim}
a, b, c = 8, 7, 6
def g(a, b, c=9):
    print(a, b, c, end=" ")

g(2, 5, 0)
print(a, b, c)
\end{verbatim}

{\begin {center}\textbf{Завдання 9}\end {center}}

Що буде виведено за виконання (повідомлення інтерпретатора про помилку позначати error)?\par
\begin{verbatim}
a,b,c=8,4,0
def f(b):
    a=3
    b+=5
    c=1
    return a+b+c

a,b,c=6,3,1
print(f(b),a,b,c)
\end{verbatim}

{\begin {center}\textbf{Завдання 10}\end {center}}


Написати функцiю f, що отримує щось та повертає щось. Невелика загадка все-таки має залишатися. Але нічого принципово нового відносно задач, що наявні у pdf-ках не планується.. Стиль запису умови також оцінюється.\par
\end {large}}
\newpage

{\Large{17.10.2024 Бета-версія: Контрольна робота \No 1, варіант  \No \textbf { 149 }\par}}
{
\begin {large}
{\begin {center}\textbf{Завдання 1}\end {center}}

Що буде виведено за виконання? (повідомлення інтерпретатора про помилку позначати error)\par
\begin{verbatim}
print(3/8/6*9)
\end{verbatim}

{\begin {center}\textbf{Завдання 2}\end {center}}

Що буде виведено за виконання? (повідомлення інтерпретатора про помилку позначати error)\par
\begin{verbatim}
print(-2**-4+False)
\end{verbatim}

{\begin {center}\textbf{Завдання 3}\end {center}}

Що буде виведено за виконання? (повідомлення інтерпретатора про помилку позначати error)\par
\begin{verbatim}
print(not 6>=7 and 6.0<=8.0 or False>2)
\end{verbatim}

{\begin {center}\textbf{Завдання 4}\end {center}}

Що буде виведено за виконання? (повідомлення інтерпретатора про помилку позначати error)\par
\begin{verbatim}
print(1 < "5.0j")
\end{verbatim}

{\begin {center}\textbf{Завдання 5}\end {center}}

Що буде виведено за виконання? (повідомлення інтерпретатора про помилку позначати error)\par
\begin{verbatim}
print(8.0 <= 8 == 3 > False)
\end{verbatim}

{\begin {center}\textbf{Завдання 6}\end {center}}

Що буде виведено за виконання? (повідомлення інтерпретатора про помилку позначати error)\par
\begin{verbatim}

def f(n):
    print("f", end="");
    return n

print(f(-1) and f(8))
\end{verbatim}

{\begin {center}\textbf{Завдання 7}\end {center}}

Що буде виведено за виконання (повідомлення інтерпретатора про помилку позначати error)?\par
\begin{verbatim}
def f(a,b):
    c=27
    if b:
        return 7
    elif a==5:
         c=1
    else: 
        return 6
    return c

print(f(-9,8))
\end{verbatim}

{\begin {center}\textbf{Завдання 8}\end {center}}

Що буде виведено за виконання (повідомлення інтерпретатора про помилку позначати error)?\par
\begin{verbatim}
a, b, c = 6, 9, 8
def f(a, b=7, c=9):
    print(a, b, c, end=" ")

f(1, 5, c=1)
print(a, b, c)
\end{verbatim}

{\begin {center}\textbf{Завдання 9}\end {center}}

Що буде виведено за виконання (повідомлення інтерпретатора про помилку позначати error)?\par
\begin{verbatim}
a,b,c=9,5,6
def h(b):
    a=5
    b-=3
    c=2
    return a+b+c

a,b,c=2,4,7
print(h(a),a,b,c)
\end{verbatim}

{\begin {center}\textbf{Завдання 10}\end {center}}


Написати функцiю f, що отримує щось та повертає щось. Невелика загадка все-таки має залишатися. Але нічого принципово нового відносно задач, що наявні у pdf-ках не планується.. Стиль запису умови також оцінюється.\par
\end {large}}
\newpage

{\Large{17.10.2024 Бета-версія: Контрольна робота \No 1, варіант  \No \textbf { 150 }\par}}
{
\begin {large}
{\begin {center}\textbf{Завдання 1}\end {center}}

Що буде виведено за виконання? (повідомлення інтерпретатора про помилку позначати error)\par
\begin{verbatim}
print(5//3*8*5)
\end{verbatim}

{\begin {center}\textbf{Завдання 2}\end {center}}

Що буде виведено за виконання? (повідомлення інтерпретатора про помилку позначати error)\par
\begin{verbatim}
print(4**-1.0*True)
\end{verbatim}

{\begin {center}\textbf{Завдання 3}\end {center}}

Що буде виведено за виконання? (повідомлення інтерпретатора про помилку позначати error)\par
\begin{verbatim}
print(9.0!=8 or not True==5 and 4<4.0)
\end{verbatim}

{\begin {center}\textbf{Завдання 4}\end {center}}

Що буде виведено за виконання? (повідомлення інтерпретатора про помилку позначати error)\par
\begin{verbatim}
print("9" >= True)
\end{verbatim}

{\begin {center}\textbf{Завдання 5}\end {center}}

Що буде виведено за виконання? (повідомлення інтерпретатора про помилку позначати error)\par
\begin{verbatim}
print(2.0 != 4 < 3.0 == 7)
\end{verbatim}

{\begin {center}\textbf{Завдання 6}\end {center}}

Що буде виведено за виконання? (повідомлення інтерпретатора про помилку позначати error)\par
\begin{verbatim}

def f(n):
    print("f", end="");
    return n>-3

print(f(-2) or f(-6))
\end{verbatim}

{\begin {center}\textbf{Завдання 7}\end {center}}

Що буде виведено за виконання (повідомлення інтерпретатора про помилку позначати error)?\par
\begin{verbatim}
def h(b):
    x=77
    if b>=1: 
        x=6
    if b==4:
         return 5
    else:
         x=0
    return x

print(h(-4))
\end{verbatim}

{\begin {center}\textbf{Завдання 8}\end {center}}

Що буде виведено за виконання (повідомлення інтерпретатора про помилку позначати error)?\par
\begin{verbatim}
a, b, c = 7, 8, 6
def h(a, b=7, c):
    print(a, b, c, end=" ")

h(a=4, 0, b=3)
print(a, b, c)
\end{verbatim}

{\begin {center}\textbf{Завдання 9}\end {center}}

Що буде виведено за виконання (повідомлення інтерпретатора про помилку позначати error)?\par
\begin{verbatim}
a,b,c=8,5,9
def f(b):
    a=2
    b=1
    c=3
    return a+b+c

a,b,c=4,9,5
print(f(b),a,b,c)
\end{verbatim}

{\begin {center}\textbf{Завдання 10}\end {center}}


Написати функцiю f, що отримує щось та повертає щось. Невелика загадка все-таки має залишатися. Але нічого принципово нового відносно задач, що наявні у pdf-ках не планується.. Стиль запису умови також оцінюється.\par
\end {large}}
\newpage

{\Large{17.10.2024 Бета-версія: Контрольна робота \No 1, варіант  \No \textbf { 151 }\par}}
{
\begin {large}
{\begin {center}\textbf{Завдання 1}\end {center}}

Що буде виведено за виконання? (повідомлення інтерпретатора про помилку позначати error)\par
\begin{verbatim}
print(6/7//5*3)
\end{verbatim}

{\begin {center}\textbf{Завдання 2}\end {center}}

Що буде виведено за виконання? (повідомлення інтерпретатора про помилку позначати error)\par
\begin{verbatim}
print(1**1.0*-2)
\end{verbatim}

{\begin {center}\textbf{Завдання 3}\end {center}}

Що буде виведено за виконання? (повідомлення інтерпретатора про помилку позначати error)\par
\begin{verbatim}
print(7<=6 and not 8!=5.0 or 3.0>False)
\end{verbatim}

{\begin {center}\textbf{Завдання 4}\end {center}}

Що буде виведено за виконання? (повідомлення інтерпретатора про помилку позначати error)\par
\begin{verbatim}
print("2.0j" > 8.0)
\end{verbatim}

{\begin {center}\textbf{Завдання 5}\end {center}}

Що буде виведено за виконання? (повідомлення інтерпретатора про помилку позначати error)\par
\begin{verbatim}
print(2 != 3.0 <= True >= 5)
\end{verbatim}

{\begin {center}\textbf{Завдання 6}\end {center}}

Що буде виведено за виконання? (повідомлення інтерпретатора про помилку позначати error)\par
\begin{verbatim}

def f(n):
    print("f", end="");
    return n

print(f(5) or f(9))
\end{verbatim}

{\begin {center}\textbf{Завдання 7}\end {center}}

Що буде виведено за виконання (повідомлення інтерпретатора про помилку позначати error)?\par
\begin{verbatim}
def h(a,b):
    c=46
    if b>-5:
        c=8
    elif a<2:
         return 9
    else: 
        return 1
    return c

print(h(-4,6))
\end{verbatim}

{\begin {center}\textbf{Завдання 8}\end {center}}

Що буде виведено за виконання (повідомлення інтерпретатора про помилку позначати error)?\par
\begin{verbatim}
a, b, c = 6, 8, 9
def g(a, b, c=6):
    print(a, b, c, end=" ")

g(a=2, b=2, c=1)
print(a, b, c)
\end{verbatim}

{\begin {center}\textbf{Завдання 9}\end {center}}

Що буде виведено за виконання (повідомлення інтерпретатора про помилку позначати error)?\par
\begin{verbatim}
a,b,c=0,8,2
def g(a):
    a=4
    b*=4
    c=5
    return a+b+c

a,b,c=5,4,1
print(g(a),a,b,c)
\end{verbatim}

{\begin {center}\textbf{Завдання 10}\end {center}}


Написати функцiю f, що отримує щось та повертає щось. Невелика загадка все-таки має залишатися. Але нічого принципово нового відносно задач, що наявні у pdf-ках не планується.. Стиль запису умови також оцінюється.\par
\end {large}}
\newpage

{\Large{17.10.2024 Бета-версія: Контрольна робота \No 1, варіант  \No \textbf { 152 }\par}}
{
\begin {large}
{\begin {center}\textbf{Завдання 1}\end {center}}

Що буде виведено за виконання? (повідомлення інтерпретатора про помилку позначати error)\par
\begin{verbatim}
print(7//4%9*8)
\end{verbatim}

{\begin {center}\textbf{Завдання 2}\end {center}}

Що буде виведено за виконання? (повідомлення інтерпретатора про помилку позначати error)\par
\begin{verbatim}
print(4**0.5*True)
\end{verbatim}

{\begin {center}\textbf{Завдання 3}\end {center}}

Що буде виведено за виконання? (повідомлення інтерпретатора про помилку позначати error)\par
\begin{verbatim}
print(not 2.0<2 or 7.0>=9 and 3==True)
\end{verbatim}

{\begin {center}\textbf{Завдання 4}\end {center}}

Що буде виведено за виконання? (повідомлення інтерпретатора про помилку позначати error)\par
\begin{verbatim}
print(2 >= "1j")
\end{verbatim}

{\begin {center}\textbf{Завдання 5}\end {center}}

Що буде виведено за виконання? (повідомлення інтерпретатора про помилку позначати error)\par
\begin{verbatim}
print(9 > 8.0  is  6.0 == 7.0)
\end{verbatim}

{\begin {center}\textbf{Завдання 6}\end {center}}

Що буде виведено за виконання? (повідомлення інтерпретатора про помилку позначати error)\par
\begin{verbatim}

def f(n):
    print("f", end="");
    return n!=3

print(f(-9) and f(0))
\end{verbatim}

{\begin {center}\textbf{Завдання 7}\end {center}}

Що буде виведено за виконання (повідомлення інтерпретатора про помилку позначати error)?\par
\begin{verbatim}
def g(a,b):
    c=95
    if a<=3:
        c=6
    if b==1:
         return 4
    else: 
        c=7
    return c

print(g(2,-7))
\end{verbatim}

{\begin {center}\textbf{Завдання 8}\end {center}}

Що буде виведено за виконання (повідомлення інтерпретатора про помилку позначати error)?\par
\begin{verbatim}
a, b, c = 9, 7, 8
def h(a, b, c):
    print(a, b, c, end=" ")

h(4, a=0)
print(a, b, c)
\end{verbatim}

{\begin {center}\textbf{Завдання 9}\end {center}}

Що буде виведено за виконання (повідомлення інтерпретатора про помилку позначати error)?\par
\begin{verbatim}
a,b,c=0,9,6
def h(a):
    a=1
    b-=2
    c=3
    return a+b+c

a,b,c=7,3,8
print(h(b),a,b,c)
\end{verbatim}

{\begin {center}\textbf{Завдання 10}\end {center}}


Написати функцiю f, що отримує щось та повертає щось. Невелика загадка все-таки має залишатися. Але нічого принципово нового відносно задач, що наявні у pdf-ках не планується.. Стиль запису умови також оцінюється.\par
\end {large}}
\newpage

{\Large{17.10.2024 Бета-версія: Контрольна робота \No 1, варіант  \No \textbf { 153 }\par}}
{
\begin {large}
{\begin {center}\textbf{Завдання 1}\end {center}}

Що буде виведено за виконання? (повідомлення інтерпретатора про помилку позначати error)\par
\begin{verbatim}
print(4/5%4*7)
\end{verbatim}

{\begin {center}\textbf{Завдання 2}\end {center}}

Що буде виведено за виконання? (повідомлення інтерпретатора про помилку позначати error)\par
\begin{verbatim}
print(9**-1.0*2)
\end{verbatim}

{\begin {center}\textbf{Завдання 3}\end {center}}

Що буде виведено за виконання? (повідомлення інтерпретатора про помилку позначати error)\par
\begin{verbatim}
print(True<=9.0 or not 6<3.0 or 2==3)
\end{verbatim}

{\begin {center}\textbf{Завдання 4}\end {center}}

Що буде виведено за виконання? (повідомлення інтерпретатора про помилку позначати error)\par
\begin{verbatim}
print("True" <= False)
\end{verbatim}

{\begin {center}\textbf{Завдання 5}\end {center}}

Що буде виведено за виконання? (повідомлення інтерпретатора про помилку позначати error)\par
\begin{verbatim}
print(6 < 7 != 8 <= False)
\end{verbatim}

{\begin {center}\textbf{Завдання 6}\end {center}}

Що буде виведено за виконання? (повідомлення інтерпретатора про помилку позначати error)\par
\begin{verbatim}

def f(n):
    print("f", end="");
    return n>=0

print(f(6) and f(-4))
\end{verbatim}

{\begin {center}\textbf{Завдання 7}\end {center}}

Що буде виведено за виконання (повідомлення інтерпретатора про помилку позначати error)?\par
\begin{verbatim}
def f(d):
    w=98
    if d: 
        return 8
    elif d<3:
         w=4
    else:
         return 3
    return w

print(f(-2))
\end{verbatim}

{\begin {center}\textbf{Завдання 8}\end {center}}

Що буде виведено за виконання (повідомлення інтерпретатора про помилку позначати error)?\par
\begin{verbatim}
a, b, c = 8, 7, 9
def f(a, b, c=6):
    print(a, b, c, end=" ")

f(3, 5)
print(a, b, c)
\end{verbatim}

{\begin {center}\textbf{Завдання 9}\end {center}}

Що буде виведено за виконання (повідомлення інтерпретатора про помилку позначати error)?\par
\begin{verbatim}
a,b,c=7,2,8
def h(a):
    a=1
    b=5
    c=2
    return a+b+c

a,b,c=3,1,0
print(h(a),a,b,c)
\end{verbatim}

{\begin {center}\textbf{Завдання 10}\end {center}}


Написати функцiю f, що отримує щось та повертає щось. Невелика загадка все-таки має залишатися. Але нічого принципово нового відносно задач, що наявні у pdf-ках не планується.. Стиль запису умови також оцінюється.\par
\end {large}}
\newpage

{\Large{17.10.2024 Бета-версія: Контрольна робота \No 1, варіант  \No \textbf { 154 }\par}}
{
\begin {large}
{\begin {center}\textbf{Завдання 1}\end {center}}

Що буде виведено за виконання? (повідомлення інтерпретатора про помилку позначати error)\par
\begin{verbatim}
print(9//9//3*6)
\end{verbatim}

{\begin {center}\textbf{Завдання 2}\end {center}}

Що буде виведено за виконання? (повідомлення інтерпретатора про помилку позначати error)\par
\begin{verbatim}
print(3**1.0--1)
\end{verbatim}

{\begin {center}\textbf{Завдання 3}\end {center}}

Що буде виведено за виконання? (повідомлення інтерпретатора про помилку позначати error)\par
\begin{verbatim}
print(not 8>=9 and False>6.0 and 5.0!=4)
\end{verbatim}

{\begin {center}\textbf{Завдання 4}\end {center}}

Що буде виведено за виконання? (повідомлення інтерпретатора про помилку позначати error)\par
\begin{verbatim}
print(True == "6")
\end{verbatim}

{\begin {center}\textbf{Завдання 5}\end {center}}

Що буде виведено за виконання? (повідомлення інтерпретатора про помилку позначати error)\par
\begin{verbatim}
print(9 > 3  is  5.0 >= 4.0)
\end{verbatim}

{\begin {center}\textbf{Завдання 6}\end {center}}

Що буде виведено за виконання? (повідомлення інтерпретатора про помилку позначати error)\par
\begin{verbatim}

def f(n):
    print("f", end="");
    return n

print(f(7) or f(4))
\end{verbatim}

{\begin {center}\textbf{Завдання 7}\end {center}}

Що буде виведено за виконання (повідомлення інтерпретатора про помилку позначати error)?\par
\begin{verbatim}
def h(a):
    v=20
    if a: 
        v=9
    if a!=0:
         v=1
    else:
         return 7
    return v

print(h(8))
\end{verbatim}

{\begin {center}\textbf{Завдання 8}\end {center}}

Що буде виведено за виконання (повідомлення інтерпретатора про помилку позначати error)?\par
\begin{verbatim}
a, b, c = 6, 9, 7
def g(a, b=8, c):
    print(a, b, c, end=" ")

g(a=4, 3, a=1)
print(a, b, c)
\end{verbatim}

{\begin {center}\textbf{Завдання 9}\end {center}}

Що буде виведено за виконання (повідомлення інтерпретатора про помилку позначати error)?\par
\begin{verbatim}
a,b,c=4,6,2
def g(a):
    a=4
    b=3
    c=2
    return a+b+c

a,b,c=3,6,5
print(g(a),a,b,c)
\end{verbatim}

{\begin {center}\textbf{Завдання 10}\end {center}}


Написати функцiю f, що отримує щось та повертає щось. Невелика загадка все-таки має залишатися. Але нічого принципово нового відносно задач, що наявні у pdf-ках не планується.. Стиль запису умови також оцінюється.\par
\end {large}}
\newpage

{\Large{17.10.2024 Бета-версія: Контрольна робота \No 1, варіант  \No \textbf { 155 }\par}}
{
\begin {large}
{\begin {center}\textbf{Завдання 1}\end {center}}

Що буде виведено за виконання? (повідомлення інтерпретатора про помилку позначати error)\par
\begin{verbatim}
print(5/3/9*8)
\end{verbatim}

{\begin {center}\textbf{Завдання 2}\end {center}}

Що буде виведено за виконання? (повідомлення інтерпретатора про помилку позначати error)\par
\begin{verbatim}
print(-3**-2+1)
\end{verbatim}

{\begin {center}\textbf{Завдання 3}\end {center}}

Що буде виведено за виконання? (повідомлення інтерпретатора про помилку позначати error)\par
\begin{verbatim}
print(not 7<5 and False>8.0 or 2.0<=8)
\end{verbatim}

{\begin {center}\textbf{Завдання 4}\end {center}}

Що буде виведено за виконання? (повідомлення інтерпретатора про помилку позначати error)\par
\begin{verbatim}
print("False" != 3)
\end{verbatim}

{\begin {center}\textbf{Завдання 5}\end {center}}

Що буде виведено за виконання? (повідомлення інтерпретатора про помилку позначати error)\par
\begin{verbatim}
print(6 < False >= 9.0 == 5)
\end{verbatim}

{\begin {center}\textbf{Завдання 6}\end {center}}

Що буде виведено за виконання? (повідомлення інтерпретатора про помилку позначати error)\par
\begin{verbatim}

def f(n):
    print("f", end="");
    return n

print(f(3) or f(-1))
\end{verbatim}

{\begin {center}\textbf{Завдання 7}\end {center}}

Що буде виведено за виконання (повідомлення інтерпретатора про помилку позначати error)?\par
\begin{verbatim}
def f(a,b):
    c=87
    if b>=-1:
        return 0
    elif a!=-3:
         c=8
    else: 
        c=2
    return c

print(f(-7,9))
\end{verbatim}

{\begin {center}\textbf{Завдання 8}\end {center}}

Що буде виведено за виконання (повідомлення інтерпретатора про помилку позначати error)?\par
\begin{verbatim}
a, b, c = 8, 9, 7
def f(a, b=6, c=7):
    print(a, b, c, end=" ")

f(c=0, b=5, b=2)
print(a, b, c)
\end{verbatim}

{\begin {center}\textbf{Завдання 9}\end {center}}

Що буде виведено за виконання (повідомлення інтерпретатора про помилку позначати error)?\par
\begin{verbatim}
a,b,c=6,9,5
def g(b):
    a+=2
    b=4
    c=1
    return a+b+c

a,b,c=2,7,3
print(g(a),a,b,c)
\end{verbatim}

{\begin {center}\textbf{Завдання 10}\end {center}}


Написати функцiю f, що отримує щось та повертає щось. Невелика загадка все-таки має залишатися. Але нічого принципово нового відносно задач, що наявні у pdf-ках не планується.. Стиль запису умови також оцінюється.\par
\end {large}}
\newpage

{\Large{17.10.2024 Бета-версія: Контрольна робота \No 1, варіант  \No \textbf { 156 }\par}}
{
\begin {large}
{\begin {center}\textbf{Завдання 1}\end {center}}

Що буде виведено за виконання? (повідомлення інтерпретатора про помилку позначати error)\par
\begin{verbatim}
print(7//7*6*6)
\end{verbatim}

{\begin {center}\textbf{Завдання 2}\end {center}}

Що буде виведено за виконання? (повідомлення інтерпретатора про помилку позначати error)\par
\begin{verbatim}
print(1**-0.5*False)
\end{verbatim}

{\begin {center}\textbf{Завдання 3}\end {center}}

Що буде виведено за виконання? (повідомлення інтерпретатора про помилку позначати error)\par
\begin{verbatim}
print(7>=4.0 or not 9!=7.0 and True==6)
\end{verbatim}

{\begin {center}\textbf{Завдання 4}\end {center}}

Що буде виведено за виконання? (повідомлення інтерпретатора про помилку позначати error)\par
\begin{verbatim}
print(6.0 < "3.0")
\end{verbatim}

{\begin {center}\textbf{Завдання 5}\end {center}}

Що буде виведено за виконання? (повідомлення інтерпретатора про помилку позначати error)\par
\begin{verbatim}
print(2.0 != 4 <= 2  is  True)
\end{verbatim}

{\begin {center}\textbf{Завдання 6}\end {center}}

Що буде виведено за виконання? (повідомлення інтерпретатора про помилку позначати error)\par
\begin{verbatim}

def f(n):
    print("f", end="");
    return n<=2

print(f(8) and f(-6))
\end{verbatim}

{\begin {center}\textbf{Завдання 7}\end {center}}

Що буде виведено за виконання (повідомлення інтерпретатора про помилку позначати error)?\par
\begin{verbatim}
def g(a,b):
    c=83
    if b!=2:
        return 3
    if a==-2:
         c=5
    else: 
        return 9
    return c

print(g(7,3))
\end{verbatim}

{\begin {center}\textbf{Завдання 8}\end {center}}

Що буде виведено за виконання (повідомлення інтерпретатора про помилку позначати error)?\par
\begin{verbatim}
a, b, c = 8, 9, 6
def h(a, b, c):
    print(a, b, c, end=" ")

h(2, c=4, b=3)
print(a, b, c)
\end{verbatim}

{\begin {center}\textbf{Завдання 9}\end {center}}

Що буде виведено за виконання (повідомлення інтерпретатора про помилку позначати error)?\par
\begin{verbatim}
a,b,c=8,4,1
def f(b):
    a=3
    b*=5
    c=2
    return a+b+c

a,b,c=0,0,1
print(f(b),a,b,c)
\end{verbatim}

{\begin {center}\textbf{Завдання 10}\end {center}}


Написати функцiю f, що отримує щось та повертає щось. Невелика загадка все-таки має залишатися. Але нічого принципово нового відносно задач, що наявні у pdf-ках не планується.. Стиль запису умови також оцінюється.\par
\end {large}}
\newpage

{\Large{17.10.2024 Бета-версія: Контрольна робота \No 1, варіант  \No \textbf { 157 }\par}}
{
\begin {large}
{\begin {center}\textbf{Завдання 1}\end {center}}

Що буде виведено за виконання? (повідомлення інтерпретатора про помилку позначати error)\par
\begin{verbatim}
print(4//6//5*5)
\end{verbatim}

{\begin {center}\textbf{Завдання 2}\end {center}}

Що буде виведено за виконання? (повідомлення інтерпретатора про помилку позначати error)\par
\begin{verbatim}
print(4**1.0*2)
\end{verbatim}

{\begin {center}\textbf{Завдання 3}\end {center}}

Що буде виведено за виконання? (повідомлення інтерпретатора про помилку позначати error)\par
\begin{verbatim}
print(not 5>=4 and True<2.0 and 3!=3.0)
\end{verbatim}

{\begin {center}\textbf{Завдання 4}\end {center}}

Що буде виведено за виконання? (повідомлення інтерпретатора про помилку позначати error)\par
\begin{verbatim}
print(9.0 <= "8j")
\end{verbatim}

{\begin {center}\textbf{Завдання 5}\end {center}}

Що буде виведено за виконання? (повідомлення інтерпретатора про помилку позначати error)\par
\begin{verbatim}
print(False > 6.0 == 7 > 4)
\end{verbatim}

{\begin {center}\textbf{Завдання 6}\end {center}}

Що буде виведено за виконання? (повідомлення інтерпретатора про помилку позначати error)\par
\begin{verbatim}

def f(n):
    print("f", end="");
    return n==-2

print(f(1) or f(-3))
\end{verbatim}

{\begin {center}\textbf{Завдання 7}\end {center}}

Що буде виведено за виконання (повідомлення інтерпретатора про помилку позначати error)?\par
\begin{verbatim}
def f(b):
    z=89
    if b>-2: 
        z=2
    elif b>=-4:
         return 6
    else:
         z=9
    return z

print(f(1))
\end{verbatim}

{\begin {center}\textbf{Завдання 8}\end {center}}

Що буде виведено за виконання (повідомлення інтерпретатора про помилку позначати error)?\par
\begin{verbatim}
a, b, c = 8, 9, 6
def h(a, b, c=7):
    print(a, b, c, end=" ")

h(1, 0, c=5)
print(a, b, c)
\end{verbatim}

{\begin {center}\textbf{Завдання 9}\end {center}}

Що буде виведено за виконання (повідомлення інтерпретатора про помилку позначати error)?\par
\begin{verbatim}
a,b,c=6,3,7
def g(a):
    a+=4
    b=5
    c=1
    return a+b+c

a,b,c=4,5,8
print(g(b),a,b,c)
\end{verbatim}

{\begin {center}\textbf{Завдання 10}\end {center}}


Написати функцiю f, що отримує щось та повертає щось. Невелика загадка все-таки має залишатися. Але нічого принципово нового відносно задач, що наявні у pdf-ках не планується.. Стиль запису умови також оцінюється.\par
\end {large}}
\newpage

{\Large{17.10.2024 Бета-версія: Контрольна робота \No 1, варіант  \No \textbf { 158 }\par}}
{
\begin {large}
{\begin {center}\textbf{Завдання 1}\end {center}}

Що буде виведено за виконання? (повідомлення інтерпретатора про помилку позначати error)\par
\begin{verbatim}
print(3/9*3*9)
\end{verbatim}

{\begin {center}\textbf{Завдання 2}\end {center}}

Що буде виведено за виконання? (повідомлення інтерпретатора про помилку позначати error)\par
\begin{verbatim}
print(-1**1.0--2)
\end{verbatim}

{\begin {center}\textbf{Завдання 3}\end {center}}

Що буде виведено за виконання? (повідомлення інтерпретатора про помилку позначати error)\par
\begin{verbatim}
print(not 9.0<=2 or 9==3 or False>8.0)
\end{verbatim}

{\begin {center}\textbf{Завдання 4}\end {center}}

Що буде виведено за виконання? (повідомлення інтерпретатора про помилку позначати error)\par
\begin{verbatim}
print("1.0" > "True")
\end{verbatim}

{\begin {center}\textbf{Завдання 5}\end {center}}

Що буде виведено за виконання? (повідомлення інтерпретатора про помилку позначати error)\par
\begin{verbatim}
print(7.0 <= 9 != 4.0 < 6)
\end{verbatim}

{\begin {center}\textbf{Завдання 6}\end {center}}

Що буде виведено за виконання? (повідомлення інтерпретатора про помилку позначати error)\par
\begin{verbatim}

def f(n):
    print("f", end="");
    return n

print(f(-5) and f(-9))
\end{verbatim}

{\begin {center}\textbf{Завдання 7}\end {center}}

Що буде виведено за виконання (повідомлення інтерпретатора про помилку позначати error)?\par
\begin{verbatim}
def g(a):
    u=34
    if a==-5: 
        return 6
    if a<-3:
         u=8
    else:
         return 2
    return u

print(g(-6))
\end{verbatim}

{\begin {center}\textbf{Завдання 8}\end {center}}

Що буде виведено за виконання (повідомлення інтерпретатора про помилку позначати error)?\par
\begin{verbatim}
a, b, c = 7, 8, 9
def f(a, b, c):
    print(a, b, c, end=" ")

f(b=0, c=2, 5)
print(a, b, c)
\end{verbatim}

{\begin {center}\textbf{Завдання 9}\end {center}}

Що буде виведено за виконання (повідомлення інтерпретатора про помилку позначати error)?\par
\begin{verbatim}
a,b,c=4,8,9
def g(b):
    a=1
    b=5
    c=4
    return a+b+c

a,b,c=1,7,0
print(g(a),a,b,c)
\end{verbatim}

{\begin {center}\textbf{Завдання 10}\end {center}}


Написати функцiю f, що отримує щось та повертає щось. Невелика загадка все-таки має залишатися. Але нічого принципово нового відносно задач, що наявні у pdf-ках не планується.. Стиль запису умови також оцінюється.\par
\end {large}}
\newpage

{\Large{17.10.2024 Бета-версія: Контрольна робота \No 1, варіант  \No \textbf { 159 }\par}}
{
\begin {large}
{\begin {center}\textbf{Завдання 1}\end {center}}

Що буде виведено за виконання? (повідомлення інтерпретатора про помилку позначати error)\par
\begin{verbatim}
print(8/5/4*3)
\end{verbatim}

{\begin {center}\textbf{Завдання 2}\end {center}}

Що буде виведено за виконання? (повідомлення інтерпретатора про помилку позначати error)\par
\begin{verbatim}
print(9**0.5*-1)
\end{verbatim}

{\begin {center}\textbf{Завдання 3}\end {center}}

Що буде виведено за виконання? (повідомлення інтерпретатора про помилку позначати error)\par
\begin{verbatim}
print(5>5.0 and not 4<=True or 7.0==6)
\end{verbatim}

{\begin {center}\textbf{Завдання 4}\end {center}}

Що буде виведено за виконання? (повідомлення інтерпретатора про помилку позначати error)\par
\begin{verbatim}
print(False < 7)
\end{verbatim}

{\begin {center}\textbf{Завдання 5}\end {center}}

Що буде виведено за виконання? (повідомлення інтерпретатора про помилку позначати error)\par
\begin{verbatim}
print(5 >= 3  is  True <= 9.0)
\end{verbatim}

{\begin {center}\textbf{Завдання 6}\end {center}}

Що буде виведено за виконання? (повідомлення інтерпретатора про помилку позначати error)\par
\begin{verbatim}

def f(n):
    print("f", end="");
    return n

print(f(0) and f(-2))
\end{verbatim}

{\begin {center}\textbf{Завдання 7}\end {center}}

Що буде виведено за виконання (повідомлення інтерпретатора про помилку позначати error)?\par
\begin{verbatim}
def f(c):
    z=94
    if c: 
        return 1
    elif c!=-1:
         z=3
    else:
         z=0
    return z

print(f(4))
\end{verbatim}

{\begin {center}\textbf{Завдання 8}\end {center}}

Що буде виведено за виконання (повідомлення інтерпретатора про помилку позначати error)?\par
\begin{verbatim}
a, b, c = 6, 7, 9
def g(a, b=6, c):
    print(a, b, c, end=" ")

g(3, a=1)
print(a, b, c)
\end{verbatim}

{\begin {center}\textbf{Завдання 9}\end {center}}

Що буде виведено за виконання (повідомлення інтерпретатора про помилку позначати error)?\par
\begin{verbatim}
a,b,c=9,2,2
def f(a):
    a*=3
    b=1
    c=3
    return a+b+c

a,b,c=0,9,7
print(f(a),a,b,c)
\end{verbatim}

{\begin {center}\textbf{Завдання 10}\end {center}}


Написати функцiю f, що отримує щось та повертає щось. Невелика загадка все-таки має залишатися. Але нічого принципово нового відносно задач, що наявні у pdf-ках не планується.. Стиль запису умови також оцінюється.\par
\end {large}}
\newpage

{\Large{17.10.2024 Бета-версія: Контрольна робота \No 1, варіант  \No \textbf { 160 }\par}}
{
\begin {large}
{\begin {center}\textbf{Завдання 1}\end {center}}

Що буде виведено за виконання? (повідомлення інтерпретатора про помилку позначати error)\par
\begin{verbatim}
print(9//4%8*4)
\end{verbatim}

{\begin {center}\textbf{Завдання 2}\end {center}}

Що буде виведено за виконання? (повідомлення інтерпретатора про помилку позначати error)\par
\begin{verbatim}
print(4**-1+True)
\end{verbatim}

{\begin {center}\textbf{Завдання 3}\end {center}}

Що буде виведено за виконання? (повідомлення інтерпретатора про помилку позначати error)\par
\begin{verbatim}
print(4.0>=False or not 7!=8 and 6.0<2)
\end{verbatim}

{\begin {center}\textbf{Завдання 4}\end {center}}

Що буде виведено за виконання? (повідомлення інтерпретатора про помилку позначати error)\par
\begin{verbatim}
print(4 == "4.0j")
\end{verbatim}

{\begin {center}\textbf{Завдання 5}\end {center}}

Що буде виведено за виконання? (повідомлення інтерпретатора про помилку позначати error)\par
\begin{verbatim}
print(False > 3.0 != 2.0 >= 2)
\end{verbatim}

{\begin {center}\textbf{Завдання 6}\end {center}}

Що буде виведено за виконання? (повідомлення інтерпретатора про помилку позначати error)\par
\begin{verbatim}

def f(n):
    print("f", end="");
    return n!=-1

print(f(5) or f(-8))
\end{verbatim}

{\begin {center}\textbf{Завдання 7}\end {center}}

Що буде виведено за виконання (повідомлення інтерпретатора про помилку позначати error)?\par
\begin{verbatim}
def h(d):
    v=93
    if d<=2: 
        return 5
    if d!=4:
         return 4
    else:
         v=7
    return v

print(h(6))
\end{verbatim}

{\begin {center}\textbf{Завдання 8}\end {center}}

Що буде виведено за виконання (повідомлення інтерпретатора про помилку позначати error)?\par
\begin{verbatim}
a, b, c = 8, 6, 8
def f(a, b=9, c=7):
    print(a, b, c, end=" ")

f(4, 5, 2)
print(a, b, c)
\end{verbatim}

{\begin {center}\textbf{Завдання 9}\end {center}}

Що буде виведено за виконання (повідомлення інтерпретатора про помилку позначати error)?\par
\begin{verbatim}
a,b,c=3,4,5
def h(b):
    a=5
    b-=4
    c=2
    return a+b+c

a,b,c=6,1,8
print(h(a),a,b,c)
\end{verbatim}

{\begin {center}\textbf{Завдання 10}\end {center}}


Написати функцiю f, що отримує щось та повертає щось. Невелика загадка все-таки має залишатися. Але нічого принципово нового відносно задач, що наявні у pdf-ках не планується.. Стиль запису умови також оцінюється.\par
\end {large}}
\newpage

{\Large{17.10.2024 Бета-версія: Контрольна робота \No 1, варіант  \No \textbf { 161 }\par}}
{
\begin {large}
{\begin {center}\textbf{Завдання 1}\end {center}}

Що буде виведено за виконання? (повідомлення інтерпретатора про помилку позначати error)\par
\begin{verbatim}
print(6/8%7*7)
\end{verbatim}

{\begin {center}\textbf{Завдання 2}\end {center}}

Що буде виведено за виконання? (повідомлення інтерпретатора про помилку позначати error)\par
\begin{verbatim}
print(1**-0.5*1)
\end{verbatim}

{\begin {center}\textbf{Завдання 3}\end {center}}

Що буде виведено за виконання? (повідомлення інтерпретатора про помилку позначати error)\par
\begin{verbatim}
print(not 9==3.0 and True<2 or 4.0>5)
\end{verbatim}

{\begin {center}\textbf{Завдання 4}\end {center}}

Що буде виведено за виконання? (повідомлення інтерпретатора про помилку позначати error)\par
\begin{verbatim}
print("False" >= 7.0)
\end{verbatim}

{\begin {center}\textbf{Завдання 5}\end {center}}

Що буде виведено за виконання? (повідомлення інтерпретатора про помилку позначати error)\par
\begin{verbatim}
print(8 == 7  is  3 < 8)
\end{verbatim}

{\begin {center}\textbf{Завдання 6}\end {center}}

Що буде виведено за виконання? (повідомлення інтерпретатора про помилку позначати error)\par
\begin{verbatim}

def f(n):
    print("f", end="");
    return n

print(f(2) and f(9))
\end{verbatim}

{\begin {center}\textbf{Завдання 7}\end {center}}

Що буде виведено за виконання (повідомлення інтерпретатора про помилку позначати error)?\par
\begin{verbatim}
def g(c):
    w=49
    if c>=-1: 
        return 1
    if c==-4:
         w=8
    else:
         return 2
    return w

print(g(-7))
\end{verbatim}

{\begin {center}\textbf{Завдання 8}\end {center}}

Що буде виведено за виконання (повідомлення інтерпретатора про помилку позначати error)?\par
\begin{verbatim}
a, b, c = 7, 9, 8
def g(a, b=6, c):
    print(a, b, c, end=" ")

g(0, c=3, b=1)
print(a, b, c)
\end{verbatim}

{\begin {center}\textbf{Завдання 9}\end {center}}

Що буде виведено за виконання (повідомлення інтерпретатора про помилку позначати error)?\par
\begin{verbatim}
a,b,c=7,2,1
def h(b):
    a=4
    b*=2
    c=1
    return a+b+c

a,b,c=3,9,8
print(h(a),a,b,c)
\end{verbatim}

{\begin {center}\textbf{Завдання 10}\end {center}}


Написати функцiю f, що отримує щось та повертає щось. Невелика загадка все-таки має залишатися. Але нічого принципово нового відносно задач, що наявні у pdf-ках не планується.. Стиль запису умови також оцінюється.\par
\end {large}}
\newpage

{\Large{17.10.2024 Бета-версія: Контрольна робота \No 1, варіант  \No \textbf { 162 }\par}}
{
\begin {large}
{\begin {center}\textbf{Завдання 1}\end {center}}

Що буде виведено за виконання? (повідомлення інтерпретатора про помилку позначати error)\par
\begin{verbatim}
print(8//4*9*3)
\end{verbatim}

{\begin {center}\textbf{Завдання 2}\end {center}}

Що буде виведено за виконання? (повідомлення інтерпретатора про помилку позначати error)\par
\begin{verbatim}
print(1**-2+False)
\end{verbatim}

{\begin {center}\textbf{Завдання 3}\end {center}}

Що буде виведено за виконання? (повідомлення інтерпретатора про помилку позначати error)\par
\begin{verbatim}
print(6<=False or not 9.0>=3 and 8!=5.0)
\end{verbatim}

{\begin {center}\textbf{Завдання 4}\end {center}}

Що буде виведено за виконання? (повідомлення інтерпретатора про помилку позначати error)\par
\begin{verbatim}
print("5" != True)
\end{verbatim}

{\begin {center}\textbf{Завдання 5}\end {center}}

Що буде виведено за виконання? (повідомлення інтерпретатора про помилку позначати error)\par
\begin{verbatim}
print(5.0 == True <= 4 > 8.0)
\end{verbatim}

{\begin {center}\textbf{Завдання 6}\end {center}}

Що буде виведено за виконання? (повідомлення інтерпретатора про помилку позначати error)\par
\begin{verbatim}

def f(n):
    print("f", end="");
    return n>=4

print(f(-7) or f(-4))
\end{verbatim}

{\begin {center}\textbf{Завдання 7}\end {center}}

Що буде виведено за виконання (повідомлення інтерпретатора про помилку позначати error)?\par
\begin{verbatim}
def h(a,b):
    c=37
    if a:
        c=7
    if b>=-5:
         return 8
    else: 
        c=5
    return c

print(h(-5,-8))
\end{verbatim}

{\begin {center}\textbf{Завдання 8}\end {center}}

Що буде виведено за виконання (повідомлення інтерпретатора про помилку позначати error)?\par
\begin{verbatim}
a, b, c = 7, 8, 6
def h(a, b=9, c=8):
    print(a, b, c, end=" ")

h(4, 4)
print(a, b, c)
\end{verbatim}

{\begin {center}\textbf{Завдання 9}\end {center}}

Що буде виведено за виконання (повідомлення інтерпретатора про помилку позначати error)?\par
\begin{verbatim}
a,b,c=0,5,4
def g(a):
    a=3
    b-=5
    c=5
    return a+b+c

a,b,c=6,3,7
print(g(a),a,b,c)
\end{verbatim}

{\begin {center}\textbf{Завдання 10}\end {center}}


Написати функцiю f, що отримує щось та повертає щось. Невелика загадка все-таки має залишатися. Але нічого принципово нового відносно задач, що наявні у pdf-ках не планується.. Стиль запису умови також оцінюється.\par
\end {large}}
\newpage

{\Large{17.10.2024 Бета-версія: Контрольна робота \No 1, варіант  \No \textbf { 163 }\par}}
{
\begin {large}
{\begin {center}\textbf{Завдання 1}\end {center}}

Що буде виведено за виконання? (повідомлення інтерпретатора про помилку позначати error)\par
\begin{verbatim}
print(4/8//7*8)
\end{verbatim}

{\begin {center}\textbf{Завдання 2}\end {center}}

Що буде виведено за виконання? (повідомлення інтерпретатора про помилку позначати error)\par
\begin{verbatim}
print(4**0.5*-2)
\end{verbatim}

{\begin {center}\textbf{Завдання 3}\end {center}}

Що буде виведено за виконання? (повідомлення інтерпретатора про помилку позначати error)\par
\begin{verbatim}
print(4!=7.0 and not 7==4 or False>6.0)
\end{verbatim}

{\begin {center}\textbf{Завдання 4}\end {center}}

Що буде виведено за виконання? (повідомлення інтерпретатора про помилку позначати error)\par
\begin{verbatim}
print("False" < True)
\end{verbatim}

{\begin {center}\textbf{Завдання 5}\end {center}}

Що буде виведено за виконання? (повідомлення інтерпретатора про помилку позначати error)\par
\begin{verbatim}
print(5  is  3.0 != 6 < False)
\end{verbatim}

{\begin {center}\textbf{Завдання 6}\end {center}}

Що буде виведено за виконання? (повідомлення інтерпретатора про помилку позначати error)\par
\begin{verbatim}

def f(n):
    print("f", end="");
    return n

print(f(-2) or f(6))
\end{verbatim}

{\begin {center}\textbf{Завдання 7}\end {center}}

Що буде виведено за виконання (повідомлення інтерпретатора про помилку позначати error)?\par
\begin{verbatim}
def g(d):
    x=87
    if d<=5: 
        x=0
    elif d>-5:
         x=9
    else:
         return 4
    return x

print(g(-4))
\end{verbatim}

{\begin {center}\textbf{Завдання 8}\end {center}}

Що буде виведено за виконання (повідомлення інтерпретатора про помилку позначати error)?\par
\begin{verbatim}
a, b, c = 7, 9, 6
def f(a, b, c):
    print(a, b, c, end=" ")

f(a=5, b=3, c=1)
print(a, b, c)
\end{verbatim}

{\begin {center}\textbf{Завдання 9}\end {center}}

Що буде виведено за виконання (повідомлення інтерпретатора про помилку позначати error)?\par
\begin{verbatim}
a,b,c=2,0,8
def f(b):
    a+=4
    b=2
    c=3
    return a+b+c

a,b,c=4,1,6
print(f(b),a,b,c)
\end{verbatim}

{\begin {center}\textbf{Завдання 10}\end {center}}


Написати функцiю f, що отримує щось та повертає щось. Невелика загадка все-таки має залишатися. Але нічого принципово нового відносно задач, що наявні у pdf-ках не планується.. Стиль запису умови також оцінюється.\par
\end {large}}
\newpage

{\Large{17.10.2024 Бета-версія: Контрольна робота \No 1, варіант  \No \textbf { 164 }\par}}
{
\begin {large}
{\begin {center}\textbf{Завдання 1}\end {center}}

Що буде виведено за виконання? (повідомлення інтерпретатора про помилку позначати error)\par
\begin{verbatim}
print(3//9/5*4)
\end{verbatim}

{\begin {center}\textbf{Завдання 2}\end {center}}

Що буде виведено за виконання? (повідомлення інтерпретатора про помилку позначати error)\par
\begin{verbatim}
print(2**2.0-2)
\end{verbatim}

{\begin {center}\textbf{Завдання 3}\end {center}}

Що буде виведено за виконання? (повідомлення інтерпретатора про помилку позначати error)\par
\begin{verbatim}
print(not 5<=8 or 2>=True and 8.0<2.0)
\end{verbatim}

{\begin {center}\textbf{Завдання 4}\end {center}}

Що буде виведено за виконання? (повідомлення інтерпретатора про помилку позначати error)\par
\begin{verbatim}
print("3j" != 7.0)
\end{verbatim}

{\begin {center}\textbf{Завдання 5}\end {center}}

Що буде виведено за виконання? (повідомлення інтерпретатора про помилку позначати error)\par
\begin{verbatim}
print(2.0 >= 9 > True <= 2)
\end{verbatim}

{\begin {center}\textbf{Завдання 6}\end {center}}

Що буде виведено за виконання? (повідомлення інтерпретатора про помилку позначати error)\par
\begin{verbatim}

def f(n):
    print("f", end="");
    return n>0

print(f(3) and f(2))
\end{verbatim}

{\begin {center}\textbf{Завдання 7}\end {center}}

Що буде виведено за виконання (повідомлення інтерпретатора про помилку позначати error)?\par
\begin{verbatim}
def h(a):
    u=95
    if a<-2: 
        return 5
    if a<=1:
         u=7
    else:
         return 6
    return u

print(h(-5))
\end{verbatim}

{\begin {center}\textbf{Завдання 8}\end {center}}

Що буде виведено за виконання (повідомлення інтерпретатора про помилку позначати error)?\par
\begin{verbatim}
a, b, c = 8, 6, 7
def h(a, b, c=9):
    print(a, b, c, end=" ")

h(2, 0, c=3)
print(a, b, c)
\end{verbatim}

{\begin {center}\textbf{Завдання 9}\end {center}}

Що буде виведено за виконання (повідомлення інтерпретатора про помилку позначати error)?\par
\begin{verbatim}
a,b,c=9,5,2
def f(a):
    a-=1
    b=3
    c=4
    return a+b+c

a,b,c=1,8,9
print(f(a),a,b,c)
\end{verbatim}

{\begin {center}\textbf{Завдання 10}\end {center}}


Написати функцiю f, що отримує щось та повертає щось. Невелика загадка все-таки має залишатися. Але нічого принципово нового відносно задач, що наявні у pdf-ках не планується.. Стиль запису умови також оцінюється.\par
\end {large}}
\newpage

{\Large{17.10.2024 Бета-версія: Контрольна робота \No 1, варіант  \No \textbf { 165 }\par}}
{
\begin {large}
{\begin {center}\textbf{Завдання 1}\end {center}}

Що буде виведено за виконання? (повідомлення інтерпретатора про помилку позначати error)\par
\begin{verbatim}
print(9/6/4*6)
\end{verbatim}

{\begin {center}\textbf{Завдання 2}\end {center}}

Що буде виведено за виконання? (повідомлення інтерпретатора про помилку позначати error)\par
\begin{verbatim}
print(2**-2-True)
\end{verbatim}

{\begin {center}\textbf{Завдання 3}\end {center}}

Що буде виведено за виконання? (повідомлення інтерпретатора про помилку позначати error)\par
\begin{verbatim}
print(not True<=7 or 3!=6 and 3.0>=4.0)
\end{verbatim}

{\begin {center}\textbf{Завдання 4}\end {center}}

Що буде виведено за виконання? (повідомлення інтерпретатора про помилку позначати error)\par
\begin{verbatim}
print(1 <= "1.0")
\end{verbatim}

{\begin {center}\textbf{Завдання 5}\end {center}}

Що буде виведено за виконання? (повідомлення інтерпретатора про помилку позначати error)\par
\begin{verbatim}
print(5 == 8.0  is  8 >= 4.0)
\end{verbatim}

{\begin {center}\textbf{Завдання 6}\end {center}}

Що буде виведено за виконання? (повідомлення інтерпретатора про помилку позначати error)\par
\begin{verbatim}

def f(n):
    print("f", end="");
    return n

print(f(7) and f(-5))
\end{verbatim}

{\begin {center}\textbf{Завдання 7}\end {center}}

Що буде виведено за виконання (повідомлення інтерпретатора про помилку позначати error)?\par
\begin{verbatim}
def f(b):
    y=66
    if b: 
        y=3
    elif b<3:
         return 9
    else:
         y=6
    return y

print(f(-9))
\end{verbatim}

{\begin {center}\textbf{Завдання 8}\end {center}}

Що буде виведено за виконання (повідомлення інтерпретатора про помилку позначати error)?\par
\begin{verbatim}
a, b, c = 9, 6, 8
def g(a, b=7, c=7):
    print(a, b, c, end=" ")

g(2, a=0)
print(a, b, c)
\end{verbatim}

{\begin {center}\textbf{Завдання 9}\end {center}}

Що буде виведено за виконання (повідомлення інтерпретатора про помилку позначати error)?\par
\begin{verbatim}
a,b,c=4,3,7
def g(b):
    a+=5
    b=1
    c=2
    return a+b+c

a,b,c=6,0,5
print(g(b),a,b,c)
\end{verbatim}

{\begin {center}\textbf{Завдання 10}\end {center}}


Написати функцiю f, що отримує щось та повертає щось. Невелика загадка все-таки має залишатися. Але нічого принципово нового відносно задач, що наявні у pdf-ках не планується.. Стиль запису умови також оцінюється.\par
\end {large}}
\newpage

{\Large{17.10.2024 Бета-версія: Контрольна робота \No 1, варіант  \No \textbf { 166 }\par}}
{
\begin {large}
{\begin {center}\textbf{Завдання 1}\end {center}}

Що буде виведено за виконання? (повідомлення інтерпретатора про помилку позначати error)\par
\begin{verbatim}
print(6//7%8*7)
\end{verbatim}

{\begin {center}\textbf{Завдання 2}\end {center}}

Що буде виведено за виконання? (повідомлення інтерпретатора про помилку позначати error)\par
\begin{verbatim}
print(9**1.0*False)
\end{verbatim}

{\begin {center}\textbf{Завдання 3}\end {center}}

Що буде виведено за виконання? (повідомлення інтерпретатора про помилку позначати error)\par
\begin{verbatim}
print(False>9 and not 6<6.0 or 7.0==8)
\end{verbatim}

{\begin {center}\textbf{Завдання 4}\end {center}}

Що буде виведено за виконання? (повідомлення інтерпретатора про помилку позначати error)\par
\begin{verbatim}
print(False == 2)
\end{verbatim}

{\begin {center}\textbf{Завдання 5}\end {center}}

Що буде виведено за виконання? (повідомлення інтерпретатора про помилку позначати error)\par
\begin{verbatim}
print(6.0 < 9 != 7 >= 9.0)
\end{verbatim}

{\begin {center}\textbf{Завдання 6}\end {center}}

Що буде виведено за виконання? (повідомлення інтерпретатора про помилку позначати error)\par
\begin{verbatim}

def f(n):
    print("f", end="");
    return n<=-3

print(f(4) or f(-1))
\end{verbatim}

{\begin {center}\textbf{Завдання 7}\end {center}}

Що буде виведено за виконання (повідомлення інтерпретатора про помилку позначати error)?\par
\begin{verbatim}
def h(a,b):
    c=43
    if b>0:
        return 1
    elif a<4:
         return 3
    else: 
        c=4
    return c

print(h(-3,-4))
\end{verbatim}

{\begin {center}\textbf{Завдання 8}\end {center}}

Що буде виведено за виконання (повідомлення інтерпретатора про помилку позначати error)?\par
\begin{verbatim}
a, b, c = 9, 6, 8
def f(a, b, c=7):
    print(a, b, c, end=" ")

f(b=4, c=1, 5)
print(a, b, c)
\end{verbatim}

{\begin {center}\textbf{Завдання 9}\end {center}}

Що буде виведено за виконання (повідомлення інтерпретатора про помилку позначати error)?\par
\begin{verbatim}
a,b,c=1,0,2
def h(a):
    a=3
    b*=1
    c=5
    return a+b+c

a,b,c=6,5,9
print(h(b),a,b,c)
\end{verbatim}

{\begin {center}\textbf{Завдання 10}\end {center}}


Написати функцiю f, що отримує щось та повертає щось. Невелика загадка все-таки має залишатися. Але нічого принципово нового відносно задач, що наявні у pdf-ках не планується.. Стиль запису умови також оцінюється.\par
\end {large}}
\newpage

{\Large{17.10.2024 Бета-версія: Контрольна робота \No 1, варіант  \No \textbf { 167 }\par}}
{
\begin {large}
{\begin {center}\textbf{Завдання 1}\end {center}}

Що буде виведено за виконання? (повідомлення інтерпретатора про помилку позначати error)\par
\begin{verbatim}
print(5//3//6*5)
\end{verbatim}

{\begin {center}\textbf{Завдання 2}\end {center}}

Що буде виведено за виконання? (повідомлення інтерпретатора про помилку позначати error)\par
\begin{verbatim}
print(1**-1.0*1)
\end{verbatim}

{\begin {center}\textbf{Завдання 3}\end {center}}

Що буде виведено за виконання? (повідомлення інтерпретатора про помилку позначати error)\par
\begin{verbatim}
print(not True<7 and 8.0<=3 and 2.0==2)
\end{verbatim}

{\begin {center}\textbf{Завдання 4}\end {center}}

Що буде виведено за виконання? (повідомлення інтерпретатора про помилку позначати error)\par
\begin{verbatim}
print("9.0j" >= "6")
\end{verbatim}

{\begin {center}\textbf{Завдання 5}\end {center}}

Що буде виведено за виконання? (повідомлення інтерпретатора про помилку позначати error)\par
\begin{verbatim}
print(6 < False == 5.0 != 7.0)
\end{verbatim}

{\begin {center}\textbf{Завдання 6}\end {center}}

Що буде виведено за виконання? (повідомлення інтерпретатора про помилку позначати error)\par
\begin{verbatim}

def f(n):
    print("f", end="");
    return n

print(f(-8) or f(-3))
\end{verbatim}

{\begin {center}\textbf{Завдання 7}\end {center}}

Що буде виведено за виконання (повідомлення інтерпретатора про помилку позначати error)?\par
\begin{verbatim}
def g(c):
    v=14
    if c: 
        return 4
    elif c>0:
         v=8
    else:
         return 2
    return v

print(g(5))
\end{verbatim}

{\begin {center}\textbf{Завдання 8}\end {center}}

Що буде виведено за виконання (повідомлення інтерпретатора про помилку позначати error)?\par
\begin{verbatim}
a, b, c = 9, 6, 8
def g(a, b=9, c):
    print(a, b, c, end=" ")

g(4, 5, 2)
print(a, b, c)
\end{verbatim}

{\begin {center}\textbf{Завдання 9}\end {center}}

Що буде виведено за виконання (повідомлення інтерпретатора про помилку позначати error)?\par
\begin{verbatim}
a,b,c=4,3,8
def f(a):
    a=2
    b-=4
    c=1
    return a+b+c

a,b,c=7,6,1
print(f(b),a,b,c)
\end{verbatim}

{\begin {center}\textbf{Завдання 10}\end {center}}


Написати функцiю f, що отримує щось та повертає щось. Невелика загадка все-таки має залишатися. Але нічого принципово нового відносно задач, що наявні у pdf-ках не планується.. Стиль запису умови також оцінюється.\par
\end {large}}
\newpage

{\Large{17.10.2024 Бета-версія: Контрольна робота \No 1, варіант  \No \textbf { 168 }\par}}
{
\begin {large}
{\begin {center}\textbf{Завдання 1}\end {center}}

Що буде виведено за виконання? (повідомлення інтерпретатора про помилку позначати error)\par
\begin{verbatim}
print(7/5*3*9)
\end{verbatim}

{\begin {center}\textbf{Завдання 2}\end {center}}

Що буде виведено за виконання? (повідомлення інтерпретатора про помилку позначати error)\par
\begin{verbatim}
print(-2**2.0+-1)
\end{verbatim}

{\begin {center}\textbf{Завдання 3}\end {center}}

Що буде виведено за виконання? (повідомлення інтерпретатора про помилку позначати error)\par
\begin{verbatim}
print(not 5>4 or False!=5.0 or 9>=9.0)
\end{verbatim}

{\begin {center}\textbf{Завдання 4}\end {center}}

Що буде виведено за виконання? (повідомлення інтерпретатора про помилку позначати error)\par
\begin{verbatim}
print("True" > 8.0)
\end{verbatim}

{\begin {center}\textbf{Завдання 5}\end {center}}

Що буде виведено за виконання? (повідомлення інтерпретатора про помилку позначати error)\par
\begin{verbatim}
print(2 <= 3 > True  is  4)
\end{verbatim}

{\begin {center}\textbf{Завдання 6}\end {center}}

Що буде виведено за виконання? (повідомлення інтерпретатора про помилку позначати error)\par
\begin{verbatim}

def f(n):
    print("f", end="");
    return n<2

print(f(1) and f(9))
\end{verbatim}

{\begin {center}\textbf{Завдання 7}\end {center}}

Що буде виведено за виконання (повідомлення інтерпретатора про помилку позначати error)?\par
\begin{verbatim}
def h(d):
    y=13
    if d>=-3: 
        y=1
    if d!=2:
         y=7
    else:
         return 0
    return y

print(h(3))
\end{verbatim}

{\begin {center}\textbf{Завдання 8}\end {center}}

Що буде виведено за виконання (повідомлення інтерпретатора про помилку позначати error)?\par
\begin{verbatim}
a, b, c = 7, 8, 6
def h(a, b, c):
    print(a, b, c, end=" ")

h(a=1, 3, b=0)
print(a, b, c)
\end{verbatim}

{\begin {center}\textbf{Завдання 9}\end {center}}

Що буде виведено за виконання (повідомлення інтерпретатора про помилку позначати error)?\par
\begin{verbatim}
a,b,c=5,4,8
def g(b):
    a+=4
    b=5
    c=2
    return a+b+c

a,b,c=2,9,3
print(g(b),a,b,c)
\end{verbatim}

{\begin {center}\textbf{Завдання 10}\end {center}}


Написати функцiю f, що отримує щось та повертає щось. Невелика загадка все-таки має залишатися. Але нічого принципово нового відносно задач, що наявні у pdf-ках не планується.. Стиль запису умови також оцінюється.\par
\end {large}}
\newpage

{\Large{17.10.2024 Бета-версія: Контрольна робота \No 1, варіант  \No \textbf { 169 }\par}}
{
\begin {large}
{\begin {center}\textbf{Завдання 1}\end {center}}

Що буде виведено за виконання? (повідомлення інтерпретатора про помилку позначати error)\par
\begin{verbatim}
print(6/7/6*7)
\end{verbatim}

{\begin {center}\textbf{Завдання 2}\end {center}}

Що буде виведено за виконання? (повідомлення інтерпретатора про помилку позначати error)\par
\begin{verbatim}
print(1**0.5*-2)
\end{verbatim}

{\begin {center}\textbf{Завдання 3}\end {center}}

Що буде виведено за виконання? (повідомлення інтерпретатора про помилку позначати error)\par
\begin{verbatim}
print(not 5.0!=7 and 9>9.0 and 5==True)
\end{verbatim}

{\begin {center}\textbf{Завдання 4}\end {center}}

Що буде виведено за виконання? (повідомлення інтерпретатора про помилку позначати error)\par
\begin{verbatim}
print("6.0" != True)
\end{verbatim}

{\begin {center}\textbf{Завдання 5}\end {center}}

Що буде виведено за виконання? (повідомлення інтерпретатора про помилку позначати error)\par
\begin{verbatim}
print(4 != 6 >= 6.0  is  False)
\end{verbatim}

{\begin {center}\textbf{Завдання 6}\end {center}}

Що буде виведено за виконання? (повідомлення інтерпретатора про помилку позначати error)\par
\begin{verbatim}

def f(n):
    print("f", end="");
    return n==4

print(f(-6) or f(0))
\end{verbatim}

{\begin {center}\textbf{Завдання 7}\end {center}}

Що буде виведено за виконання (повідомлення інтерпретатора про помилку позначати error)?\par
\begin{verbatim}
def f(a,b):
    c=73
    if a:
        c=2
    elif b<=-4:
         return 9
    else: 
        return 0
    return c

print(f(-8,4))
\end{verbatim}

{\begin {center}\textbf{Завдання 8}\end {center}}

Що буде виведено за виконання (повідомлення інтерпретатора про помилку позначати error)?\par
\begin{verbatim}
a, b, c = 8, 9, 6
def g(a, b=7, c):
    print(a, b, c, end=" ")

g(a=1, 3, b=0)
print(a, b, c)
\end{verbatim}

{\begin {center}\textbf{Завдання 9}\end {center}}

Що буде виведено за виконання (повідомлення інтерпретатора про помилку позначати error)?\par
\begin{verbatim}
a,b,c=2,8,9
def h(a):
    a=4
    b=1
    c=5
    return a+b+c

a,b,c=4,1,7
print(h(b),a,b,c)
\end{verbatim}

{\begin {center}\textbf{Завдання 10}\end {center}}


Написати функцiю f, що отримує щось та повертає щось. Невелика загадка все-таки має залишатися. Але нічого принципово нового відносно задач, що наявні у pdf-ках не планується.. Стиль запису умови також оцінюється.\par
\end {large}}
\newpage

{\Large{17.10.2024 Бета-версія: Контрольна робота \No 1, варіант  \No \textbf { 170 }\par}}
{
\begin {large}
{\begin {center}\textbf{Завдання 1}\end {center}}

Що буде виведено за виконання? (повідомлення інтерпретатора про помилку позначати error)\par
\begin{verbatim}
print(7//4%4*4)
\end{verbatim}

{\begin {center}\textbf{Завдання 2}\end {center}}

Що буде виведено за виконання? (повідомлення інтерпретатора про помилку позначати error)\par
\begin{verbatim}
print(9**-0.5*-1)
\end{verbatim}

{\begin {center}\textbf{Завдання 3}\end {center}}

Що буде виведено за виконання? (повідомлення інтерпретатора про помилку позначати error)\par
\begin{verbatim}
print(2<=8 or not 4>=4.0 or 6.0<False)
\end{verbatim}

{\begin {center}\textbf{Завдання 4}\end {center}}

Що буде виведено за виконання? (повідомлення інтерпретатора про помилку позначати error)\par
\begin{verbatim}
print("False" >= 9)
\end{verbatim}

{\begin {center}\textbf{Завдання 5}\end {center}}

Що буде виведено за виконання? (повідомлення інтерпретатора про помилку позначати error)\par
\begin{verbatim}
print(5 <= 4.0 > 3.0 < 3)
\end{verbatim}

{\begin {center}\textbf{Завдання 6}\end {center}}

Що буде виведено за виконання? (повідомлення інтерпретатора про помилку позначати error)\par
\begin{verbatim}

def f(n):
    print("f", end="");
    return n

print(f(8) and f(-7))
\end{verbatim}

{\begin {center}\textbf{Завдання 7}\end {center}}

Що буде виведено за виконання (повідомлення інтерпретатора про помилку позначати error)?\par
\begin{verbatim}
def g(a,b):
    c=93
    if a<5:
        c=6
    if b>3:
         c=8
    else: 
        return 6
    return c

print(g(1,-5))
\end{verbatim}

{\begin {center}\textbf{Завдання 8}\end {center}}

Що буде виведено за виконання (повідомлення інтерпретатора про помилку позначати error)?\par
\begin{verbatim}
a, b, c = 6, 9, 8
def h(a, b, c=7):
    print(a, b, c, end=" ")

h(2, 5, a=4)
print(a, b, c)
\end{verbatim}

{\begin {center}\textbf{Завдання 9}\end {center}}

Що буде виведено за виконання (повідомлення інтерпретатора про помилку позначати error)?\par
\begin{verbatim}
a,b,c=7,0,5
def h(a):
    a=3
    b*=1
    c=4
    return a+b+c

a,b,c=3,1,6
print(h(b),a,b,c)
\end{verbatim}

{\begin {center}\textbf{Завдання 10}\end {center}}


Написати функцiю f, що отримує щось та повертає щось. Невелика загадка все-таки має залишатися. Але нічого принципово нового відносно задач, що наявні у pdf-ках не планується.. Стиль запису умови також оцінюється.\par
\end {large}}
\newpage

{\Large{17.10.2024 Бета-версія: Контрольна робота \No 1, варіант  \No \textbf { 171 }\par}}
{
\begin {large}
{\begin {center}\textbf{Завдання 1}\end {center}}

Що буде виведено за виконання? (повідомлення інтерпретатора про помилку позначати error)\par
\begin{verbatim}
print(8/9//9*3)
\end{verbatim}

{\begin {center}\textbf{Завдання 2}\end {center}}

Що буде виведено за виконання? (повідомлення інтерпретатора про помилку позначати error)\par
\begin{verbatim}
print(-2**False+False)
\end{verbatim}

{\begin {center}\textbf{Завдання 3}\end {center}}

Що буде виведено за виконання? (повідомлення інтерпретатора про помилку позначати error)\par
\begin{verbatim}
print(6>=3 or not 9==8.0 or True!=7.0)
\end{verbatim}

{\begin {center}\textbf{Завдання 4}\end {center}}

Що буде виведено за виконання? (повідомлення інтерпретатора про помилку позначати error)\par
\begin{verbatim}
print("5j" <= 4.0)
\end{verbatim}

{\begin {center}\textbf{Завдання 5}\end {center}}

Що буде виведено за виконання? (повідомлення інтерпретатора про помилку позначати error)\par
\begin{verbatim}
print(5.0 == 7 != True >= 8)
\end{verbatim}

{\begin {center}\textbf{Завдання 6}\end {center}}

Що буде виведено за виконання? (повідомлення інтерпретатора про помилку позначати error)\par
\begin{verbatim}

def f(n):
    print("f", end="");
    return n

print(f(5) or f(-9))
\end{verbatim}

{\begin {center}\textbf{Завдання 7}\end {center}}

Що буде виведено за виконання (повідомлення інтерпретатора про помилку позначати error)?\par
\begin{verbatim}
def g(a,b):
    c=59
    if b>=-3:
        return 5
    if a==1:
         c=1
    else: 
        c=2
    return c

print(g(5,8))
\end{verbatim}

{\begin {center}\textbf{Завдання 8}\end {center}}

Що буде виведено за виконання (повідомлення інтерпретатора про помилку позначати error)?\par
\begin{verbatim}
a, b, c = 6, 7, 8
def f(a, b=9, c=8):
    print(a, b, c, end=" ")

f(3, c=2, b=0)
print(a, b, c)
\end{verbatim}

{\begin {center}\textbf{Завдання 9}\end {center}}

Що буде виведено за виконання (повідомлення інтерпретатора про помилку позначати error)?\par
\begin{verbatim}
a,b,c=0,3,5
def f(b):
    a=3
    b=2
    c=3
    return a+b+c

a,b,c=6,8,1
print(f(b),a,b,c)
\end{verbatim}

{\begin {center}\textbf{Завдання 10}\end {center}}


Написати функцiю f, що отримує щось та повертає щось. Невелика загадка все-таки має залишатися. Але нічого принципово нового відносно задач, що наявні у pdf-ках не планується.. Стиль запису умови також оцінюється.\par
\end {large}}
\newpage

{\Large{17.10.2024 Бета-версія: Контрольна робота \No 1, варіант  \No \textbf { 172 }\par}}
{
\begin {large}
{\begin {center}\textbf{Завдання 1}\end {center}}

Що буде виведено за виконання? (повідомлення інтерпретатора про помилку позначати error)\par
\begin{verbatim}
print(3//8*8*6)
\end{verbatim}

{\begin {center}\textbf{Завдання 2}\end {center}}

Що буде виведено за виконання? (повідомлення інтерпретатора про помилку позначати error)\par
\begin{verbatim}
print(-3**-4-1)
\end{verbatim}

{\begin {center}\textbf{Завдання 3}\end {center}}

Що буде виведено за виконання? (повідомлення інтерпретатора про помилку позначати error)\par
\begin{verbatim}
print(not 3.0<=7 and 5>3 and 2.0<False)
\end{verbatim}

{\begin {center}\textbf{Завдання 4}\end {center}}

Що буде виведено за виконання? (повідомлення інтерпретатора про помилку позначати error)\par
\begin{verbatim}
print("True" == "8")
\end{verbatim}

{\begin {center}\textbf{Завдання 5}\end {center}}

Що буде виведено за виконання? (повідомлення інтерпретатора про помилку позначати error)\par
\begin{verbatim}
print(False > 9 < 2.0  is  2)
\end{verbatim}

{\begin {center}\textbf{Завдання 6}\end {center}}

Що буде виведено за виконання? (повідомлення інтерпретатора про помилку позначати error)\par
\begin{verbatim}

def f(n):
    print("f", end="");
    return n<1

print(f(6) and f(-9))
\end{verbatim}

{\begin {center}\textbf{Завдання 7}\end {center}}

Що буде виведено за виконання (повідомлення інтерпретатора про помилку позначати error)?\par
\begin{verbatim}
def h(a,b):
    c=92
    if a:
        return 0
    elif b!=-1:
         return 7
    else: 
        c=9
    return c

print(h(4,-9))
\end{verbatim}

{\begin {center}\textbf{Завдання 8}\end {center}}

Що буде виведено за виконання (повідомлення інтерпретатора про помилку позначати error)?\par
\begin{verbatim}
a, b, c = 6, 7, 9
def g(a, b, c):
    print(a, b, c, end=" ")

g(1, 4, 5)
print(a, b, c)
\end{verbatim}

{\begin {center}\textbf{Завдання 9}\end {center}}

Що буде виведено за виконання (повідомлення інтерпретатора про помилку позначати error)?\par
\begin{verbatim}
a,b,c=8,7,2
def g(b):
    a*=3
    b=5
    c=2
    return a+b+c

a,b,c=4,0,9
print(g(a),a,b,c)
\end{verbatim}

{\begin {center}\textbf{Завдання 10}\end {center}}


Написати функцiю f, що отримує щось та повертає щось. Невелика загадка все-таки має залишатися. Але нічого принципово нового відносно задач, що наявні у pdf-ках не планується.. Стиль запису умови також оцінюється.\par
\end {large}}
\newpage

{\Large{17.10.2024 Бета-версія: Контрольна робота \No 1, варіант  \No \textbf { 173 }\par}}
{
\begin {large}
{\begin {center}\textbf{Завдання 1}\end {center}}

Що буде виведено за виконання? (повідомлення інтерпретатора про помилку позначати error)\par
\begin{verbatim}
print(9/3//5*9)
\end{verbatim}

{\begin {center}\textbf{Завдання 2}\end {center}}

Що буде виведено за виконання? (повідомлення інтерпретатора про помилку позначати error)\par
\begin{verbatim}
print(4**1.0*True)
\end{verbatim}

{\begin {center}\textbf{Завдання 3}\end {center}}

Що буде виведено за виконання? (повідомлення інтерпретатора про помилку позначати error)\par
\begin{verbatim}
print(2>=8 or not 4>8.0 and True==4.0)
\end{verbatim}

{\begin {center}\textbf{Завдання 4}\end {center}}

Що буде виведено за виконання? (повідомлення інтерпретатора про помилку позначати error)\par
\begin{verbatim}
print(4 > "3.0j")
\end{verbatim}

{\begin {center}\textbf{Завдання 5}\end {center}}

Що буде виведено за виконання? (повідомлення інтерпретатора про помилку позначати error)\par
\begin{verbatim}
print(9 <= 7.0 == True < 8)
\end{verbatim}

{\begin {center}\textbf{Завдання 6}\end {center}}

Що буде виведено за виконання? (повідомлення інтерпретатора про помилку позначати error)\par
\begin{verbatim}

def f(n):
    print("f", end="");
    return n

print(f(5) and f(-1))
\end{verbatim}

{\begin {center}\textbf{Завдання 7}\end {center}}

Що буде виведено за виконання (повідомлення інтерпретатора про помилку позначати error)?\par
\begin{verbatim}
def f(b):
    z=51
    if b==-5: 
        z=5
    if b>-4:
         return 3
    else:
         return 1
    return z

print(f(-3))
\end{verbatim}

{\begin {center}\textbf{Завдання 8}\end {center}}

Що буде виведено за виконання (повідомлення інтерпретатора про помилку позначати error)?\par
\begin{verbatim}
a, b, c = 6, 9, 7
def f(a, b=8, c):
    print(a, b, c, end=" ")

f(c=4, b=2, c=1)
print(a, b, c)
\end{verbatim}

{\begin {center}\textbf{Завдання 9}\end {center}}

Що буде виведено за виконання (повідомлення інтерпретатора про помилку позначати error)?\par
\begin{verbatim}
a,b,c=5,2,4
def f(a):
    a-=2
    b=5
    c=1
    return a+b+c

a,b,c=7,0,9
print(f(a),a,b,c)
\end{verbatim}

{\begin {center}\textbf{Завдання 10}\end {center}}


Написати функцiю f, що отримує щось та повертає щось. Невелика загадка все-таки має залишатися. Але нічого принципово нового відносно задач, що наявні у pdf-ках не планується.. Стиль запису умови також оцінюється.\par
\end {large}}
\newpage

{\Large{17.10.2024 Бета-версія: Контрольна робота \No 1, варіант  \No \textbf { 174 }\par}}
{
\begin {large}
{\begin {center}\textbf{Завдання 1}\end {center}}

Що буде виведено за виконання? (повідомлення інтерпретатора про помилку позначати error)\par
\begin{verbatim}
print(5//5%7*8)
\end{verbatim}

{\begin {center}\textbf{Завдання 2}\end {center}}

Що буде виведено за виконання? (повідомлення інтерпретатора про помилку позначати error)\par
\begin{verbatim}
print(3**-1-2)
\end{verbatim}

{\begin {center}\textbf{Завдання 3}\end {center}}

Що буде виведено за виконання? (повідомлення інтерпретатора про помилку позначати error)\par
\begin{verbatim}
print(False<6 and not 2!=2.0 or 8<=3.0)
\end{verbatim}

{\begin {center}\textbf{Завдання 4}\end {center}}

Що буде виведено за виконання? (повідомлення інтерпретатора про помилку позначати error)\par
\begin{verbatim}
print(2.0 < False)
\end{verbatim}

{\begin {center}\textbf{Завдання 5}\end {center}}

Що буде виведено за виконання? (повідомлення інтерпретатора про помилку позначати error)\par
\begin{verbatim}
print(8.0 != 2 == 9.0 >= 4)
\end{verbatim}

{\begin {center}\textbf{Завдання 6}\end {center}}

Що буде виведено за виконання? (повідомлення інтерпретатора про помилку позначати error)\par
\begin{verbatim}

def f(n):
    print("f", end="");
    return n>=3

print(f(4) or f(-7))
\end{verbatim}

{\begin {center}\textbf{Завдання 7}\end {center}}

Що буде виведено за виконання (повідомлення інтерпретатора про помилку позначати error)?\par
\begin{verbatim}
def f(a):
    w=43
    if a<=2: 
        w=6
    elif a==3:
         w=9
    else:
         return 8
    return w

print(f(-8))
\end{verbatim}

{\begin {center}\textbf{Завдання 8}\end {center}}

Що буде виведено за виконання (повідомлення інтерпретатора про помилку позначати error)?\par
\begin{verbatim}
a, b, c = 8, 7, 9
def h(a, b=6, c=9):
    print(a, b, c, end=" ")

h(b=5, c=0, 3)
print(a, b, c)
\end{verbatim}

{\begin {center}\textbf{Завдання 9}\end {center}}

Що буде виведено за виконання (повідомлення інтерпретатора про помилку позначати error)?\par
\begin{verbatim}
a,b,c=8,1,6
def h(b):
    a=4
    b+=3
    c=3
    return a+b+c

a,b,c=3,4,7
print(h(a),a,b,c)
\end{verbatim}

{\begin {center}\textbf{Завдання 10}\end {center}}


Написати функцiю f, що отримує щось та повертає щось. Невелика загадка все-таки має залишатися. Але нічого принципово нового відносно задач, що наявні у pdf-ках не планується.. Стиль запису умови також оцінюється.\par
\end {large}}
\newpage

{\Large{17.10.2024 Бета-версія: Контрольна робота \No 1, варіант  \No \textbf { 175 }\par}}
{
\begin {large}
{\begin {center}\textbf{Завдання 1}\end {center}}

Що буде виведено за виконання? (повідомлення інтерпретатора про помилку позначати error)\par
\begin{verbatim}
print(4//6/3*5)
\end{verbatim}

{\begin {center}\textbf{Завдання 2}\end {center}}

Що буде виведено за виконання? (повідомлення інтерпретатора про помилку позначати error)\par
\begin{verbatim}
print(3**4+1)
\end{verbatim}

{\begin {center}\textbf{Завдання 3}\end {center}}

Що буде виведено за виконання? (повідомлення інтерпретатора про помилку позначати error)\par
\begin{verbatim}
print(not True<5.0 or 3!=7.0 and 7==5)
\end{verbatim}

{\begin {center}\textbf{Завдання 4}\end {center}}

Що буде виведено за виконання? (повідомлення інтерпретатора про помилку позначати error)\par
\begin{verbatim}
print(7 < "5.0j")
\end{verbatim}

{\begin {center}\textbf{Завдання 5}\end {center}}

Що буде виведено за виконання? (повідомлення інтерпретатора про помилку позначати error)\par
\begin{verbatim}
print(6 <= 5.0 > 5  is  False)
\end{verbatim}

{\begin {center}\textbf{Завдання 6}\end {center}}

Що буде виведено за виконання? (повідомлення інтерпретатора про помилку позначати error)\par
\begin{verbatim}

def f(n):
    print("f", end="");
    return n<=-1

print(f(-3) and f(-4))
\end{verbatim}

{\begin {center}\textbf{Завдання 7}\end {center}}

Що буде виведено за виконання (повідомлення інтерпретатора про помилку позначати error)?\par
\begin{verbatim}
def f(a,b):
    c=67
    if b<=3:
        return 4
    if a<=5:
         c=3
    else: 
        c=2
    return c

print(f(0,0))
\end{verbatim}

{\begin {center}\textbf{Завдання 8}\end {center}}

Що буде виведено за виконання (повідомлення інтерпретатора про помилку позначати error)?\par
\begin{verbatim}
a, b, c = 7, 8, 6
def h(a, b, c):
    print(a, b, c, end=" ")

h(1, a=3)
print(a, b, c)
\end{verbatim}

{\begin {center}\textbf{Завдання 9}\end {center}}

Що буде виведено за виконання (повідомлення інтерпретатора про помилку позначати error)?\par
\begin{verbatim}
a,b,c=6,0,9
def f(a):
    a=4
    b+=1
    c=2
    return a+b+c

a,b,c=5,1,3
print(f(a),a,b,c)
\end{verbatim}

{\begin {center}\textbf{Завдання 10}\end {center}}


Написати функцiю f, що отримує щось та повертає щось. Невелика загадка все-таки має залишатися. Але нічого принципово нового відносно задач, що наявні у pdf-ках не планується.. Стиль запису умови також оцінюється.\par
\end {large}}
\newpage

{\Large{17.10.2024 Бета-версія: Контрольна робота \No 1, варіант  \No \textbf { 176 }\par}}
{
\begin {large}
{\begin {center}\textbf{Завдання 1}\end {center}}

Що буде виведено за виконання? (повідомлення інтерпретатора про помилку позначати error)\par
\begin{verbatim}
print(9/8*5*9)
\end{verbatim}

{\begin {center}\textbf{Завдання 2}\end {center}}

Що буде виведено за виконання? (повідомлення інтерпретатора про помилку позначати error)\par
\begin{verbatim}
print(4**-1.0*-2)
\end{verbatim}

{\begin {center}\textbf{Завдання 3}\end {center}}

Що буде виведено за виконання? (повідомлення інтерпретатора про помилку позначати error)\par
\begin{verbatim}
print(not 9.0<=6.0 and False>=4 or 6>9)
\end{verbatim}

{\begin {center}\textbf{Завдання 4}\end {center}}

Що буде виведено за виконання? (повідомлення інтерпретатора про помилку позначати error)\par
\begin{verbatim}
print("True" > "6")
\end{verbatim}

{\begin {center}\textbf{Завдання 5}\end {center}}

Що буде виведено за виконання? (повідомлення інтерпретатора про помилку позначати error)\par
\begin{verbatim}
print(3 != 7.0 < 7 > 4.0)
\end{verbatim}

{\begin {center}\textbf{Завдання 6}\end {center}}

Що буде виведено за виконання? (повідомлення інтерпретатора про помилку позначати error)\par
\begin{verbatim}

def f(n):
    print("f", end="");
    return n

print(f(0) or f(1))
\end{verbatim}

{\begin {center}\textbf{Завдання 7}\end {center}}

Що буде виведено за виконання (повідомлення інтерпретатора про помилку позначати error)?\par
\begin{verbatim}
def f(a,b):
    c=49
    if a>-3:
        return 1
    elif b<2:
         return 7
    else: 
        c=0
    return c

print(f(6,-2))
\end{verbatim}

{\begin {center}\textbf{Завдання 8}\end {center}}

Що буде виведено за виконання (повідомлення інтерпретатора про помилку позначати error)?\par
\begin{verbatim}
a, b, c = 6, 7, 8
def g(a, b, c=9):
    print(a, b, c, end=" ")

g(4, 2)
print(a, b, c)
\end{verbatim}

{\begin {center}\textbf{Завдання 9}\end {center}}

Що буде виведено за виконання (повідомлення інтерпретатора про помилку позначати error)?\par
\begin{verbatim}
a,b,c=9,0,2
def g(b):
    a=1
    b=5
    c=4
    return a+b+c

a,b,c=7,5,3
print(g(b),a,b,c)
\end{verbatim}

{\begin {center}\textbf{Завдання 10}\end {center}}


Написати функцiю f, що отримує щось та повертає щось. Невелика загадка все-таки має залишатися. Але нічого принципово нового відносно задач, що наявні у pdf-ках не планується.. Стиль запису умови також оцінюється.\par
\end {large}}
\newpage

{\Large{17.10.2024 Бета-версія: Контрольна робота \No 1, варіант  \No \textbf { 177 }\par}}
{
\begin {large}
{\begin {center}\textbf{Завдання 1}\end {center}}

Що буде виведено за виконання? (повідомлення інтерпретатора про помилку позначати error)\par
\begin{verbatim}
print(4//9/7*5)
\end{verbatim}

{\begin {center}\textbf{Завдання 2}\end {center}}

Що буде виведено за виконання? (повідомлення інтерпретатора про помилку позначати error)\par
\begin{verbatim}
print(2**1.0--1)
\end{verbatim}

{\begin {center}\textbf{Завдання 3}\end {center}}

Що буде виведено за виконання? (повідомлення інтерпретатора про помилку позначати error)\par
\begin{verbatim}
print(not 9.0>=True and 2!=7 or 9>5.0)
\end{verbatim}

{\begin {center}\textbf{Завдання 4}\end {center}}

Що буде виведено за виконання? (повідомлення інтерпретатора про помилку позначати error)\par
\begin{verbatim}
print(False == 5.0)
\end{verbatim}

{\begin {center}\textbf{Завдання 5}\end {center}}

Що буде виведено за виконання? (повідомлення інтерпретатора про помилку позначати error)\par
\begin{verbatim}
print(2 <= 6.0  is  True == 4)
\end{verbatim}

{\begin {center}\textbf{Завдання 6}\end {center}}

Що буде виведено за виконання? (повідомлення інтерпретатора про помилку позначати error)\par
\begin{verbatim}

def f(n):
    print("f", end="");
    return n!=-4

print(f(9) or f(-2))
\end{verbatim}

{\begin {center}\textbf{Завдання 7}\end {center}}

Що буде виведено за виконання (повідомлення інтерпретатора про помилку позначати error)?\par
\begin{verbatim}
def g(a,b):
    c=88
    if a==1:
        c=9
    if b!=4:
         return 5
    else: 
        return 6
    return c

print(g(8,1))
\end{verbatim}

{\begin {center}\textbf{Завдання 8}\end {center}}

Що буде виведено за виконання (повідомлення інтерпретатора про помилку позначати error)?\par
\begin{verbatim}
a, b, c = 6, 7, 9
def f(a, b, c=8):
    print(a, b, c, end=" ")

f(b=0, c=5, 1)
print(a, b, c)
\end{verbatim}

{\begin {center}\textbf{Завдання 9}\end {center}}

Що буде виведено за виконання (повідомлення інтерпретатора про помилку позначати error)?\par
\begin{verbatim}
a,b,c=8,2,4
def h(b):
    a=5
    b*=5
    c=3
    return a+b+c

a,b,c=1,8,2
print(h(a),a,b,c)
\end{verbatim}

{\begin {center}\textbf{Завдання 10}\end {center}}


Написати функцiю f, що отримує щось та повертає щось. Невелика загадка все-таки має залишатися. Але нічого принципово нового відносно задач, що наявні у pdf-ках не планується.. Стиль запису умови також оцінюється.\par
\end {large}}
\newpage

{\Large{17.10.2024 Бета-версія: Контрольна робота \No 1, варіант  \No \textbf { 178 }\par}}
{
\begin {large}
{\begin {center}\textbf{Завдання 1}\end {center}}

Що буде виведено за виконання? (повідомлення інтерпретатора про помилку позначати error)\par
\begin{verbatim}
print(3/4*4*3)
\end{verbatim}

{\begin {center}\textbf{Завдання 2}\end {center}}

Що буде виведено за виконання? (повідомлення інтерпретатора про помилку позначати error)\par
\begin{verbatim}
print(-3**True+2)
\end{verbatim}

{\begin {center}\textbf{Завдання 3}\end {center}}

Що буде виведено за виконання? (повідомлення інтерпретатора про помилку позначати error)\par
\begin{verbatim}
print(False<6.0 or not 5<=3 and 2.0==8)
\end{verbatim}

{\begin {center}\textbf{Завдання 4}\end {center}}

Що буде виведено за виконання? (повідомлення інтерпретатора про помилку позначати error)\par
\begin{verbatim}
print("2.0" >= "9j")
\end{verbatim}

{\begin {center}\textbf{Завдання 5}\end {center}}

Що буде виведено за виконання? (повідомлення інтерпретатора про помилку позначати error)\par
\begin{verbatim}
print(8.0 >= 9 > 5  is  7)
\end{verbatim}

{\begin {center}\textbf{Завдання 6}\end {center}}

Що буде виведено за виконання? (повідомлення інтерпретатора про помилку позначати error)\par
\begin{verbatim}

def f(n):
    print("f", end="");
    return n

print(f(2) and f(3))
\end{verbatim}

{\begin {center}\textbf{Завдання 7}\end {center}}

Що буде виведено за виконання (повідомлення інтерпретатора про помилку позначати error)?\par
\begin{verbatim}
def g(d):
    u=40
    if d>=4: 
        u=4
    if d<0:
         return 7
    else:
         return 2
    return u

print(g(-1))
\end{verbatim}

{\begin {center}\textbf{Завдання 8}\end {center}}

Що буде виведено за виконання (повідомлення інтерпретатора про помилку позначати error)?\par
\begin{verbatim}
a, b, c = 7, 8, 6
def f(a, b, c):
    print(a, b, c, end=" ")

f(3, 0)
print(a, b, c)
\end{verbatim}

{\begin {center}\textbf{Завдання 9}\end {center}}

Що буде виведено за виконання (повідомлення інтерпретатора про помилку позначати error)?\par
\begin{verbatim}
a,b,c=9,6,0
def g(b):
    a-=1
    b=4
    c=2
    return a+b+c

a,b,c=5,7,3
print(g(b),a,b,c)
\end{verbatim}

{\begin {center}\textbf{Завдання 10}\end {center}}


Написати функцiю f, що отримує щось та повертає щось. Невелика загадка все-таки має залишатися. Але нічого принципово нового відносно задач, що наявні у pdf-ках не планується.. Стиль запису умови також оцінюється.\par
\end {large}}
\newpage

{\Large{17.10.2024 Бета-версія: Контрольна робота \No 1, варіант  \No \textbf { 179 }\par}}
{
\begin {large}
{\begin {center}\textbf{Завдання 1}\end {center}}

Що буде виведено за виконання? (повідомлення інтерпретатора про помилку позначати error)\par
\begin{verbatim}
print(7/7%6*8)
\end{verbatim}

{\begin {center}\textbf{Завдання 2}\end {center}}

Що буде виведено за виконання? (повідомлення інтерпретатора про помилку позначати error)\par
\begin{verbatim}
print(9**-1.0*True)
\end{verbatim}

{\begin {center}\textbf{Завдання 3}\end {center}}

Що буде виведено за виконання? (повідомлення інтерпретатора про помилку позначати error)\par
\begin{verbatim}
print(not 4==4.0 and 6!=6 and True>8.0)
\end{verbatim}

{\begin {center}\textbf{Завдання 4}\end {center}}

Що буде виведено за виконання? (повідомлення інтерпретатора про помилку позначати error)\par
\begin{verbatim}
print(3.0 != "False")
\end{verbatim}

{\begin {center}\textbf{Завдання 5}\end {center}}

Що буде виведено за виконання? (повідомлення інтерпретатора про помилку позначати error)\par
\begin{verbatim}
print(False != 9.0 <= 2.0 == True)
\end{verbatim}

{\begin {center}\textbf{Завдання 6}\end {center}}

Що буде виведено за виконання? (повідомлення інтерпретатора про помилку позначати error)\par
\begin{verbatim}

def f(n):
    print("f", end="");
    return n

print(f(-6) and f(8))
\end{verbatim}

{\begin {center}\textbf{Завдання 7}\end {center}}

Що буде виведено за виконання (повідомлення інтерпретатора про помилку позначати error)?\par
\begin{verbatim}
def h(c):
    x=55
    if c: 
        x=3
    elif c!=1:
         return 5
    else:
         x=0
    return x

print(h(7))
\end{verbatim}

{\begin {center}\textbf{Завдання 8}\end {center}}

Що буде виведено за виконання (повідомлення інтерпретатора про помилку позначати error)?\par
\begin{verbatim}
a, b, c = 9, 8, 9
def h(a, b=7, c=6):
    print(a, b, c, end=" ")

h(4, 2, 5)
print(a, b, c)
\end{verbatim}

{\begin {center}\textbf{Завдання 9}\end {center}}

Що буде виведено за виконання (повідомлення інтерпретатора про помилку позначати error)?\par
\begin{verbatim}
a,b,c=6,4,2
def f(a):
    a=2
    b=4
    c=3
    return a+b+c

a,b,c=4,8,1
print(f(a),a,b,c)
\end{verbatim}

{\begin {center}\textbf{Завдання 10}\end {center}}


Написати функцiю f, що отримує щось та повертає щось. Невелика загадка все-таки має залишатися. Але нічого принципово нового відносно задач, що наявні у pdf-ках не планується.. Стиль запису умови також оцінюється.\par
\end {large}}
\newpage

{\Large{17.10.2024 Бета-версія: Контрольна робота \No 1, варіант  \No \textbf { 180 }\par}}
{
\begin {large}
{\begin {center}\textbf{Завдання 1}\end {center}}

Що буде виведено за виконання? (повідомлення інтерпретатора про помилку позначати error)\par
\begin{verbatim}
print(5//6//8*4)
\end{verbatim}

{\begin {center}\textbf{Завдання 2}\end {center}}

Що буде виведено за виконання? (повідомлення інтерпретатора про помилку позначати error)\par
\begin{verbatim}
print(1**1.0*False)
\end{verbatim}

{\begin {center}\textbf{Завдання 3}\end {center}}

Що буде виведено за виконання? (повідомлення інтерпретатора про помилку позначати error)\par
\begin{verbatim}
print(7.0>=2 or not 5<7 or 3.0<=False)
\end{verbatim}

{\begin {center}\textbf{Завдання 4}\end {center}}

Що буде виведено за виконання? (повідомлення інтерпретатора про помилку позначати error)\par
\begin{verbatim}
print(5 <= True)
\end{verbatim}

{\begin {center}\textbf{Завдання 5}\end {center}}

Що буде виведено за виконання? (повідомлення інтерпретатора про помилку позначати error)\par
\begin{verbatim}
print(8 < 3 >= 6 != 3.0)
\end{verbatim}

{\begin {center}\textbf{Завдання 6}\end {center}}

Що буде виведено за виконання? (повідомлення інтерпретатора про помилку позначати error)\par
\begin{verbatim}

def f(n):
    print("f", end="");
    return n>-2

print(f(-5) or f(7))
\end{verbatim}

{\begin {center}\textbf{Завдання 7}\end {center}}

Що буде виведено за виконання (повідомлення інтерпретатора про помилку позначати error)?\par
\begin{verbatim}
def h(a,b):
    c=62
    if a:
        c=8
    elif b>=0:
         c=3
    else: 
        return 4
    return c

print(h(3,-1))
\end{verbatim}

{\begin {center}\textbf{Завдання 8}\end {center}}

Що буде виведено за виконання (повідомлення інтерпретатора про помилку позначати error)?\par
\begin{verbatim}
a, b, c = 9, 7, 8
def g(a, b=6, c):
    print(a, b, c, end=" ")

g(3, c=4, b=1)
print(a, b, c)
\end{verbatim}

{\begin {center}\textbf{Завдання 9}\end {center}}

Що буде виведено за виконання (повідомлення інтерпретатора про помилку позначати error)?\par
\begin{verbatim}
a,b,c=7,3,5
def h(b):
    a=2
    b=5
    c=1
    return a+b+c

a,b,c=9,6,0
print(h(b),a,b,c)
\end{verbatim}

{\begin {center}\textbf{Завдання 10}\end {center}}


Написати функцiю f, що отримує щось та повертає щось. Невелика загадка все-таки має залишатися. Але нічого принципово нового відносно задач, що наявні у pdf-ках не планується.. Стиль запису умови також оцінюється.\par
\end {large}}
\newpage

{\Large{17.10.2024 Бета-версія: Контрольна робота \No 1, варіант  \No \textbf { 181 }\par}}
{
\begin {large}
{\begin {center}\textbf{Завдання 1}\end {center}}

Що буде виведено за виконання? (повідомлення інтерпретатора про помилку позначати error)\par
\begin{verbatim}
print(6/3%3*7)
\end{verbatim}

{\begin {center}\textbf{Завдання 2}\end {center}}

Що буде виведено за виконання? (повідомлення інтерпретатора про помилку позначати error)\par
\begin{verbatim}
print(2**-3-1)
\end{verbatim}

{\begin {center}\textbf{Завдання 3}\end {center}}

Що буде виведено за виконання? (повідомлення інтерпретатора про помилку позначати error)\par
\begin{verbatim}
print(8<3 or not 8.0<=False and 9==9.0)
\end{verbatim}

{\begin {center}\textbf{Завдання 4}\end {center}}

Що буде виведено за виконання? (повідомлення інтерпретатора про помилку позначати error)\par
\begin{verbatim}
print(6.0 != "8j")
\end{verbatim}

{\begin {center}\textbf{Завдання 5}\end {center}}

Що буде виведено за виконання? (повідомлення інтерпретатора про помилку позначати error)\par
\begin{verbatim}
print(True > 5.0 <= 3 < 3.0)
\end{verbatim}

{\begin {center}\textbf{Завдання 6}\end {center}}

Що буде виведено за виконання? (повідомлення інтерпретатора про помилку позначати error)\par
\begin{verbatim}

def f(n):
    print("f", end="");
    return n==-1

print(f(-8) and f(9))
\end{verbatim}

{\begin {center}\textbf{Завдання 7}\end {center}}

Що буде виведено за виконання (повідомлення інтерпретатора про помилку позначати error)?\par
\begin{verbatim}
def h(a,b):
    c=89
    if b:
        return 7
    if a>=-2:
         c=0
    else: 
        c=5
    return c

print(h(-2,5))
\end{verbatim}

{\begin {center}\textbf{Завдання 8}\end {center}}

Що буде виведено за виконання (повідомлення інтерпретатора про помилку позначати error)?\par
\begin{verbatim}
a, b, c = 7, 8, 6
def h(a, b, c=9):
    print(a, b, c, end=" ")

h(5, 2, a=0)
print(a, b, c)
\end{verbatim}

{\begin {center}\textbf{Завдання 9}\end {center}}

Що буде виведено за виконання (повідомлення інтерпретатора про помилку позначати error)?\par
\begin{verbatim}
a,b,c=0,4,3
def f(a):
    a=5
    b+=2
    c=4
    return a+b+c

a,b,c=7,6,8
print(f(b),a,b,c)
\end{verbatim}

{\begin {center}\textbf{Завдання 10}\end {center}}


Написати функцiю f, що отримує щось та повертає щось. Невелика загадка все-таки має залишатися. Але нічого принципово нового відносно задач, що наявні у pdf-ках не планується.. Стиль запису умови також оцінюється.\par
\end {large}}
\newpage

{\Large{17.10.2024 Бета-версія: Контрольна робота \No 1, варіант  \No \textbf { 182 }\par}}
{
\begin {large}
{\begin {center}\textbf{Завдання 1}\end {center}}

Що буде виведено за виконання? (повідомлення інтерпретатора про помилку позначати error)\par
\begin{verbatim}
print(8//5//9*6)
\end{verbatim}

{\begin {center}\textbf{Завдання 2}\end {center}}

Що буде виведено за виконання? (повідомлення інтерпретатора про помилку позначати error)\par
\begin{verbatim}
print(1**-0.5*False)
\end{verbatim}

{\begin {center}\textbf{Завдання 3}\end {center}}

Що буде виведено за виконання? (повідомлення інтерпретатора про помилку позначати error)\par
\begin{verbatim}
print(not 4>4.0 and 6!=True or 2>=6.0)
\end{verbatim}

{\begin {center}\textbf{Завдання 4}\end {center}}

Що буде виведено за виконання? (повідомлення інтерпретатора про помилку позначати error)\par
\begin{verbatim}
print(False < "1.0j")
\end{verbatim}

{\begin {center}\textbf{Завдання 5}\end {center}}

Що буде виведено за виконання? (повідомлення інтерпретатора про помилку позначати error)\par
\begin{verbatim}
print(4  is  9 == 6.0 >= False)
\end{verbatim}

{\begin {center}\textbf{Завдання 6}\end {center}}

Що буде виведено за виконання? (повідомлення інтерпретатора про помилку позначати error)\par
\begin{verbatim}

def f(n):
    print("f", end="");
    return n

print(f(-9) or f(7))
\end{verbatim}

{\begin {center}\textbf{Завдання 7}\end {center}}

Що буде виведено за виконання (повідомлення інтерпретатора про помилку позначати error)?\par
\begin{verbatim}
def g(b):
    u=14
    if b==5: 
        return 0
    elif b!=-2:
         u=2
    else:
         return 9
    return u

print(g(9))
\end{verbatim}

{\begin {center}\textbf{Завдання 8}\end {center}}

Що буде виведено за виконання (повідомлення інтерпретатора про помилку позначати error)?\par
\begin{verbatim}
a, b, c = 8, 6, 7
def f(a, b=9, c=7):
    print(a, b, c, end=" ")

f(2, c=0)
print(a, b, c)
\end{verbatim}

{\begin {center}\textbf{Завдання 9}\end {center}}

Що буде виведено за виконання (повідомлення інтерпретатора про помилку позначати error)?\par
\begin{verbatim}
a,b,c=9,5,2
def h(b):
    a*=3
    b=1
    c=2
    return a+b+c

a,b,c=1,8,0
print(h(b),a,b,c)
\end{verbatim}

{\begin {center}\textbf{Завдання 10}\end {center}}


Написати функцiю f, що отримує щось та повертає щось. Невелика загадка все-таки має залишатися. Але нічого принципово нового відносно задач, що наявні у pdf-ках не планується.. Стиль запису умови також оцінюється.\par
\end {large}}
\newpage

{\Large{17.10.2024 Бета-версія: Контрольна робота \No 1, варіант  \No \textbf { 183 }\par}}
{
\begin {large}
{\begin {center}\textbf{Завдання 1}\end {center}}

Що буде виведено за виконання? (повідомлення інтерпретатора про помилку позначати error)\par
\begin{verbatim}
print(9//5*8*8)
\end{verbatim}

{\begin {center}\textbf{Завдання 2}\end {center}}

Що буде виведено за виконання? (повідомлення інтерпретатора про помилку позначати error)\par
\begin{verbatim}
print(-2**-4+True)
\end{verbatim}

{\begin {center}\textbf{Завдання 3}\end {center}}

Що буде виведено за виконання? (повідомлення інтерпретатора про помилку позначати error)\par
\begin{verbatim}
print(not 8>7 or False==2.0 and 3.0<4)
\end{verbatim}

{\begin {center}\textbf{Завдання 4}\end {center}}

Що буде виведено за виконання? (повідомлення інтерпретатора про помилку позначати error)\par
\begin{verbatim}
print("True" >= 3)
\end{verbatim}

{\begin {center}\textbf{Завдання 5}\end {center}}

Що буде виведено за виконання? (повідомлення інтерпретатора про помилку позначати error)\par
\begin{verbatim}
print(8 <= 2.0 >= 6  is  2)
\end{verbatim}

{\begin {center}\textbf{Завдання 6}\end {center}}

Що буде виведено за виконання? (повідомлення інтерпретатора про помилку позначати error)\par
\begin{verbatim}

def f(n):
    print("f", end="");
    return n

print(f(-6) and f(2))
\end{verbatim}

{\begin {center}\textbf{Завдання 7}\end {center}}

Що буде виведено за виконання (повідомлення інтерпретатора про помилку позначати error)?\par
\begin{verbatim}
def f(a):
    y=47
    if a: 
        y=7
    if a>=-1:
         y=1
    else:
         return 5
    return y

print(f(0))
\end{verbatim}

{\begin {center}\textbf{Завдання 8}\end {center}}

Що буде виведено за виконання (повідомлення інтерпретатора про помилку позначати error)?\par
\begin{verbatim}
a, b, c = 8, 9, 6
def g(a, b, c):
    print(a, b, c, end=" ")

g(b=4, a=1, b=5)
print(a, b, c)
\end{verbatim}

{\begin {center}\textbf{Завдання 9}\end {center}}

Що буде виведено за виконання (повідомлення інтерпретатора про помилку позначати error)?\par
\begin{verbatim}
a,b,c=3,4,5
def g(a):
    a=3
    b-=4
    c=5
    return a+b+c

a,b,c=2,9,7
print(g(a),a,b,c)
\end{verbatim}

{\begin {center}\textbf{Завдання 10}\end {center}}


Написати функцiю f, що отримує щось та повертає щось. Невелика загадка все-таки має залишатися. Але нічого принципово нового відносно задач, що наявні у pdf-ках не планується.. Стиль запису умови також оцінюється.\par
\end {large}}
\newpage

{\Large{17.10.2024 Бета-версія: Контрольна робота \No 1, варіант  \No \textbf { 184 }\par}}
{
\begin {large}
{\begin {center}\textbf{Завдання 1}\end {center}}

Що буде виведено за виконання? (повідомлення інтерпретатора про помилку позначати error)\par
\begin{verbatim}
print(6/3/7*9)
\end{verbatim}

{\begin {center}\textbf{Завдання 2}\end {center}}

Що буде виведено за виконання? (повідомлення інтерпретатора про помилку позначати error)\par
\begin{verbatim}
print(9**0.5*2)
\end{verbatim}

{\begin {center}\textbf{Завдання 3}\end {center}}

Що буде виведено за виконання? (повідомлення інтерпретатора про помилку позначати error)\par
\begin{verbatim}
print(9>=3 and not True<=5 or 5.0!=7.0)
\end{verbatim}

{\begin {center}\textbf{Завдання 4}\end {center}}

Що буде виведено за виконання? (повідомлення інтерпретатора про помилку позначати error)\par
\begin{verbatim}
print("9.0" > 4.0)
\end{verbatim}

{\begin {center}\textbf{Завдання 5}\end {center}}

Що буде виведено за виконання? (повідомлення інтерпретатора про помилку позначати error)\par
\begin{verbatim}
print(7 > True != 7.0 == 9.0)
\end{verbatim}

{\begin {center}\textbf{Завдання 6}\end {center}}

Що буде виведено за виконання? (повідомлення інтерпретатора про помилку позначати error)\par
\begin{verbatim}

def f(n):
    print("f", end="");
    return n>4

print(f(3) or f(1))
\end{verbatim}

{\begin {center}\textbf{Завдання 7}\end {center}}

Що буде виведено за виконання (повідомлення інтерпретатора про помилку позначати error)?\par
\begin{verbatim}
def h(c):
    z=57
    if c<-3: 
        return 4
    if c<=3:
         z=6
    else:
         z=8
    return z

print(h(8))
\end{verbatim}

{\begin {center}\textbf{Завдання 8}\end {center}}

Що буде виведено за виконання (повідомлення інтерпретатора про помилку позначати error)?\par
\begin{verbatim}
a, b, c = 7, 8, 6
def h(a, b=9, c):
    print(a, b, c, end=" ")

h(a=3, 4, c=3)
print(a, b, c)
\end{verbatim}

{\begin {center}\textbf{Завдання 9}\end {center}}

Що буде виведено за виконання (повідомлення інтерпретатора про помилку позначати error)?\par
\begin{verbatim}
a,b,c=1,0,4
def h(a):
    a=3
    b=4
    c=1
    return a+b+c

a,b,c=6,9,3
print(h(b),a,b,c)
\end{verbatim}

{\begin {center}\textbf{Завдання 10}\end {center}}


Написати функцiю f, що отримує щось та повертає щось. Невелика загадка все-таки має залишатися. Але нічого принципово нового відносно задач, що наявні у pdf-ках не планується.. Стиль запису умови також оцінюється.\par
\end {large}}
\newpage

{\Large{17.10.2024 Бета-версія: Контрольна робота \No 1, варіант  \No \textbf { 185 }\par}}
{
\begin {large}
{\begin {center}\textbf{Завдання 1}\end {center}}

Що буде виведено за виконання? (повідомлення інтерпретатора про помилку позначати error)\par
\begin{verbatim}
print(3//9//6*4)
\end{verbatim}

{\begin {center}\textbf{Завдання 2}\end {center}}

Що буде виведено за виконання? (повідомлення інтерпретатора про помилку позначати error)\par
\begin{verbatim}
print(4**False+-1)
\end{verbatim}

{\begin {center}\textbf{Завдання 3}\end {center}}

Що буде виведено за виконання? (повідомлення інтерпретатора про помилку позначати error)\par
\begin{verbatim}
print(8<=4 and not 5>=True and 4.0>3.0)
\end{verbatim}

{\begin {center}\textbf{Завдання 4}\end {center}}

Що буде виведено за виконання? (повідомлення інтерпретатора про помилку позначати error)\par
\begin{verbatim}
print(False == "4")
\end{verbatim}

{\begin {center}\textbf{Завдання 5}\end {center}}

Що буде виведено за виконання? (повідомлення інтерпретатора про помилку позначати error)\par
\begin{verbatim}
print(5 < 4  is  False >= 3)
\end{verbatim}

{\begin {center}\textbf{Завдання 6}\end {center}}

Що буде виведено за виконання? (повідомлення інтерпретатора про помилку позначати error)\par
\begin{verbatim}

def f(n):
    print("f", end="");
    return n

print(f(-5) and f(-1))
\end{verbatim}

{\begin {center}\textbf{Завдання 7}\end {center}}

Що буде виведено за виконання (повідомлення інтерпретатора про помилку позначати error)?\par
\begin{verbatim}
def f(a,b):
    c=47
    if a!=-1:
        return 2
    elif a>-4:
         return 3
    else: 
        c=9
    return c

print(f(9,2))
\end{verbatim}

{\begin {center}\textbf{Завдання 8}\end {center}}

Що буде виведено за виконання (повідомлення інтерпретатора про помилку позначати error)?\par
\begin{verbatim}
a, b, c = 9, 8, 7
def g(a, b=6, c=9):
    print(a, b, c, end=" ")

g(0, 2, c=5)
print(a, b, c)
\end{verbatim}

{\begin {center}\textbf{Завдання 9}\end {center}}

Що буде виведено за виконання (повідомлення інтерпретатора про помилку позначати error)?\par
\begin{verbatim}
a,b,c=6,1,4
def g(b):
    a*=1
    b=1
    c=2
    return a+b+c

a,b,c=2,9,7
print(g(a),a,b,c)
\end{verbatim}

{\begin {center}\textbf{Завдання 10}\end {center}}


Написати функцiю f, що отримує щось та повертає щось. Невелика загадка все-таки має залишатися. Але нічого принципово нового відносно задач, що наявні у pdf-ках не планується.. Стиль запису умови також оцінюється.\par
\end {large}}
\newpage

{\Large{17.10.2024 Бета-версія: Контрольна робота \No 1, варіант  \No \textbf { 186 }\par}}
{
\begin {large}
{\begin {center}\textbf{Завдання 1}\end {center}}

Що буде виведено за виконання? (повідомлення інтерпретатора про помилку позначати error)\par
\begin{verbatim}
print(5/4%5*6)
\end{verbatim}

{\begin {center}\textbf{Завдання 2}\end {center}}

Що буде виведено за виконання? (повідомлення інтерпретатора про помилку позначати error)\par
\begin{verbatim}
print(4**-0.5*-2)
\end{verbatim}

{\begin {center}\textbf{Завдання 3}\end {center}}

Що буде виведено за виконання? (повідомлення інтерпретатора про помилку позначати error)\par
\begin{verbatim}
print(not 6.0==3 or 2<7 or 7.0!=False)
\end{verbatim}

{\begin {center}\textbf{Завдання 4}\end {center}}

Що буде виведено за виконання? (повідомлення інтерпретатора про помилку позначати error)\par
\begin{verbatim}
print("True" <= 2)
\end{verbatim}

{\begin {center}\textbf{Завдання 5}\end {center}}

Що буде виведено за виконання? (повідомлення інтерпретатора про помилку позначати error)\par
\begin{verbatim}
print(8.0 <= 2 > 4.0 < 9)
\end{verbatim}

{\begin {center}\textbf{Завдання 6}\end {center}}

Що буде виведено за виконання? (повідомлення інтерпретатора про помилку позначати error)\par
\begin{verbatim}

def f(n):
    print("f", end="");
    return n>=0

print(f(-4) or f(-8))
\end{verbatim}

{\begin {center}\textbf{Завдання 7}\end {center}}

Що буде виведено за виконання (повідомлення інтерпретатора про помилку позначати error)?\par
\begin{verbatim}
def g(a,b):
    c=22
    if b<-5:
        return 1
    elif b<=5:
         c=6
    else: 
        return 4
    return c

print(g(-1,-3))
\end{verbatim}

{\begin {center}\textbf{Завдання 8}\end {center}}

Що буде виведено за виконання (повідомлення інтерпретатора про помилку позначати error)?\par
\begin{verbatim}
a, b, c = 6, 7, 8
def f(a, b=8, c):
    print(a, b, c, end=" ")

f(1, a=5)
print(a, b, c)
\end{verbatim}

{\begin {center}\textbf{Завдання 9}\end {center}}

Що буде виведено за виконання (повідомлення інтерпретатора про помилку позначати error)?\par
\begin{verbatim}
a,b,c=8,5,6
def h(a):
    a=3
    b+=5
    c=4
    return a+b+c

a,b,c=0,3,1
print(h(b),a,b,c)
\end{verbatim}

{\begin {center}\textbf{Завдання 10}\end {center}}


Написати функцiю f, що отримує щось та повертає щось. Невелика загадка все-таки має залишатися. Але нічого принципово нового відносно задач, що наявні у pdf-ках не планується.. Стиль запису умови також оцінюється.\par
\end {large}}
\newpage

{\Large{17.10.2024 Бета-версія: Контрольна робота \No 1, варіант  \No \textbf { 187 }\par}}
{
\begin {large}
{\begin {center}\textbf{Завдання 1}\end {center}}

Що буде виведено за виконання? (повідомлення інтерпретатора про помилку позначати error)\par
\begin{verbatim}
print(4//8/9*3)
\end{verbatim}

{\begin {center}\textbf{Завдання 2}\end {center}}

Що буде виведено за виконання? (повідомлення інтерпретатора про помилку позначати error)\par
\begin{verbatim}
print(1**-1.0*True)
\end{verbatim}

{\begin {center}\textbf{Завдання 3}\end {center}}

Що буде виведено за виконання? (повідомлення інтерпретатора про помилку позначати error)\par
\begin{verbatim}
print(5.0<=6 and not 9>False or 6>=9.0)
\end{verbatim}

{\begin {center}\textbf{Завдання 4}\end {center}}

Що буде виведено за виконання? (повідомлення інтерпретатора про помилку позначати error)\par
\begin{verbatim}
print(7.0 > False)
\end{verbatim}

{\begin {center}\textbf{Завдання 5}\end {center}}

Що буде виведено за виконання? (повідомлення інтерпретатора про помилку позначати error)\par
\begin{verbatim}
print(5 != 6.0 == True <= 3.0)
\end{verbatim}

{\begin {center}\textbf{Завдання 6}\end {center}}

Що буде виведено за виконання? (повідомлення інтерпретатора про помилку позначати error)\par
\begin{verbatim}

def f(n):
    print("f", end="");
    return n

print(f(-2) or f(-7))
\end{verbatim}

{\begin {center}\textbf{Завдання 7}\end {center}}

Що буде виведено за виконання (повідомлення інтерпретатора про помилку позначати error)?\par
\begin{verbatim}
def f(d):
    w=18
    if d>-2: 
        return 3
    elif d>1:
         return 6
    else:
         w=4
    return w

print(f(-2))
\end{verbatim}

{\begin {center}\textbf{Завдання 8}\end {center}}

Що буде виведено за виконання (повідомлення інтерпретатора про помилку позначати error)?\par
\begin{verbatim}
a, b, c = 6, 9, 7
def h(a, b, c=9):
    print(a, b, c, end=" ")

h(a=1, 0, b=2)
print(a, b, c)
\end{verbatim}

{\begin {center}\textbf{Завдання 9}\end {center}}

Що буде виведено за виконання (повідомлення інтерпретатора про помилку позначати error)?\par
\begin{verbatim}
a,b,c=3,9,7
def f(a):
    a=3
    b-=2
    c=1
    return a+b+c

a,b,c=0,5,8
print(f(a),a,b,c)
\end{verbatim}

{\begin {center}\textbf{Завдання 10}\end {center}}


Написати функцiю f, що отримує щось та повертає щось. Невелика загадка все-таки має залишатися. Але нічого принципово нового відносно задач, що наявні у pdf-ках не планується.. Стиль запису умови також оцінюється.\par
\end {large}}
\newpage

{\Large{17.10.2024 Бета-версія: Контрольна робота \No 1, варіант  \No \textbf { 188 }\par}}
{
\begin {large}
{\begin {center}\textbf{Завдання 1}\end {center}}

Що буде виведено за виконання? (повідомлення інтерпретатора про помилку позначати error)\par
\begin{verbatim}
print(7/7*3*5)
\end{verbatim}

{\begin {center}\textbf{Завдання 2}\end {center}}

Що буде виведено за виконання? (повідомлення інтерпретатора про помилку позначати error)\par
\begin{verbatim}
print(9**0.5*2)
\end{verbatim}

{\begin {center}\textbf{Завдання 3}\end {center}}

Що буде виведено за виконання? (повідомлення інтерпретатора про помилку позначати error)\par
\begin{verbatim}
print(True<9 or not 3!=8.0 and 2.0==4)
\end{verbatim}

{\begin {center}\textbf{Завдання 4}\end {center}}

Що буде виведено за виконання? (повідомлення інтерпретатора про помилку позначати error)\par
\begin{verbatim}
print("True" != "1j")
\end{verbatim}

{\begin {center}\textbf{Завдання 5}\end {center}}

Що буде виведено за виконання? (повідомлення інтерпретатора про помилку позначати error)\par
\begin{verbatim}
print(6 > 7 != 2.0 == 8)
\end{verbatim}

{\begin {center}\textbf{Завдання 6}\end {center}}

Що буде виведено за виконання? (повідомлення інтерпретатора про помилку позначати error)\par
\begin{verbatim}

def f(n):
    print("f", end="");
    return n!=-2

print(f(4) and f(8))
\end{verbatim}

{\begin {center}\textbf{Завдання 7}\end {center}}

Що буде виведено за виконання (повідомлення інтерпретатора про помилку позначати error)?\par
\begin{verbatim}
def h(a):
    x=24
    if a!=-3: 
        return 5
    elif a==0:
         x=2
    else:
         return 3
    return x

print(h(2))
\end{verbatim}

{\begin {center}\textbf{Завдання 8}\end {center}}

Що буде виведено за виконання (повідомлення інтерпретатора про помилку позначати error)?\par
\begin{verbatim}
a, b, c = 7, 8, 6
def f(a, b, c):
    print(a, b, c, end=" ")

f(4, c=3, b=5)
print(a, b, c)
\end{verbatim}

{\begin {center}\textbf{Завдання 9}\end {center}}

Що буде виведено за виконання (повідомлення інтерпретатора про помилку позначати error)?\par
\begin{verbatim}
a,b,c=2,7,8
def f(a):
    a=5
    b=2
    c=5
    return a+b+c

a,b,c=5,3,2
print(f(a),a,b,c)
\end{verbatim}

{\begin {center}\textbf{Завдання 10}\end {center}}


Написати функцiю f, що отримує щось та повертає щось. Невелика загадка все-таки має залишатися. Але нічого принципово нового відносно задач, що наявні у pdf-ках не планується.. Стиль запису умови також оцінюється.\par
\end {large}}
\newpage

{\Large{17.10.2024 Бета-версія: Контрольна робота \No 1, варіант  \No \textbf { 189 }\par}}
{
\begin {large}
{\begin {center}\textbf{Завдання 1}\end {center}}

Що буде виведено за виконання? (повідомлення інтерпретатора про помилку позначати error)\par
\begin{verbatim}
print(8//6*4*7)
\end{verbatim}

{\begin {center}\textbf{Завдання 2}\end {center}}

Що буде виведено за виконання? (повідомлення інтерпретатора про помилку позначати error)\par
\begin{verbatim}
print(-1**2-1)
\end{verbatim}

{\begin {center}\textbf{Завдання 3}\end {center}}

Що буде виведено за виконання? (повідомлення інтерпретатора про помилку позначати error)\par
\begin{verbatim}
print(not 2==5 or 8.0>=True and 8!=9.0)
\end{verbatim}

{\begin {center}\textbf{Завдання 4}\end {center}}

Що буде виведено за виконання? (повідомлення інтерпретатора про помилку позначати error)\par
\begin{verbatim}
print("8.0" == 7)
\end{verbatim}

{\begin {center}\textbf{Завдання 5}\end {center}}

Що буде виведено за виконання? (повідомлення інтерпретатора про помилку позначати error)\par
\begin{verbatim}
print(7 < 5 >= 8.0  is  False)
\end{verbatim}

{\begin {center}\textbf{Завдання 6}\end {center}}

Що буде виведено за виконання? (повідомлення інтерпретатора про помилку позначати error)\par
\begin{verbatim}

def f(n):
    print("f", end="");
    return n<=-4

print(f(5) or f(-3))
\end{verbatim}

{\begin {center}\textbf{Завдання 7}\end {center}}

Що буде виведено за виконання (повідомлення інтерпретатора про помилку позначати error)?\par
\begin{verbatim}
def g(b):
    v=58
    if b: 
        v=1
    if b>=-4:
         return 0
    else:
         v=9
    return v

print(g(-8))
\end{verbatim}

{\begin {center}\textbf{Завдання 8}\end {center}}

Що буде виведено за виконання (повідомлення інтерпретатора про помилку позначати error)?\par
\begin{verbatim}
a, b, c = 9, 6, 8
def g(a, b, c=7):
    print(a, b, c, end=" ")

g(1, 4, 2)
print(a, b, c)
\end{verbatim}

{\begin {center}\textbf{Завдання 9}\end {center}}

Що буде виведено за виконання (повідомлення інтерпретатора про помилку позначати error)?\par
\begin{verbatim}
a,b,c=1,6,4
def f(b):
    a=5
    b+=4
    c=1
    return a+b+c

a,b,c=2,6,9
print(f(b),a,b,c)
\end{verbatim}

{\begin {center}\textbf{Завдання 10}\end {center}}


Написати функцiю f, що отримує щось та повертає щось. Невелика загадка все-таки має залишатися. Але нічого принципово нового відносно задач, що наявні у pdf-ках не планується.. Стиль запису умови також оцінюється.\par
\end {large}}
\newpage

{\Large{17.10.2024 Бета-версія: Контрольна робота \No 1, варіант  \No \textbf { 190 }\par}}
{
\begin {large}
{\begin {center}\textbf{Завдання 1}\end {center}}

Що буде виведено за виконання? (повідомлення інтерпретатора про помилку позначати error)\par
\begin{verbatim}
print(5/9//9*6)
\end{verbatim}

{\begin {center}\textbf{Завдання 2}\end {center}}

Що буде виведено за виконання? (повідомлення інтерпретатора про помилку позначати error)\par
\begin{verbatim}
print(4**-2+-1)
\end{verbatim}

{\begin {center}\textbf{Завдання 3}\end {center}}

Що буде виведено за виконання? (повідомлення інтерпретатора про помилку позначати error)\par
\begin{verbatim}
print(False<=4.0 and not 3.0<7 or 2>6)
\end{verbatim}

{\begin {center}\textbf{Завдання 4}\end {center}}

Що буде виведено за виконання? (повідомлення інтерпретатора про помилку позначати error)\par
\begin{verbatim}
print(5 >= False)
\end{verbatim}

{\begin {center}\textbf{Завдання 5}\end {center}}

Що буде виведено за виконання? (повідомлення інтерпретатора про помилку позначати error)\par
\begin{verbatim}
print(9  is  6 <= 5.0 == 7.0)
\end{verbatim}

{\begin {center}\textbf{Завдання 6}\end {center}}

Що буде виведено за виконання? (повідомлення інтерпретатора про помилку позначати error)\par
\begin{verbatim}

def f(n):
    print("f", end="");
    return n

print(f(0) and f(6))
\end{verbatim}

{\begin {center}\textbf{Завдання 7}\end {center}}

Що буде виведено за виконання (повідомлення інтерпретатора про помилку позначати error)?\par
\begin{verbatim}
def h(d):
    v=68
    if d: 
        v=7
    if d<=4:
         return 8
    else:
         return 8
    return v

print(h(4))
\end{verbatim}

{\begin {center}\textbf{Завдання 8}\end {center}}

Що буде виведено за виконання (повідомлення інтерпретатора про помилку позначати error)?\par
\begin{verbatim}
a, b, c = 6, 7, 9
def f(a, b, c):
    print(a, b, c, end=" ")

f(a=3, c=0, b=3)
print(a, b, c)
\end{verbatim}

{\begin {center}\textbf{Завдання 9}\end {center}}

Що буде виведено за виконання (повідомлення інтерпретатора про помилку позначати error)?\par
\begin{verbatim}
a,b,c=8,1,5
def g(b):
    a=1
    b=2
    c=4
    return a+b+c

a,b,c=4,0,9
print(g(a),a,b,c)
\end{verbatim}

{\begin {center}\textbf{Завдання 10}\end {center}}


Написати функцiю f, що отримує щось та повертає щось. Невелика загадка все-таки має залишатися. Але нічого принципово нового відносно задач, що наявні у pdf-ках не планується.. Стиль запису умови також оцінюється.\par
\end {large}}
\newpage

{\Large{17.10.2024 Бета-версія: Контрольна робота \No 1, варіант  \No \textbf { 191 }\par}}
{
\begin {large}
{\begin {center}\textbf{Завдання 1}\end {center}}

Що буде виведено за виконання? (повідомлення інтерпретатора про помилку позначати error)\par
\begin{verbatim}
print(6//6/6*7)
\end{verbatim}

{\begin {center}\textbf{Завдання 2}\end {center}}

Що буде виведено за виконання? (повідомлення інтерпретатора про помилку позначати error)\par
\begin{verbatim}
print(4**1.0*False)
\end{verbatim}

{\begin {center}\textbf{Завдання 3}\end {center}}

Що буде виведено за виконання? (повідомлення інтерпретатора про помилку позначати error)\par
\begin{verbatim}
print(7.0!=6.0 and not 8>=9 or True==5)
\end{verbatim}

{\begin {center}\textbf{Завдання 4}\end {center}}

Що буде виведено за виконання? (повідомлення інтерпретатора про помилку позначати error)\par
\begin{verbatim}
print("True" <= "1.0j")
\end{verbatim}

{\begin {center}\textbf{Завдання 5}\end {center}}

Що буде виведено за виконання? (повідомлення інтерпретатора про помилку позначати error)\par
\begin{verbatim}
print(3 > True >= 8 != 4)
\end{verbatim}

{\begin {center}\textbf{Завдання 6}\end {center}}

Що буде виведено за виконання? (повідомлення інтерпретатора про помилку позначати error)\par
\begin{verbatim}

def f(n):
    print("f", end="");
    return n

print(f(-6) or f(-1))
\end{verbatim}

{\begin {center}\textbf{Завдання 7}\end {center}}

Що буде виведено за виконання (повідомлення інтерпретатора про помилку позначати error)?\par
\begin{verbatim}
def f(c):
    w=23
    if c<2: 
        w=5
    elif c>=-1:
         return 9
    else:
         w=1
    return w

print(f(-2))
\end{verbatim}

{\begin {center}\textbf{Завдання 8}\end {center}}

Що буде виведено за виконання (повідомлення інтерпретатора про помилку позначати error)?\par
\begin{verbatim}
a, b, c = 8, 9, 8
def h(a, b=7, c=6):
    print(a, b, c, end=" ")

h(5, 4)
print(a, b, c)
\end{verbatim}

{\begin {center}\textbf{Завдання 9}\end {center}}

Що буде виведено за виконання (повідомлення інтерпретатора про помилку позначати error)?\par
\begin{verbatim}
a,b,c=7,1,0
def h(a):
    a-=2
    b=5
    c=4
    return a+b+c

a,b,c=2,8,5
print(h(b),a,b,c)
\end{verbatim}

{\begin {center}\textbf{Завдання 10}\end {center}}


Написати функцiю f, що отримує щось та повертає щось. Невелика загадка все-таки має залишатися. Але нічого принципово нового відносно задач, що наявні у pdf-ках не планується.. Стиль запису умови також оцінюється.\par
\end {large}}
\newpage

{\Large{17.10.2024 Бета-версія: Контрольна робота \No 1, варіант  \No \textbf { 192 }\par}}
{
\begin {large}
{\begin {center}\textbf{Завдання 1}\end {center}}

Що буде виведено за виконання? (повідомлення інтерпретатора про помилку позначати error)\par
\begin{verbatim}
print(7/8%5*5)
\end{verbatim}

{\begin {center}\textbf{Завдання 2}\end {center}}

Що буде виведено за виконання? (повідомлення інтерпретатора про помилку позначати error)\par
\begin{verbatim}
print(-1**-2--2)
\end{verbatim}

{\begin {center}\textbf{Завдання 3}\end {center}}

Що буде виведено за виконання? (повідомлення інтерпретатора про помилку позначати error)\par
\begin{verbatim}
print(False<4 or not 7<=5.0 and 2.0>3)
\end{verbatim}

{\begin {center}\textbf{Завдання 4}\end {center}}

Що буде виведено за виконання? (повідомлення інтерпретатора про помилку позначати error)\par
\begin{verbatim}
print(8.0 < "4")
\end{verbatim}

{\begin {center}\textbf{Завдання 5}\end {center}}

Що буде виведено за виконання? (повідомлення інтерпретатора про помилку позначати error)\par
\begin{verbatim}
print(9.0 < 2 != 4.0 <= False)
\end{verbatim}

{\begin {center}\textbf{Завдання 6}\end {center}}

Що буде виведено за виконання? (повідомлення інтерпретатора про помилку позначати error)\par
\begin{verbatim}

def f(n):
    print("f", end="");
    return n<3

print(f(-8) and f(-4))
\end{verbatim}

{\begin {center}\textbf{Завдання 7}\end {center}}

Що буде виведено за виконання (повідомлення інтерпретатора про помилку позначати error)?\par
\begin{verbatim}
def f(a,b):
    c=78
    if a:
        c=8
    if a==-2:
         c=4
    else: 
        return 5
    return c

print(f(-9,7))
\end{verbatim}

{\begin {center}\textbf{Завдання 8}\end {center}}

Що буде виведено за виконання (повідомлення інтерпретатора про помилку позначати error)?\par
\begin{verbatim}
a, b, c = 9, 7, 6
def g(a, b=8, c):
    print(a, b, c, end=" ")

g(b=0, c=2, 1)
print(a, b, c)
\end{verbatim}

{\begin {center}\textbf{Завдання 9}\end {center}}

Що буде виведено за виконання (повідомлення інтерпретатора про помилку позначати error)?\par
\begin{verbatim}
a,b,c=3,4,8
def g(b):
    a=3
    b*=3
    c=5
    return a+b+c

a,b,c=7,2,0
print(g(a),a,b,c)
\end{verbatim}

{\begin {center}\textbf{Завдання 10}\end {center}}


Написати функцiю f, що отримує щось та повертає щось. Невелика загадка все-таки має залишатися. Але нічого принципово нового відносно задач, що наявні у pdf-ках не планується.. Стиль запису умови також оцінюється.\par
\end {large}}
\newpage

{\Large{17.10.2024 Бета-версія: Контрольна робота \No 1, варіант  \No \textbf { 193 }\par}}
{
\begin {large}
{\begin {center}\textbf{Завдання 1}\end {center}}

Що буде виведено за виконання? (повідомлення інтерпретатора про помилку позначати error)\par
\begin{verbatim}
print(9//7//4*9)
\end{verbatim}

{\begin {center}\textbf{Завдання 2}\end {center}}

Що буде виведено за виконання? (повідомлення інтерпретатора про помилку позначати error)\par
\begin{verbatim}
print(4**1.0*False)
\end{verbatim}

{\begin {center}\textbf{Завдання 3}\end {center}}

Що буде виведено за виконання? (повідомлення інтерпретатора про помилку позначати error)\par
\begin{verbatim}
print(False==2 and not 8.0>=4.0 or 7>4)
\end{verbatim}

{\begin {center}\textbf{Завдання 4}\end {center}}

Що буде виведено за виконання? (повідомлення інтерпретатора про помилку позначати error)\par
\begin{verbatim}
print(3 > 3.0)
\end{verbatim}

{\begin {center}\textbf{Завдання 5}\end {center}}

Що буде виведено за виконання? (повідомлення інтерпретатора про помилку позначати error)\par
\begin{verbatim}
print(6 >= 4 == 2.0  is  5.0)
\end{verbatim}

{\begin {center}\textbf{Завдання 6}\end {center}}

Що буде виведено за виконання? (повідомлення інтерпретатора про помилку позначати error)\par
\begin{verbatim}

def f(n):
    print("f", end="");
    return n==-3

print(f(5) or f(-9))
\end{verbatim}

{\begin {center}\textbf{Завдання 7}\end {center}}

Що буде виведено за виконання (повідомлення інтерпретатора про помилку позначати error)?\par
\begin{verbatim}
def g(b):
    x=35
    if b<5: 
        return 4
    if b<=-5:
         x=0
    else:
         x=7
    return x

print(g(9))
\end{verbatim}

{\begin {center}\textbf{Завдання 8}\end {center}}

Що буде виведено за виконання (повідомлення інтерпретатора про помилку позначати error)?\par
\begin{verbatim}
a, b, c = 7, 9, 8
def g(a, b=6, c=8):
    print(a, b, c, end=" ")

g(b=5, c=2, 4)
print(a, b, c)
\end{verbatim}

{\begin {center}\textbf{Завдання 9}\end {center}}

Що буде виведено за виконання (повідомлення інтерпретатора про помилку позначати error)?\par
\begin{verbatim}
a,b,c=1,3,6
def f(b):
    a=4
    b+=2
    c=1
    return a+b+c

a,b,c=4,5,9
print(f(b),a,b,c)
\end{verbatim}

{\begin {center}\textbf{Завдання 10}\end {center}}


Написати функцiю f, що отримує щось та повертає щось. Невелика загадка все-таки має залишатися. Але нічого принципово нового відносно задач, що наявні у pdf-ках не планується.. Стиль запису умови також оцінюється.\par
\end {large}}
\newpage

{\Large{17.10.2024 Бета-версія: Контрольна робота \No 1, варіант  \No \textbf { 194 }\par}}
{
\begin {large}
{\begin {center}\textbf{Завдання 1}\end {center}}

Що буде виведено за виконання? (повідомлення інтерпретатора про помилку позначати error)\par
\begin{verbatim}
print(3/4/8*4)
\end{verbatim}

{\begin {center}\textbf{Завдання 2}\end {center}}

Що буде виведено за виконання? (повідомлення інтерпретатора про помилку позначати error)\par
\begin{verbatim}
print(2**-3-True)
\end{verbatim}

{\begin {center}\textbf{Завдання 3}\end {center}}

Що буде виведено за виконання? (повідомлення інтерпретатора про помилку позначати error)\par
\begin{verbatim}
print(not True!=8 or 9<3.0 and 5<=5.0)
\end{verbatim}

{\begin {center}\textbf{Завдання 4}\end {center}}

Що буде виведено за виконання? (повідомлення інтерпретатора про помилку позначати error)\par
\begin{verbatim}
print("5.0j" != "9")
\end{verbatim}

{\begin {center}\textbf{Завдання 5}\end {center}}

Що буде виведено за виконання? (повідомлення інтерпретатора про помилку позначати error)\par
\begin{verbatim}
print(9 < False > 3 != 8.0)
\end{verbatim}

{\begin {center}\textbf{Завдання 6}\end {center}}

Що буде виведено за виконання? (повідомлення інтерпретатора про помилку позначати error)\par
\begin{verbatim}

def f(n):
    print("f", end="");
    return n

print(f(-2) and f(4))
\end{verbatim}

{\begin {center}\textbf{Завдання 7}\end {center}}

Що буде виведено за виконання (повідомлення інтерпретатора про помилку позначати error)?\par
\begin{verbatim}
def g(a,b):
    c=93
    if a>=-3:
        c=6
    if b!=-1:
         return 7
    else: 
        return 3
    return c

print(g(-6,-6))
\end{verbatim}

{\begin {center}\textbf{Завдання 8}\end {center}}

Що буде виведено за виконання (повідомлення інтерпретатора про помилку позначати error)?\par
\begin{verbatim}
a, b, c = 7, 9, 6
def h(a, b, c=8):
    print(a, b, c, end=" ")

h(a=1, 0, b=3)
print(a, b, c)
\end{verbatim}

{\begin {center}\textbf{Завдання 9}\end {center}}

Що буде виведено за виконання (повідомлення інтерпретатора про помилку позначати error)?\par
\begin{verbatim}
a,b,c=6,7,2
def h(b):
    a=3
    b=2
    c=5
    return a+b+c

a,b,c=7,4,5
print(h(a),a,b,c)
\end{verbatim}

{\begin {center}\textbf{Завдання 10}\end {center}}


Написати функцiю f, що отримує щось та повертає щось. Невелика загадка все-таки має залишатися. Але нічого принципово нового відносно задач, що наявні у pdf-ках не планується.. Стиль запису умови також оцінюється.\par
\end {large}}
\newpage

{\Large{17.10.2024 Бета-версія: Контрольна робота \No 1, варіант  \No \textbf { 195 }\par}}
{
\begin {large}
{\begin {center}\textbf{Завдання 1}\end {center}}

Що буде виведено за виконання? (повідомлення інтерпретатора про помилку позначати error)\par
\begin{verbatim}
print(4/5%3*3)
\end{verbatim}

{\begin {center}\textbf{Завдання 2}\end {center}}

Що буде виведено за виконання? (повідомлення інтерпретатора про помилку позначати error)\par
\begin{verbatim}
print(1**-0.5*-2)
\end{verbatim}

{\begin {center}\textbf{Завдання 3}\end {center}}

Що буде виведено за виконання? (повідомлення інтерпретатора про помилку позначати error)\par
\begin{verbatim}
print(False<=9.0 and not 6>=3 or 2.0<8)
\end{verbatim}

{\begin {center}\textbf{Завдання 4}\end {center}}

Що буде виведено за виконання? (повідомлення інтерпретатора про помилку позначати error)\par
\begin{verbatim}
print(False >= "True")
\end{verbatim}

{\begin {center}\textbf{Завдання 5}\end {center}}

Що буде виведено за виконання? (повідомлення інтерпретатора про помилку позначати error)\par
\begin{verbatim}
print(True > 6.0 < 7 <= 8)
\end{verbatim}

{\begin {center}\textbf{Завдання 6}\end {center}}

Що буде виведено за виконання? (повідомлення інтерпретатора про помилку позначати error)\par
\begin{verbatim}

def f(n):
    print("f", end="");
    return n

print(f(8) or f(6))
\end{verbatim}

{\begin {center}\textbf{Завдання 7}\end {center}}

Що буде виведено за виконання (повідомлення інтерпретатора про помилку позначати error)?\par
\begin{verbatim}
def h(a,b):
    c=79
    if b<1:
        c=8
    elif b<=-4:
         return 0
    else: 
        c=2
    return c

print(h(-7,-5))
\end{verbatim}

{\begin {center}\textbf{Завдання 8}\end {center}}

Що буде виведено за виконання (повідомлення інтерпретатора про помилку позначати error)?\par
\begin{verbatim}
a, b, c = 9, 6, 7
def f(a, b, c):
    print(a, b, c, end=" ")

f(c=0, a=1, a=5)
print(a, b, c)
\end{verbatim}

{\begin {center}\textbf{Завдання 9}\end {center}}

Що буде виведено за виконання (повідомлення інтерпретатора про помилку позначати error)?\par
\begin{verbatim}
a,b,c=2,0,9
def h(a):
    a=5
    b-=4
    c=3
    return a+b+c

a,b,c=7,5,8
print(h(b),a,b,c)
\end{verbatim}

{\begin {center}\textbf{Завдання 10}\end {center}}


Написати функцiю f, що отримує щось та повертає щось. Невелика загадка все-таки має залишатися. Але нічого принципово нового відносно задач, що наявні у pdf-ках не планується.. Стиль запису умови також оцінюється.\par
\end {large}}
\newpage

{\Large{17.10.2024 Бета-версія: Контрольна робота \No 1, варіант  \No \textbf { 196 }\par}}
{
\begin {large}
{\begin {center}\textbf{Завдання 1}\end {center}}

Що буде виведено за виконання? (повідомлення інтерпретатора про помилку позначати error)\par
\begin{verbatim}
print(8//3*7*8)
\end{verbatim}

{\begin {center}\textbf{Завдання 2}\end {center}}

Що буде виведено за виконання? (повідомлення інтерпретатора про помилку позначати error)\par
\begin{verbatim}
print(9**0.5*1)
\end{verbatim}

{\begin {center}\textbf{Завдання 3}\end {center}}

Що буде виведено за виконання? (повідомлення інтерпретатора про помилку позначати error)\par
\begin{verbatim}
print(not 4>6.0 or 7.0!=7 and 3==True)
\end{verbatim}

{\begin {center}\textbf{Завдання 4}\end {center}}

Що буде виведено за виконання? (повідомлення інтерпретатора про помилку позначати error)\par
\begin{verbatim}
print(True <= "7j")
\end{verbatim}

{\begin {center}\textbf{Завдання 5}\end {center}}

Що буде виведено за виконання? (повідомлення інтерпретатора про помилку позначати error)\par
\begin{verbatim}
print(2 >= 7.0 == False  is  5)
\end{verbatim}

{\begin {center}\textbf{Завдання 6}\end {center}}

Що буде виведено за виконання? (повідомлення інтерпретатора про помилку позначати error)\par
\begin{verbatim}

def f(n):
    print("f", end="");
    return n>=2

print(f(9) and f(0))
\end{verbatim}

{\begin {center}\textbf{Завдання 7}\end {center}}

Що буде виведено за виконання (повідомлення інтерпретатора про помилку позначати error)?\par
\begin{verbatim}
def h(a,b):
    c=41
    if b>3:
        return 9
    elif a==2:
         c=1
    else: 
        c=2
    return c

print(h(8,3))
\end{verbatim}

{\begin {center}\textbf{Завдання 8}\end {center}}

Що буде виведено за виконання (повідомлення інтерпретатора про помилку позначати error)?\par
\begin{verbatim}
a, b, c = 7, 6, 9
def g(a, b=8, c):
    print(a, b, c, end=" ")

g(3, 4, 2)
print(a, b, c)
\end{verbatim}

{\begin {center}\textbf{Завдання 9}\end {center}}

Що буде виведено за виконання (повідомлення інтерпретатора про помилку позначати error)?\par
\begin{verbatim}
a,b,c=4,1,3
def g(a):
    a*=1
    b=2
    c=1
    return a+b+c

a,b,c=6,0,4
print(g(b),a,b,c)
\end{verbatim}

{\begin {center}\textbf{Завдання 10}\end {center}}


Написати функцiю f, що отримує щось та повертає щось. Невелика загадка все-таки має залишатися. Але нічого принципово нового відносно задач, що наявні у pdf-ках не планується.. Стиль запису умови також оцінюється.\par
\end {large}}
\newpage

{\Large{17.10.2024 Бета-версія: Контрольна робота \No 1, варіант  \No \textbf { 197 }\par}}
{
\begin {large}
{\begin {center}\textbf{Завдання 1}\end {center}}

Що буде виведено за виконання? (повідомлення інтерпретатора про помилку позначати error)\par
\begin{verbatim}
print(8//7/7*9)
\end{verbatim}

{\begin {center}\textbf{Завдання 2}\end {center}}

Що буде виведено за виконання? (повідомлення інтерпретатора про помилку позначати error)\par
\begin{verbatim}
print(-2**4+-1)
\end{verbatim}

{\begin {center}\textbf{Завдання 3}\end {center}}

Що буде виведено за виконання? (повідомлення інтерпретатора про помилку позначати error)\par
\begin{verbatim}
print(not 6>=5 and 9==True or 8.0!=6.0)
\end{verbatim}

{\begin {center}\textbf{Завдання 4}\end {center}}

Що буде виведено за виконання? (повідомлення інтерпретатора про помилку позначати error)\par
\begin{verbatim}
print(1 < "7.0")
\end{verbatim}

{\begin {center}\textbf{Завдання 5}\end {center}}

Що буде виведено за виконання? (повідомлення інтерпретатора про помилку позначати error)\par
\begin{verbatim}
print(2 < 4.0 == 4 >= 3.0)
\end{verbatim}

{\begin {center}\textbf{Завдання 6}\end {center}}

Що буде виведено за виконання? (повідомлення інтерпретатора про помилку позначати error)\par
\begin{verbatim}

def f(n):
    print("f", end="");
    return n<1

print(f(-3) and f(3))
\end{verbatim}

{\begin {center}\textbf{Завдання 7}\end {center}}

Що буде виведено за виконання (повідомлення інтерпретатора про помилку позначати error)?\par
\begin{verbatim}
def f(a,b):
    c=75
    if a>=4:
        return 3
    if b!=0:
         return 8
    else: 
        c=4
    return c

print(f(-9,-2))
\end{verbatim}

{\begin {center}\textbf{Завдання 8}\end {center}}

Що буде виведено за виконання (повідомлення інтерпретатора про помилку позначати error)?\par
\begin{verbatim}
a, b, c = 9, 8, 7
def f(a, b, c=6):
    print(a, b, c, end=" ")

f(0, c=2)
print(a, b, c)
\end{verbatim}

{\begin {center}\textbf{Завдання 9}\end {center}}

Що буде виведено за виконання (повідомлення інтерпретатора про помилку позначати error)?\par
\begin{verbatim}
a,b,c=9,6,1
def f(a):
    a=4
    b=3
    c=1
    return a+b+c

a,b,c=3,0,8
print(f(b),a,b,c)
\end{verbatim}

{\begin {center}\textbf{Завдання 10}\end {center}}


Написати функцiю f, що отримує щось та повертає щось. Невелика загадка все-таки має залишатися. Але нічого принципово нового відносно задач, що наявні у pdf-ках не планується.. Стиль запису умови також оцінюється.\par
\end {large}}
\newpage

{\Large{17.10.2024 Бета-версія: Контрольна робота \No 1, варіант  \No \textbf { 198 }\par}}
{
\begin {large}
{\begin {center}\textbf{Завдання 1}\end {center}}

Що буде виведено за виконання? (повідомлення інтерпретатора про помилку позначати error)\par
\begin{verbatim}
print(9/4%4*6)
\end{verbatim}

{\begin {center}\textbf{Завдання 2}\end {center}}

Що буде виведено за виконання? (повідомлення інтерпретатора про помилку позначати error)\par
\begin{verbatim}
print(1**-2+2)
\end{verbatim}

{\begin {center}\textbf{Завдання 3}\end {center}}

Що буде виведено за виконання? (повідомлення інтерпретатора про помилку позначати error)\par
\begin{verbatim}
print(7.0>False or not 2<=9.0 and 7<3)
\end{verbatim}

{\begin {center}\textbf{Завдання 4}\end {center}}

Що буде виведено за виконання? (повідомлення інтерпретатора про помилку позначати error)\par
\begin{verbatim}
print("False" == 4.0)
\end{verbatim}

{\begin {center}\textbf{Завдання 5}\end {center}}

Що буде виведено за виконання? (повідомлення інтерпретатора про помилку позначати error)\par
\begin{verbatim}
print(9.0 <= 9 != True  is  8)
\end{verbatim}

{\begin {center}\textbf{Завдання 6}\end {center}}

Що буде виведено за виконання? (повідомлення інтерпретатора про помилку позначати error)\par
\begin{verbatim}

def f(n):
    print("f", end="");
    return n

print(f(-7) or f(7))
\end{verbatim}

{\begin {center}\textbf{Завдання 7}\end {center}}

Що буде виведено за виконання (повідомлення інтерпретатора про помилку позначати error)?\par
\begin{verbatim}
def h(a):
    z=94
    if a: 
        return 6
    elif a>-4:
         return 2
    else:
         z=3
    return z

print(h(8))
\end{verbatim}

{\begin {center}\textbf{Завдання 8}\end {center}}

Що буде виведено за виконання (повідомлення інтерпретатора про помилку позначати error)?\par
\begin{verbatim}
a, b, c = 8, 9, 7
def h(a, b=6, c):
    print(a, b, c, end=" ")

h(3, 1)
print(a, b, c)
\end{verbatim}

{\begin {center}\textbf{Завдання 9}\end {center}}

Що буде виведено за виконання (повідомлення інтерпретатора про помилку позначати error)?\par
\begin{verbatim}
a,b,c=9,3,2
def g(b):
    a*=3
    b=5
    c=2
    return a+b+c

a,b,c=8,7,5
print(g(a),a,b,c)
\end{verbatim}

{\begin {center}\textbf{Завдання 10}\end {center}}


Написати функцiю f, що отримує щось та повертає щось. Невелика загадка все-таки має залишатися. Але нічого принципово нового відносно задач, що наявні у pdf-ках не планується.. Стиль запису умови також оцінюється.\par
\end {large}}
\newpage

{\Large{17.10.2024 Бета-версія: Контрольна робота \No 1, варіант  \No \textbf { 199 }\par}}
{
\begin {large}
{\begin {center}\textbf{Завдання 1}\end {center}}

Що буде виведено за виконання? (повідомлення інтерпретатора про помилку позначати error)\par
\begin{verbatim}
print(4//9*3*4)
\end{verbatim}

{\begin {center}\textbf{Завдання 2}\end {center}}

Що буде виведено за виконання? (повідомлення інтерпретатора про помилку позначати error)\par
\begin{verbatim}
print(2**True--1)
\end{verbatim}

{\begin {center}\textbf{Завдання 3}\end {center}}

Що буде виведено за виконання? (повідомлення інтерпретатора про помилку позначати error)\par
\begin{verbatim}
print(not 3.0<=4 or False<5.0 and 2>5)
\end{verbatim}

{\begin {center}\textbf{Завдання 4}\end {center}}

Що буде виведено за виконання? (повідомлення інтерпретатора про помилку позначати error)\par
\begin{verbatim}
print(8 >= 9.0)
\end{verbatim}

{\begin {center}\textbf{Завдання 5}\end {center}}

Що буде виведено за виконання? (повідомлення інтерпретатора про помилку позначати error)\par
\begin{verbatim}
print(6 > 3 >= True == 7.0)
\end{verbatim}

{\begin {center}\textbf{Завдання 6}\end {center}}

Що буде виведено за виконання? (повідомлення інтерпретатора про помилку позначати error)\par
\begin{verbatim}

def f(n):
    print("f", end="");
    return n

print(f(-5) and f(1))
\end{verbatim}

{\begin {center}\textbf{Завдання 7}\end {center}}

Що буде виведено за виконання (повідомлення інтерпретатора про помилку позначати error)?\par
\begin{verbatim}
def g(a,b):
    c=13
    if b:
        return 7
    elif a<-5:
         c=9
    else: 
        c=6
    return c

print(g(2,2))
\end{verbatim}

{\begin {center}\textbf{Завдання 8}\end {center}}

Що буде виведено за виконання (повідомлення інтерпретатора про помилку позначати error)?\par
\begin{verbatim}
a, b, c = 6, 7, 9
def f(a, b, c):
    print(a, b, c, end=" ")

f(4, c=5, b=3)
print(a, b, c)
\end{verbatim}

{\begin {center}\textbf{Завдання 9}\end {center}}

Що буде виведено за виконання (повідомлення інтерпретатора про помилку позначати error)?\par
\begin{verbatim}
a,b,c=6,1,4
def h(b):
    a+=4
    b=3
    c=5
    return a+b+c

a,b,c=6,0,2
print(h(a),a,b,c)
\end{verbatim}

{\begin {center}\textbf{Завдання 10}\end {center}}


Написати функцiю f, що отримує щось та повертає щось. Невелика загадка все-таки має залишатися. Але нічого принципово нового відносно задач, що наявні у pdf-ках не планується.. Стиль запису умови також оцінюється.\par
\end {large}}
\newpage

{\Large{17.10.2024 Бета-версія: Контрольна робота \No 1, варіант  \No \textbf { 200 }\par}}
{
\begin {large}
{\begin {center}\textbf{Завдання 1}\end {center}}

Що буде виведено за виконання? (повідомлення інтерпретатора про помилку позначати error)\par
\begin{verbatim}
print(5/3//6*7)
\end{verbatim}

{\begin {center}\textbf{Завдання 2}\end {center}}

Що буде виведено за виконання? (повідомлення інтерпретатора про помилку позначати error)\par
\begin{verbatim}
print(1**-1.0*True)
\end{verbatim}

{\begin {center}\textbf{Завдання 3}\end {center}}

Що буде виведено за виконання? (повідомлення інтерпретатора про помилку позначати error)\par
\begin{verbatim}
print(not 2.0>=6 and 9!=8 or True==4.0)
\end{verbatim}

{\begin {center}\textbf{Завдання 4}\end {center}}

Що буде виведено за виконання? (повідомлення інтерпретатора про помилку позначати error)\par
\begin{verbatim}
print("2.0" > "True")
\end{verbatim}

{\begin {center}\textbf{Завдання 5}\end {center}}

Що буде виведено за виконання? (повідомлення інтерпретатора про помилку позначати error)\par
\begin{verbatim}
print(9.0 > 7 != 2.0  is  5)
\end{verbatim}

{\begin {center}\textbf{Завдання 6}\end {center}}

Що буде виведено за виконання? (повідомлення інтерпретатора про помилку позначати error)\par
\begin{verbatim}

def f(n):
    print("f", end="");
    return n<=0

print(f(2) or f(-6))
\end{verbatim}

{\begin {center}\textbf{Завдання 7}\end {center}}

Що буде виведено за виконання (повідомлення інтерпретатора про помилку позначати error)?\par
\begin{verbatim}
def h(a,b):
    c=72
    if b==2:
        return 5
    if a<=4:
         c=1
    else: 
        return 0
    return c

print(h(9,-1))
\end{verbatim}

{\begin {center}\textbf{Завдання 8}\end {center}}

Що буде виведено за виконання (повідомлення інтерпретатора про помилку позначати error)?\par
\begin{verbatim}
a, b, c = 8, 6, 7
def h(a, b=8, c=9):
    print(a, b, c, end=" ")

h(5, 4, b=2)
print(a, b, c)
\end{verbatim}

{\begin {center}\textbf{Завдання 9}\end {center}}

Що буде виведено за виконання (повідомлення інтерпретатора про помилку позначати error)?\par
\begin{verbatim}
a,b,c=5,3,7
def f(a):
    a-=2
    b=1
    c=4
    return a+b+c

a,b,c=8,9,1
print(f(a),a,b,c)
\end{verbatim}

{\begin {center}\textbf{Завдання 10}\end {center}}


Написати функцiю f, що отримує щось та повертає щось. Невелика загадка все-таки має залишатися. Але нічого принципово нового відносно задач, що наявні у pdf-ках не планується.. Стиль запису умови також оцінюється.\par
\end {large}}
\newpage

%end
\end{document}

